SV
00101 oM namo buddhAya / 
00102 	vidhUtakalpanAjAlagambhIraudAramUrtaye /  
00103 	namaH samantabhadrAya samantaspharaNatviSe //
00104 	prAyaH prAkRtasaktir apratibalaprajJo janaH kevalaM / 
00105 	na anarthy eva subhASitaiH parigato vidveSTy api IrSyAmalaiH //
00106 	tena ayaM na paraupakAra iti naz cintA api cetaz ciraM /
00107 	sUktAbhyAsavivardhitavyasanam ity atra anubaddhaspRham //  
00108 arthAnarthavivecanasya anumAnAzrayatvAt tadvipratipattes tadvyavasthApanAya
00109 Aha /  
00110 	pakSadharmas tadaMzena vyApto hetus tridhA eva saH / 
00111 	avinAbhAvaniyamAd hetvAbhAsAs tato apare //1//
00112 pakSo dharmI / avayave samudAyaupacArAt / prayojanAbhAvAd anupacAra
00113 iti cet / na / sarvadharmidharmapratiSedhArthatvAt / tadekadezatvAt
00114 tadupacArayogyadharmidharmapratipattyartham / tathA
00115 ca cAkSuSatvAdiparihAraH / dharmavacanena api dharmyAzraya
00201 siddhau dharmigrahaNasAmarthyAt pratyAsattyA sAdhyadharmisiddhir
00202 iti cet / na / dRSTAntadharmiNo api pratyAsatteH / tadaMza
00203 vyAptyA dRSTAntadharmiNi sattvasiddher dharmivacanAt sAdhyadharmi
00204 parigrahaH / siddhe punar vacanaM niyamArtham AzaGkyeta /  
00205 sajAtIya eva sattvam iti siddhe api vijAtIyavyatireke sAdhya 
00206 abhAve asattvavacanavat / sAmarthyAd arthagatau pratipatti   
00207 gauravaparihArArthaM ca pakSavacanam / pakSasya dharmatve tad 
00208 vizeSaNApekSasya anyatra ananuvRtter asAdhAraNatA iti cet / na / ayogavyavacchedena
00209 vizeSaNAt / yathA caitro dhanurdharaH / na anyayogavyavacchedena /
00210 yathA pArtho dhanurdhara ity AkSepsyAmaH / 
00211 tadaMzas taddharmaH / vaktur abhiprAyavazAt / na tadekadezaH /
00212 pakSazabdena samudAyAvacanAt / vyAptir vyApakasya tatra bhAva
00213 eva / vyApyasya vA tatra eva bhAvaH etena avayavavyatirekau yathAsvaM
00214 pramANena nizcitAv uktau pakSadharmaz ca / ta ete kArya  
00215 svabhAvAnupalabdhilakSaNAs trayo hetavaH / yathA agnir atra dhUmAt /
00216 vRkSo ayaM ziMzapAtvAt / pradezavizeSe kvacin na ghaTa upalabdhi
00217 lakSaNaprAptasya anupalabdheH / yadi syAd upalabdhyasattva
00218 eva syAn na anyathA / tena upalabdhilakSaNaprAptasattvasya ity uktaM
00219 bhavati / tatra dvau vastusAdhanau ekaH pratiSedhahetuH /  svabhAvapratibandhe
00220 hi saty artho arthaM na vyabhicarati / sa ca
00221 tadAtmatvAt / tadAtmatve sAdhyasAdhanabhedAbhAva iti cet /
00222 na / dharmabhedaparikalpanAd iti vakSyAmaH / tathA ca Aha / sarva
00223 eva ayam anumAnAnumeyavyavahAro buddhyArUDhena dharmadharmi
00301 bhedena iti / bhedo dharmadharmitayA buddhyAkArakRto na  artho
00302 api / vikalpabhedAnAM svatantrANAm anarthAzrayatvAt / tat  
00303 kalpitaviSayAd arthapratItAv anarthapratilambha eva syAt / kAryasya
00304 api svabhAvapratibandhaH / tatsvabhAvasya tadutpatter iti /
00305 etau dvAv anumeyapratyayau sAkSAd anutpatter atatpratibhAsitve
00306 api tadutpattes tadvyabhicAriNAv iti pramANaM pratyakSavat /
00307 pratyakSasya api hy arthAvyabhicAra eva prAmANyam / tad abhAve 
00308 bhAvinas tadvipralambhAt / avyabhicAraz ca anyasya ko anyas tadutpatteH /
00309 anAyattarUpANAm sahabhAvaniyamAbhAvAt / yadi tadutpatteH 
00310 kAryaM gamakam sarvathA gamyagamakabhAvaH sarvathA
00311 janyajanakabhAvAt / na / tadabhAve bhavatas tadutpatti
00312 niyamAbhAvAt / tasmAt / 
00313 	kAryaM svabhAvair yAvadbhir avinAbhAvi kAraNe /
00314 teSAM  
00315 	hetuH / 
00316 tatkAryatvaniyamAt tair eva dharmair ye tair vinA na bhavati / 
00317 aMzena janyajanakatvaprasaGga iti cet / na / tajjanyavizeSagrahaNe
00318 abhimatatvAl liGgavizeSaupAdhInAM ca sAmAnyAnAm / aviziSTa  
00319 sAmAnyavivakSAyAM vyabhicArAn na iSyate / 
00401 	svabhAve bhAvo api bhAvamAtrAnurodhini //2//
00402 hetur iti vartate / tAdAtmyaM hi arthasya tanmAtrAnurodhini eva  
00403 na anyA Ayatte / tadbhAve abhUtasya pazcAd bhAvaniyamAbhAvAt / kAraNAnAM
00404 kAryavyabhicArAt / 
00405 	apravRttiH pramANAnAm /
00406 anupalabdhiH  
00407 	apravRttiphalA asati /  
00408 sajjJAnazabdavyavahAra pratiSedhaphalA / upalabdhipUrvakatvAt teSAm
00409 iti idaM sadasatpratiSedhavidhihetvos tulyaM rUpam / tathA 
00410 hi sattvam upalabdhir eva vastuyogyatAlakSaNA tadAzrayA vA  
00411 jJAnapravRttiH / tataH sajjJAnazabdavyavahAravRtteH / asatAM  
00412 ca asattvam anupalabdhiH  
00413 	asajjJAnaphalA kAcid hetubhedavyapekSayA //3//
00414 hetur anupalabdhiH / bhedo asyA vizeSaNam upalabdhilakSaNa 
00415 prAptasattvam / atra anupalabdher liGgAd asattAyAm upalabdher  
00416 abhAvo api anyayA anupalabdhyA sAdhya iti anavasthAnAd  apratipattiH
00417 syAt / atha upalabdhyabhAvo vinA anupalabdhyA syAt / tathA
00418 sattA abhAvo api syAt / apArthikAnupalabdhiH / atha anyaupalabdhyA
00419 anupalabdhisiddhir iti pratyakSasiddhA anupalabdhiH / tathA anyasattayA
00420 asattA kiM na sidhyati iti / yadA punar evaMvidhA anupalabdhir
00421 eva asatAm asattA tadA siddhe api viSaye mohAd viSayiNo
00422 asajjJAnazabdavyavahArAn apratipadyamAno viSayapradarzanena 
00423 samaye pravartyate / yathA gaur ayam sAsnAAdisamudAyAtmakatvAd 
00424 iti / tathA ca dRSTAntAsiddhicodanA api prativyUDhA
00501 viSayapratipattAv apy apratipannaviSayINAM darzanAt /
00502 evam anayor anupalabdhyoH svaviparyayahetvabhAvabhAvAbhyAM 
00503 sadvyavahArapratiSedhaphalatvaM tulyam / ekatra saMzayAd anyatra  
00504 viparyayAt / tatra AdyA sadvyavahAraniSedhaupayogAt pramANam  
00505 uktA / na tu vyatirekadarzanAdAv upayujyate / saMzayAt / dvitIyA 
00506 tv atra pramANaM nizcayaphalatvAt / sA ca prayogabhedAt 
00507 	viruddhakAryayoH siddhir asiddhir hetubhAvayoH /  
00508 	dRzyAtmanor abhAvArthAnupalabdhiz caturvidhA //4//
00509 yAvAn kazcit pratiSedhaH sa sarvo anupalabdheH / tathA hi sa dvidhA
00510 kriyeta kasyacid vidhinA niSedhena vA / vidhau viruddho vA  vidhIyeta 
00511 aviruddho vA / aviruddhasya vidhau sahabhAvavirodhAbhAvAd 
00512 apratiSedhaH / viruddhasya apy anupalabdhyabhAvena virodhApratipattiH 
00513 / tathA hy aparyantakAraNasya bhavato anyabhAve 
00514 abhAvAd virodhagatiH / sa ca anupalabdheH / anyonyaupalabdhiparihAra
00515 sthitalakSaNatA vA virodho nityAnityatvavat / tatra apy  ekaupalabdhyA
00516 anyAnupalabdhir eva ucyate / anyathA aniSiddhaupalabdher abhAva  
00517 asiddheH / ekasya niSedhena anyAbhAvasAdhane siddhA eva anupalabdhiH /
00518 niSedhasya anupalabdhirUpatvAt / tatra apy arthAntaraniSedhe  
00519 kAryakAraNayor anubhavasya vA / tatra anubhavasya apratibandhAt tad 
00520 abhAve anyena na bhavitavyam iti kuta etat / kAryAnupalabdhAv  
00521 api na avazyaM kAraNA api tadvanti bhavanti iti tadabhAvaH kutaH / 
00522 tasmAt kAraNAnupalabdhir eva abhAvaM gamayati iti / svabhAvAnupalabdhis
00523 tu svayam asattA eva / tatra kevalaM viSayI sAdhyate /  
00524 asyAm api yadA vyApakadharmAnupalabdhyA vyApyAbhAvam Aha
00525 tadA abhAvo api iti / iyaM pratiSedhaviSayAnupalabdhiH prayoga
00601 bhedena caturdhA bhavati / viruddhasiddhyA yathA na zItasparzo
00602 atra agneH / etena vyApakaviruddhasiddhir uktA veditavyA yathA
00603 na tuSArasparzo atra agneH / viruddhakAryasiddhyA yathA na zIta  
00604 sparzo atra dhUmAt / hetvasiddhyA yathA na atra dhUmo anagneH /  
00605 svabhAvAsiddhyA yathA na atra dhUmo anupalabdheH / etena vyApaka  
00606 svabhAvAsiddher uktA yathA na atra ziMzapA vRkSAbhAvAt / 
00607 sarvatra ca asyAm abhAvasAdhanyAm anupalabdhau dRzyAtmanAm eva
00608 teSAM tadviruddhAnAM ca siddhir asiddhiz ca veditavyA / anyeSAm  
00609 abhAvavirodhAsiddheH / yadi viruddhakAryaupalabdhyA apy  abhAva 
00610 siddhiH tatkAraNaupalabdhyA kiM na sidhyati / 
00611 	tadviruddhanimittasya yA upalabdhiH prayujyate /
00612 	nimittayor viruddhatvAbhAve sA vyabhicAriNI //5//  
00613 yathA na zItasparzo atra kASThAd iti / nimittayoH punar virodhe
00614 gamikA eva yathA na asya romaharSAdivizeSAH santi saMnihitadahana
00615 vizeSatvAt etena tat kAryAd api tadviruddhakAryAbhAvagatir
00616 uktA veditavyA yathA na romaharSAdivizeSayuktapuruSavAn ayaM
00617 pradezo dhUmAt / iyaM ca hetvasiddhyA eva tadviruddhasiddhiH  
00618 prAg eva nirdiSTA iti iyaM prayogabhedAd aSTadhAnupalabdhiH / tatra
00619 yA iyaM viruddhakAryaupalabdhir uktA tatra
00620 	iSTaM viruddhakArye api dezakAlAdyapekSaNam / 
00621 	anyathA vyabhicAri syAd bhasmA iva azItasAdhane //6//  
00622 yas tarhi samagreNa hetunA kAryautpAdo anumIyate sa kathaM  
00623 trividhe hetAv antarbhavati /  
00624 	hetunA yaH samagreNa kAryautpAdo anumIyate /  
00625 	arthAntarAnapekSatvAt sa svabhAvo anuvarNitaH //7//
00626 asAv api yathA asaMnihitAn na anyam apekSata iti tanmAtrAnubandhI 
00627 svabhAvo bhAvasya / tatra hi kevalaM samagrAt kAraNAt  kAryautpatti
00628 sambhavo anumIyate samagrANAM kAryautpAdanayogyatAnumAnAt /
00629 yogyatA ca sAmagrImAtrAnubandhinI iti svabhAvabhUtA eva anumIyate /
00701 kiM punaH kAraNaM sAmagryAH kAryam eva na anumIyate / 
00703 	sAmagrIphalazaktInAM pariNAmAnubandhini / 
00704 	anaikAntikatA kArye pratibandhasya sambhavAt //8//
00705 na hi samagrANi ity eva kAraNadravyANi svakAryaM janayanti sAmagrI
00706 janmanAM zaktInAM pariNAmApekSatvAt kAryautpAdasya / atra 
00707 antare ca pratibandhasambhavAn na kAryAnumAnam / yogyatAyAs 
00708 tu dravyAntarAnapekSatvAn na virudhyate anumAnam / uttarottara  
00709 zaktipariNAmena kAryautpAdanasamarthA iyaM kAraNasAmagrI 
00710 zaktipariNAmapratyayasya anyasya apekSaNIyasya abhAvAd iti / pUrva
00711 sajAtimAtrahetutvAt zaktiprasUteH sAmagryA yogyatA ananyApekSiNI
00712 ity ucyate / yA tarhy akAryakAraNabhUtena anyena rasAdinA rUpa  
00713 AdigatiH sA katham / sA api 
00714 	ekasAmagryadhInasya rUpAder asato gatiH /  
00715 	hetudharmAnumAnena dhUmaindhanavikAravat //9//  
00716 tatra hetur eva tathAbhUto anumIyate / pravRttazaktirUpaupAdAna 
00717 kAraNasahakAripratyayo hi rasahetU rasaM janayati / indhanavikAra
00718 vizeSaupAdAnahetusahakAripratyayAgnidhUmajananavat / tathA  
00719 hi / 
00720 	zaktipravRttyA /  
00721 svakAraNasya phalautpAdanaM praty Abhimukhyena  
00801 	na vinA rasaH sa eva anyakAraNam /  
00802 rUpaupAdAnahetUnAM pravRttikAraNam / sA api rasaupAdAnakAraNapravRttI
00803 rUpaupAdAnakAraNapravRttisahakAriNI / tasmAd yathAbhUtAd heto
00804 rasa utpannas tathAbhUtam anumApayan rUpam anumApayati
00805 	ity /
00806 tatra api  
00807 	atItaekaikAlAnAM gatis /
00808 na anAgatAnAM vyabhicArAt  
00809 	tat /
00810 tasmAd iyam  
00811 	kAryaliGgajA //10//  
00812 tena na anyo hetur gamako asti / apratibaddhasvabhAvasya avinAbhAva  
00813 niyamAbhAvAt / etena pipIlikautsaraNamatsyavikArAder 
00814 varSAdyanumAnam uktam / tatra api bhUtapariNAma eva  varSahetuH  
00815 pipIlikAdi saMkSobhAdihetur iti /
00816 	hetunA tv asamagreNa yat kAryam anumIyate /  
00817 	zeSavat tad asAmarthyAd dehAd rAgAnumAnavat //11// 
00818 samagrANy eva hi kAraNAni yogyatAm apy anumApayanty asamagrasya
00819 ekAntAsAmarthyAt / yathA dehaindriyabuddhibhyo rAgAdyanumAnam /
00820 AtmAtmIyAbhinivezapUrvakA hi rAgAdayo ayonizomanaskAra
00821 pUrvakatvAt sarvadoSautpatteH / dehAdInAM hetutve api
00822 na kevalAnAM sAmarthyam asti iti / vipakSavRtter adRSTAv api
00823 zeSavadanumAnAt saMzayaH / tathA  
00901 	vipakSe adRSTimAtreNa kAryasAmAnyadarzanAt /  
00902 	hetujJAnaM pramANAbhaM vacanAd rAgitAdivat //12//  
00903 na hi rAgAdInAm eva kAryaM spandanavacanAdayaH  vaktukAmatA 
00904 sAmAnyahetutvAt / sa eva rAga iti cet iSTatvAn na kiMcid bAdhitaM  
00905 syAt / nityasukhAtmAtmIyadarzanAkSiptaM sAsravadharmaviSayaM
00906 cetaso abhiSvaGgaM rAgam AhuH / na evaM karuNAAdayo anyathA
00907 api sambhavAd iti nivedayiSyAmaH / atra yathA rakto bravIti  
00908 tathA virakto api iti vacanamAtrAd apratipattiH / na api vizeSAt /
00909 abhiprAyasya durbodhatvAt / vyavahArasaMkareNa sarveSAM vyabhicArAt /
00910 prayojanAbhAvAd avyAhAra iti cet na parArthatvAt / 
00911 na yukto vItarAgatvAd iti cet na karuNA api vRtteH /  sa eva
00912 rAga iti cet / iSTam / aviparyAsasamudbhavAn na doSaH / asaty
00913 apy Atmagrahe duHkhavizeSadarzanamAtreNa abhyAsabalautpAdinI
00914 bhavaty eva karuNA / tathA hi / sattvadharmAdyAlambanA maitryAdaya
00915 iSyante / etAz ca sajAtIyAbhyAsavRttayo na rAgApekSiNyaH /
00916 na evaM rAgAdayo viparyAsAbhAve abhAvAt / kAruNikasya api
00917 niSphala Arambho aviparyAsAd iti cet / na / parArthasya eva  
00918 phalatvena iSTatvAt / icchAlakSaNatvAt phalasya / sarvathA abhUta  
00919 asamAropAn nirdoSaH / tadanyena doSavattvasAdhane na kiMcid 
00920 aniSTam / vaktary Atmani rAgAdidarzanena anyatra tadanumAne 
00921 atiprasaGgaH / vyabhicArAt / ananyAnumAne iha avyabhicAra iti 
01001 ko nizcayaH / karaNaguNavaktukAmate hi vacanam anumApayet /  
01002 rAgautpAdanayogyatArahite vacanAdarzanAt tadanumAne
01003 atiprasaGgaH uktaH / rAgasya anupayoge kathaM taczaktir upayujyate / 
01004 zaktyupayoge hi sa eva upayuktaH syAt tac ca na asti ity uktam /  
01005 tasmAn nAntarIyakam eva kAryaM kAraNam anumApayati  tat
01006 pratibandhAt na anyad vipakSe adarzane api / sarvadarzino hi darzana
01007 vyAvRttiH sarvatra abhAvaM gamayet / kvacit tathA dRSTAnAm  
01008 api dezakAlasaMskArabhedena anyathA adarzanAd / yathA AmalakyaH  
01009 kSIrAvasekena madhuraphalA bhavanti / na ca evaM bahulaM dRzyante /  
01010 tena evaM syAd yuktaM vaktum mAdRzo vaktA rAgI iti 
01011 rAgautpattipratyayavizeSeNa AtmadarzanAyonizomanaskAreNa yogAt /
01012 tadA apy apArthako vacanaudAhAraH / tasmAd vipakSe adRSTir ahetuH / 
01013 	na ca adarzanamAtreNa vipakSe avyabhicAritA /  
01014 	sambhAvyavyabhicAratvAt sthAlItaNDulapAkavat //13//  
01015 na hi bahulaM pakvadarzane api sthAlyantargamanamAtreNa pAkaH 
01016 sidhyati / vyabhicAradarzanAt / evaM tu syAd evaM svabhAvA 
01017 etatsamAnapAkahetavaH pakvA iti / anyathA tu zeSavad etad anumAnaM
01018 vyabhicAri / kiM punar etac zeSavat /
01019 	yasya adarzanamAtreNa vyatirekaH pradarzyate /
01020 	tasya saMzayahetutvAc zeSavat tad udAhRtam //14//  
01021 sa tasya vyatireko na nizcita iti vipakSe vRttir AzaGkyeta / vyatireka  
01022 sAdhanasya adarzanamAtrasya saMzayahetutvAt / na sarvAnupalabdhir  
01023 gamikA / tasmAd ekanivRttyA anyanivRttim icchatA tayoH  
01024 kazcit svabhAvapratibandho apy eSTavyaH / anyathA agamako
01025 hetuH syAt / 
01026 	hetos triSu api rUpeSu nizcayas tena varNitaH /
01027 	asiddhaviparItArthavyabhicArivipakSataH //15// 
01028 na hy asati pratibandhe anvayavyatirekanizcayo asti / tena tam 
01101 eva darzayan nizcayam Aha / tatra anvayasya nizcayena viruddha 
01102 tat pakSyANAM nirAsaH / vyatirekasya anaikAntikasya tatpakSasya ca
01103 zeSavadAdeH / dvayor ity ekasiddhapratiSedhaH / prasiddhavacanena 
01104 sandigdhayoH zeSavadasAdhAraNayoH sapakSavipakSayor 
01105 api / anyathA hy asati pratibandhe adarzanamAtreNa vyatireke /
01106 	vyabhicArivipakSeNa vaidharmyavacanaM ca yat /  
01107 yad Aha / eSa tAvan nyAyo yad ubhayaM vaktavyam viruddha anaikAntika
01108 pratipakSeNa iti / vaidharmyavacanam anaikAntikapratipakSeNa /
01110 	yady adRSTiphalaM tac ca /  
01111 yadi tena vipakSe adarzanaM khyApyate / 
01112 	tad anukte api gamyate //16// 
01113 na hi tasya prAg darzanabhrAntir yA vacanena nivartyeta / smRtir
01201 vAcAdarzane kriyata iti cet / darzanaM khalv apratIyamAnam  
01202 anaGgam iti yuktaM tatra smaraNAdhAnam / adarzanaM tu darzana  
01203 abhAvaH / sa darzanena bAdhyate / tadabhAve tu siddha eva ity apArthakaM
01204 tatsiddhaye vacanaM / na vai anupalambhamAnasya 
01205 tAvatA na asti iti bhavati tadarthaM vacanam iti cet /  
01206 	na ca na asti iti vacanAt tan na asty eva yathA yadi / 
01207 	na asti sa khyApyate nyAyas tadA na asti iti gamyate //17//
01208 yady anupalabhamAno na asti iti na pratyeti vacanAd api na eva pratyeSyati /
01209 tad api hy anupalambham eva khyApayati / na ca ekAnupalambho  
01210 anyAbhAvaM sAdhayaty atiprasaGgAt / na ca tena  
01211 na asti iti vacanAt tathA bhavaty atiprasaGgAt / tat kathaM vaidharmyavacanena
01212 anaikAntikaparihAraH / tasmAd vyAvRttim icchatA
01213 tatra nyAyo vaktavyaH yato asya vyAvRttam iti bhavati /  
01214 nanu tadabhAve anupalambhAt siddhA vyAvRttiH /  
01215 	yady adRSTyA nivRttiH syAc zeSavad vyabhicAri kim /
01216 yathA pakvAny etAni phalAny evaMrasAni vA rUpAvizeSAd ekazAkhA
01217 prabhavatvAd vA upayuktavad iti / atra api vivakSitAzeSapakSIkaraNe
01218 hetoH sAdhyAbhAve anupalambho asti iti kathaM vyabhicAraH / 
01219 pratyakSabAdhAAzaGkA vyabhicAra ity eke / na / pakSIkRta
01220 viSaye abhAvAt / kadAcid bhaved iti cet / tathA zaGkAyAm atiprasaGgaH /
01221 anyatra apy abhAvaniyamAbhAvAt / vRttaM pramANaM bAdhakam / 
01222 avRttabAdhane sarvatra anAzvAsaH / vyatirekas tu siddha
01223 eva sAdhanam iti tathA abhAvanizcayam apekSate / anupalambhAt tu  
01224 kvacid abhAvasiddhAv apy apratibaddhasya tadabhAve sarvatra 
01225 abhAvAsiddheH saMzayAd avyatireko vyabhicAraH zeSavataH /
01226 kiM ca /
01301 	vyatireky api hetuH syAn /  
01302 na idaM nirAtmakaM jIvaccharIram aprANAdimattvaprasaGgAd iti /
01303 nirAtmakeSu ghaTAdiSu dRSTeSu prANAdyadarzanAt tannivRtty  
01304 AtmagatiH syAt / adRzyAnupalambhAd abhAvAsiddhau  ghaTAdInAM
01305 nairAtmyAsiddheH prANAder anivRttiH / abhyupagamAt siddham  
01306 iti cet / katham idAnIm AtmasiddhiH / parasya apy apramANikA  
01307 kathaM nairAtmyasiddhiH / abhyupagamena ca sAtmakAnAtmakau 
01308 vibhajya tatra abhAvena gamakatvaM kathayatA Agamikatvam
01309 Atmani pratipannaM syAt / na anumeyatvam / tasmAd adarzane apy
01310 AtmanonivRttyasiddheH / tannivRttau kvacin nivRttAv api prANAdInAm
01311 apratibandhAt / sarvatra nivRttyasiddher agamakatvam /
01312 yA apy asiddhiyojanA tathA sapakSe sann asann ity evamAdiSv api  
01313 yathAyogam udAharyam ity evamAdikA / sA api 
01314 	na vAcyA asiddhiyojanA //18//
01315 anupalambha eva saMzayAt / upalambhe tadabhAvAt / anupalambhAc 
01316 ca vyatireka iti saMzayito anivAryaH syAt / yathAyogavacanAt
01317 anivArita eva iti cet / na / ya eva tu ubhayanizcitavAcI ity AdivacanAt /
01318 tena anupalambhe api saMzayAd anivRttiM manyamAnaH tatpratiSedham
01319 Aha / kiM ca /  
01401 	vizeSasya vyavacchedahetutA syAd adarzanAt /
01402 zrAvaNatvasya api nityAnityayor adarzanAd vyAvRttir iti tadvyavaccheda
01403 hetutA syAt / na hi tadvyAvRtter anyad vyavacchedanam /
01404 avyavacchedas tu kutazcid vyAvRtter eva anizcayAt / yo hi yatra
01405 na asti iti nizcitaH sa bhavaMs tadabhAvaM kathaM na gamayet / 
01406 	pramANAntarabAdhA cen /
01407 atha api syAd ubhayavyavacchede pramANAntaraM bAdhakam asti / 
01408 anyonyavyavacchedarUpANAm ekavyavacchedena anyavidhAnAd 
01409 apratiSedhaH / vidhipratiSedhayor virodhAt / 
01410 	na idAnIM nAstitA adRzaH //19// 
01411 na evam adarzanaM pramANaM bAdhAsambhavAt / 
01412 	tathA anyatra api sambhAvyaM pramANAntarabAdhanam /  
01413 lakSaNayukte bAdhAsambhave tallakSaNam eva dUSitaM syAd iti 
01414 sarvatra anAzvAsaH / anumAnaviSaye api pratyakSAnumAnavirodha 
01415 darzanAd anAzvAsaprasaGga iti cet / na / yathokte asambhavAt / 
01416 sambhavinaz ca atallakSaNatvAt / viruddhAvyabhicAryavacanam iti
01417 cet / anumAnaviSaye avacanAd iSTam / viSayaM ca asya  nivedayiSyAmaH /
01418 kiM ca /
01419 	dRSTA ayuktir adRSTez ca syAt sparzasya avirodhinI //20//
01420 yadi hy anupalambhena abhAvaH sidhyet / yad Aha / yady adarzana  
01421 mAtreNa dRSTebhyaH pratiSedhaH kriyate / na ca so api yukta  
01501 iti / katham ayuktaH / anupalambhAd abhAvasiddheH / nanu upalabdhi 
01502 lakSaNaprApteH sparzasya yukta eva pratiSedhaH / na yuktaH /
01503 dRzyatatsvabhAvaviSayamAtrApratiSedhAt / pRthivyAdi sAmAnyena
01504 gRhItvA ayaM pratiSedham Aha / tatra ca tUlaupalapallava  
01505 AdiSu tadbhAve api sparzabhedadarzanAt asya api kvacid vizeSe 
01506 sambhavAzaGkayA bhavitavyam iti sarvatra adarzanamAtreNa ayuktaH 
01507 pratiSedha iti / evam AcAryIyaH kazcid anupalambhAd abhAvaM 
01508 bruvANa upAlabdhaH / api ca / 
01509 	dezAdibhedAd dRzyante bhinnA dravyeSu zaktayaH /
01510 	tatra ekadRSTyA na anyatra yuktas tadbhAvanizcayaH //21// 
01511 yadi kathaMcid vipakSe adarzanamAtreNa apratibaddhasya api tad
01512 avyabhicAraH / kvacid deze kAnicid dravyANi kathaMcid dRSTAni
01513 punar anyathA anyatra dRzyante / yathA kAzcid oSadhayaH kSetra  
01514 vizeSe viziSTarasavIryavipAkA bhavanti / na anyatra / tathA kAla 
01515 saMskArabhedAt / na ca taddezais tathA dRSTA iti sarvAs tattvena
01516 tathAbhUtAH sidhyanti / guNAntarANAM kAraNAntarApekSatvAt /  
01517 vizeSahetvabhAve tu syAd anumAnam / yathA adRSTakartRkam api 
01518 vAkyaM puruSasaMskArapUrvakam iti / vAkyeSu vizeSAbhAvAt /
01519 sarvaprakArANAM puruSaiH karaNasya darzanAt /  na evam asambhavad
01520 vizeSahetavaH puruSA yena vacanAdeH kiMcinmAtrasAdharmyAt 
01521 sarvAkArasAmyam anumIyeta / sarvaguNeSu vizeSadarzanAt /  
01522 saMskArabhedena vizeSapratipatteH / tadvad anyasya api  sambhavAd /
01523 asambhavAnumAne ca bAdhakahetvabhAvAt vairAgyAdRSTeH  
01524 adRSTena ca bAdhyabAdhakabhAvAsiddheH / rAgAdyavyabhicAri  
01525 kAryAbhAvAt / sambhave api vizeSANAM draSTum azakyatvAt /
01601 tAdRzAM ca apratikSepArthatvAt / na evaM vAkyAni dRzyavizeSatvAt /
01602 adRzyatve apy adRSTavizeSANAM vijAtIyatvaupagamavirodhAt / tad
01603 vizeSANAm anyatra api zakyakriyatvAt / pratyakSANAM zabdAnAm
01604 apratyakSasvabhAvAbhAvAt / bhrAntinimittAbhAvAt / bAdhakAbhAvAd
01605 bhrAntyasiddheH / puruSeSu vizeSadarzanasya bAdhakatvAd  
01606 asamAnam / parabhAvabhUtasya ca vAkyavizeSasya atadvizeSatvAt /
01607 tad abhinnasvabhAvAnAM sarveSAM puruSakriyA na vA kasyacit /
01608 kiM ca /
01609 	AtmamRccetanAAdInAM yo abhAvasya aprasAdhakaH /
01610 	sa eva anupalambhaH kiM hetvabhAvasya sAdhakaH //22//  
01611 anupalambhaM ca asya pramANayata AtmavAdo nirAlambaH syAt /
01612 apratyakSatvAd Atmanas tatkAryAsiddheH / indriyAdInAM tu 
01613 vijJAnakAryasya kAdAcitkatvAt sApekSyasiddhyA prasiddhir ucyate /
01614 kim apy asya kAraNam asti iti / ca tv evaMbhUtam iti /
01615 na evaM sukhAdikAryaM prasAdhitaM kaMcid arthaM puSNAti / 
01616 yena kenacit kAraNavattvAbhyupagamAt / tathA ca anupalambha 
01617 eva AtmanaH syAt / taM tena pratyAcakSANaH kim iti prativyUDho 
01618 anupalambhasya asAdhanatvAd iti / katham asAdhanaM vyatirekaM 
01619 sAdhayet / mRdaH khalv api kazcic caitanyam anupalabhyamAnam
01620 api icchann adarzanAd vacanAder vyAvRttim Aha / dadhyAdikaM 
01621 ca aparaH kSIrAdiSv aparArtheSu saMghAtatvAdarzanAd vyatirekam /
01622 ko hy atra niyamaH saMhatair avazyaM parArthair bhavitavyam iti /  
01623 asty eva upalambho dadhyAdInAM kSIrAdiSv anumAnam /  
01624 azaktAd anutpatteH / atha kA iyaM zaktiH sa eva bhAvaH  
01625 uta anyad eva kiMcit / sa eva cet tathA eva upalabhyeta vizeSAbhAvAt /
01626 anyac cet katham anyabhAve tad asti / upacAramAtraM tu syAd
01627 ity ayam eSAM parasparavyAghAtaH /
01628 	tasmAt tanmAtrasambandhaH svabhAvo bhAvam eva vA /
01629 	nivartayet /  
01630 yathA vRkSaH ziMzapAM / zAkhAAdimadvizeSasya eva kasyacit tathA  
01701 prasiddheH sa tasya svabhAvaH / svaM ca svabhAvaM parityajya
01702 kathaM bhAvo bhavet / svabhAvasya eva bhAvatvAd iti tasya svabhAvapratibandhAd
01703 avyabhicAraH /  
01704 	kAraNaM vA kAryam avyabhicArataH //23//  
01705 kAraNaM nivartamAnaM kAryaM nivartayati / anyathA tat tasya  
01706 kAryam eva na syAt / siddhas tu kAryakAraNabhAvaH svabhAvaM
01707 niyamayati ity ubhayathA svabhAvapratibandhAd eva nivRttiH /  
01708 	anyathA ekanivRttyA anyavinivRttiH kathaM bhavet /  
01709 	na azvavAn iti martyena na bhAvyaM gamatA api kim //24//
01710 	saMnidhAnAt tathA ekasya katham anyasya saMnidhiH /  
01711 	gomAn ity eva martyena bhAvyam azvavatA api kim //25//
01712 tasmAt svabhAvapratibandhAd eva hetuH sAdhyaM gamayati /  
01713 sa ca tadbhAvalakSaNas tadutpattilakSaNo vA / sa eva  avinAbhAvo
01714 dRSTAntAbhyAM pradarzyate /  
01715 	tasmAd vaidharmyadRSTAnte na iSTo avazyam iha AzrayaH /
01716 	tadabhAve ca tan na iti vacanAd api tadgateH //26//  
01717 yataH / 
01718 	tadbhAvahetubhAvau hi dRSTAnte tadavedinaH /
01719 	khyApyete /  
01720 dRSTAnte hi sAdhyadharmasya tadbhAvas tanmAtrAnubandhena
01721 tatsvabhAvatayA khyApyate / yaH kRtakaM svabhAvaM janayati  
01722 so anityasvabhAvaM santaM janayati iti pramANaM dRSTAntena upadarzyate /
01723 anyathA ekadharmasadbhAvAt tadanyena api bhavitavyam
01724 iti niyamAbhAvAt sAdhanasya sAdhyavyabhicArAzaGkA syAt /  
01801 tena ca pramANena sAdhyadharmasya tanmAtrAnubandhaH khyApyate /
01802 svakAraNAd eva kRtakas tathAbhUto jAto yo nazvaraH
01803 kSaNasthitodharmA / anyatas tasya tadbhAvaniSedhAt /  hetubhAvo  
01804 vA tasmin saty eva bhAvAd iti dRSTAntena pradarzyate arthAntarasya  
01805 tathA prasiddhe tadbhAve hetubhAve vA anityatvAbhAve  kRtakatvaM 
01806 na bhavati dahanAbhAve ca dhUmaH / tathA hi sa tasya  
01807 svabhAvo hetur vA / kathaM svaM svabhAvaM hetuM vA antareNa bhaved 
01808 ity Azrayam antareNa api vaidharmyadRSTAnte prasidhyati vyatirekaH /
01809 yeSAM punaH prasiddhAv eva tadbhAvahetubhAvau teSAm
01811 	viduSAM vAcyo hetur eva hi kevalaH //27//  
01812 yadarthe dRSTAnta ucyate so arthaH siddha iti kiM tadvacanena 
01813 tadA / tatpradarzane api kiM vaidharmyadRSTAntAzrayeNa iti manyamAna
01814 AzrayaM pratikSipti sma /
01815 	tena eva jJAtasambandhe dvayor anyatarauktitaH /
01816 	arthApattyA dvitIye api smRtiH samupajAyate //28//
01817 yad Aha arthApattyA vA anyatareNa ubhayapradarzanAd iti / tatra api 
01818 dRSTAntena tadbhAvahetubhAvapradarzanaM manyamAno arthApattyA
01819 ekavacanena dvitIyasiddhim Aha / tathA hi yat kRtakaM tad
01820 anityam ity ukte anarthAntarabhAve vyaktam ayam asya svabhAvas 
01821 tanmAtrAnubandhI pramANadRSTas tadbhAvaniyamAd iti / jJAta
01822 tadbhAvasya arthApattyA anityatvAbhAve kRtakatvaM na bhavati iti
01823 bhavati / na hi svabhAvasya abhAve bhAvo bhavaty abhedAt /  anyathA  
01824 tadbhAve bhavati ity eva na syAt / tathA tadabhAve na bhavati 
01825 ity ukte / tata eva tadbhAvatAvedinaH / tathA hy ayam asya  
01826 svabhAvo yena tadabhAve na bhavati / anyathA ayogAd iti tat
01901 tatsvabhAvatApratipattyA anvayasmRtir bhavati / tathA yatra  
01902 dhUmas tatra agnir ity ukte kAryaM dhUmo dahanasya / yena  
01903 dhUme avazyam agnir bhavati / anyathA arthAntarasya tadanubandha 
01904 niyamAbhAvAt svAtantryaM bhAvasya syAt / atas tadabhAve 
01905 api svabhAvAvaikalyAn na abhAvaH / kArye tv avazyaM kAraNaM  
01906 bhavati / idam eva hi kAraNasya kAraNatvaM yad arthAntarabhAve
01907 svabhAvaupadhAnaM / kAryasya api tadbhAva eva bhAvaH
01908 tac ca asti dhUme / tasmAt kAryaM dhUma ity anvayena vidhita
01909 tatkAryatvasya dahanAbhAve dhUmo na bhavati ity arthAd vyatirekapratipattir
01910 bhavati / tathA asaty agnau dhUmo na asti ity
01911 ukte agnir dhUme bhavaty avazyam ity arthAd anvayapratipattiH /  
01912 anyathA hi tadabhAve kiM na bhaved iti / nanu ca nityAnitya
01913 arthakAryatvAbhAve api zravaNajJAnaM na bhavati tadabhAve /
01914 na vai na bhavati / tayor eva tataH saMzayAt / anyathA
01915 abhAvena nizcitAt kathaM tadbhAvaparAmarzena saMzayaH syAt / 
01916 kevalaM tu bhAvanizcayAbhAvAn na asti ity ucyate / yadA punar
01917 dRSTAntena na agnidhUmayoH kAryakAraNabhAvaH  pradarzyate / tadA
01918 yatra dhUmas tatra agnir ity eva na syAt / pratibandhAbhAvAt /
01919 kuto agnyabhAve dhUmo na asti ity arthAd vyatirekasiddhiH / tathA  
01920 vaidharmyeNa abhAvAsiddher anvayasmRtiH tasmAd dRSTAntena ayam
01921 eva yathoktasvabhAvapratibandhaH pradarzyate / ekasadbhAve  
01922 anyasya prasiddhyartham / tadabhAve asambhavAt / 
01923 	hetusvabhAvAbhAvo ataH pratiSedhe ca kasyacit /  
01924 	hetur / 
01925 tAv eva hi nivartamAnau svapratibaddhaM nivartayata iti kasyacid 
01926 arthasya pratiSedham api sAdhayitukAmena hetor vyApakasya  
02001 vA svabhAvasya nivRttir hetutvena AkhyeyA / apratibandhe hi katham 
02002 ekasya nivRttir anyasya nivRttiM sAdhayet /
02003 	yuktaupalabhasya tasya ca anupalambhanam //29//
02004 pratiSedhahetuH / pratiSedhaviSayavyavahArahetus tadhetur ity 
02005 uktaH / svayaM tathAbhUtAnupalambhasya pratiSedharUpatvAt /  
02006 hetuvyApakAnupalabdhir ubhayasya api hetuH /  
02007 	iti iyaM trividhA apy uktAnupalabdhir anekadhA /  
02008 	tat tadviruddhAdyagatigatibhedaprayogataH //30// 
02009 trividha eva hi pratiSedhahetuH / upalabhyasattvasya hetos tathAbhAva
02010 nizcaye vyApakasya svAtmanaz ca anupalabdhir iti / sa 
02011 prayogavazena tattadviruddhAdyagatigatibhedaprayogato  aneka
02012 prakAra uktaH / tasya agatyA tadviruddhagatyA viruddhakAryagatyA 
02013 ity Adibhedaprayogair yathoktaM prAk /  
02014 	kAryakAraNabhAvAd vA svabhAvAd vA niyAmakAt /  
02015 	avinAbhAvaniyamo adarzanAn na na darzanAt //31// 
02016 	avazyaMbhAvaniyamaH kaH parasya anyathA paraiH /  
02017 	arthAntaranimitte vA dharme vAsasi rAgavat //32// 
02018 ity antarazlokau / api ca /  
02019 	arthAntaranimitto hi dharmaH syAd anya eva saH /  
02020 na hi tasmin niSpanne aniSpanno bhinnahetuko vA tatsvabhAvo  
02021 yuktaH / ayam eva khalu bhedo bhedahetur vA bhAvAnAM viruddhadharma
02022 adhyAsaH kAraNabhedaz ca / tau cen na bhedakau tadA 
02023 na kasyacit kutazcid bheda ity ekaM dravyaM vizvaM syAt / 
02024 tataz ca saha utpattivinAzau sarvasya ca sarvatra upayogaH syAt /
02101 anyathA ekam ity eva na syAt / nAmAntaraM vA / arthabhedam abhyupagamya
02102 tathA abhidhAnAt / nanv anarthAntarahetutve api bhAva
02103 kAle anityatA aniSpattes tulyA atatsvabhAvatA / na vai kAcid anyA 
02104 anityatA nAma yA pazcAn niSpadyeta / sa eva hi bhAvaH kSaNa
02105 sthitidharmA anityatA / vacanabhede api dharmidharmatayA nimittaM
02106 vakSyAmaH / tAM punar asya kSaNasthitodharmatAM  svabhAvaM 
02107 svahetor eva tathAutpatteH pazyann api mandabuddhiH sattA  
02108 upalambhena sarvadA tadbhAvazaGkAvipralabdho na vyavasyati  
02109 sadRzAparautpattivipralabdho vA / antyakSaNadarzinAM nizcayAt /  
02110 pazcAd asya anupalabdhyA asthitopratipatter nizcayakAla iti 
02111 tadA anityatA vyavasthApyate / kAryautpAdanazakteH  kAraNasvabhAvatve
02112 api adRSTatatkAryasya kAraNadarzane api apratipanna  
02113 tadbhAvasya kAryadarzanAt tatpratipattivat / anyathA  arthAntaram
02114 eva anityatA syAt / anyanimittatve animittatve vA / tathA ca bhAvas 
02115 tadvAn na syAt / tadanupayogAt / upayoge vA sa eva asya  
02116 svAtmabhUto anityatA iti kim anyayA / svabhAvena vA acalasya arthAntara
02117 yoge api tadbhAvAnutpatteH / sa ca arthAntarAd bhavann anityatA
02118 anyo vA dharmo hetuH phalaM vA syAt / ahetuphalasya asambandhAt / 
02119 tatra bhAvAnumAnasya asambhavAt / tatra / 
02120 	pazcAd bhAvAn na hetutvaM phale apy ekAntatA kutaH //33// 
02121 sa hi niSpanne bhAve arthAntarataH pazcAd bhavan kathaM
02122 tasya hetuH syAt / phalasya api na avazyaM hetau bhAva iti tad
02123 bhAvahetor anaikAntikatvam / tan na arthAntaranimitto dharmo  
02124 bhAve avazyaM bhAvi ity ananumAnam / yadi tarhi darzana adarzane
02125 na anvayavyatirekagater AzrayaH kathaM dhUmo agniM na  vyabhicarati
02126 iti gamyate / yasmAt /
02201 kAryaM dhUmo hutabhujaH kAryadharmAnuvRttitaH /
02202 yeSAm upalambhe tallakSaNam anupalabdhaM yad upalabhyate /  
02203 tatra ekAbhAve api na upalabhyate / tat tasya kAryaM / tac ca  dhUme 
02204 asti /  
02205 	sa bhavaMs tadabhAve tu hetumattAM vilaGghayet //34// 
02206 sakRd api tathA adarzanAt kAryaM siddhaH / akAryatve akAraNAt 
02207 sakRd apy abhAvAt / kAryasya ca svakAraNam antareNa bhAve  
02208 ahetumattA eva syAt / na hi yasya yam antareNa bhAvaH sa tasya
02209 hetur bhavati / bhavati ca dhUmo agnim antareNa tan na tadhetuH
02210 anyahetukatvAn na ahetukatvam iti cet / na / tatra api
02211 tulyatvAt / tadabhAve apy agnau bhavati iti / kathaM vA tato 
02212 anyato vA atajjananasvabhAvAd bhavet / svayam atatsvabhAvasya
02213 ajananAt tasya ahetutA syAt / na vai sa eva bhavati tAdRzasya
02214 bhAvAt / anyAdRzAd bhavan kathaM tAdRzaH syAt / tAdRzAd hi 
02215 bhavan tAdRzaH syAt / anyAdRzAd api tAdRzo bhAve taczaktiniyama 
02216 abhAvAn na hetubhedo bhedaka iti akAraNaM vizvasya vaizvarUpyaM
02217 syAt / sarvaM vA sarvasmAj jAyeta / tasmAt kAraNa
02218 bhedAbhedAbhyAM kAryabhedAbhedau / tan na dhUmo arthAd dRSTa
02219 AkAravijAtIyAd bhavaty ahetukatvaprasaGgAt / tathA ca / 
02220 	nityaM sattvam asattvaM vA ahetor anyAnapekSaNAt /
02221 	apekSAto hi bhAvAnAM kAdAcitkatvasambhavaH //35// 
02222 sa hi dhUmo ahetur bhavan nirapekSatvAn na kadAcin na bhavet /
02223 tadbhAve vaikalyAbhAvAd iSTakAlavat / tadA api vA na bhavet /  
02224 abhAvakAlAvizeSAt / apekSayA hi bhAvAH kAdAcitkA bhavanti /
02301 bhAvAbhAvakAlayos tadbhAvayogyatA ayogyatAyogAt tulyayogyatA
02302 yogyatayor dezakAlayos tadvattA itarayor niyamAyogAt / sA ca 
02303 yogyatA hetubhAvAt kim anyat / tasmAd ekadezakAlaparihAreNa anya
02304 dezakAlayor vartamAno bhAvas tatsApekSo nAma bhavati / 
02305 tathA hi / tathA vRttir eva apekSA tatkRtaupakArAnapekSasya tanniyama
02306 ayogAt / tan niyatadezakAlatvAd dhUmo atra dRSTaH sakRd 
02307 vaikalye ca punar na dRSTaH tajjanyo asya svabhAvaH / anyathA 
02308 sakRd apy abhAvAt / sa tatpratiniyato anyatra kathaM bhavet / 
02309 bhavan vA na dhUmaH syAt / tajjanito hi svabhAvavizeSo dhUma 
02310 iti / tathA hetur api tathAbhUtakAryajananasvabhAvaH / tasya anyato
02311 api bhAve na sa tasya svabhAva iti sakRd api na janayet /
02312 na vA sa dhUmo adhUmajananasvabhAvAd bhAvAt / tatsvabhAvatve
02313 ca sa eva agnir ity avyabhicAraH /
02314 	agnisvabhAvaH zakrasya mUrddhA yady agnir eva saH / 
02315 	atha anagnisvabhAvo asau dhUmas tatra kathaM bhavet //36//  
02316 	dhUmahetusvabhAvo hi vahnis taczaktibhedavAn / 
02317 	adhUmahetor dhUmasya bhAve sa syAd ahetukaH //37//  
MatangaV vidya, p. 190
02318 iti saMgrahazlokau / kathaM tarhi idAnIM bhinnAt sahakAriNaH kArya  
02319 utpatter yathA cakSUrUpAder vijJAnasya / na vai kiMcid ekaM 
02320 janakaM tatsvabhAvam / kiM tu sAmagrI janikA tatsvabhAvA / sa eva 
02321 anumIyate / sa eva ca sAmagrI svabhAvasthityAzrayaH kAryasya / 
02322 ata eva sahakAriNAm apy aparyAyeNa jananam / yad api kiMcid
02323 vijAtIyAd bhavad dRSTaM gomayAdeH zAlukAdi / tatra api tathA abhidhAne
02324 apy asty eva svabIjaprabhavAt svabhAvabhedaH / hetusvabhAvabhedAt /
02325 yathA kadalI bIjakandaudbhavA / sphuTam eva 
02326 tAdRzaM loko vivecayaty AkArabhedAt / tasmAn na suvivecitAkAraM 
02327 kAryaM kAraNaM vyabhicarati / 
02401 	anvayavyatirekAd yo yasya dRSTo anuvartakaH / 
02402 	svabhAvas tasya tadhetur ato bhinnAn na sambhavaH //38// 
02403 iti saMgrahazlokaH / tasmAt sakRd api darzanAdarzanAbhyAM
02404 kAryakAraNabhAvasiddher bhavati tatas tatpratipattiH / na anyathA / 
02405 anvayavyatirekayor niHzeSadarzanAdarzanAyattatvAt / kvacid  
02406 amUrtatve nityatvadarzane apy anyatra anyathA adRSTeH / kvacin
02407 nityatvAbhAve apy adRSTasya punar dRSTer iti bhavatu kAryasya
02408 kAraNena avinAbhAvas tadutpatteH / svabhAva idAnIM katham avinAbhAvaH /
02410 	svabhAve apy avinAbhAvo bhAvamAtrAnurodhini / 
02411 yo hi bhAvamAtrAnurodhI svabhAvas tatra avinAbhAvo bhAvasya iSyate /
02413 	tadabhAve svayaM bhAvasya abhAvaH syAd abhedataH //39//
02414 ya eva bhAvo bhAvamAtrAnurodhI svabhAva ity ucyate / sa eva 
02415 svayaM vastuto bhAvaH / sa ca AtmAnaM parityajya kathaM bhavet / 
02416 ya eva tarhi kRtakaH sa eva anityo bhedAbhAvAt / pratijJArthaikadezo 
02417 hetuH syAt / na eSa doSaH / yasmAt /
02418 	sarve bhAvAH svabhAvena svasvabhAvavyavasthiteH /  
02419 	svabhAvaparabhAvAbhyAM yasmAd vyAvRttibhAginaH //40// 
02420 	tasmAd yato yato arthAnAM vyAvRttis tannibandhanAH /  
02421 	jAtibhedAH prakalpyante tadvizeSAvagAhinaH //41// 
02422 	tasmAd yo yena dharmeNa vizeSaH sampratIyate /
02423 	na sa zakyas tato anyena tena bhinnA vyavasthitiH //42//
02424 sarva eva hi bhAvAH svarUpasthitayaH / te na AtmAnaM pareNa
02425 mizrayanti / tasya aparatvaprasaGgAt / yad apy eSAm abhinnam  
02426 AtmabhUtaM rUpaM na tat teSAm / tadAnIM teSAm abhAvAt / 
02501 tad eva hi syAd abhinnasya bhAvAt / tadvyatiriktasya bhinnasya  
02502 ca abhAvAt / tasya eva ca punar bhedavirodhAt / tac ca Atmani vyavasthitam
02503 amizram eva / arthAntaram apy anekasambandhe api  
02504 na tat teSAM sAmAnyam atadrUpatvAt / dvitvAdisaMyogakArya 
02505 dravyeSv api prasaGgAt / na hi sambandhinA apy anyena anye samAnA  
02506 nAma / tadvanto nAma syuH / bhUtavat / kaNTheguNena / 
02507 na abhinnapratyayaviSayAH / bhUtavat / tadAtmAnam eva hi buddhiH
02508 saMsRjantI sAmAnyaviSayA pratibhAsate / na ekasambandhinAv iti 
02509 bhUtavat / taddarzinyAH sA bhrAntir iti cet / taddarzinI iti kutaH /
02510 nirbIjabhrAntyayogAd iti cet / ta eva tadekakAryA bIjam / saMkhyA
02511 saMyogakAryadravyAdimatsu bhUtAdiSv abhAvAc ca / tan na tathA
02512 sAmAnyabuddhau niveSAbhAvAt sAmAnyam anyat / sati vA tasya api  
02513 svAtmani vyavasthAnAd amizraNam anyena / tasmAd ime bhAvAH 
02514 sajAtIyAbhimatAd anyasmAc ca vyatiriktAH svabhAvena  ekarUpatvAt /
02515 yato yato bhinnAs tadbhedapratyAyanAya kRtasaMnivezaiH  
02516 zabdais tatas tato bhedam upAdAya svabhAvAbhede apy aneka  
02517 dharmANaH pratIyante / te api zabdAH sarvabhedAnAkSepe apy eka 
02518 bhedacodanAt tatsvalakSaNaniSThA eva bhavanti / tad ekasmAd
02519 api tasya bhedo asti iti / tasmAd ekasya bhAvasya yAvanti pararUpANi 
02520 tAvat yas tadapekSayA vyAvRttayaH / tadasambhavikAryakAraNasya
02521 tadbhedAt / yAvat yaz ca vyAvRttayas tAvat yaH zrutayo
02522 atatkAryakAraNaparihAreNa vyavahArArthAH / yathA prayatnAnantarIyakaH
02523 zabdaH zrAvaNa ity atatkAryakAraNaparihArArthaH / 
02524 tasmAt svabhAvAbhede api yena yena dharmeNa nAmnA yo vizeSo
02525 bhedaH pratIyate na sa zakyo anyena pratyAyayitum iti na ekArthAH 
02526 sarvazabdAH / tan na pratijJArthaikadezo hetur iti / kathaM
02527 punar etad gamyate vyavacchedaH zabdaliGgAbhyAM pratipAdyate 
02528 vidhinA na vasturUpam eva iti / pramANAntarasya zabda antarasya 
02601 ca pravRtteH / tathA hi /
02602 	ekasya arthasvabhAvasya pratyakSasya sataH svayam / 
02603 	ko anyo na dRSTo bhAgaH syAd yaH pramANaiH parIkSyate //43//
02604 eko hy arthAtmA / sa pratyakSaH asiddhe dharmiNi sAdhanAsambhavAt /
02605 yathA anityatve sAdhye zabdaH / tasya pratyakSeNa eva siddheH  
02606 sarvAkArasiddhiH / tadanyasya asiddhasya abhAvAt / bhAve vA  
02607 atatsvabhAvatvam / na hi yo yad ekayogakSemo na bhavati sa  
02608 tatsvabhAvo yuktaH / tanmAtranibandhanatvAt / bhedavyavahArasya / 
02609 anyathA abhAvaprasaGgAd ity uktam / tasmAt pratyakSe dharmiNi  
02610 tatsvabhAvasAkalyaparicchedAt tatra anavakAzA pramANAntara 
02611 vRttiH syAt /
02612 	no ced bhrAntinimittena saMyojyeta guNAntaram /
02613 	zuktau vA rajatAkAro rUpasAdharmyadarzanAt //44// 
02614 yadi dRSTasarvatattvasya api bhAvasya tathA nizcayapratirodhinA  
02615 bhrAntinimittena guNAntaraM na saMyojyate / yathA zuktau rajata  
02616 AkAraH / na hi zuktau dve rUpe samAnaM viziSTaM ca tathA 
02617 pratipattiprasaGgAt / apratipattau vA vivekena dvitvavikalpAyogAt /
02618 atiprasaGgAc ca / tasmAt pazyan zuktirUpaM viziSTam eva
02619 pazyati / nizcayapratyayavaikalyAt tv anizcinvan tatsAmAnyaM  
02620 pazyAmi iti manyate / tato asya rajatasamAropaH / tathA / sadRzAparApara
02621 utpattyA alakSitanAnAtvasya tadbhAvasamAropAt sthitobhrAntiH / 
02622 yAvanto asya parabhAvAs tAvanta eva yathAsvaM nimitta
02623 bhAvinaH samAropA iti tadvyavacchedakAni bhavanti pramANAni  
02624 saphalAni syuH / teSAM tu vyavacchedaphalAnAM na apratItavastvaMza
02701 pratyAyane pravRttis tasya dRSTatvAt / anaMzasya ca ekadezena 
02702 darzanAyogAt /
02703 	tasmAd dRSTasya bhAvasya dRSTa eva akhilo guNaH /
02704 	bhrAnter nizcIyate na iti sAdhanaM sampravartate //45//  
02705 iti saMgrahazlokaH / tasmAn na adRSTagrahaNAya dRSTe pramANAntaravRttiH /
02707 	vastugrahe anumAnAc ca dharmasya ekasya nizcaye /  
02708 	sarvadharmagraho apohe na ayam doSaH prasajyate //46// 
02709 na kevalaM pratyakSadRSTe pramANAntarAvRttiH kvacit / yadA anumAnam
02710 api vastu vidhinA pratyAyayati na vyavacchedakRt tadA eka
02711 dharmanizcaye tadavyatirekAt sarvadharmanizcaya iti pramANa  
02712 antarAvRttiH na hi tasmin nizcite tadAtmA anizcito yuktaH /  
02713 yadA punar anumAnena samAropavyavacchedaH kriyate tadA na eka 
02714 samAropavyavacchedAd anyavyavacchedaH kRto bhavati iti tadartham  
02715 anyat pravartate / nanu na avazyaM viparyAsapUrvaka eva apratIta  
02716 nizcayo bhavati / yathA akasmAd dhUmAd agnipratipattiH / na hi 
02717 tatra anagnisamAropaH sambhAvyate / tan na sarvatra vyavacchedaH 
02718 kriyate / uktam atra / dharmipratipattAv abhedAt sarvapratipattiH /
02719 bhede vA asambaddhasya tatra apratipattir iti / tasmAt
02720 tatra api taddarzinas tatsvabhAvAnizcayaH kutaH viparyAsAt / sa 
02721 ca taM pradezaM tadviviktena rUpeNa nizcinvann agnisattAbhAvanA
02722 vimuktayA buddhyA katham aviparyasto nAma / tadAkAra
02801 samAropasaMzayarahitaz ca tatpratipattau na liGgam anusaret /
02802 na ca tasya anvayavyatirekayor Adriyeta / 
02803 	tasmAd apohaviSayam iti liGgaM prakIrtitam / 
02804 	anyathA dharmiNaH siddhAv asiddhaM kim ataH param //47//  
02805 iti saMgrahazlokaH /
02806 	kvacid dRSTe api yaj jJAnaM sAmAnyArthaM vikalpakam /  
02807 	asamAropitAnyAMze tanmAtrApohagocaram //48//
02808 yad rUpAdidarzanAnantaram aliGgaM nizcayajJAnaM bhavati /  
02809 tat katham asati samArope bhavad vyavacchedaviSayaM bhavati /  
02810 samAropaviSaye tasya abhAvAt / yatra hy asya samAropo yathA
02811 sthiraH sAtmaka iti vA na tatra bhede nizcayo bhavati /  
02812 	nizcayAropamanasor bAdhyabAdhakabhAvataH /  
02813 na hi sarvato bhinno dRSTo api bhAvas tathA eva pratyabhijJAyate /  
02814 kvacid bhede vyavadhAnasambhavAt / yathA zukteH zuktitve /
02815 yatra tu pratipattur bhrAntinimittaM na asti tatra eva asya taddarzana
02816 avizeSe api smArto nizcayo bhavati / samAropanizcayayor bAdhyabAdhakabhAvAt
02817 nizcayasya 
02818 	samAropaviveke asya pravRttir iti gamyate //49// 
02819 tadviveka eva ca anyApohaH / tasmAt tad api tanmAtrApohagocaram / 
02820 na vastusvabhAvanizcayAtmakam / tathA hi kasyacin nizcaye 
02821 apy anyasya apratipattidarzanAt / tatsvabhAvanizcaye ca tasya ayogAt /
02901 	yAvanto aMzasamAropAs tannirAse vinizcayAH /
02902 	tAvantA eva zabdAz ca tena te bhinnagocarAH //50//  
02903 	anyathA ekena zabdena vyApta ekatra vastuni / 
02904 	buddhyA vA na anyaviSaya iti paryAyatA bhavet //51//
02905 ity antarazlokaH /  
02906 	yasyApi nAnopAdher dhIr grAhikArthasya bhedinaH /
02907 yo api manyate bhinnA eva upAdhayaH parasparam AzrayAc ca /  
02908 tannibandhanAH zrutayas tadAdhAreSu vartante / tatra eva vA
02909 tad ayam aprasaGga iti / tasya api /  
02910 	nAnopAdhyupakArAGgazaktyabhinnAtmano grahe //52//
02911 	sarvAtmanopakArasya ko bhedakaH syAd anizcitaH /
02912 yady apy upAdhayo bhinnA eva zabdajJAnAntarANAM nimittam 
02913 arthe / sa tu tais tadvAn eka eva upalIyate / tasya nAnAupAdhInAm 
02914 upakArAzrayazaktisvabhAvasya svAtmany abhedAt sarvAtmanA 
02915 grahaNe ka eva upAdhibhedo anizcitaH syAt / sarvaupAdhyupakArakatvena
02916 grahaNAt / na hy upakArakatvaM anyad eva tasya svarUpeNa
02917 gRhyamANasya agRhItaM nAma / ato yad eva asya svabhAvena  
02918 grahaNaM tad eva upakAratvena api iti / 
02919 	tayor Atmani sambandhAd ekajJAne dvayagrahaH //53// 
02920 AtmabhUtasya upAdhitadvator upakAryopakArakabhAvasya grahaNAt  
02921 ekajJAne dvayor api grahaNam iti ekaupAdhiviziSTe api tasmin
02922 gRhyamANe sarvaupAdhInAM grahaNam / tadgrahaNa nAntarIyakatvAd  
02923 upAdhimadgrahaNasya / anyathA tathA api na gRhyeta /
02924 na hy anya eva anyaupakArako yo na gRhItaH syAt / na ca apy upakArake
02925 tathA gRhIte upakAryAgrahaNaM tasya apy agrahaNaprasaGgAt /  
02926 svasvAmitvavat / tasmAd arthAntaraupAdhivAde api samAnaH prasaGgaH /
02927 atha api syAd bhinnA eva zaktayaH zaktimato yAbhir upAdhIn  
02928 upakaroti / tato na ayaM prasaGga iti /  
03001 	dharmaupakArazaktInAM bhede tAs tasya kiM yadi /
03002 	na upakAras tatas tAsAM tathA syAd anavasthitiH //54// 
03003 yadi pratyupAdhy upakArakatvAni tasya na svAtmabhUtAny eva 
03004 na api tata upakAram anubhavanti / kiM tasya iti tA ucyante / upakAre 
03005 vA svAtmabhUtAbhir ayam ekaH zaktibhiH zaktIr upakurvan  
03006 eka upAdhinA api gRhyamANaH sarvAtmanA gRhyata eva / tathA 
03007 hi / ekaupAdhigrahaNe tadupakAriNyAH zakter grahaNam / tadgrahaNe  
03008 tadupakArI bhAvaH svAtmabhUtasakalazaktyupakAro gRhItaH  
03009 sarvAH zaktIr grAhayati tAz ca svaupAdhIn iti tadavasthaH
03010 prasaGgaH / atha tA api zaktyupakAriNyaH zaktayo bhinnA eva  
03011 bhAvAt / evam upAdhInAM taczaktInAM ca aparA parAsv eva zaktiSv 
03012 aparyavasAnena ghaTanAt sa ekas tAbhiH kadAcid apy agRhItas 
03013 tadupakArAtmA tadvattvena na gRhyate / yadi punaH kevalAn 
03014 eva upAdhIn zabdajJAnAny upalIyeran / tadA tasya asamAveSAn na
03015 bhavati tatpratipattimukhena sarvapratipattiH / tadA api tasya  
03016 zabdair anAkSepAn na syAt tatra pravRttir iti vyarthaH zabdaprayogaH  
03017 syAt / arthakriyAAzrayo hi sarvo vidhipratiSedhAbhyAM vyavahAraH  
03018 upAdhayaz ca tatra asamarthAH samarthaz ca na eva ucyata 
03019 iti kiM zabdaprayogaiH / tataz ca upAdhayo na upAdhayaH syuH /
03020 kvacit pravRttau hi kasyacit pradhAnasya aGgabhAvAt tadapekSayA  
03021 tathA ucyate / tasya zabdair anAkSepAn na te kasyacid aGga
03022 bhUtA iti kim upAdhayaH / lakSitalakSaNAd adoSa iti cet samAnaH  
03023 prasaGgaH / sa tAvat tair upAdhibhir nAntarIyakatayA  upalakSyamANa 
03024 ekena apy upalakSaNe sarvAtmanA upalakSita iti tadavasthaH 
03025 prasaGgaH / ko hy atra vizeSaH zabdA vA enam upalakSayeyus tallakSitA 
03101 vA upAdhayaH / sa tAvat tadAnIM nizcIyate sarvaupakArakaH
03102 iti na kiMcid etat / tasmAd /  
03103 	ekaupakArake grAhye na upakAras tato apare /
03104 	dRSTe tasminn adRSTA ye tadgrahe sakalagrahaH //55//  
03105 iti saMgrahazlokaH /
03106 	yadi bhrAntinivRttyarthaM gRhIte apy anyad iSyate /  
03107 syAd etat nirbhAgasya vastuno grahaNe ko anyas tadA na gRhIto
03108 nAma / sa tu bhrAntyA na avadhAryata iti pramANAntaraM pravartate /
03109 yady evam /
03110 	tad vyavacchedaviSayaM siddhaM tadvat tato aparam //56//
03111 	asamAropaviSaye vRtter /
03112 tat tarhi bhrAntinivRttyarthaM pravRttaM pramANam / anyasamAropa
03113 vyavacchedaphalam iti siddham anyApohaviSayam / tadvat  
03114 anyad api / asamAropaviSaye vRtteH / yatra asya samAropo na tatra
03115 nizcaya iti samAropAbhAve vartamAno anyApohaviSayaH siddhaH /
03116 	api ca nizcayaiH /
03117 	yan na nizcIyate rUpaM tat teSAM viSayaH katham //57//
03118 iyam eva khalu nizcayAnAM svArthapratipattir yat tannizcayanam /  
03119 tac ced AkArAntaravad anizcitaM kathaM tair gRhItaM /
03120 katham idAnIm anizcIyamAnaM pratyakSeNa api gRhItaM nAma / 
03121 na pratyakSaM kasyacin nizcAyakam / tad ayam api gRhNAti tan
03122 na nizcayena / kiM tarhi / tatpratibhAsena / tan na nizcayAnizcaya
03123 vazAt pratyakSasya grahaNAgrahaNe / na evaM nizcayAnAM kiMcin  
03124 nizcinvato apy anyatra anizcayena pravRttibhedAd grahaNa agrahaNam /
03125 tasmAt tad eva asya grahaNaM yo nizcayaH / anyathA ekAkAre 
03126 api tan na syAt / kiM punaH kAraNaM sarvato bhinne vasturUpe 
03127 anubhavautpattAv api tathA eva na smArto nizcayo bhavati /  sahakAri
03128 vaikalyAt tataz ca /  
03201 	pratyakSeNa gRhIte api vizeSe aMzavivarjite / 
03202 	yad vizeSAvasAye asti pratyakSaH sa pratIyate //58// 
03203 yady apy aMzarahitaH sarvato bhinnasvabhAvo bhAvo anubhUtas 
03204 tathA api na sarvabhedeSu tAvatA nizcayo bhavati / kAraNAntara 
03205 apekSyatvAt / anubhavo hi yathA avikalpAbhyAsaM nizcayapratyayAn
03206 janayati / yathA rUpadarzanAvizeSe api kuNapakAminIbhakSya
03207 vikalpAH / tatra buddhipATavaM tadvAsanAbhyAsaH prakaraNam 
03208 ity Adayo anubhavAd bhedanizcayautpattisahakAriNaH / teSAm
03209 eva ca pratyAsattitAratamyAdibhedAt paurvAparyam / yathA  
03210 janakatvAdhyApakatvAvizeSe api pitaram AyAntaM dRSTvA pitA me 
03211 Agacchati na upAdhyAya iti / so api bhavan nizcayo asati bhrAnti  
03212 kAraNe bhavati / tasmAn na abhUta ity eva sarvAkAranizcayaH /  
03213 	tatra api ca anyavyAvRttir anyavyAvRtta ity api /
03214 	zabdAz ca nizcayAz ca eva saMketam anurundhate //59//  
03215 tatra apy anyApohe na vyAvRttir anyA anya eva vyAvRttas  tadvyAvRtter
03216 nivartamAnasya tadbhAvaprasaGgAt / tathA ca vyAvRtter  
03217 abhAvaH / tasmAd yA eva vyAvRttiH sa eva vyAvRttaH / zabda 
03218 pratipattibhedas tu saMketabhedAt / na vAcyabhedo asti / 
03219 nanu ca vAcyavizeSAbhAvAt saMketabhedo apy ayukto dvayor  
03220 ekAbhidhAnAt / tathA ca vyatirekiNyA vibhakter ayogas tasyA  
03221 bhedAzrayatvAt /
03222 	dvayor ekAbhidhAne api vibhaktir vyatirekiNI / 
03223 	bhinnam artham iva anveti vAcye lezavizeSataH //60//
03224 na vai zabdAnAM kAcid viSayasvabhAvAyattA vRttir icchAto vRtty 
03225 abhAvaprasaGgAt / te yathA vyatirikte avyatirikte vA prayoktum 
03226 iSyante tathA niyuktAs tam artham pratibandhena prakAzayanti /
03227 tena gaur gotvam ity ekArthAbhidhAne api kasyacid  
03228 vizeSasya pratyAyanArthaM kRte saMketabhede vyatiriktArthA
03301 vibhaktir arthAntaram iva adarzayantI pratibhAty anarthAntare 
03302 api tathA prayogadarzanAbhyAsAt / na tAvatA sarvatra bhedaH /  
03303 anyatra api puruSaicchAvazAt pravRttasya pratibandhAbhAvAt / yathA
03304 ekaM kvacid ekavacanena khyApyate tadavizeSe api gaurava  
03305 AdikhyApanArthaM bahuvacanena / prayojanAbhAvAt tu  saMketa 
03306 bhedo na syAt / tad apy asty eva / tathA hi /
03307 	bhedAntarapratikSepApratikSepau tayor dvayoH /
03308 	saMketabhedasya padaM jJAtRvAJcAnurodhinaH //61// 
03309 yadA ayaM pratipattA tadanyavyavacchedabhAvAnapekSaH piNDa 
03310 vizeSe azvavyavacchedamAtraM jijJAsate tathA abhUtajJApanArthaM 
03311 tathA AkRtasaMketena zabdena prabodhyata anazvatvam asya asti 
03312 iti / yadA punar vyavacchedAntarAnirAkAGkSas taM jJAtum icchati
03313 tadA aparityaktavyavacchedAntare tatra eva azvavyavacchede  
03314 tathA aprakAzanAya prayuJjate anazvo ayam iti / ata eva pUrvatra
03315 pratikSiptabhedAntaratvAc zabdavRtter na sAmAnAdhikaraNyaM  
03316 vizeSaNavizeSyabhAvo vA / gotvam asya zuklam iti / 
03317 tanmAtravizeSeNa buddhes tadAzrayabhUtAyA ekatvena apratibhAsanAt /
03318 nirAkAGkSatvAc ca / dvitIye tu bhavati / tathA saMketa  
03319 anusAreNa saMhRtasakalavyavacchedadharmair vibhAgavata ekasya  
03320 iva sandarzanena pratibhAsanAt / vyavacchedAntarasAkAGkSatvAc  
03321 ca / 
03322 	bhedo ayam eva sarvatra dravyabhAvAbhidhAyinoH / 
03323 	zabdayor na tayor vAcye vizeSas tena kazcana //62//  
03401 tasmAn na sarvatra dharmadharmivAcinoH zabdayor vAcye arthe 
03402 nizcayapratyayaviSayatvena kazcid vizeSaH / ekas tam eva pratyAyayan
03403 pratikSiptabhedAntaraH pratyAyayati / anyo pratikSepeNa
03404 ity ayaM vizeSaH /  
03405 jijJApayiSur arthaM taM taddhitena kRtA api vA /  
03406 anyena vA yadi brUyAd bhedo na asti tato aparaH //63 /  
03407 etAvantam eva ca bhedaM darzayaMs taddhitena vA darzayet 
03408 pAcakatvam iti kRtA api vA pAka ity anyena vA tathAbhUtajJApanAya
03409 svayaM kRtena samayena / na punas tathA abhidhAnamAtreNa 
03410 arthAntaram eva tad bhavati / tathAbhUtasya eva jJApanAya  
03411 zabdasya kRtasaMketatvAt / nanu ca pAcakatvam iti sambandha
03412 ucyate na pAka eva / na vai pAkena anya eva kazcit pAcako nAma abhidhIyate
03413 yAdRzo varNyate / yat punar asya abhidheyaM tat kathitaM  
03414 tad eva pAcakatvena api ity alam apratiSThair mithyAvikalpaiH /  
03415 yathA ca na arthAntarabhUtA kriyA asti tatsamavAyo vA tat 
03416 pratyAyayiSyAmaH /
03417 	tena anyApohaviSaye tadvat pakSaupavarNanam /
03418 	pratyAkhyAtaM pRthaktve hi syAd doSo jAtitadvatoH //64//  
03419 yad AhuH / anyApohe api zabdArthe tadviziSTasya abhidhAnAt tadvat 
03420 pakSauditaH sarvaH prasaGgaH samAna iti tad apy anena  prativyUDham /  
03421 tatra hy arthAntaram upAdAya anyatra vartamAno dhvanir
03422 asvAtantryAdidoSair upadrUyate / na ca arthAntaram anyasmAd
03423 vyAvRttir vyAvRttAd dvayor ekAbhidhAnAd ity uktam / katham
03424 idAnIm ekasya vyAvRttasya anyAnanugamAd anyavyAvRttiH  
03425 sAmAnyam / tadbuddhau tathA pratibhAsanAt / na vai kiMcit  
03426 sAmAnyaM nAma asti / zabdAzrayA buddhir anAdivAsanAsAmarthyAd
03427 asaMsRSTAn api dharmAn saMsRjantI jAyate / tasyAH pratibhAsa 
03501 vazena sAmAnyaM sAmAnAdhikaraNyaM ca vyavasthApyate  
03502 asadartho api / arthAnAM saMsargabhedAbhAvAt / tasya sarvasya  
03503 tatkAryakAraNatayA anyebhyo bhidyamAnA arthAH samAzrayo  
03504 dhvaniz ca aniSTaparihAreNa pravartayati ity anyApohaviSaya uktaH /
03505 tatra anapekSitabAhyatattvo buddhipratibhAsavazAd eko aneka 
03506 vyAvRttaH zabdair viSayIkriyate tadanubhavAhitavAsanAprabodha  
03507 janmabhir vikalpair adhyavasitatadbhAvArthaiH / tatra eva
03508 ca ayaM dharmadharmivyavahAraH parasparaM tattvAnyatvAbhyAm
03509 avAcyaH pratanyate / na hy anyo dharmo dharmiNo anartha  
03510 antarAbhidhAnAt na api sa eva / tadvAcinAm iva dharmavAcinAm 
03511 api vyavacchedAntarAkSepaprasaGgAt / tathA ca iSTapratyAyanAt  
03512 saMketabhedAkaraNam iti / etac zabdArthe avAcyatvaM  
03513 dharmadharmiNoH / vastuni tu svalakSaNe sAmAnyalakSaNam avAcyam 
03514 abhAvAt / nanu ca dharmadharmiNor abhede bhede vA
03515 dRSTAH SaSThyAdivibhaktayo dharmabahutvAt tatra dRSTo vacana 
03516 bhedaz ca dharmiNi na syAt / uktam atra zabdAnAM svAtantrya 
03517 abhAvAd iti / api ca /  
03518 	yeSAM vastuvazA vAco na vivakSAparAzrayAH /
03519 	SaSThIvacanabhedAdicodyaM tAn prati yuktimat //65// 
03520 yadi nAma ete zabdAH puruSaiH kvacit praNinISitA api na zakyante
03601 praNetuM vastupratibandhAd dhUmAdivat / tadA ayam upAlambhaH  
03602 syAt kathaM SaSThyAdaya iti / yadA punaH / 
03603 	yad yathA vAcakatvena vaktRbhir viniyamyate / 
03604 	anapekSitabAhyArthaM tat tathA vAcakaM vacaH //66//
03605 na hi vyatireke SaSThI bAhulye jasAdaya ity etad api  puruSAbhiprAya
03606 nirapekSaM vastusaMnidhimAtreNa svayaM pravRttam / te
03607 tu tatra tathA prayuJjata iti tatas tathA pratipattir bhavati /
03608 evam anyatra api kathaMcit taiH prayuktAs tathA eva pratItihetavo  
03609 bhavanti / tatra vAcyeSu puruSAyattavRttInAM zabdAnAm avastu 
03610 sandarzinAM yathA abhyAsaM vikalpaprabodhahetUnAM pravRtticintA  
03611 tadvazAd vastuvyavasthApanaM ca kevalaM jADya  
03612 khyApanam / tathAkRtavyavasthAH zabdA dharmadharmyAdiSu na  
03613 punar vyatirekavastubhedAd iti kuta etat / tathA vyavahAra ayogAt /
03614 na hi dharmadharmiNor bhede tattvarUpatve vA sAmAnya
03615 tatsambandhasAmAnAdhikaraNyavizeSaNavizeSyabhAvA yujyante
03616 zabdAnAM vA yathAvastuvRttAv iti vakSyAmaH /  yaz ca ayaM sarvatra
03617 vastukRtam eva zabdapravRttibhedam icchati tasya  
03618 	dArAH SaNNagarI ity Adau bhedAbhedavyavasthiteH / 
03619 	khasya svabhAvaH khatvaM ca ity atra vA kiM nibandhanam //67 //
03621 yadA ekA api strI dArA ekam api sikatAdravyaM sikatA iti vyavahAras
03622 tatra kiM bAhulyaM yena evaM bhavati / zaktibheda iti cet / 
03623 sarvatra utsannam idAnIM ekavacanam ekazakter abhAvAt / yatnaz
03624 ca vyarthaH / vastvabhedAd anyatra ekavacanam iti cet / 
03625 iha apy astu / tad ayaM nirvastuko niyamaH kriyamANaH svAtantryam 
03626 icchAyAH zabdaprayoge khyApayati / SaNNagarI iti ca
03627 kathaM bahuSv ekavacanam / na hi nagarANy eva kiMcit /
03701 kutas teSAm samAhAraH / prAsAdapuruSAdInAM vijAtIyAnAm  
03702 anArambhAt kutas tatsamudAyo dravyam / asaMyogAz ca / 
03703 na saMyogaH / prAsAdasya svayaM saMyogAtmakasya pareNa asaMyogAc 
03704 ca / tata eva saMkhyAbhAvaH / tatsaMyogapuruSaviziSTA 
03705 sattA nagaram iti cet / kim asyA niratizayAyA vizeSaNam / 
03706 sattAyAz ca ekatvAt / nagarabahutve api nagarANi iti bahuvacanaM 
03707 na syAt / dvayasya parasparasahitatA iti cet / anupakAryaupakArakaH 
03708 kaH sahAyIbhAvaH / puruSasaMyogasattAnAM ca bahutvAn  
03709 nagaram iti katham ekavacanam / tathAbhUtAnAM kvacid  
03710 arthe abhinnA zaktir asti sA nimittam iti cet / na / zakter vastu
03711 rUpAvyatirekAt / vyatireke vA anupakAryasya pAratantryAyogAt /  
03712 upakAre vA zaktyupakAriNyA api zakter vyatireka ity anavasthiter  
03713 apratipattiH / tadavyatireke vA AdyAyAm api prasaGga iti yat
03714 kiMcid etat / khasya svabhAvaH khatvam iti vyatirekAzrayA
03715 SaSThI na syAt / na hi tatra sAmAnyam asti na api vibhutvAdayo  
03716 guNAs tathA ucyante / arthAntarasya tatsvabhAvatvAyogAt / teSAM 
03717 ca niHsvabhAvatvaprasaGgAt / tasya apy arthAntarasvabhAvatve 
03718 atiprasaGgaH / tathA ca apratipattiH / evaM SaTpadArthavargAdayo
03719 api vAcyAH / na hi tatra sAmAnyaM saMkhyA saMyogo vA sambhavati /
03720 katham idAnIm asaty atizaye khasya svabhAva iti na tu
03721 kham ity eva / khasya arthAntarasAdhAraNarUpAparAmarzena kha
03722 zabdapravRttinibandhanaM rUpaM tathA jijJAsAyAm evam ucyate /
03723 nanu sarvato vyAvRttasya rUpasya abhidhAnaM na sambhavati /
03724 na vai tad eva rUpaM buddhau samarpyate / anatIndriyatva
03725 prasaGgAt / kevalam ayaM tathAbhUtaM pratyAyayiSyAmi iti zabdena  
03726 zrotary asaMsRSTatatsvabhAvaM vikalpapratibimbam arpayati /  
03727 yad Aha / adRSTArthe arthavikalpamAtram iti / na evaM
03801 pratipAdyapratipAdakAbhyAM svalakSaNaM pratipannaM pratipAditaM  
03802 vA bhavati / svargAdizravaNe api tadanubhAvinAm iva
03803 pratibhAsAbhedaprasaGgAt / tasmAd ayam apratipadyamAno api
03804 bhAvasvabhAvaM tathAbhUta eva vikalpapratibimbe tadadhyavasAyI
03805 saMtuSyati / tathAbhUtatvAd eva zabdArthapratipattes 
03806 tena etad evam ucyate zabdaH svarUpam Aha iti / na punaH  svarUpa  
03807 pratibhAsasya eva vijJAnasya jananAt / kathaM tarhi idAnIm ekAnta  
03808 vyAvRttarUpeSu bhAveSu sAmAnyaM nAma / teSAm asaMsargAd 
03809 anyasya ca abhAvAt / uktaM yAdRzaM sAmAnyam asaMsRSTAnAm  
03810 ekAsaMsargas tadvyatirekiNAM samAnatA iti / api ca / 
03811 	pararUpaM svarUpeNa yayA saMvriyate dhiyA// 
03812 	ekArthapratibhAsinyA bhAvAn Azritya bhedinaH //68 // 
03813 	tayA saMvRtanAnArthAH saMvRtyA bhedinaH svayam / 
03814 	abhedina iva AbhAnti bhAvA rUpeNa kenacit //69// 
03815 	tasyA abhiprAyavazAt sAmAnyaM sat prakIrtitam /
03816 	tad asat paramArthena yathA saMkalpitaM tayA //70//
03817 buddhiH khalu tadanyavyatirekiNaH padArthAn Azritya utpadyamAnA 
03818 vikalpikA svavAsanAprakRtiM anuvidadhatI bhinnam eSAM  
03819 rUpaM tirodhAya pratibhAsam abhinnam AtmIyam adhyasya 
03820 tAn saMsRjantI sandarzayati / sA ca ekasAdhyasAdhanatayA anya  
03821 vivekinAM bhAvAnAM tadvikalpavAsanAyAz ca prakRtir yad evam
03822 eSA pratibhAti tadudbhavA / sA ca saMvRtiH saMvriyate anayA
03823 svarUpeNa pararUpam iti / te ca tayA saMvRtabhedAH svayaM
03824 bhedino apy abhedina iva kenacid rUpeNa pratibhAnti / tad eSAM
03825 buddhipratibhAsam anurundhAnaiH buddhiparivartinAm eva bhAvAnAm
03826 AkAravizeSaparigrahAd bahir iva parisphuratAM sAmAnyam
03901 ity ucyate / katham idAnIm anyApohaH sAmAnyaM / sa eva khalv
03902 anyApohas / tam eva gRhNatI sA prakRtivibhramAd vikalpAnAM
03903 vastugrAhiNI iva pratibhAti / sA hi tadanyavivekiSv eva bhAveSu
03904 bhavantI vivekaviSayA iti gamyate / nanu bAhyA vivekino na ca  
03905 teSu vikalpapravRttir iti kathaM teSu bhavati / vyAkhyAtAraH 
03906 khalv evaM vivecayanti na vyavahartAraH / te tu svAlambanam  
03907 eva arthakriyAyogyaM manyamAnA dRzyavikalpyAv arthAv ekIkRtya  
03908 pravartante / tadabhiprAyavazAd evam ucyate / tatkAritayA 
03909 atatkAribhyo bhinnAMs tathA zabdena pratipAdayanti iti / pratibhAsabhedAdibhyas
03910 tu tattvacintakA na abhedam anumanyante / 
03911 yadi pratipattRabhiprAyo anuvidhIyate anyApoho api sAmAnyaM
03912 mA bhUt na hy evaM pratipattir iti / na vai kevalam
03913 evam apratipattiH / vyaktivyatiriktAvyatiriktaekanityavyApitAdy
03914 AkArair api na eva pratipattiH / kevalam abhinnAkArA buddhir 
03915 utpadyate / tasyAH ka Azraya ity anyApoha ucyate / tasya vastuSu 
03916 bhAvAt / avirodhAt / vyavahArasya ca zabdAzrayasya tathA
03917 darzanAt / na punar vastubhUtaM kiMcit sAmAnyaM nAma asti  
03918 yathA iyaM buddhiH pratibhAti / yasmAt / 
03919 	vyaktayo na anuyanty anyad anuyAyi na bhAsate //  
03920 na hi imA vyaktayaH parasparam anvAvizanti / bhedAbhAve na sAmAnyasya
03921 eva abhAvaprasaGgAt / anyac ca na tAbhyo vyatiriktaM
03922 kiMcit tathA buddhau pratibhAty apratibhAsamAnaM ca katham
03923 AtmanA anyaM grAhayed vyapadezayed vA / na ca taiH sambaddham
03924 ekam ity eva sAmAnyaM bhavaty atiprasaGgAd ity uktam /  
03925 abhinnAbhidhAnapratyayanimittam ekaM sAmAnyaM na sarvam
03926 iti cet / katham anyato anyatra pratyayavRttiH / tatsambandhAt / 
03927 saMkhyAkAryadravyAdiSv api prasaGgaH asAmAnya AtmakatvAn
03928 na iti cet / nanu sa eva ayaM vicAryate / ko ayaM sAmAnyAtmA 
03929 iti / tatra sati sambandhe pratyayavRttis tataH sAmAnyam ity  
03930 atra ucyate / anekasambandhebhyaH kAryadravyAdibhyo nimitta  
04001 sambhavAt pratyayavRttis tataz ca sAmAnyAtmatA / anyathA anyatra 
04002 api mA bhUt / vizeSAbhAvAt / tathA ca dravyaguNasAmAnyAnAM
04003 rUpasaMkara iti / evaM tarhi buddher eva pratibhAso jJAnarUpatvAt  
04004 san eva sAmAnyam / tan na / yasmAt 
04005 	jJAnAd avyatiriktaM ca katham arthAntaraM vrajet /  71//
04006 jJAnasya rUpaM katham arthAnAM sAmAnyam / tasya teSv abhAvAt / 
04007 tadbhAvAdhyavasAyAt tathAbhrAntyA vyavahAra iti cet / tatra 
04008 tathAjJAnautpatteH kiM nibandhanam anAzrayasya ca utpattau sarvatra  
04009 syAt / athavA jJAnAd avyatiriktam ekasmAt katham anyasya 
04010 punar jJAnasya rUpaM syAt vyaktyantarabhAvinaH / ta taz ca 
04011 jJAnAntaraM vyaktyantaraM vA avyApnuvat kathaM sAmAnyaM  
04012 syAt / 
04013 	tasmAn mithyAvikalpo ayaM artheSv ekAtmatAgrahaH / 
04014 na hy arthA vyatiriktena avyatiriktena vA kenacid AtmanA samAnAH 
04015 tathA eSAM grahaNaM mithyAvikalpa eva / 
04016 	itaretarabhedo asya bIjaM saMjJA yadarthikA //72//  
04017 yasya pratyAyanArthaM saMketaH kriyate abhinnasAdhyAn bhAvAn
04018 atatsAdhyebhyo bhedena jJAtvA tatparihAreNa pravarteta iti so 
04019 ayam itaretarabhedas tasya ekAtmatApratibhAsino mithyAvikalpasya
04020 bIjaM / tam eva gRhNan eSa vikalpaH svavAsanAprakRter evaM  
04021 pratibhAti / kathaM punar bhinnAnAm abhinnaM kAryaM yena  
04022 tadanyebhyo bhedAd abheda ity ucyate / prakRtir eSA bhavAnAM
04023 yad  
04024 	ekapratyavamarzArthajJAnAdyekArthasAdhane / 
04025 	bhede api niyatAH kecit svabhAvena indriyAdivat //73//
04101 yathA indriyaviSayAlokamanaskArA AtmA indriyamano arthatatsaMnikarSA
04102 vA asaty api tadbhAvaniyate sAmAnye rUpavijJAnam ekaM
04103 janayanti evaM ziMzapAAdayo api bhedAH parasparAnanvaye api
04104 prakRtyA eva ekaM ekAkAraM pratyabhijJAnaM janayanti anyAM vA  
04105 yathApratyayaM dahanagRhAdikAM kASThasAdhyAm arthakriyAm  
04106 na tu bhedAvizeSe api jalAdayaH zrotRAdivad rUpAdijJAne /
04107 	jvarAdizamane kAzcit saha pratyekam eva vA / 
04108 	dRSTA yathA vA oSadhayo nAnAtve api na ca aparAH //74//  
04109 yathA vA guDUcIvyaktyAdayaH saha pratyekaM vA jvarAdizamana  
04110 lakSaNam ekaM kAryaM kurvanti / na ca tatra sAmAnyam apekSante / 
04111 bhede api tatprakRtitvAt / na tadavizeSe api dadhitrapusAAdayaH / 
04112 syAd etat sAmAnyam eva kiMcit tAsu tathAbhUtAsu  
04113 vidyate tata eva tad ekaM kAryam iti / tad ayuktam /
04114 	avizeSAn /
04115 sAmAnyasya  
04116 	na sAmAnyam /  
04117 tatkAryakRt / tasya api
04118 	avizeSaprasaGgataH /
04119 	tAsAM kSetrAdibhede api / 
04120 yadi hi sAmAnyAj jvarAdizamanaM kAryaM syAt / tasya avizeSAd 
04121 vyaktInAM kSetrAdibhede api cirazIghraprazamanAdayo vizeSA
04122 guNatAratamyaM ca na syAt / vizeSe vA sAmAnyasya svabhAva
04123 bhedAt svarUpahAnam / 
04124 	dhrauvyAc ca /  
04125 sAmAnyasya  
04201 	anupakArataH //75// 
04202 yadi hy upakuryAd anAdheyavizeSasya ananyApekSaNAt sakRt sarvaM
04203 svakAryaM janayet / na vA tajjananasvabhAvam / vyaktayas
04204 tu kAladezasaMskAravazena viziSTautpattayo vizeSavat kAryaM 
04205 kuryur ity avirodhaH / tadvad arthA api kecit svabhAvabhede api  
04206 ekapratyabhijJAnAdikAm arthakriyAM kurvantas tadakAribhyo  
04207 bhedAd abhinnA ity ucyante ekena vA aneko janito atajjanyebhyo  
04208 bhedAt / kiM punar anena bhedalakSaNena sAmAnyena svalakSaNaM
04209 samAnam iti pratyeyam atha anyad eva / kiM ca ataH / yadi 
04210 svalakSaNaM kathaM vikalpasya viSayaH / anyato vA katham  
04211 arthakriyA / svalakSaNe ca anityatvAdipratIter atAdrUpyam teSAM  
04212 ca avastudharmatA / na eSa doSaH / jJAnapratibhAsinyarthe sAmAnya
04213 sAmAnAdhikaraNyadharmadharmivyavahArAH / yad etaj jJAnaM
04214 vastusvabhAvagrAhiNAnubhavena AhitAM vAsanAm Azritya  
04215 vikalpakam utpadyate atadviSayam api tadviSayam iva tadanubhava 
04216 AhitavAsanAprabhavaprakRter adhyavasitatadbhAvasvarUpam  
04217 abhinnakAryapadArthaprasUter abhinnArthagrAhI iva tadanyabheda
04218 paramArthasamAnAkAram tatra yo arthAkAraH pratibhAti bAhya
04219 iva eka iva anarthakriyAkAry api tatkArI iva vyavahAriNAM tathA adhyavasAya
04220 pravRtteH anyathA pravRttyayogAt tad arthakriyAkAritayA 
04221 pratibhAsanAt tadakAribhyo bhinnam iva na ca tat 
04222 tattvaM parIkSAnaGgatvAd iti pratipAdayiSyAmaH / te arthA buddhi
04223 nivezinas tena samAnA iti gRhyante kutazcid vyAvRttyA pratibhAsanAt 
04224 na svalakSaNam tatra apratibhAsanAt / ta eva ca kutazcid 
04225 vyAvRttAH punar anyato api vyAvRttimanto abhinnAz ca pratibhAnti  
04226 iti / svayaM asatAm api tathA buddhyA upadarzanAn
04301 mithyArtha eva sAmAnyasAmAnAdhikaraNyavyavahAraH kriyate /  
04302 sarvaz ca ayaM svalakSaNAnAm eva darzanAhitavAsanAkRto viplava  
04303 iti tatpratibaddhajanmanAM vikalpAnAm atatpratibhAsitve api  
04304 vastuny avisaMvAdo maNiprabhAyAm iva maNibhrAnteH na anyeSAm 
04305 tadbhedaprabhave saty api yathA adRSTavizeSAnusaraNaM parityajya 
04306 kiMcit sAmAnyagrahaNena vizeSAntarasamAropAd dIpa
04307 prabhAyAm iva maNibuddheH / tena na vikalpaviSayeSv artheSv  
04308 arthakriyAkAritvam / na api svalakSaNasya anityatvAdyabhAvaH / yasmAn
04309 na anityatvaM nAma kiMcid anyac calAd vastunaH / kSaNa  
04310 pratyupasthAnadharmatayA tasya tathAbhUtasya grahaNAd etad 
04311 evaM bhavaty anityo ayam anityatvam asya iti vA / taddharmatAm
04312 eva avataranto vikalpA nAnAekadharmavyatirekAn sandarzayanti /  
04313 na ca te nirAzrayAs tadbhedadarzanAzrayatvAt / na avastu
04314 dharmatA tatsvabhAvasya eva tathA khyAteH / vastunas tu  
04315 nAnAekavyatirekagraho vibhramaH syAt / tasya ekAnekakAryakAriNas
04316 tathAbhAvajijJAsAsu tathAbhAvakhyApanAya tathAkRtasthititvAt /
04317 na vastubhedAt / tasya ekasya anekatvAyogAt / anekasya ca
04318 ekatvAyogAt / vyatiriktasya ca niSedhAt / teSAM prakRtibhedAd
04319 yathAvastu zabdArthAbhyupagame sAmAnAdhikaraNyAyogAt / tad
04320 upAdher ekasya dvAbhyAm abhidhAnAd adoSa iti cet / anupakAriNi
04321 pAratantryAyogAd anupAdhiH / pAratantrye ca janyajanaka bhAvAt 
04322 saha anavasthiter dvayor anabhidhAnam / ekasya buddhyA
04323 adhyAhAre na vastuviSayaH zabdaH syAt / buddhipratibhAsa 
04324 viSayatve ca sarvaM tathA eva astu / tathA bhinnaupAdhim ata ekasya
04325 grahaNe apratibhAsanAt / upakAryaupakAriNor apy upAdhitadvatoH  
04326 saha avasthAnAd adoSa iti cet / na / niSpannasya pAratantryAbhAvAd 
04327 anupAdhitvam / na aniSpannasya svarUpAsiddheH / sarvathA asat  
04328 pAratantryam iti kalpanAAropitaM kRtvA vyavahAre sarvathA 
04329 sa eva kiM na buddhir anuvidhIyate / ekazabdapramANena viSayIkaraNe  
04330 vastubalAd azeSAkSepAt tadanyavaiyarthyaM ca na syAt /  
04401 buddhipratibhAsasya nirvastukatvAt vastusAmarthyAbhAvinAm  
04402 doSANAm aprasaGgaH / tad abhinnam ekAkAraviSayIkaraNe apy anizcita
04403 anyAkAram AkArAntarasAkAGkSabuddhigrAhyaM bhinnazabdArtha
04404 upasaMhAre apy abhinnaM buddhau pratibhAti iti sAmAnya vizeSaNavizeSyabhAva
04405 sAmAnAdhikaraNyAni yathApratIti na virudhyante
04406 dharmadharmibhedo apy asya / anekArthabhedasambhave  
04407 tadekArthabhedavidhipratiSedhajijJAsAyAM tad eva vastu 
04408 pratikSiptabhedAntareNa dharmazabdena saMcodya buddhes tathA  
04409 pratibhAsanAd vyatiriktaM dharmam iva avizeSeNa aparam asya svabhAvaM  
04410 dharmitayA vyavasthApya pradarzyate / tAvatA ca aMzena 
04411 dharmadharmiNor bhedAd bhedavatI iva buddhiH pratibhAti / na 
04412 vastubhedAt / yathoktadoSAt / tathAbhUtabhedabAhulyacodanayA
04413 vacanabhedaH sAdhyasAdhanabhedaz ca tatsvabhAvasamAzrayair 
04414 dharmapratibhAsabhedais tatsvabhAvapratipattaye kriyata iti / 
04415 	tatsvabhAvagrahaNAt yA dhIs tadarthA iva apy anarthikA //
04416 	vikalpikA atatkAryArthabhedaniSThA prajAyate //76//
04417 	tasyAM yad rUpam AbhAti bAhyam ekam iva anyataH /
04418 	vyAvRttam iva nistattvaM parIkSAnaGgabhAvataH //77// 
04419 	arthA jJAnaniviSTAs te yato vyAvRttirUpiNaH /  
04420 	tena abhinnA iva AbhAnti vyAvRttAH punar anyataH //78//
04421 	ta eva teSAM sAmAnyasamAnAdhAragocaraiH /
04422 	jJAnAbhidhAnair mithyArtho vyavahAraH pratanyate //79//
04423 	sa ca sarvaH padArthAnAm anyonyAbhAvasaMzrayaH /  
04424 	tena anyApohaviSayo vastulAbhasya ca AzrayaH //80// 
04425 	yatra asti vastusambandho yathoktAnumitau yathA / 
04426 	na anyatra bhrAntisAmye api dIpatejo maNau yathA //81//  
04427 	tatra ekakAryo aneko api tadakAryAnyatAAzrayaiH /  
04501 	ekatvena abhidhAjJAnair vyavahAraM pratAryate //82//
04502 	tathA anekakRd eko api tadbhAvaparidIpane /  
04503 	atatkAryArthabhedena nAnAdharmapratIyate //83//
04504 	yathApratItikathitaH zabdArtho asAv asann api /
04505 	sAmAnAdhikaraNyaM ca vastuny asya na sambhavaH //84// 
04506 	dharmadharmivyavasthAnaM bhedo abhedaz ca yAdRzaH /  
04507 	asamIkSitatattvArtho yathA loke pratIyate //85//
04508 	taM tathA eva samAzritya sAdhyasAdhanasaMsthitiH / 
04509 	paramArthAvatArAya vidvadbhir avakalpyate //86// 
04510 	saMsRjyante na bhidyate svato arthAH pAramArthikAH / 
04511 	rUpam ekam anekaM ca teSu buddher upaplavaH //87//
04512 	bhedas tato ayaM bauddhe arthe sAmAnyaM bheda ity api /
04513 	tasya eva ca anyavyAvRttyA dharmabhedaH prakalpyate //88//  
04514 	sAdhyasAdhanasaMkalpe vastudarzanahAnitaH / 
04515 	bhedaH sAmAnyasaMsRSTo grAhyo na atra svalakSaNam //89//
04516 	samAnabhinnAdyAkArair na tad grAhyaM kathaMcana /  
04517 	bhedAnAM bahubhedAnAM tatra ekasminn ayogataH //90//
04518 	tadrUpaM sarvato bhinnaM tathA tatpratipAdikA / 
04519 	na zrutiH kalpanA vA asti sAmAnyena eva vRttitaH //91// 
04520 iti saMgrahazlokAH / kiM punaH kAraNaM svalakSaNe zabdA na 
04521 prayujyante / yasmAt /
04522 	zabdAH saMketitaM prAhur vyavahArAya sa smRtaH /  
04523 	tadA svalakSaNaM na asti saMketas tena tatra na //92 //
04524 na hi zabdA asaMketitam arthaM prakAzayanti / saMketaz ca
04525 vyavahArArthaM kriyate api nAma itaH zabdAt kRtasaMketAd  
04526 uttarakAlaM imam arthaM pratipadyeta iti / na ca prAkkRtasambandhasya
04527 ekatra svalakSaNe zabdasya pazcAd prayogo yuktaH /  
04528 tasya dezakAlavyaktibhedAnAskandhanAt / tasmAn na svalakSaNe  
04529 samayaH / sAmAnyaM tarhi vyatiriktam avyatiriktaM vA vyApi
04530 zabdair abhidhIyate / tan na vyavahArakAlAbhAvadoSaH / na etad 
04531 asti / yasmAt
04532 	api pravarteta pumAn vijJAyArthakriyAkSamAn//  
04601 	tat sAdhanAya ity artheSu saMyojyante abhidhAyakAH //93//  
04602 na khalu vai vyasanam eva etal lokasya yad ayam asaMketayann  
04603 aprayuJjAno vA zabdAn duHkham AsIta / kiM tarhi sarva eva asya avadheya
04604 ArambhaH phalArthaH / niSphalArambhasya upekSanIyatvAt / 
04605 tad ayaM zabdAn api kvacin niyuJjAnaH phalam eva kiMcid  
04606 IhituM yuktaH / tac ca sarvaM tyAgAptilakSaNam iSTAniSTayoH /  
04607 tena ayam iSTAniSTayoH sAdhanam asAdhanaM ca jJAtvA tatra  
04608 pravRttinivRttI kuryAM kArayeyaM vA iti zabdAn niyuJjIta niyoge 
04609 vA Adriyeta / anyathA apekSaNIyatvAt / 
04610 	tatra anarthakriyAyogyA jAtiH /
04611 na hi jAtiH kvacid vAhadohaupatiSThate / na ca tAdRzaM prakaraNam 
04612 antareNa loke zabdaprayogo vyavahAreSu / vyakter azakya 
04613 codanatvAl lakSitalakSaNArthaM jAticodanA iti cet / azabdacodite  
04614 saty api sambandhe kathaM pravarteta / na hi kazcid daNDaM 
04615 chindhi ity ukte daNDinaM chinatti / na apy asambhavAd vyaktau pravRttiH /
04616 evaM hy asambaddhapralApI syAt / na tato anyatra pravRttir
04617 balIvardadohacodanAvat / na ca arthAntaracodanena arthAntarasya lakSaNam /
04618 na hi sambandhe saty api daNDazabdAd daNDini pratipattiH /  
04619 aniyatasambandhatvAt tatra na iti cet / tat tulyaM jAtAv  
04620 api / vyaktInAm apAye kevalAyA jAter avasthAnAt / bhrAtRAdi
04621 zabdAs tu sambandhizabdatvAd AkSipeyuH param / na tathA  
04622 gotvAdizrutayaH sambandhivAcinyaH / apetavyaktInAm api jAtInAM
04623 taczrutibhyo nityam anugamanaprasaGgAt / sarvadA tatsambandha 
04624 yogyatApratIter iSTam eva iti cet / sarvadA tarhi gozabdAd 
04625 apravRttiH / sahitAsahitAvasthayor vizeSeNa anAkSepAt / vyakti  
04626 sambandhinyA jAtez codanAd adoSa iti cet / sarvadA tarhi tadvizeSaNatvena
04627 avasthitA vyaktir AkSiptA eva iti tadvAn abhidheyaH syAt /  
04628 na ca jAtivyaktyoH kazcit sambandho anyonyam ajanyajanakatvena 
04629 anupakArAt / tato lakSaNam apy ayuktam / tasmAn na jAtau  
04630 zabdanivezanaM phalAbhAvAt / evaM tarhi / 
04631 tadvAn alaM  
04701 arthakriyAsv iti tatra zabdo niyojyate /
04702 	sa ca / 
04703 	sAkSAn na yojyate kasmAd /
04704 yadi vyaktau zabdanivezanaM phalavat / sa ca zabdaH kasmAt sAkSAd 
04705 vyaktiSv eva na niyujyate / kiM tatra anyena vyavadhinA / 
04706 	AnantyAc ced idaM samam //94// 
04707 syAd etad AnantyAd vyaktInAm azakyaH zabdena sambandhaH  
04708 kartum / evaM sati idam AnantyaM tadvaty api samAnam / jAtyA api
04709 hi vyaktasya eva vaktavyA ity akRtasambandhasya anabhidhAnAd 
04710 avazyaM tatra sambandhaH karaNIyaH / sa ca na zakyate /
04711 tatsambandhini karaNAt tatra api kRta eva iti cet / uktam
04712 atra sambandhe apy ekatra kRtAd anyatra apratItiH na ca sambandho 
04713 asti iti / api ca / 
04714 	tatkAriNAm atatkAribhedasAmye na kiM kRtaH /  
04715 yAm arthakriyAm adhikRtya ayam artheSu zabdAn niyuGkte tat
04716 kAriNAm arthAnAm anyebhyo bhedAt tatra eva ca eSAm abhede kiM 
04717 na zabdaH prayujyate / 
04718 	tadvaddoSasya sAmyAc ced astu jAtir alaM parA //95//
04719 syAd etat anyavyAvRtte api zabdArthe vyAvRttiviziSTasya tadvato  
04720 abhidhAnAn na tadvatpakSAd vizeSaH / ko hy atra vizeSo vyAvRttir  
04801 jAtir vyAvRttimAn jAtimAn iti / astu nAma tadvaddoSaH / jAtir  
04802 anyA mA bhUt / jAtim api hy abhyupagacchatA avazyaM bhAvAnAM
04803 bhedo abhyupagantavyaH / tadabhAve tasyA apy abhAva
04804 prasaGgAt sa ca ekasmAd bhedas tadanyeSAm abhedas tadviziSTeSv 
04805 artheSu pratipattir astu / sarvathA doSaparihArasya kartum 
04806 azakyatvAt / arthAntarAbhyupagame prayojanAbhAvAt / tad
04807 arthasya anyena sAdhanAt / tadabhyupagamasya ca avazyaM bhAvitvAt /  
04808 api ca /
04809 	tadanyaparihAreNa pravarteta iti ca dhvaniH /
04810 	ucyate tena tebhyo asya avyavacchede kathaM ca saH //96//
04811 zabdaM hy eSa prayuJjAno artheSv aniSTaparihAreNa pravarteta iti 
04812 ca prayuGkte / tatra anyatra ca pravRttyanujJAyAM tan nAmagrahaNa  
04813 vaiyarthyAt pravRttinivRttyanujJAyAM ca ekacodanA anAdarAd 
04814 avacanam eva syAd anyavyAvRttyanabhidhAne / tasmAd avazyaM  
04815 zabdena vyavacchedaz codanIyaH / sa ca abhinnas tadanyeSv
04816 iti jAtidharmo apy asti / tan niyatAbhyupagamaM niyatacodanaM  
04817 jAtyarthaprasAdhanaM ca parityajyArthAntarakalpanaM  kevalam  
04818 anarthanirbandha eva / yathAkalpanam asya ayogAt / na vai vyavacchedo 
04819 na kriyate / pravRttiviSayaM tu kathayadbhir jAtir uktA /  
04820 	vyavacchedo asti ced asya nanv etAvat prayojanam /
04821 	zabdAnAm iti kiM tatra sAmAnyena apareNa vaH //97//  
04822 nanu uktaM pravRttiviSayaH pradarzyate iti / uktam idam / ayuktaM  
04823 tu uktam / tathA hi na sA pravRttiyogyA iti niveditam etat / 
04824 taddvAreNa acodite pravRttir api pratyuktA / tadvaccodane ca 
04825 vyavadhAnaM / jAtitadvatoH pravRttiviSayatve vyAvRttitadvantau
04826 kiM na iSyete / vyAvRtter avastutvena asAdhanatvAc cet / tat  
04827 tulyaM jAteH / tadvataH sAdhanAd adoSa iti cet / tulyaM tad vyAvRttimataH /
04828 avastugrAhI ca vyAvRttivAdinAM zAbdaH pratyayaH /
04901 sa vibhramavazAd akArake api kArakAdhyavasAyI pravartayati /
04902 vastusaMvAdas tu vastuutpattyA tatpratibandhe sati bhavati /
04903 anyathA na eva asti / vastuutpatter abhrAntir iti cet / na / atatpratibhAsinas
04904 tadadhyavasAyAt / maNiprabhAyAM maNibhrAntidarzanena 
04905 vyabhicArAc ca / bhrAnter avastusaMvAda iti cet / na /
04906 yathoktena eva vyabhicArAt / vitathapratibhAso hi bhrAnti lakSaNam /
04907 tannAntarIyakatayA tu saMvAdo na pratibhAsApekSI / vastuni 
04908 tu yathAbhAvam arpitacetasaH pravRttau grAhyasya sAmAnyasya
04909 anarthakriyAyogyatvAd apravRttiH / anyatra ca pravRttAv
04910 atiprasaGgaH / tadvadgrahaNe ca sAmAnyavaiyarthyAdayaH proktAH /  
04911 jAtigrahaNe api sambandhAc zliSTAbhAsA buddhiH pravartayati 
04912 iti cet / tadA na jAtir na tadvAn ekasya api svabhAvasthiter
04913 agrahaNAd iti paravAda eva AzritaH syAt / evaM tarhy anvayinaH 
04914 kasyacid arthasya abhAvAt prakRtibhinneSv artheSu tad eva idam iti  
04915 pratyabhijJAnaM na syAt / na eSa doSaH / yasmAt / 
04916 	jJAnAdyarthakriyAM tAM tAM dRSTvA  bhede api kurvataH /  
04917 	arthAMs tadanyavizleSaviSayair dhvanibhiH saha //98//
04918 	saMyojya pratyabhijJAnaM kuryAd apy anyadarzane /  
04919 uktam etat bhede api bhAvAs tulyArthakriyAkAriNaz cakSurAdivad
04920 iti / tAm ekAm jJAnAdikAm arthakriyAM teSu pazyato vastudharmatayA
04921 eva anyebhyo bhidyamAnA bhAvAs tadvyAvRttiviSayadhvani 
04922 saMsRSTaM tad eva idam iti svAnubhavavAsanAprabodhena saMsRSTa  
04923 bhedaM mithyApratyayaM janayanti / anyathA na bheda 
04924 saMsargavatI buddhiH syAt / yathA daNDiSu / na hi tatra ekadaNDa  
04925 yoge apy anyatra sa eva ayam iti bhavati / kiM tarhi / tad iha iti /
04926 na ca evaM pratyabhijJAnam / kiM tarhi / tad eva idam iti / tan na 
04927 tad ekaM anekatra pazyato api bhedasaMsargavad yuktam / vibhrama 
05001 balAt tu tathA jJAne na virodhaH / nimittAbhAvAd vibhramo 
05002 ayukta iti cet / ta eva bhAvas tadekArthakAriNo anubhava dvAreNa  
05003 prakRtyA vibhramaphalAyA vikalpavAsanAyA hetutvAn nimittam /  
05004 marIcikAdiSv api hi jalAdibhrAntes tAv ekAbhinna AkAraparAmarza
05005 pratyayanimittAnubhavajananau bhAvau kAraNaM bhinnAv
05006 api / na hi tatra anyad eva kiMcit sAmAnyam asti yat tathA 
05007 pratIyeta / sattve vA sadarthagrAhiNI buddhir bhrAntir na syAt /  
05008 abhUtAkArasamAropAd bhrAntir iti cet / na tarhi sA  tatsAmAnya
05009 grAhiNI / yam eva khalv AkAram iyam Aropayati sa eva asyA  
05010 viSaya iti / aviSayIkRtasya azakyasamAropAt / AkArAntaravat / sa 
05011 ca tatra na asti ity asAmAnyaM / sati sAmAnyagrahaNe tadAropo
05012 na anyathA atiprasaGgAd iti cet / saty ekakAryakArigrahaNa iti  
05013 kiM na iSyate / avazyaM ca icchatA api sAmAnyaM vyaktInAm eka 
05014 kAryajananazaktir eSTavyA / tatas ta eva anyebhyo bhidyamAnAs 
05015 tAdRzaM pratyayaM janayanti iti kim atra sAmAnyena / yathAbhAvam  
05016 eva asaMsRSTabhedaM kiM na pratyeti iti cet / azaktir eSA
05017 vikalpAnAm avidyAprabhavAd / na vai bAhyApekSA eva bhrAntayo 
05018 bhavanti / kiM tu viplavAd AntarAd api kezAdivibhramavat / avidyA
05019 udbhavAd viplavatve cakSurvijJAnAdiSv api prasaGgaH / na /  
05020 tasyA vikalpalakSaNatvAt / vikalpa eva hy avidyA / sA svabhAvena 
05101 eva viparyasyati / na evam indriyajJAnAni vikalpakAni / na
05102 vA teSv apy eSa doSo advayAnAM dvayanirbhAsAd iti vakSyAmaH /
05103 sarveSAM viplave api pramANatadAbhAsavyavasthA A AzrayaparAvRtter
05104 arthakriyAyogyAbhimatasaMvAdanAt / mithyAtve api
05105 prazamAnukUlatvAn mAtRsaMjJAAdivat / marIcikAyAM jalajJAnasya anyasya
05106 ca bhinnabhAvautpatter vibhramasya ca avizeSe apy abhimata
05107 arthakriyAyogyAyogyautpatter arthasaMvAdA itarau / ayogyAt /  
05108 katham utpattir iti cet / vikalpAnAm arthapratibandhaniyamAbhAvAt /
05109 na hi vikalpA yathArtham eva jAyante / sati marIcikAdarzane
05110 jalabhrAntir iti tadudbhavA ity ucyate na tu yathAsvabhAvam  
05111 ajalavivekinArthena svabhAvAnukArapratyarpaNena jananAt /  
05112 sA tu vizeSalakSaNApATavAt pratyayApekSiNA svavAsanA
05113 prabodhena janyate / tasmAd bhinnabhAvajanmano / vikalpavibhramAt
05114 tad eva idam iti pratyabhijJAnaM na vyatiriktasya sAmAnyasya  
05115 darzanAt na avyatiriktasya vyaktivad ananvayAt / api 
05116 ca / 
05117 	parasya api na sA buddhiH sAmAnyAd eva kevalAt //99//  
05118 na hi paro api enAM buddhiM kevalasAmAnyabhAvinIM  vaktum 
05119 arhati /
05120 	nityaM tanmAtravijJAne vyaktyajJAnaprasaGgataH /  
05121 yadi hi nityam anayA buddhyA sAmAnyam eva gRhyeta apratItA 
05122 eva vyaktiH syAd anena jJAnena / 
05123 	tadA kadAcit sambaddhasya agRhItasya tadvataH //100// 
05124 	tadvattA anizcayo na syAd vyavahAras tataH katham /
05125 yadA sAmAnyagrAhiNo vijJAnasya na bheda AlambanabhAvena upayujyate
05201 tadA na tau kadAcid api zliSTau gRhItAv  iti idam asya 
05202 sAmAnyam ayaM vA tadvAn iti na syAt / tathA ca tatpratipattyA  
05203 tadvati pratipattir na syAd arthAntaravat /
05204 	ekavastusahAyAz ced vyaktayo jJAnakAraNam //101// 
05205 syAd etad bhavanti vyaktayas tasya AlambanabhAvena kAraNaM na  
05206 tu kevalAH / yadA punar AsAm ekaM sahakAry asti tadA tatsahitA 
05207 gRhyanta iti /
05208 	tad ekaM vastu kiM tAsAM nAnAtvaM samapohati / 
05209 	nAnAtvAc ca ekavijJAnahetutA tAsu na iSyate //102// 
05210 kiM vai teSAM bhedAnAM tena ekena nAnAtvaM nirAkriyate /
05211 nAnAtvaM hi teSv ekavijJAnAkAraNatve kAraNam ucyate / 
05212 	anekam api yady ekam apekSyAbhinnabuddhikRt /  
05213 na brUmo anekam ekaM pratyayaM na janayati bhedAd iti /  
05214 na bhedo jananavirodhI / kiM tarhi / kaivalyam / tena ekena 
05215 sahitA janayanty eva / evaM tarhi /  
05216 	tAbhir vinA api pratyekaM kriyamANAM dhiyaM prati //103//  
05217 	tena ekena api sAmarthyaM tAsAM na ity agraho dhiyA /
05218 katham idAnIM vyaktInAM tatra jJAne sAmarthyagatiH / pratyekaM  
05219 tAsAm abhAve api tadbhAvAt / asati sAmAnye abhAvAd itarathA 
05220 ca bhAvAt / na eSa doSaH / yathA nIlAdiSv ekApAye api cakSurvijJAnaM
05221  bhavati iti / na samUhe api teSAm asAmarthyam /  
05222 tathA iha api pratyekam ekApAye api bhavati iti na sarvadA sAmarthyaM /
05223 viSama upanyAsaH / tathA hi / 
05301 	nIlAder netravijJAne pRthak sAmarthyadarzanAt //104// 
05302 	zaktisiddhiH samUhe api na evaM vyakteH kathaMcana /
05303 nIlAdInAM hi cakSurvijJAne pratyekam api sAmarthyaM dRSTam  
05304 iti samUhe api zaktir aviruddhA / tathA na kadAcid vyaktayaH 
05305 sAmAnyanirapekSA anvayi vijJAnaM janayanti / tasmAd asamarthA
05306 eva vyaktayas tatra iti na tena gRhyeran /  
05307 	tAsAm anyatamApekSaM tac cec zaktaM na kevalam //105//  
05308 atha api syAt na vemarahitaH kuvindaH paTaM karoti praty ekam /
05309 vemAbhAve api kuvindaH karoti iti na tata eva paTautpattiH / tathA
05310 na kevalam ekaikavyaktyapAye vijJAnautpattAv api sAmAnyaM
05311 tadhetuH / kiM tarhi / vyaktInAm ekAM kAMcid apekSya vijJAnam 
05312 utpAdayati / evaM sati / 
05313 	tad ekam upakuryus tAH kathaM ekAM dhiyaM ca na /
05314 bhinnAnAM hy arthAnAm ekArthaupakriyA virodhinI iti sarvo ayam
05315 ArambhaH / tAz ced vyaktayo bhinnA apy ekaM sAmAnyam upakurvanti  
05316 kaH punar AsAM vijJAnena aparAdhaH kRto yat tan na upakurvanti /
05317 kim antargaDunA sAmAnyena / yathA asambhinnAnAm
05318 apy ekasAmAnyaupakaraNazaktiH evaM tad eva ekaM vijJAnaM kurvantu / 
05319 kiM ca /
05320 	kAryaz ca tAsAM prApto asau jananaM yad upakriyA //106//  
05321 na hy anatizayam AtmAnam asya pUrvavad bibhrataH kazcid 
05322 upakArako nAma / atiprasaGgAt / arthAntarajanane api  tasya
05323 kiM tena / tasya tadAzrayatve anupakAriNaH ko ayam AzrayAzrayibhAvaH /
05324 atiprasaGgo vA / upakAre api tatra eva tatpratibandha
05325 iti kim anyas tatkaraNAt tadupakArI / tadapekSasya Azrayasya 
05326 tadupayoge anupakAryatve kA iyam apekSA nAma / tadutpattidharmA 
05327 dharmA bhAvaH svabhAvapratibandhAd apekSate nAma anAdheyAtizaya
05401 AtmA parair apekSate ca iti vyAhatam etat / tasmAd yaH 
05402 kazcit kasyacit kvacit pratibandhaH sa sarvo janyatAyAm eva antarbhavati /
05403 parabhAvautpAdane tadanupakArAt / akiMcitkarasya  
05404 ca anupakArAt / tasmAd vyaktyupakRtasya sAmAnyasya vijJAnajanane  
05405 janane vyaktam asya tatkAryatAnuyujyate / kevalasya sAmarthye 
05406 api vyaktInAM kvacid apy atra sAmarthyAsiddheH agrAhyatvaM /  
05407 samarthA vyaktayo vijJAne pratibhAsanAd iti cet /
05408 katham asiddhaupakArANAM pratibhAsa iti sa eva sAmAnyAbhyupagame 
05409 cintyate / yasmAn na anupakArako viSayo atiprasaGgAt /  
05410 na aviSayasya vijJAne pratibhAsaH / anupakArakasya aviSayatve 
05411 atItAnAgatAdInAm aviSayatvam asatAm upakArAsAmarthyAd
05412 iti cet / bhavantu nAma tadviSayANi nirviSayANi / nirviSayatve
05413 api tadanukArI pratibhAsas tadrUpAnubhavAhitavAsanAutpatter
05414 AtmabhUta eva vijJAnasya / bhAvAbhAvAnuvidhAnAc ca sAmarthyaM 
05415 na pratibhAsanAt / apratibhAsino api vyaktivyatirekeNa sAmAnyasya
05416 bhAvAt / pratibhAsinAm api kezAdiviplavAnAm abhAvAt /
05417 	abhinnapratibhAsA dhIr na bhinneSv iti cen matam /
05418 na brUmo anekam ekakAryakRn na bhavati iti / kiM tarhi / na  
05419 bhinneSv artheSv arpitatadAkArA buddhir abhinnapratibhAsinI 
05420 syAt / na vai sAmAnyagrAhiNISu svalakSaNapratibhAsaH tadabhAve  
05421 api tAsAM bhAvAt AkArAntareNa ca svajJAne pratibhAsanAt
05501 anekAkArAyogAd ekasya atiprasaGgAc ca / tasmAn na iyaM bhinnArtha
05502 grAhiNy abhinnA pratibhAti tadudbhavA / atatpratibhAsiny  
05503 apy adhyavasAyavibhramAd vyavahArayati lokam / sa tu
05504 tasyAM pratibhAsamAna AkAro na artheSv asti / anyatra bhedAd abhedinaH /
05505 sa ca arUpaH / tam eva eSA gRhNatI tathA viplavata ity 
05506 uktaM prAk / api ca / vastusAmAnyavAdino api hi bhinnA eva vyaktayaH /
05507 kathaM tAsv abhinnAkArA buddhir iti tulyaM codyam /  
05508 na tulyaM tatra abhinnasya sAmAnyasya sadbhAvAt / nanu tatra 
05509 tasya AbhAsaH sato api na lakSyate / sA hi varNasaMsthAnapratibhAsavatI
05510 vibhAvyate / na ca IdRzaM sAmAnyam / na ca tato vyatiriktaH  
05511 kazcid abhinna AkAro asti / AkRtisAmAnyavAdino api vizeSavat  
05512 tasya avyatirekAd arthAntare avRttir iti bhedAn na abhinnapratibhAso
05513 yujyate / athavA astu 
05514 	pratibhAso dhiyaM bhinnaH samAnA iti tadgrahAt //107//  
05515 na eva tAsv abhinnaH pratibhAso asti samAnA iti grahAt / na hy 
05516 ekasmin pratibhAse samAnA iti yuktaM / kiM tarhi / tad eva iti / 
05517 dvayasya grahaNAd adoSa iti cet / tathA api tad iha iti syAt / na samAna
05518 iti / tad eva tAsAM sAmyam iti cet / katham anyonyasya  
05519 sAmyam / tatsambandhAd iti cet / na / pratibaddhasya sambandha  
05520 ayogAt /  atiprasaGgAc ca /
05521 	kathaM tA bhinnadhIgrAhyAH samAz ced /  
05601 nanu samAnA iti grahaNAd eva Asv abhinnapratibhAsaH / na vai  
05602 taddarzane bhinnAbhinnayoH pratibhAsanAt samAnA iti pratItiH / 
05603 kiM tarhi /
05604 	ekakAryatA /  
05605 	sAdRzyaM /
05606 na hi vayam arthajJAne dvAv AkArau pazyAmaH / apazyantaz ca
05607 katham arthadvayakalpanena AtmAnam eva vipralabhAmahe / eka
05608 kAryAs tu vyaktayaH kalpanAviSayatAm upayAntyas tathA anayA 
05609 vibhramAn mizrIkriyanta ity anavadyam etat / 
05610 	nanu dhIH kAryaM tAsAM sA ca vibhidyate //108//  
05611 pratibhAvaM / tadvat tatpratibhAsino vijJAnasya api bhedAt / katham  
05612 ekakAryAH / tad hi tAsAM kAryaM tac ca bhidyate / yad
05613 apy udakAharaNAdikam ekaM ghaTAdikAryaM tad api pratidravyaM
05614 bhedAd bhidyata eva iti na ekaM bhedAnAM kAryam asti / 
05615 na eSa doSaH / yasmAt /  
05616 	ekapratyavamarzasya hetutvAd dhIr abhedinI /
05617 	ekadhIhetubhAvena vyaktInAm apy abhinnatA //109//
05618 niveditam etad yathA na bhAvAnAM svabhAvasaMsargo asti iti /  
05619 tatra saMsRSTAkArA buddhir bhrAntir eva / tAM tu bhedinaH 
05620 padArthAH krameNa vikalpahetavo bhavanto janayanti svabhAvata
05621 iti ca / sa tv eSAm abhinno bheda ity ucyate jJAnAdeH 
05701 kasyacid ekasya karaNAt atatkArisvabhAvavivekaH / tad api  
05702 pratidravyaM bhidyamAnam api prakRtyA ekapratyavamarzasya abheda
05703 avaskandino hetur bhavad abhinnaM khyAti / tathAbhUtapratyavamarzahetor
05704 abhedAvabhAsino jJAnAder arthasya hetutvAd
05705 vyaktayo api saMsRSTAkAraM svabhAvabhedaparamArthaM svabhAvata
05706 ekaM pratyayaM janayanti ity asakRd uktam etat / tasmAd
05707 ekakAryatA eva bhAvAnAm abhedaH /  
05708 	sA ca atatkAryavizleSas /
05709 eva  
05710 	tadanyasya anuvartinaH /  
05711 vastunaH
05712 	adRSTeH pratiSedhAc ca / 
05713 na hi dRzyaM vibhAgena apratibhAsamAnam asti ity uktam etat
05714 sati vA kvacid anAzritaM kathaM jJAnahetur iti / pratiSedhasya  
05715 ca vidhAnAt tatkalpanA ayuktA iti / tasmAt  
05716 	saMketas / 
05717 api  
05718 	tadvidarthikaH //110//
05719 eva yuktaH / yo ayam anyonyaM viveko bhAvAnAM tatpratItaya  
05720 eva saMketo api kriyamANaH zobheta atatkArivivekena pravRttyarthatayA /
05721 yadi hi na tatpratItyarthaH saMketas tasya vyavahAra  
05722 kAle apy asaMsparzAn na anyaparihAreNa pravarteta / na hi sa 
05723 teSAM tebhyo vivekaH zabdena codita iti / sA ca zrutiH / 
05801 	akAryakRtitatkAritulyarUpAvabhAsinIm / 
05802 	dhiyaM vastupRthagbhAvamAtrabIjAm anarthikAm //111//
05803 	janayanty apy atatkAriparihArAGgabhAvataH /
05804 	vastubhedAzrayAc ca arthe na visaMvAdikA matA //112//  
05805 	tato anyApohaviSayA tatkartRAzritabhAvataH /  
05806 ekasvabhAvarahiteSv artheSu tam adhyAropyautpadyamAnAM mithyA 
05807 pratibhAsitvAd akAryakAriNam api tatkAryakAriNam iva adhyavasyantIM
05808 vastupRthagbhAvamAtrabIjAM samAnAdhyavasAyAM
05809 mithyAbuddhiM zrutir janayanty api tadanyaparihArAGgAbhAvAt 
05810 paramArthatas tadvyatirekiSu padArtheSu na visaMvAdikA ity ucyate /  
05811 tathA hi sa teSu vyatireko bhUtaH / sarvathA avyatikramaNIyatvAt /
05812 na eko vyatirikto avyatirikto vA sarvathA ayogAt / tasya
05813 samAvezane vastuni dUrautsRSTam eva vastu syAc zabdajJAnAbhyAm /  
05814 tadviSayAbhimatasya tasya abhAvAt / anyasya ca vastu 
05815 dharmasya kasyacid asaMsparzAt / tata eva ca sA zrutir anyApoha 
05816 viSayayA ity ucyate / anyavyAvRtteSv artheSu vyAvRttibhedam  
05817 upAdAya avizeSeNa nivezanAt / vyavahAre apy anyaparihAreNa 
05818 pravartanAt /
05819 	avRkSavyatirekeNa vRkSArthagrahaNe dvayam //113// 
05820 	anyonyAzrayam ity ekagrahAbhAve dvayAgrahaH /
05821 	saMketAsambhavas tasmAd iti kecit pracakSate //114//  
05822 yady avRkSebhyo bhedo vRkSas tasya avRkSagrahaNam antareNa
05823 tathA grahItum azakyatvAt avijJAtavRkSeNa avRkSasya api tadvyavaccheda
05824 rUpasya aparijJAnAt buddhAv anArUDhe arthe na saMketaH  
05825 zakyata ity eke /
05901 	teSAm avRkSAH saMkete vyavacchinnA na vA / 
05902 ya evam ekaM vastusAmAnyam abhyupagamya itaretarAzrayam 
05903 anyavyavacchedena saMkete codayanti / teSAM tatra api saMketa 
05904 karaNe avRkSA vyavacchinnA na vA
05905 	yadi /
05906 	vyavacchinnAH kathaM jJAtAH prAg vRkSagrahaNAd Rte //115//  
05907 na hi tadA pratipattA vRkSaM vetti na avRkSaM tajjJAnAya eva
05908 tadarthitayA upagamAt / so ajAnAnaH katham avRkSavyavacchedaM 
05909 pratipadyeta saMkete / apratipattau ca aparihRtatadanya 
05910 nivezinaH zabdAd
05911 	anirAkaraNe teSAM saMkete vyavahAriNAm / 
05912 	na syAt tatparihAreNa pravRttir vRkSabhedavat //116//  
05913 na hi saMkete parAvyavacchedena nivezitAc zabdAd vyavahAre  
05914 tatparihAreNa pravRttir yuktA / ziMzapAAdibhedavat / atha api syAt /
05915 	avidhAya niSidhyAnyat pradarzya ekaM puraH sthitam /  
05916 	vRkSo ayam iti saMketaH kriyate tat prapadyate //117//
05917 	vyavahAre api tena ayam adoSa iti cet / 
05918 na vai vastusat sAmAnyavAdinA kasyacid vyavacchedena kiMcid
05919 vidhIyate / kiM tarhi / ekam agrato vyavasthitaM vastu sandarzya
05920 vRkSo ayam iti saMketaH kriyate / saMketakAle tathAdRSTam 
05921 eva ayam arthaM tatsambandhinaM vA vyavahAre api pratipadyata 
06001 iti na samAnaH prasaGgaH / na samAnaH / yasmAt tatra api / 
06002 	taruH /  
06003 	ayam apy ayam eva iti prasaGgo na nivartate //118//
06004 ekaM pradarzya ayaM vRkSa ity api bruvANo ayam apy ayam eva ity 
06005 ubhayIM gatiM na ativartate / tayoz ca sa eva doSaH / na doSaH  
06006 dRSTaviparItasya sujJAnatvAt / ekaM hi kiMcit pazyato anyatra
06007 tadAkAravivekinIM buddhim anubhavatas tato anyad iti yathAnubhavaM
06008 tadvivecano vaidharmyanizcaya utpadyate / sa hy 
06009 ayam eva vRkSa iti pradarzya vyutpAditaH / yatra eva taM na  
06010 pazyati tam eva avRkSaM svayam eva pratipadyate / na idaM  vyavaccheda
06011 vAdinaH sambhavati / ekatra dRSTasya rUpasya kvacid
06012 ananvayAd darzanena pratipattau vyaktyantare api na syAt tathA  
06013 pratItiH / evaM tarhi tatra api tulyam etat / yasmAt /
06014 	ekapratyavamarzAkhye jJAna ekatra hi sthitaH /
06015 	prapattA tadatadhetUn arthAn vibhajate svayam //119//  
06016 niveditam etat prAg yathA ete bhAvAH prakRtibhedino api jJAnAdikam
06017 ekaM kAryaM kecit kurvanti na anya iti / tAn ayaM tatra svayam 
06018 eva tadhetUn atadhetUMz ca vibhajya pratyeti / tasya / 
06019 	tadbuddhivartino bhAvAn bhAto hetutayA dhiyaH / 
06020 	aheturUpavikalAn ekarUpAn iva svayam //120//
06021 	bhedena pratipadyeta ity uktir bhede niyujyate /
06022 	taM tasyAH pratiyatI dhIr bhrAntyA ekaM vastv iva IkSyate //121//
06023 teSAM prakRtyA eva pratyayavazAt tathAbhUtavikalpakAraNAnAm 
06024 anvayAt taddraSTur buddhau viparivartamAnAn tajjJAna
06025 hetutayA tadanyavyAvRttyA ca atathAbhUtAn api tathAdhyavasitAn 
06026 avibhaktabAhyAdhyAtmikabhedAn pratipattA pratipattim anusRtya  
06027 ete vRkSA iti svaparavikalpeSv ekapratibhAsAn Adarzya vikalpa  
06028 vijJAne vyavasthitas tadvijJAnahetUn bhedena pratipadyeta
06101 ity uktim atadhetubhyo bhede niyuGkte / taM tasyAH pratipadyamAnA  
06102 buddhir vikalpikA bhrAntivazAd eva ekavastugrahiNI iva 
06103 pratibhAti / na punar ekaM vastu tatra dRzyam asti yasya darzana  
06104 adarzanAbhyAM bhinnadarzane apy eSa vRkSAvRkSavibhAgaM 
06105 kurvIta / tasya zAkhAAdipratibhAsavibhAgena daNDavad daNDiny 
06106 agrahaNAt / agRhItasya ca aparapravibhAgena anupalakSaNAt / AkRter  
06107 apy ekatra dRSTAyA anyatra draSTum azakyatvAt / tadatadvator 
06108 vRkSAvRkSatve vyaktir eka eva vRkSaH syAt / bhavatu nAma ghaTAdi  
06109 zabdeSv arthAntaravyavacchedaH / atha jJeyAdipadeSu katham / 
06110 na hy ajJeyaM kiMcid asti yato bhedaH syAt / tato bhedena 
06111 viSayIkaraNa eva tasya jJeyatvAt / na eSa doSaH / yasmAt / 
06112 	kvacin nivezanAya arthe vinivartya kutazcana / 
06113 	buddheH prayujyate zabdas tadarthasya avadhAraNAt //122//
06114 	vyartho anyathA prayogaH syAt taj jJeyAdipadeSv api / 
06115 	vyavahAraupanIteSu vyavacchedyo asti kazcana //123//  
06116 zabdaM hi prayuJjAnaH sarvo anvayavyatirekau na ativartate /
06117 tasya pravRttinivRttyarthatvAt / yadi hy ayaM na kasyacit  kutazcin 
06118 nivartayet pravartayed vA buddhiM yathAbhUtAnujJAnAt
06119 sarvavyavahAreSu na kiMcid vyAharet / vyAhArasya avadhAraNa  
06120 nAntarIyakatvAt / yathA ghaTena udakam Anayeti / yadi ghaTena aJjalinA
06121 vA udakAnayanaM yathAkathaMcid abhimataM syAt udakam 
06122 Anayety eva vaktavyaM syAt / na ghaTena iti / tathA paMzunA  
06123 vA anena vA yena kenacid AnItena artha Anayety eva syAd  anAkSipta
06124 karaNakarmakam / evam Anayanam anyad vA yat kiMcid
06125 anuSThAnaM vA abhimataM tadA Anayety api na
06126 brUyAt / vyarthatvAd vacanasya / tathA vyavahAraupanItAnAM jJeyAdi
06201 padAnAm api kenacid vyavacchedyena bhavitavyam / ananyAzaGkAyAM
06202 prayogAyogAt tatra hi yad eva mUDhamater AzaGkAsthAnaM  
06203 tad eva nivartyam / anAzaGkamAno vA kiM parasmAd 
06204 upadezam apekSate / azrotRsaMskAraM ca bruvANaH kathaM  
06205 na unmattaH / tatsaMskArAya eva zabdAnAM kRtasaMketatvAt / avyavahAra
06206 upanItaz ca na eva kazcit jJeyAdizabdo asti / vAkyagatasya
06207 padasya arthacintanAt / kva punar ete zabdAH prayujyanta iti  
06208 prayogaviSayacintAyAm anyApoha ucyate / anirdiSTaprayogaM tu  
06209 jJeyazabdasya ko artha iti prazne na kazcid arthaH / tataH kvacid  
06210 apratipatteH / tathA ghaTAdizabdAnAm api / yA api kvacit prakaraNe
06211 kevalazabdazravaNAt pratipattir dRSTaprayogAnusAreNa sAkAGkSatvAt
06212 sA aparisamAptatadviplava eva ghaTAdizabdeSu /  
06213 tAdRzo jJeyAdizabdeSv api yathAdarzanam asty eva / tasmAt
06214 sarva eva zabdaprayogaH kutazcid buddhiM nivartya kvacin nivezana 
06215 arthaH / tatsAphalyAt /
06216 	nivezanaM ca yo yasmAd bhidyate vinivartya tam /  
06217 	tadbhede bhidyamAnAnAM samAnAkArabhAsini //124// 
06218 	sa ca ayam anyavyAvRttyA gamyate tasya vastunaH /
06219 	kazcid bhAga iti prokto rUpaM na asya api kiMcana //125//  
06220 	tadgatAv eva zabdebhyo gamyate anyanivartanam / 
06221 	na tatra gamyate kazcid viziSTaH kenacit paraH //126//
06222 	na ca api zabdo dvayakRd anyonyAbhAva ity asan /  
06223 	arUpo rUpavattvena darzanaM buddhiviplavaH //127// 
06224 nivezyamAno apy eSa zabdo yasmAd bhidyate taM vinivartya 
06225 bhidyamAnAnAM bhede samAnarUpapratibhAsiny AkSiptatadanyavyAvRttir
06226 nivezyata iti / sa eva ayam arthAntaravyAvRttyA 
06227 tasya vastunaH kazcid bhAgo gamyate zabdo arthAntaranivRtti 
06301 viziSTAn eva bhAvAn Aha ity AdinA nirdiSTaH /  sa hi taM bhedaM  
06302 kathayann arthAntaravyavacchedam AkSipan eva vartate / 
06303 ekagatabhedacodanAyAs tadanyavyAvRttyAkSepanAntarIyakatvAt / 
06304 sa eva bhedas tadvyAvRttyA gato bhAgas tadgates tadupAdhitvAt  
06305 tadviziSTo gata ity ucyate / na punar arthAntaranivRttir vizeSaNa 
06306 bhUtA keSAMcid arthAnAM yayA viziSTAH zabdaiz codyante / daNDivat /
06307 dvayor hi bhidyamAnayor bhedasya ubhayagatatvAd
06308 ekabhedAbhidhAne api nAntarIyakas tadanyAkSepo bhavati iti tayor
06309 na vizeSaNavizeSyabhAvaH / ekabhedAbhidhAne apy anyavyAvRtti 
06310 gater anvayavyatirekacodanayA vyavahArAGgatAM zabdAnAM  
06311 darzayaMs tadvyAvRttyA gamyate tadviziSTo vA ity Aha / ata 
06312 eva ca zabdasya na dvau vyApArau tadanyavyAvartanaM svArtha
06313 abhidhAnaM ca / svArthAbhidhAnAd eva tadanyavyAvRttigateH /
06314 svArthasya bhedarUpatvAt / na hy anvayo avyatireko ananvayo  
06315 vA vyatirekaH / ekAnvayasya parihAryAbhAve niSphalacodanatvAt 
06316 tathA ekaparihArasya kvacit sthityabhAve / sa ca ayaM bhedo arUpaH /
06317 rUpavattvena tv asya darzanaM kevalaM buddhiviplava
06318 eva /
06401 	tena eva aparamArtho asAv anyathA na hi vastunaH / 
06402 	vyAvRttir vastu bhavati bhedo asya asmAd iti IraNAt //128//
06403 rUpaM hi paramArthaH / bhedaz ced rUpaM syAt / tadrUpaM vA
06404 syAd atadrUpaM vA / tAdrUpye tad eva iti na anyas tato bhidyeta /
06405 na hi tasya rUpam anyasya syAt / na tad eva bhedasya rUpaM /  
06406 rUpaM ca anyad eva syAt / tataz ca bhAvas tasmAd vyAvarteta /  
06407 tato asmAt tasya bheda iti na syAt / yat khalu yadbhedAd vyAvartate  
06408 tad eva tad bhavati iti so asya bheda iti ca na syAt / na 
06409 hy anyonyasya bhedo bhavati / sambandhAbhAvAt / sati vA sa  
06410 kAryakAraNabhAva iti rUpaM tajjanitaM bheda ity avizeSAt / 
06411 sarvakAryANi svakAraNAnAM vyAvRttayaH syuH / rUpAntaratve
06412 ca bhedasya tato apy asya bheda iti bhedaupAdhitvAt dravyAntaravan  
06413 na bhedaH syAt / na hi bhedaupAdhir eva bhedo ayam ata  
06414 iti vizeSanirdezAt / tataz ca upAdhyabhAve bhedasya eva abhAvaH syAt /
06415 tasmAn na vyatiriktaH / tadanyagatyabhAvAc ca vastuno na  paramArthaH /
06416 kathaM tarhy abhinnasya vastunaH zabdena codane
06417 tasya eva anyato api bhedAd anaMzayaekabhedacodane sarvabheda 
06418 gates tatra zabdapramANAntarANi vyarthAni na syuH / yasmAt / 
06419 	ekArthazleSaviccheda eko vyApriyate dhvaniH / 
06420 	liGgaM vA tatra vicchinnaM vAcyaM vastu na kiMcana //129//
06421 	yasya abhidhAnato vastusAmarthyAd akhile gatiH /  
06422 	bhaven nAnAphalaH zabda ekAdhAro bhavaty ataH //130//
06423 uktaM prAg yathA saMsRSTabAhyAdhyAtmikabhedA buddhiH svam  
06424 eva AbhAsaM vyavahAraviSayam arthakriyAyogyam adhyavasAya 
06425 zabdArtham upanayati iti / tatra eva ca te zabdAs tais tair bhrAnti 
06426 kAraNaiH saMsRSTarUpa iva AbhAti yathAsaMketaM vicchedAya vyApriyante/
06427 na ca ekasAdhyaM vyavacchedam anyaH karoti / saMketa 
06501 pratiniyamAt na ca vicchinnaM kiMcid vastv AkSipyate  
06502 yasya abhidhAnAd vastubalena akhile gatiH syAt / zabdAnAM buddhi
06503 viplavaviSayatvAt / tatra ca avastuni vastusAmarthyAbhAvAt / tathAbhUta
06504 arthadarzanadvAreNa ayaM nAnAekadharmabhedAbhedapratibhAsa
06505 viplavAnusArI vyavahAra iti tasya tatpratibandhe sati
06506 tadavyabhicAraH / tato vitathAd apy ante tathAbhUta eva vastuni
06507 jJAnasaMvAdAt / na punar bhinnAkAragrAhiNAM jJAnazabdAnAM  
06508 ekavastuviSayatvAt nAnAphalaH zabda ekAdhAro vyAghAtAt / yathA varNite
06509 tu buddhipratibhAsAzraye na doSa iti / 
06510 	vicchedaM sUcayan ekam apratikSipya vartate /  
06511 	yadA anyaM tena sa vyApta ekatvena ca bhAsate //131// 
06512 	sAmAnAdhikaraNyaM syAt tadA buddhyanurodhataH / 
06513 	vastudharmasya saMsparzo vicchedakaraNe dhvaneH //132//
06514 	syAt satyaM sa hi tatra iti na ekavastvabhidhAyini / 
06515 	buddhAv abhAsamAnasya dRzyasya abhAvanizcayAt //133// 
06516 ity antarazlokAH / 
06517 	tena anyApohaviSayAH proktAH sAmAnyagocarAH / 
06518 	zabdAz ca buddhayaz ca eva vastuny eSAm asambhavAt //134//  
06519 yadi hi vidhirUpeNa vastv eva zabdair vikalpair vA api viSayIkriyeta  
06520 so ayaM sarvArthasarvAkArapratItiprasaGgo asAmAnAdhikaraNya 
06521 Adayaz ca iti manyamAnaH praNetA nyAyazAstrasya anyApoha 
06522 viSayAv etau prAha / tathA hi / 
06523 	ekatvAd vasturUpasya bhinnarUpA matiH kutaH /
06524 	anvayavyatirekau vA na ekasya ekArthagocarau //135//
06525 tad ekam anaMzaM vastu kathaM bhinnAkArAbhir buddhibhir  
06526 viSayIkriyate / AkArabhedAzrayatvAd bhedasya / tasya ca abhedAt /
06527 tadAtmano api sAmAnyasya tadekayogakSematvAt / tad ayam  
06528 anyonyArthaparihAreNa ekaviSayayor vRttyabhAvAt sAmAnAdhikaraNya
06601 Adir na ca vastvAtmana ekasya tatra eva vRttir 
06602 avRttiz ca yuktA vyAghAtAt / na ca anyatra avartamAnaM sAmAnyaM 
06603 syAt / sAmAnyasya vRttir na vizeSasya iti cet / na / bhedAbhAvAt / tad
06604 hi ekarUpaM sAmAnyaM vA bhavet vizeSo vA / na hy asati rUpabhede  
06605 ayaM pravibhAgo yuktaH sati vA avyatireko na syAd ity uktam / 
06606 tad ayam avibhAgo anviyAd vA na vA / na punar ananvayo anvayI  
06607 ca / yo api bhinnam eva sAmAnyam Aha / tasya api /
06608 	abhedavyavahArAz ca bhede syur anibandhanAH /  
06609 yathAsvaM zabdA bhinnam artham abhidadhAnAH katham ekArtha
06610 buddhyAzrayAH syuH / arthAntarAbhidhAyinaz ca anAkSepakAH 
06611 nirAkAGkSatvAt kathaM vizeSaNavizeSyabhAvAzrayAH / 
06612 	sarvatra bhAvAd vyAvRtter na ete doSAH prasaGginaH //136// 
06613 yathA hy ekas tasmAd bhinnas tathA anyo api iti bhedasya asAmAnya
06614 doSo api na asti / pariziSTAbhAvas tu prAg eva uktaH / api ca / 
06615 	ekakAryeSu bhedeSu tatkAryaparicodane / 
06616 	gauravAzaktivaiphalyAd bhedAkhyAyAH samA zrutiH //137//  
06617 	kRtA vRddhair atatkAryavyAvRttivinibandhanA /  
06618 	na bhAve sarvabhAvAnAM svabhAvasya vyavasthiteH //138//  
06619 	yad rUpaM zAbaleyasya bAhuleyasya na asti tat / 
06620 	atatkAryaparAvRttir dvayor api ca vidyate //139// 
06621 	arthAbhedena ca vinA zabdAbhedo na yujyate /  
06622 	tasmAt tatkAryatA api iSTA atatkAryAd eva bhinnatA //140//
06623 	cakSurAdau yathA rUpavijJAnaekaphale kvacit// 
06624 	avizeSeNa tatkAryacodanAsambhave sati //141// 
06625 	sakRt sarvapratItyarthaM kazcit sAMketikIM zrutim / 
06626 	kuryAd Rte api tadrUpasAmAnyAd vyatirekiNaH //142//
06701 yo api manyate katham abhinnam artham antareNa bahuSv ekA 
06702 zrutiH teSAm asAmyAt ekavRtter anyatra pratyayAjananAt 
06703 apratyAsattike ca pratyayautpAde atiprasaGgAt teSu ca eka
06704 zabdanivezanavaiphalyAd ekArthaniyogAbhAvAt bhinna svabhAvAnAM
06705 pRthaG niyoge ca tathAcoditAnAM vibhAgAparijJAnAd 
06706 iti / tasya apy ekam asti ity eva lokena zabdo nivezanIyaH tad  
06707 vA ekam enAM zrutiM vastuzaktyA eva dhvanayati iti / na asty etat /
06708 kiM tarhi kenacit prayojanena kecic zabdAH kvacin nivezyante /  
06709 tatra yady anekam ekatra upayujyeta tad avazyaM tatra codanIyaM /  
06710 tasya pRthak pRthak codane atigauravaM syAt / na ca asya ananya
06711 sAdhAraNaM rUpaM zakyaM codayitum / na apy asya AyAsa asya  
06712 kiMcit sAphalyaM / kevalam anena tatra yogyAs te arthAz  codanIyAH / 
06713 ta ekena vA zabdena codyeran bahubhir vA iti svAtantryam
06714 atra vaktuH / tad iyam ekA zrutir bahuSu vaktRabhiprAya
06715 vazAt pravartamAnA na upAlambham arhati / na ca iyam azakya 
06716 pravartanA / icchAdhInatvAt / yadi hi na prayoktur icchA katham 
06717 iyam ekatra api pravarteta / icchAyAM vA ka enAM bahuSv api  
06718 pratibandhuM samarthaH / prayojanAbhAvAd / eva apravartanam iti 
06719 cet / uktaM prayojanam / bhinneSv ekasmAt pratItir atatprayojana
06720 bhedena ity uktam / na punaH svabhAvasya ekatvAt / yathAsvaM  
06721 vyavasthitasvabhAvAnAm anyonyarUpAzleSAt katham ekasvabhAvanimittaH
06722 zabdo bhinneSu bhaved ity uktaM prAk / atatprayojanavyAvRttis
06723 tu bhinnAnAm aviruddhA iti sa eva arthAbhedaH zabdAbhedasya  
06724 kAraNaM bhavatu / tena ime tatprayojanA ity atatprayojanebhyo
06725 bhinnA eva uktAH / na punar eSAm anyA tatkAryatA anyatra anyato 
06726 bhedAt / yathA cakSUrUpAlokamanaskAresv Atmaindriyamano artha
06801 tatsaMnikarSeSu vA rUpavijJAnaekakAryeSu tat kAryasAmAnyacodanA
06802 sambhave kuto rUpavijJAnam iti vyavahAralAghavArthaM kazcit 
06803 sAMketikIM zrutiM nivezayet yaro rUpavijJAnahetuH zaso  
06804 vA iti / api nAma sarveSAM tadhetUnAM sakRt pratItir yathA syAd  
06805 iti / na ca atra anugAmi kiMcid rUpam asti / kevalaM tadarthatayA  
06806 te bhAvA atadarthebhyo bhinnA iti bheda eva eSAm abhedaH / evaMjAtIyAz
06807 ca sarve samUhasantAnAvasthAvizeSazabdA ye samastAH 
06808 kiMcid ekaM kAryaM kurvanti teSAM tatra vizeSAbhAvAd apArthikA 
06809 vizeSacodanA iti sakRt sarveSAM niyojanArtham ekam ayaM  
06810 lokaH zabdaM teSu niyuGkte ghaTa iti / te api sajAtIyAd anyataz ca  
06811 bhedAvizeSe api tatprayojanAGgatayA tadanyebhyo bhidyanta 
06812 ity abhedAt tato avizeSeNa pratIyante / tatra ghaTasya rUpAdayo  
06813 ity api ghaTasvabhAvA rUpAdaya udakadhAraNavizeSAdikAryasamarthA
06814 iti yAvat / sAmAnyakAryasAdhanaprasiddhena AtmanA rUpa  
06815 AdizabdaiH prasiddhA viziSTakAryasAdhanAkhyena vizeSeNa viziSTAs
06816 ta evam ucyante / na punar atra anyat kiMcid yathAvarNitalakSaNaM
06817 dravyam asti tasya tAdRzasya anupalambhAt / ekavacanam  
06818 api tadekazaktisUcanArthaM saMketaparatantraM vA / tathA
06819 ye hetuphalavizeSabhUtAH kiMcid ekaM sAdhayanti sAdhyante 
06820 vA te api sakRt pratyayArthaM vrIhyAdizabdaiH kRtasaMketAH 
06821 kathyanta iti pUrvavad vAcyam / ye api pRthak samastA vA kvacid 
06822 upayujyante ta avasthAvizeSavAcibhiH sakRd eva zabdaiH pratyaya
06823 arthaM khyApyante sanidarzanAH sapratighA vA iti tadanyebhyo
06824 bhedasAmAnyena / yathA ekakAryAs tatkAryacodanAyAM
06825 tadanyabhedena ghaTAdizabdaiH kRtasamayAH / tathA kAraNApekSayA
06901 apy aneka ekena vyavahArArtham eva / yathA zAbaleyo
06902 bAhuleyaH prayatnAnantarIyakaH zabdaH kRtako vA iti / tathA 
06903 tatkAryapratiSedhena apy acAkSuSaH zabdo anityo anAtma iti /
06904 tatkAraNapratiSedhena apy asvAmikaH zUnya iti / evaM yathAyogam  
06905 anyad api vAcyam / zUnyAnityAdizabdeSu yathAkalpanaM samIhita 
06906 AkAraM buddhAv Aropya tadvyavacchedena vyapadezaH kriyate / 
06907 buddhisamIhA sandarzitavibhAgatvAt sarvasya zabdArthasya /  
06908 apratipakSadoSaupakSepAdayo durmativispanditAni ity  upekSanIyAH /
06909 atha 
06910 	ekavRtter aneko api yady ekazrutimAn bhavet /
06911 na kevalam ekakAryAs tadanyabhedAvizeSAd ekazabdena ucyante
06912 api tv ekavRttyA apy aneka ekazabdena ucyeta / ko virodhaH syAt /  
06913 uktam atra / tasya upalabhyAbhimatasya anupalabdher abhAvAt  
06914 anupalabhyatAyAM vA taddarzanAzrayA vyapadeza pratyabhijJAna
06915 Adayo na bhaveyur ity Adi / api ca /
06916 	vRttir AdheyatA vyaktir iti tasmin na yujyate //143//  
06917 yad etad ekam anekatra vartamAnam ekAM zrutiM vartayati  
06918 tasya kA iyaM vRttiH / AdheyatA vA syAt / yathA kuNDe badarANi  
06919 vartanta iti / vyaktir vA tair abhivyakter / yady AdheyatA /  
06920 	nityasya anupakAryatvAn na AdhAraH / 
06921 nityaM hi sAmAnyam iSyate / anityatve aparAparautpatter  anekatvAd
06922 bhedavad ekapratyayAyogAt / nityasya ca kiM kurvANa  
06923 AdhAraH syAt / tasya tatra samavAyAd AdhAra iti cet ko 
07001 ayaM samavAyo nAma / apRthaksiddhAnAm AzrayAzrayibhAvaH /
07002 tad eva idam anupakArakasya AzrayatvaM na sambhAvayAmaH / 
07003 atiprasaGgabhayAt / tasmAt samavAyasaMyogAv ekArthasamavAya 
07004 Adayo api vastusambandhAH kAryakAraNabhAvAn na vyatiricyante /  
07005 parasparam anyato vA anupakAriNAm apratibandhAt / apratibaddhasya  
07006 ca asambandhAt / yady apy ekArthasamavAyinAM parasparam 
07007 anupakAraH / tata ekasmAd upakAreNa bhAvyam / abhAve  
07008 yathoktadoSaprasaGgAt / ataH svaupakAradvAreNa eva param api 
07009 buddhyA saMghaTayya khyApyate / tasmAt tatra api kAryakAraNabhAva  
07010 kRta eva pratibandhaH / tad ayam AzrayaH sAmAnyasya  
07011 svAtmany anupakurvANo anapekSasya AdhAra iti yAcitakamaNDanam 
07012 etat / kathaM tarhi idAnIm ajanakaM kuNDaM badarANAm
07013 AdhAraH /
07014 	pravisarpataH /
07015 	zaktis taddezajananaM kuNDAder badarAdiSu //144//
07016 prakRtyA eva guruNo dravyasya asamAnadezakAryautpAdana dharmaNaH
07017 samAnadezakAryautpAdanabhAva AdhArakRtaH / tasmAt tatpUrva 
07018 kSaNasahakAri kuNDaM tatra eva badarakAryaM janayad AdhAra 
07019 ity ucyate / anyathA iha kuNDe badarANi ity api na syAt / na vai 
07020 tadupakArakRto ayaM vyapadezaH / kiM tarhi / saMyogakRtaH /  
07021 kiM punaH sa tayor eva saMyogaH / tAbhyAM jananAt samavAyAd 
07022 vA / sa kim ekatra eva na samavaiti janyate vA / tasya asAmarthyAt /  
07023 tad asamarthaM pRthak tat sahitam api tAdRzam eva ity anupakArakatvAn  
07024 na saMyogena tadvat syAt / sahitasya tadanyaupakArAd  
07025 vizeSautpatteH sAmarthyam / ko ayam ajanyajanakabhUtAnAm 
07026 upakAraH / svarUpasya siddher akAryatvAt pararUpakriyAyAm
07101 api tatra anupakArAt / ubhayathA kArakasya akiMcitkaratvena  
07102 anupakArakatvAd ity uktaprAyam / tasmAt sarva eva vastu 
07103 sambandhA janakasya eva upayogavizeSavazAt pravibhAgena kAryakAraNabhAvAd
07104 vyavasthApyante / tad ayaM kuNDAdInAm apy
07105 AdhArabhAvo badarAdiSu jananazaktir eva
07106 	na sambhavati sA apy atra / 
07107 na hy ayaM jananavizeSalakSaNa AdhArabhAvaH sAmAnye sAmAnya 
07108 Azrayasya sambhavati / tasya ajanyatvAt
07109 	tadabhAve apy avasthiteH /  
07110 	na sthitiH /  
07111 atha api syAt sthApaka AzrayaH sAmAnyasya tataH sthitihetutvAd  
07112 AdhAro na jananAd iti / tad ayuktam / tasya tadabhAve api  sthAnAt /
07113 patanadharmANAM hi bhAvAnAM pAtapratibandhAd ajanako api 
07114 sthApako bhavet / atra api yadi kazcit pratibandhaM na paryanuyuJjIta / 
07115 sa hi pAtapratibandho na arthAntaram eva yaH sthApayitrA
07116 kriyeta / arthAntaratve tatra eva asya upayoga iti kaH patataH pratibandhaH /
07117 pratibandhAd apAte api tulyaH paryanuyogo anavasthA  
07118 vA / tasmAt pAtAbhAvaH pAtapratibandhaH sa kathaM kenacit
07119 kriyate / abhAvaM karoti iti hi na abhAvo nAma kazcit kAryaH / tasya  
07120 kathaMcit kAryarUpatve abhAvAyogAt / tasmAd bhAva kriyApratiSedha
07121 nirdeza eSa bhAvaM na karoti iti yAvat / tathA ca ayaM na  
07122 kiMcit karoti ity akiMcitkaraz ca kaH kasya sthApako nAma / tena  
07123 ayaM kenacid apratibaddha iti na kadAcit tiSThet / tasmAt pAta  
07124 pratibandha ity api kSaNikAnAM bhAvAnAm upAdAnasamAna  
07125 dezautpAdanam ucyate / astu nAma ajananaM pAtinAM tatpratibaddhaH / 
07126 tatkaraNAd gatimato dravyasya kazcit sthApako apy astu /
07127 sAmAnyasya akriyasya kiMlakSaNAM sthitiM kurvANaH sthApakaH
07128 syAt / sthitir hi tasya svarUpApracyutir eva / sA ca na AzrayAyattA
07129 nityatvAt /  
07130 	sA apy ayuktA eva bhedAbhedavivecane //145//  
07201 astu nAma AzrayahetukA sthitiH sAmAnyasya / sA sAmAnyAd anyA
07202 vA syAd ananyA vA / sA anyA cet tAm eva sa AzrayaH karoti / sA
07203 ca apratibaddhA sAmAnya iti kim sAmAnyasya AzrayeNa / pratibaddhe
07204 vA kaH pratibandha iti vAcyam / sthitikaraNaM cet tatra api tulyaH  
07205 prasaGgo anavasthA ca / tata upakArAnavadhAraNAd asya iyaM 
07206 sthitir ity apratItiH / jananaM cet kim AzrayeNa apekSitena anupakAriNA /
07207 apekSA iti hi tatpratibandhaH sa ca anAdheyavizeSasya  sAmAnyasya 
07208 ayukta iti / kevalaM janayed iti / na asty anyaH sthitihetuH /  
07209 abhede vA sthiteH sAmAnyAt svarUpam eva tat tasya / tac ca  
07210 nityam asti iti / na sthitir asya kenacit kriyate / tasmAn na sAmAnyasya
07211 AdhAro asti / tan na AdheyatA asya vRttiH / atha punaH sato api 
07212 sAmAnyasya avyaktasya vyaktyA jJAnakAraNatvAt tadvyaktis tatra 
07213 vRttiH syAt / na yuktam evaM bhavitum / yasmAt /
07214 	vijJAnautpattiyogyatvAya Atmany anyAnurodhi yat / 
07215 	tad vyaGgyaM yogyatAyAz ca kAraNam kArakaM matam //146//
07217 	prAg eva asya ca yogyatve tadapekSA na yujyate /
07218 	sAmAnyasya avikAryasya tat sAmAnyavataH kutaH //147//
07219 na khalu vai kArakAd vyaJjakasya kazcid bhedaH / svaviSayavijJAna
07220 utpAdanasamartham aparaM sajAtIyaupAdAnApekSam anapekSaM
07221 vA janayan bhAvam eva vyaJjaka ucyate / paratra tu jJAna  
07222 jananazaktir anAkSiptA janyasya iti jananamAtreNa kArakatvam /
07223 yo hi yato vijJAnautpAdanayogyatAM pratilabhate sa cen
07224 na tasya janyaH syAt / sA asya svabhAvabhUtA yogyatA prAg eva asti
07225 iti na vijJAnajanane tam apekSeta / parabhUtAyAM ca asyAM 
07226 sa eva tato bhavati iti sthitivat prasaGgaH / tasmAd vyaJjako na  
07227 taM karoti na apy anyam ity akiMcitkaraz ca apekSyata iti vyAhatam 
07301 etat / nanv ajanakA api kAryatvAd dhUmAdayo vyaJjakAH /  
07302 satyaM vyaJjakA na tu dhUmam apekSya agnir Atmani jJAnaM  
07303 janayati / tathAbhUtasya agneH sAkSAd ajanakatvAt / kevalam 
07304 upAdAnabalena eva tatra jJAnam utpadyate na viSayabalena asaty 
07305 api tasmin bhAvAt paraMparayA liGgAnusAreNa / na api  sAmAnyalakSaNa
07306 avabhAsinAM pratyayAnAM saMnihitaviSayatA viSaya 
07307 balena utpattir vA iti niveditam etat / nivedayiSyate ca / tasmAd ye 
07308 viSayAH sAkSAd upayogena vijJAnaM janayantas tatra param  
07309 apekSante te avazyaM tata AtmAnaM pratilabhante / na ca ayam
07310 AtmapratilambhaH sAmAnyasya nityasya kutazcit sambhavati /
07311 tasmAn na tat kenacit vyaGgyam / na vai yogyatApratilambhaM
07312 sAmAnyasya vyaktiM brUmaH / kiM tarhi / svAzrayasamavAyam /
07313 svAzrayasamavetaM hi tad Atmany anyatra vA vijJAnahetur iti / 
07314 uktam atra ajanyajanakayoH ko ayam AzrayAzrayibhAvalakSaNaH 
07315 samavAya iti / svAzrayasamavAyApekSo vijJAnahetus tena janya
07316 eva syAt / tadhetoH svabhAvasya prAgabhAvAt pazcAc ca tato  
07317 bhAvAt / nityaM tatsvabhAvasadbhAve prAg api samavAyAd vijJAna 
07318 udayaprasaGgAt / na vai vyaktiH sAmAnyasya saMskArAd 
07319 vyaJjikA / kiM tarhi / tadgrAhiNa indriyasya / so api / 
07320 	aJjanAder iva vyakteH saMskAro na indriyasya ca / 
07321 	pratipatter abhinnatvAt tadbhAvAbhAvakAlayoH //148//
07322 saMskRtam aJjanAdibhir indriyaM pratipattau kiMcid atizayam
07323 AsAdayati / spaSTAspaSTabhedAt / atatkAriNaz ca atatsaMskArakatvAt /
07324 na evaM vyakter indriyasya  kazcit saMskAras tadbhAvAbhAva
07401 kAlayoH pratipattiM praty avizeSAt / viSayasaMskAras tv indriya
07402 avizeSe api tadvizeSAdhAnAd upakArI syAt / na indriyasaMskAraH /
07403 prAg adRzye darzanazaktyAdhAnAd upakAraka iti cet / so atIndriyam 
07404 arthaM darzayan kathaM na pratipatter bhedakaH / ekapratiniyame 
07405 ca sAmAnyAntarasya darzako na syAt / vyaktyA ca indriya
07406 saMskArAt taddarzane tadvyaGgyeSu sAmAnyeSu kadAcit anizcayo 
07407 na syAd ekanizcayo vA / tasyA avibhAgAyAs teSu vizeSAbhAvAt /  
07408 vyakter yadi indriyasaMskAro yadi na indriyasaMskAraH sAmAnyasya 
07409 vijJAnajananasvabhAva iti svabhAvAd apracyuter anapekSya eva  
07410 indriyasaMskAraM vijJAnaM janayet / saMskRtaindriyasahakAritvAt 
07411 kevalam asamartham iti cet / ko ayam anAdheyAtizayasya
07412 sahakArArthaH / anityA hi bhAvAH sahakAriNo viziSTAtmalAbhAt 
07413 tam apekSeran / yo hy eSAM janaka AtmA sa tadA eva tato bhavati 
07414 iti janyatA eva eSAM parasparato apekSA / sAmAnyaM punar anAsAdya  
07415 paraM nityaM tatsvabhAvaM kim iti indriyam apekSate / 
07416 hi tasya kevalasya yo na tatsvabhAvaH sa punaH kathaMcid 
07417 bhAvI vyakter indriyasaMskAraH / tatsahakAri sAmAnyaM vijJAna
07418 hetur ity api pAraMparyeNa vyakteH kAryam eva sAmAnyam 
07419 uktaM syAt / api ca / 
07420 	vyaJjakasya ca jAtInAM jAtimattA yadi iSyate /
07421 	prApto gotvAdinA tadvAn pradIpAdiH prakAzakaH //149//  
07422 yo hi yadviSayavijJAnahetuH sa tasya vyaJjakaH / vijJAnahetutvaM 
07423 gotvAdiSu pradIpAder apy asti / tejaHsaMskArApekSiNaH 
07424 cakSuSo arthapratipatteH / tataH pradIpAdayo gotvAdinA tadvantaH 
07501 syuH / na hi vyakter api jJAnahetutAM muktvA anyA kAcid asty  
07502 abhivyaktiH sAmAnyasya svabhAvAtizayasya AdhAtum azakyatvAt /  
07503 samavAyo abhivyaktir iti cet / uktA uttaram etat / tasya samavAyAyogAd
07504 iti / samavAyamAtraM hi vyaktyA saha asya jAtaM na anyaH 
07505 kazcid vizeSa iti / pUrvavat pazcAd api na jJAnahetuH syAt /
07506 samavAyAd eva jJAnahetutve svAzrayasamavAyinAm anyeSAm api 
07507 dRzyatAApattiH / tasmAj jJAnahetutA eva vyaJjakatvaM / tac ca
07508 tulyaM pradIpAdAv iti sa eva prasaGgaH / tan na AdheyatA na vyaktir
07509 vRttiH sAmAnyasya iti / avRtter na anekatra jJAnahetuH / ata eva / 
07510 	vyakter anyA athavA ananyA yeSAM jAtis tu vidyate / 
07511 	teSAM vyaktiSv apUrvAsu kathaM sAmAnyabuddhayaH //150//  
07512 vidyata eva ity avadhAraNArthas tuzabdaH / vidyamAno hi padArthaH
07513 svasAmarthyena anyatra buddhiM janayan svarUpAnukAriNIM
07514 tatsambandham apekSate / anyathA atiprasaGgAt / sa ca sAmAnyasya  
07515 satas tattvAnyattvapakSayor na sambhavati /  
07516 	ekatra dRSTasya anyatra darzanAsambhavAt /  
07517 sA hi buddhir ekabhAvinI vyaktyantaram evam Askanded bhUta  
07518 grAhiNI yadi tatra dRSTaM kiMcid anyatra pazyet / tac ca 
07519 	sataH /  
07520 	ananyatve anvayAbhAvAd anyatve apy anapAzrayAt //151//
07521 na sambhavati / svabhAvo hi svabhAvAn na tattvam anyatvaM  
07522 vA laGghayati / rUpasya atadbhUtasya anyatvAvyatikramAt / idam
07523 eva khalu rUpasya anyatvaM yan na tad AkArAntaravad avizeSAt /  
07524 tac cet sAmAnyasya rUpam ananyat tad eva tad bhavati / atattve
07525 vastvantaravad anyatvaprasaGgAt / na ca ekavyaktyAtmano  vyakty
07526 antarAnvAvezo avyaktyantaratvaprasaGgAt / tato na avyatirekiNaH 
07527 sAmAnyAd anvayinI buddhiH syAt / na api vyatirekiNas tasya  
07601 kvacid anAzrayAd anyasya api vyaGgyavyaJjakabhAvAdeH sambandhasya
07602 kenacid anupakAryasya apratibandhena abhAvAt / asambandhAc
07603 ca jJAnautpattAv api prasaGgAt / tad ayam ekavastu
07604 darzanena ekavRtteH pratyayasya anyatra vRttim icchaMs tattvAnyattve 
07605 na atikrAmati ity ayuktam etat / tasmAd iyam artheSv ekarUpA
07606 pratItir vikalpavAsanAsamutthitA bhrAntir eva / bhAvabhedo vAsanA
07607 prakRtiz ca tasya Azraya iti nirloThitam etat / kathaM tarhi idAnIM 
07608 pradhAnaIzvarAdikAryazabdA bhAveSv atadbhUtabhedeSv
07609 abhedena vartante / te api yathAsaMketAhitavAsanAupaskRtatvAd
07610 vijJAnasantateH sarvArthadarzaneSv anapekSya api tadbhedaM tathA
07611 adhyavasAyAd atathAbhUtakalpitavyavacchedena vikalpavijJAna
07612 pratibhAsiny artha upAdAnabalaprabhavavikalpasamutthitAH 
07613 pravartante / na hi teSv atathAbhUteSu kiMcid vyatiriktam  avyatiriktaM 
07614 vA sAmAnyam asti / tathAbhAvakalpanayA tu tadanya 
07615 bhedaH pratipattRabhiprAyavazAt syAt / tadabhiprAyAd eva sAmAnyaM
07616 kiM na iti cet / tena avazyaM hi tatra bhedo nAntarIyakatvAd  
07617 eSTavyaH / sa eva sAmAnyakArye paryApta iti niHprayojanA
07618 sAmAnyakalpanA / yadi satsv asatsu vA bhAveSu sAmAnyabuddhir  
07619 na iyam arthavatI kevalaM viplava eva iti na asmAkam asyA viSaya
07620 nirUpaNaM prati kazcid AdaraH kvacid avisaMvAdo asyA vastuni 
07621 kAryakAraNabhAvapratibandhAn na tathAbhUtagrAhyasamAvezAt  
07622 pratyakSavad atathAbhAve api bhAvAd iti nivedayiSyAmo niveditaM  
07623 ca / bhedaviSayatvaM punar asyA bahulaM bhinnapadArtha
07624 darzanabalena teSu bhAvAdhyavasAyAt / tathA bhAvakalpanAyAm 
07625 eva aparatra bhAvAt / api ca ayaM sAmAnyam arthAntaraM kalpayan  
07701 svAzrayamAtragataM vA kalpayet sarvagataM vA AkAzAdivat / 
07702 tatra yadi svAzrayamAtragataM ghaTatvAdizUnyeSu pradezeSu  
07703 ghaTAdyutpattau kathaM teSu bhinnadezadravyavartinaH sAmAnyasya  
07704 sambhavaH / yasmAt tat pUrvadravyAd utpitsu dravyaM 
07705 	na yAti / 
07706 niSkriyatvaupagamAt / na hy anyadravyavRtter bhAvasya tato avicalato 
07707 bhinnadezena bhAvena tadubhayAntarAlAvyApino yogo 
07708 yuktaH / prAk sa
07709 	na ca tatra AsIt asti pazcAn /  
07710 na ca tatra utpanno na kutazcid AyAta iti ka imaM vyAghAtabhAram 
07711 udvoDhuM samartho anyatra jADyAt / api ca 
07712 	na ca aMzavat /  
07713 	jahAti pUrvaM na AdhAram /  
07714 utpitsudezAd bhinnadezam / tayoz ca vartata iti / 
07715 	aho vyasanasantatiH //152//  
07716 bhinnadezayor hi bhAvayoH sambandho dvidhA bhavet / nAnAvayava 
07717 AtmatayA anyonyAbhyAm avayavAbhyAM tatsambandhAd  Aloka
07718 rajjuvaMzadaNDAdivat / na hi sa avayavatvam antareNa bhinna  
07719 dezAbhyAM yugapat kasyacid yogo yuktaH / tasya dvitIya AtmAbhAvAt /
07720 ekAtmanaz ca tatpradezavartisambandhirUpatvAt / anyathA  
07721 tatsambandhAyogAt / ekasya Adheyasya tatra sthAnaM tadA eva  
07722 tatra tena eva AtmanA asthAnam iti tatsthitAsthitAtmanor ekasya  
07801 virodhAd ayuktam etat / sarvatra sarvadA sarvAkArasthitAtmA iti
07802 cet / tatsvabhAvadarzanAzrayaH pratyayaH sarvatra sarvAkAraH 
07803 syAt / tathA ca gAm apy azva iti pratIyAt / azvasthitAtmanA dravyatvena
07804 sambandhAt tatsvabhAvapratipattyA ca tathA nizcayAt 
07805 tasya ca ekasya adRSTAkArAntarAbhAvAt / tasmAn na anavayavam aneka
07806 deze yugapad AdhIyate / pUrvAdhAratyAge tu bhinnadeze api 
07807 varteta / sa ca na abhimataH /  
07808 	anyatra vartamAnasya tato anyasthAnajanmani /  
07809 	svasmAd acalataH sthAnAd vRttir ity atiyuktimat //153//  
07810 	yatra asau vartate bhAvas tena sambadhyate api na /  
07811 	taddezinaM ca vyApnoti kim apy etan mahAdbhutam //154//  
07812 ity antarazlokau / yasya tu sarvagataM sAmAnyaM tasya api  
07813 	vyaktA eva ekatra sA vyaktyA abhedAt sarvatragA yadi / 
07814 	jAtir dRzyeta sarvatra / 
07815 na jAter nityAyAH kAcid vyaktir iti niSiddham etat / tasmAn nityam
07816 anapekSitaparaupaskArA dRzyeta vA na vA kadAcit tasmin svabhAve
07817 vyavasthAnAt / svabhAvAntarasya kutazcid anutpatteH / abhyupagamya
07818 api vyaktiM vyApiny ekatra vyaktyA bhedAbhAvAd vyaktA  
07819 eva sarvatra iti vyaktizUnyeSv api pradezeSu dRzyeta /
07820 	na ca sA vyaktyapekSiNI //155//  
07821 yadi hi vyaktyapekSiNI syAt /
07822 	vyaJjakApratipattau hi na vyaGgyaM sampratIyate /  
07823 	viparyayaH punaH kasmAd iSTaH sAmAnyatadvatoH //156//
07824 yo api hi svAzrayaindriyasaMyogApekSapratipattikaM sAmAnyam 
07825 AzrayazUnyeSu pradezeSu na dRzyata iti pratisamAdadhIta tasya apy 
07826 asty eva AzrayaindriyasaMyoga upakAraka iti tatas taddarzI yathA  
07827 sthitAM pazyet / na hi tasyAM dRzyamAnAyAm adRSTaM tadIyaM 
07901 yuktam / vyaktivyaGgyatvAt sAmAnyasya vyaJjakarahiteSu
07902 pradezeSv adarzanam ity api mithyA / tathAbhUtasya vyaGgyavyaJjakabhAvasya
07903 tatra abhAvAt / svapratipattyA aparapratipattihetur  
07904 hi vyaJjakaH pradIpAdiH svarUpazUnye deze svavyaGgyaM na 
07905 darzayat / na evaM vyaktir viparyayAt / kathaM hi sA vyaJjikA 
07906 ca syAt sAmAnyasya / tatpratipattidvAreNa ca dRzyA syAt / vyaGgyA
07907 ca sA eva prasajyate pradIpaghaTavat / kathaMcit tatpratipattim
07908 antareNa adRzyarUpatvAt / api ca anena kim asambhavad  abhisamIkSya
07909 evaM bahvAyAsaH sAmAnyavAda AzritaH / parasparato bhedAd
07910 vyatirekiNISu vyaktiSv anvayinaH pratyayasya ayogAt / katham
07911 idAnIm
07912 	pAcakAdiSv abhinnena vinA apy arthena vAcakaH / 
07913 pAcakapAThakAdiSv anyonyam ananvayiSv api zabdapratyayAnuvRttir  
07914 asti / pAcakaH pAThaka iti / na hi teSv anyad ekam abhinnam 
07915 asti yena bhinnAs tathA pratIyeran / karma asti cet / vyaktibhya 
07916 eva tarhi sa pratyayo astu / kim idAnIM karmaNA anyena vA / bhinnam
07917 abhinnapratyayahetur na bhavati ity ekaM sAmAnyam iSTam / 
07918 tad yadi bhinnam api karmAbhinnam pratyayaM janayet / vyaktibhiH 
07919 ko aparAdhaH kRto yena tAs tathA na iSyante / tAsAm ekarUpatvAt / 
07920 asya idam iti vyatirekapratItir atadAkAravizeSavatI ca na 
07921 syAd iti cet / uktam atra yathA vyatireko vizeSapratyayAz ca 
07922 yathAsvam arthAntaravivekAd iti / tasmAd vyaktivat  
07923 	bhedAn na hetuH karma asya /
07924 pAcakAdyabhedapratyayasya / tatkarma jAtir abhedAd hetur
07925 iti cet /  
07926 	na jAtiH karmasaMzrayAt //157//  
08001 na hy arthAntarasambandhinI jAtir arthAntare pratyayahetur  
08002 gotvam iva karkAdiSu / pAcakakarmasu ca karmajAtir na ca tAni  
08003 karmANi pAcakazabdena abhidhIyante / kiM tarhi / tatkarmAzrayo 
08004 dravyam / tasya karmasAmAnyasya 
08005 	zrutyantaranimittatvAt /  
08006 ca pAkaH pAka iti hi tataH syAn na punaH pAcaka iti / tasya karma 
08007 nimittatve proktam / kiM ca  
08008 	sthityabhAvAc ca karmaNaH /
08009 na hy anityaM karma sarvadA asti / tasya ca pratyayasya karmanimittatve 
08010 niruddhe karmaNi na pAcaka ity ucyeta / pacata eva karma 
08011 sadbhAvAt / tata eva 
08012 	asambandhAn na sAmAnyaM na ayuktam zabdakAraNam //158//
08013 	atiprasaGgAt /  
08014 vinaSTe hi karmaNi tat sAmAnyaM na karmaNi na kartari iti sambaddha 
08015 sambandhe apy asya na asti ity asambandhAn na zabdajJAna  
08016 hetuH / anyathA hy atiprasaGgaH syAt / atItam anAgataM vA
08017 nimittIkRtya tayoH pravRttir iti cet /  
08018 	karma api na asajjJAnAbhidhAnayoH /  
08019 nimittam / tayoH 
08020 	anaimittikatAApatteH / 
08021 asad hi nirupAkhyaM kathaM nimittaM syAt / kAryakaraNa
08101 lakSaNatvAd  vastutvasya / tat pracyutAsamprAptarUpam atItAnAgataM
08102 karmanimittam / anyac ca vyaktyAdikaM na iSTam ity
08103 animitte te syAtAm / tathA ca na jAtisiddhis tasyA jJAnAbhidhAnayoH /
08104 nimittatve na iSTatvAt / zaktiH pAcakAdizabdanimittaM  
08105 na karma na sAmAnyam iti cet  
08106 	na ca zaktir ananvayAt //159//  
08107 na hi zaktir nAma kiMcid anyad eva pAcakAdInAm / tasyA eva  
08108 pAcakAdyarthakriyAupayogena dravyasya anupayogitvaprasaGgAt / tasyAM
08109 tasya upayoga iti cet kim idAnIM zaktyA / zaktyupayogAya 
08110 zaktyantarasya vyatirekiNo abhyupagame atiprasaGgAd dravyam 
08111 eva upayujyata iti vAcyam / tasya upayoge zaktAv arthakriyAyAm  
08112 eva upayujyata iti kiM na iSyata iti kim antarAle anarthikayA zaktyA / 
08113 tasmAc zaktir iti dravyam eva tatkAryaM tac ca na anveti iti / tato 
08114 anvayI zabdo na syAt /  
08115 	sAmAnyaM pAcakatvAdi yadi prAg eva tad bhavet /
08116 	vyaktaM sattAAdivan no cen na pazcAd avizeSataH //160//
08117 atha api pAcakatvam iti sAmAnyam eva kiMcid bhavet / saty arthe  
08118 tatsamavAyasya akAdAcitkatvAt sattAAdivat / prAg eva vyaktaM 
08119 syAt / yAvanti hi sAmAnyAny arthe samavAyadharmANi tAni saha  
08120 utpAdena asya samavayanti iti samayaH / tadvyatikrame tasya 
08121 pazcAd apy avizeSAn na tatsamavAyaH syAt / tatsambandhisvabhAva  
08122 vaiguNyAd hi sa tasya prAG na AsIt tatra eva ca asya svabhAve  
08123 sthitasya pazcAd bhavati iti duranvayam etat /
08124 	kriyAupakArApekSasya vyaJjakatve avikAriNaH /
08125 	na apekSAtizaye apy asya kSaNikatvAt kriyA kutaH //161//
08126 karmaupakAram apekSya pAcakatvaM dravyeNa vyajyata iti cet /  
08201 sthirasvabhAvasya anatizayAd avizeSAdhAyini kA apekSA / atizaye  
08202 vA kSaNikatvAt karmaNaH pratikSaNaM svabhAvabhUtasya anyAnyasya  
08203 atizayasya utpattes tad api kSaNikaM syAt / tataH svautpattisthAnavinAzinaH
08204 kriyA kutaH yadapekSaM vyaJjakaM syAt / kathaM 
08205 tarhi idAnIm asty abhinne vastuni jJAnazabdayor anvayinor vRttiH / 
08206 yathA pAcakAdiSu / nanu tad eva idaM cintyate kathaM teSv
08207 api iti / cintitam etad yathA na teSu sambhavati / tat kim idAnIm
08208 animitte te syAtAm / na animitte / kiM tarhi / na bAhyatattvanimitte /
08209 yathAsvaM vAsanAprabodhAd vikalpautpattiH / tataH zabdAH /
08210 na punar vikalpAbhidhAnayor vastusattA samAzraya ity uktaprAyam
08211 etat / yathAsvaM samayavAsanAvazAd virodhirUpasamAvezena  
08212 aparAparadarzane apy anvayinos tayor darzanAt / na ca tatra
08213 tannibandhanaH kazcit svabhAvo asti parasparavirodhinor yugapad  
08214 ekatra samAvezAyogAt / aniyamena tarhi syAt / na hy animittaM  
08215 bhavat kvacid bhavati kvacin na bhavati iti niyamam  
08216 arhati / na khalu vai tad animittaM vAsanAvizeSanimittatvAt /  
08217 bAhyaM tu tathAbhUtaM dRzyaM na asti iti brUmaH /  na ca asati tasmin 
08218 na bhavitavyam / suptataimirikaupalabdheSv artheSv abhAveSu
08219 samayavAsanAAropitarUpavizeSeSu ca tathA vikalpautpatteH / na
08220 ca te asatsu utpadyanta iti sarvatra sarvAkAraH / vibhAge na eva tathAupalabdhAnAM
08221 vikalpanAt / uktaM ca atra kiMcid asmAbhiH prakRtyA  
08222 api kecid ekajJAnakAryAH svabhAvabhedAd iti / api ca /
08223 	tulye bhede yayA jAtiH pratyAsattyA prasarpati /  
08224 	kvacin na anyatra sa eva astu zabdajJAnanibandhanam //162// 
08225 ity antarazlokaH /  
08226 	na nivRttiM vihAya asti yadi bhAvAnvayo aparaH /  
08227 	ekasya kAryam anyasya na syAd atyantabhedataH //163// 
08228 yady ete bhAvA vyAvRttiM muktvA svabhAvena kenacid anvayinA
08301 zUnyAH / na eSAM bahUnAm ekaM kAryaM syAt /  yo hi tasya svabhAvo 
08302 janakaH / na hi so anyasya asti / yo asti sa na janako vyatirekas ya
08303 niHsvabhAvatvAt / yajjanakaM tad eva vastu tajjanakaM  
08304 ca aparatra na asti iti na aparaM janayet / sa hi tasya svabhAvo yo janakaH
08305 so anyasya api yadi syAt / sa tena svabhAvena tato abhinnaH 
08306 syAd iti asti svabhAvAnvayaH / 
08307 	yady ekAtmatayA anekaH kAryasya ekasya kArakaH /  
08308 	AtmA ekatra api so asti iti vyarthAH syuH sahakAriNaH //164// 
08309 yady ekasvabhAvatvAd aneka ekasya kArakaH sa teSAm abhinnaH
08310 svabhAvaH / ekasaMnidhAne apy asti iti / avaikalyAt kAraNasya eko 
08311 api janakaH syAt / yasmAt /
08312 	na apaity abhinnaM tadrUpaM vizeSAH khalv apAyinaH / 
08313 na hi tasya abhinnasvabhAvasya arthAntare vizeSo asti / vizeSo abheda
08314 hAneH / sa ca tatra apy asti iti na ekasthitAv api tasya apAyo  asti / ye
08315 vizeSAs teSAM sahasthitiniyamAbhAvAt syAd apAyaH / na ca te  
08316 janakA iSTAH / sahakAriNAm ekasvabhAvatayA janakatvAbhyupagamAt /
08317 tato janakasya sthAnAt / asthAyinaz ca ajanakatvAd  
08318 ekasthitAv api kAryautpattiH syAt / na ca bhavati / ataH 
08319 	ekApAye phalAbhAvAd vizeSebhyas tadudbhavaH //165// 
08320 tatkAryam anekasahakArisAdhAraNam ekavizeSApAye api na bhavati /
08321 punar apy avikaleSu sarveSu vizeSeSu bhavati / na tv avikale  
08322 apy abhinne rUpe / kAryaM hi kutazcid bhAvadharmi yan na  bhavati  
08323 tat tasya eva vaikalyAt / na ca abhinnasya rUpasya ekasthitAv
08324 api vaikalyam asti / avikale api tasminn abhavat tasya ajanakAtmatAM  
08325 sUcayati / yatsAkalyavaikalyAbhyAM ca kAryaM bhAvAbhAvavat 
08401 tata eva utpattiH / tasmin sati niyamena bhavatas tadanyasmAd  utpatti 
08402 kalpanAyAm atiprasaGgAt / tasmAd vizeSA eva janakA na
08403 sAmAnyaM / tatas ta eva vastu / yasmAt / 
08404 	sa pAramArthiko bhAvo ya eva arthakriyAkSamaH /  
08405 idam eva hi vastvavastunor lakSaNaM yad arthakriyAyogyatA ayogyatA
08406 ca iti vakSyAmaH / 
08407 	sa ca / 
08408 arthakriyAyogyo arthaH  
08409 	na anveti yo anveti na tasmAt kAryasambhavaH //166//  
08410 tasmAt sarvaM sAmAnyam anarthakriyAyogyatvAd avastu / vastu
08411 tu vizeSa eva tata eva tanniSpatteH / svabhAvAnanvayAt tarhy
08412 ekasya janakaM rUpam anyasya na asti ity ajanakaH syAt / janakatve
08413 vA bhedAvizeSAt sarvo janakaH syAt / na etad asti / yasmAt /  
08414 	tena AtmanA hi bhede api hetuH kazcin na ca aparaH /
08415 	svabhAvo ayam / 
08416 ekasya janakAd Atmano bhidyamAnAH sarve samaM janakA na vA 
08417 avizeSe kazcid iti syAd etad yady eSAM na vizeSaH sambhavet / tato bheda
08418 api kutazcid AtmAtizayAt kazcij janako na aparaH / sa hi  
08419 tasya svabhAvo na aparasya / na hi svabhAvA bhAvAnAM  paryanuyogam  
08420 arhanti kim agnir dahaty uSNo vA na udakam iti / etAvat tu  
08421 syAt kuto ayaM svabhAva iti / nirhetukatve anapekSiNo niyamAbhAvena  
08422 atiprasaGgAt / tasmAt svabhAvo asya svahetor ity ucyate / 
08423 tasya api tajjananAtmatA tadanyasmAd ity anAdir hetuparaMparA
08501 bhinnAnAM hi kazcid hetur na anyaH svabhAvAd ity atra na kiMcid 
08502 bAdhakam /
08503 	abhede tu syAtAM nAzaudbhavau sakRt //167// 
08504 abhedAt / svabhAvena eva vizvasya svAtmavad vibhAgautpattisthiti  
08505 nirodhAdayo na syuH tathA upalakSaNAd abhedasya / idam eva hi  
08506 bhedAbhedalakSaNam ekAkArasya api vyatireko avyatirekaz ca / 
08507 virodhinor ekAtmany asambhavAt / 
08508 	bhedo api tena na evaM ced / 
08509 na vai sarvAkArAvyatirekaM brUmo yena evaM syAt / kazcid asya AtmA
08510 bhinno na anya iti bhedAn na sahotpattyAdayaH / evaM tarhi 
08511 	ya ekasmin vinazyati /  
08512 	tiSThaty AtmA na tasya /
08513 bhedaH / sthAnAsthAnayor ekAtmAzrayatve ko anyo dharmo bhedaka 
08514 iti nAnAtvam eva kvacin na syAt / sarvAkAraviveka avivekinor vA arthayor
08515 abhyupagamAn nAma kevalaM na iSTaM syAn na vastv  
08516 ity uktam / tad ime na ekayogakSemA bhAvAH bhinnA eva  
08517 	ato na syAt sAmAnyabhedadhIH //168//
08518 tad idam arthAntaram anAyattam ajanyatvAd asya idaM sAmAnyaM  
08519 bhedo vA iti vyapadezaM na arhati / anyApohe apy eSa tulyaH prasaGga 
08520 iti cet / na tulyaH / yataH /
08521 	nivRtter niHsvabhAvatvAn na sthAnAsthAnakalpanA /
08522 na hy anyApoho nAma kiMcit tasya ca svabhAvAnuSaGgiNyaH  
08523 svabhAvasthitipracyutikalpanA na kalpante /  
08524 	upaplavaz ca sAmAnyadhiyas tena apy adUSaNA //169// 
08525 nirviSayam eva khalv idaM mithyAjJAnaM yad anekatra ekAkAram 
08601 iti na tadviSayasya abhAvAt sthitir asthitir vA / yat punar etad uktaM
08602 tajjanako hi sa tasya svabhAvaH  
08603 	yat tasya janakaM rUpaM tato anyo janakaH katham /
08604 tatra na brUmo anyasya tajjanakaM rUpaM na asti iti / kiM tarhi /
08605 yad ekasya tajjanakaM tad anyasya na ity anyo api svarUpeNa eva
08606 janako na pararUpeNa atattvAt / te yathAsvaM bhinnAz ca tajjanakAz  
08607 ca svabhAvena iti ko atra virodhaH / ekarUpavikalas tadrUpo
08608 na syAt na atatkAryaH / tena eva ca tatkAryaM kartavyam iti ko 
08609 atra nyAyaH / api ca /
08610 	bhinnA vizeSA janakA / 
08611 ity uktam / na ca te vizeSAs tena AtmanA parasparam anuyanti /
08612 yad ekasya janakaM rUpam anyasya tan na asti / na ca tAvatA
08613 ajanakAH /
08614 	apy abhedo api teSu cet //170//  
08615 syAd etat satyaM vizeSA janakAH na punas teSAM viziSTam eva  
08616 rUpaM kiM tv abhinnam api tadekazaktiyogAj janakAH / 
08617 	tena te ajanakAH proktAH / 
08618 saty api sAmAnye rUpe na tena te janakAs tasya anapAyAd  ekasthitAv
08619 api kAryautpattiprasaGgAd ity uktaM prAk / kiM ca /
08620 	pratibhAso api bhedakaH / 
08621 	ananyabhAk /
08622 utpattisthitivinAzAdibhedaz na ity apizabdAt / yo ayam abhinnAn 
08623 sarvArthAn manyate tasya ayam artheSu buddhipratibhAsabhedo
08624 viruddhadharmAdhyAsaz ca na syAt / sati vA tasminn abhede api 
08625 na kazcid bhedaH syAt / tathA ca ayaM pravibhAgo na syAd ekAtmavat /
08626 tasmAd ayaM bhinnapratibhAsAdir vizeSa eva / na ca atra aparam
08627 abhinnaM pratibhAsaM pazyAmo yadbalena abhedapratItiH 
08628 syAt / ato vizeSa eva /  
08701 	sa eva arthas tasya vyAvRttayo apare //170//
08702 	tatkAryaM kAraNaM ca uktaM tat svalakSaNam iSyate /
08703 	tattyAgAptiphalAH sarvAH puruSANAM pravRttayaH //172//  
08704 yad arthakriyAkAri tad eva vastv ity uktam / sa ca vizeSa eva /
08705 yat punar etat sAmAnyaM nAma tat tasya eva aparasmAd bhedaH /
08706 na hi tasya arthatve dRzyasya rUpAnupalakSaNaM yuktam / tad  
08707 upalakSaNakRtatvAd bhedeSv abhinnapratyayasya / api ca /
08708 	yathA abhedAvizeSe api na sarvaM sarvasAdhanam / 
08709 	tathA bhedAvizeSe api na sarvaM sarvasAdhanam //173//  
08710 yad uktam / kathaM tajjanakasvabhAvAd bhinno asya janakaH 
08711 syAt / janakatve vA avizeSAt sarvo janakaH syAd iti / uktam  
08712 atra / yady avizeSaH syAt syAd etad iti / yathA ca asya svayam
08713 abhedavAdino abhedAvizeSe api na sarvaH sarvasya janaka iti tathA 
08714 bhedAvizeSe api bhaviSyati / athavA /
08715 	bhede hi kArakaM kiMcid vastudharmatayA bhavet /
08716 	abhede tu virudhyete tasya  ekasya kriyAkriye //174// 
08717 bhedamAtrAvizeSe api svahetupratyayaniyamitasvabhAvatvAt kecid  
08718 eva kArakAH syuH na anye atatsvabhAvatvAd ity atra na eva kiMcid  
08719 viruddham asti / ekatve tu tasya tatra eva tathA kArakatvam akArakatvaM
08720 ca iti vyAhatam etat /
08721 	bhedo apy asty akriyAtaz cen na kuryuH sahakAriNaH /
08722 na vai sarvAkArAvivekaM brUmo bhedasya api bhAvAt / tasmAt
08723 kazcid akArako api iti / tathA api kathaMcid bhedAt sahakAriNo  
08724 akArakAH syuH /
08725 	paryAyeNa atha kartRtvaM sa kiM tasya eva vastunaH //175// 
08726 atha api syAn na eva kazcid akArako asti / sarveSAM sarvatra paryAyeNa
08727 upayogAt / zakter vA vipariNatAyAs tannivezinyA rUpAntareNa
08728 upayogAt / sa eva khalv ayaM paryAyo bhedAzraya ekasya katham /
08801 pariNAmo vA avyatirekiNyAH / vizeSe vA kathaMcid ekatva  
08802 hAnir iti yat kiMcid etat / kiM ca / 
08803 	atyantabhedAbhedau ca syAtAM tadvati vastuni /
08804 	anyonyaM vA tayor bhedaH sadRzAsadRzAtmanoH //176//
08805 bhAvAz ced abhinnena AtmanA svAtmabhUtena bhedinas tadvantaH  
08806 syuH tadabhinnasvabhAvAtmatvAd bhedasya api kutaH parasparaM 
08807 bhedaH / atha na sa tasya samAna AtmA / tathA sati tadAtmanA  
08808 tena api na yuktaM tathA bhavitum / tathA abhAve hy atad  
08809 dharmA syAt / na hy ayaM pravRttinivRttimAn svabhAva eko yuktaH / 
08810 na sarvAtmanA abheda eva / 
08811 	tayor api bhaved bhedo yadi /
08812 na hi kvacid asya ekAntiko bhedo abhedo vA vivekena vyavasthApanAt / 
08813 sAmAnyaM vizeSa iti /
08814 	yena AtmanA tayoH /
08815 	bhedaH sAmAnyam ity etad yadi bhedas tadAtmanA //177//  
08816 	bheda eva /
08817 yadi sAmAnyavizeSayor yam AtmAnam Azritya sAmAnyaM vizeSa
08818 iti sthitis tena AtmanA bhedas tadA bheda eva / yasmAt tau hi
08819 tayoH svAtmAnau tau ced vyatirekiNau vyatireka eva sAmAnya  
08820 vizeSayoH svabhAvabhedAt / svabhAvo hi bhAva iti / 
08821 	tathA ca syAn niHsAmAnyavizeSatA /
08822 	bhedasAmAnyayor yadvad ghaTAdInAM parasparam //178//
08823 vyatireke ca bhedasAmAnyayor na bhedaH sAmAnyAn na sAmAnyaM 
08824 bhedavat sambandhAbhAvAt parasparaM ghaTAdivad 
08825 ity uktam / api ca /
08826 	yam AtmAnaM puraskRtya puruSo ayaM pravartate /  
08827 	tatsAdhyaphalavAJchAvAn bhedAbhedau tadAzrayau //179// 
08828 	cintyete svAtmanA bhedo vyAvRttyA ca samAnatA /
08901 	asty eva vastu na anveti pravRttyAdiprasaGgataH //180 /  
08902 sarva eva gaur azvAd bhinno abhinno vA iti bhedam abhedaM vA  
08903 pRcchan vizeSam eva bhAvasya svabhAvAkhyam adhikRtya pravartate / 
08904 sa eva hi tathA ucyate / dravyatvAdayas tu na tatra zabdacoditAH 
08905 yathAsvaM pRthagabhidhAnAt / arthasya tadavyabhicArAt  
08906 tato gatiH syAt / nirloThitaM ca etad AcAryeNa / tad ayaM gavAdizabda
08907 pratyupasthApitam arthaM bhinnam abhinnaM vA pRcchann 
08908 arthAntaraupakSepeNa tatra kim iti dvimukhabuddhiH kriyate / 
08909 tasmAd yo asya AtmA ananyasAdhAraNo yaM puraskRtya puruSo  
08910 viziSTArthakriyArthI pravartate yathA gor vAhadohAdau na anya
08911 sambhavino arthasya yathA yuddhapraveze sa eva svabhAvo 
08912 yathAsvaM zabdacodito na dravyatvAdi sAmAnyam / taccodanayA
08913 tadA prAptum anabhipretatvAt / gavAdisamAvezAt tad AtmabhUtAnAM
08914 ca ananvayena tatra anubhayarUpatvAt / tam eva ca ayaM bhAvaM
08915 prakAraiH paryanuyuGkte / tasya bhede dravyatvAdyabhedo  
08916 asya abAdhaka eva / sarvatra svabhAvena bhedasya abhyupagamAt 
08917 sAmAnyasya ca vyAvRttilakSaNasya svabhAvabhUtasya ca sAmAnyasya  
08918 abhede apy uktam / svAtmanA eva abhede tu tatsvabhAvanibandhana  
08919 arthakriyArthI  samaM dvayor api pravarteta / eko api  
08920 tAm arthakriyAM tatsvabhAvatvAd eva karoti / tadanyasya api
08921 tat tulyam iti so api kiM na karoti / 
08922 	etena eva yad ahrIkAH kiM apy azlIlam Akulam / 
08923 	pralapanti pratikSiptaM tad apy ekAntasambhavAt //181//  
08924 yad ayam ahrIkaH syAd uSTro dadhi syAn na iti kim apy azlIlam
08925 ayuktam aheyaupAdeyam apariniSThAnAd AkulaM pralapanti / tad
08926 apy anena nirastaM svabhAvena ekAntabhedAt / tadanvaye vA /
08927 	sarvasya ubhayarUpatve tadvizeSanirAkRteH /
08928 	codito dadhi khAdeti kim uSTraM na abhidhAvati //182// 
09001 tathA hy uSTro api syAd dadhi na api sa eva uSTraH ye na anyo api syAd 
09002 uSTraH / tathA dadhy api syAd uSTraH na api tad eva dadhi yena anyad 
09003 api syAd dadhi / tad anayor ekasya api kasyacit tadrUpAbhAvasya  
09004 abhAvAt svarUpasya vA atadbhAvinaH svaniyatasya abhAvAt 
09005 na kazcid vizeSa iti / dadhi khAdeti codita uSTram api khAdet /
09006 	atha asty atizayaH kazcid yena bhedena vartate /
09007 	sa eva dadhi so anyatra na asti ity anubhayaM param //183// 
09008 atha anayoH kazcid atizayo asti yena ayaM tathA coditaH  kSIravikAra  
09009 eva pravartate na anyatra / sa eva atizayo arthakriyArthipravRttiviSayo
09010 dadhi / tatphalaupAdAnabhAvalakSitasvabhAvaM hi
09011 vastu dadhi iti / sa ca tAdRzaH svabhAvo anyatra na asti iti / pravRttyabhAvAd
09012 arthinaH / tasmAt tan na ubhayarUpam ity ekAntavAdaH /  
09013 api ca /
09014 	sarvAtmatve ca sarveSAM bhinnau syAtAM na dhIdhvanI /
09015 	bhedasaMhAravAdasya tadabhAvAd asambhavaH //184// 
09016 so ayam ahrIkaH kvacid apy ekam AkAraM pratiniyatam apazyan
09017 vibhAgAbhAvAd bhAvAnAM kathaM asaMsRSTAnyAkAravatyA buddhy
09018 Adhimucyeta arthAn abhilaped vA / tato bhedAgrahAt tatsaMhAra 
09019 vAdo na syAt syAd uSTro dadhi syAn na iti / atha punar asaMsRSTAv 
09020 AkArau pratipadya saMharet / ekarUpasaMsargiNyAH buddheH
09021 kvacit pratiniyamAt tatpratibhAsabhedakRta eva tayo rUpayoH
09022 svabhAvabhedo api syAt / ekAnekavyavasthiteH pratibhAsaviSayatvAt /
09023 tathA ca na ekas tadubhayarUpaH syAd iti mithyAvAda
09024 eSaH / sthitam etat na bhAvAnAM kazcit svabhAvAnvayo asti  
09025 bhedalakSaNam eva tu sAmAnyam / atha ca prakRtyA kecid  ekajJAna
09026 AdiphalAH kecin na iti / bhavatu nAma bhAvAnAM svabhAva 
09101 bhedaH sAmAnyam / yeSAM tu nirupAkhyAnAM svabhAva eva na asti
09102 tatra kathaM svabhAvabhedaviSayAH zabdAH / teSv avazyaM  
09103 zabdapravRttyA bhAvyam / kathaMcid avyavasthApiteSu  vidhipratiSedha
09104 ayogAt / tathA ca sarvatra ayam anvayavyatirekAzrayo vyavahAro  
09105 na syAt uSNasvabhAvo agnir na anuSNa ity api / svabhAvAntarasya  
09106 asataH kathaMcid avyavasthApanAt / sarvathA apratipatter
09107 agnisvabhAvasya apratipattir iti vyAmUDhaM jagat syAt / syAd etat
09108 na tatra kasyacid asato niSedhaH anuSNaM sad eva arthAntaraM
09109 niSidhyata iti / katham idAnIM sad asan nAma / na brUmaH sarvatra  
09110 asat tatra na asti iti dezakAladharmaniSedha eva sarvabhAveSu 
09111 kriyate na dharmiNaH tanniSedhe tadviSayazabdapravRttyabhAvAt
09112 anirdiSTaviSayasya naJo aprayogAt / so api tarhi dezAdi 
09113 pratiSedhaH katham / yasmAn na tatra api dezAdInAM pratiSedho 
09114 na apy arthasya / sambandho niSedhyata iti cet / nanu tanniSedhe
09115 api tulyo doSo niSedhAd asati zabdApravRttir ity Adi / asato vA asya 
09116 niSedhe tadvad dharmiNo api niSedhaH / na vai sambandhasya 
09117 na asti iti niSedhaH / kiM tarhi / na iha ghaTo na idAnIM na evam ity uktau
09118 na anena sambandho asti na etad dharmA vA iti pratItiH / tathA  
09119 ca sambandho niSiddho bhavati iti / tathA api kathaM niSiddho yAvad  
09120 asya sambandho dharmo vA na asti iti matir na bhavati / na ca asyAH  
09121 kathaMcid bhAve sambhavo abhAveSu tathA abhAvAt / tasmAt sambandha 
09122 abhAvapratIter na ayam iha ity AdyA pratItiH / sa tadabhAve 
09201 na syAt / pratItau vA tadabhAvasya / yathA pratItimatas tatprabhavAH
09202 zabdAH kena nivAryante / sa eva hi zabdAnAM na viSayo  
09203 yo na vitarkANAm / te cet pravRttAH ko vacanasya niSeddhA /  
09204 na hy avAcyam arthaM buddhayaH samIhante / sambandhasya  
09205 tu svarUpeNa anabhidhAnam uktam / abhidhAne sambandhitvena 
09206 buddhAv upasthAnAt / yathAbhiprAyam apratItiH / tad ayaM pratIyamAno
09207 api sambandhirUpa eva iti svarUpeNa na abhidhIyate / tasmAn 
09208 na abhAvavat sambandhe api prasaGgaH / api ca / ayam  
09209 abhAvam abhidheyaM bruvANaM prati pratividadhad abruvANaH  
09210 kathaM pratividadhyAt / vacane vA asya katham abhAvo anuktaH /
09211 atha abhAvam eva na icchet / tena avacanam / tad eva idAnIM  katham 
09212 abhAvo na asti iti / yat punar etad arthaniSedhe anarthakazabdAprayogAn
09213 nirviSayasya naJo aprayoga ity atra uttaraM vakSyate / tasmAt
09214 santy abhAveSu zabdAH / teSu kathaM svabhAvabheda iti /
09215 tatra api /
09216 	rUpAbhAvAd abhAvasya zabdA rUpAbhidhAyinaH /  
09217 	na AzaGkyA eva siddhAs te vyavacchedasya vAcakAH //185//  
09218 vastuvRttInAM zabdAnAM kiM rUpam abhidheyam Ahosvid bheda 
09219 iti zaGkA syAt / abhAvas tu vivekalakSaNa eva nimittIkartavyasya  
09220 kasyacid rUpasya abhAvAt tadbhAve abhAvAyogAt / tadbhAvalakSaNatvAd
09221 bhAvasya / tasmAd ayam eva sa mukhyo vivekaH /  
09222 tasya tathA abhAvakhyApinaH zabdAH kiM vivekaviSayA ity asthAnam  
09223 eva etad AzaGkAyAH / tasmAt siddham etat sarve zabdA vivekaviSayA  
09301 vikalpAz ca / ta ete ekavastupratizaraNA api yathAsvam
09302 avadhibhedaupakalpitair bhedair bhinneSv iva pratibhAtsu buddhau 
09303 vivekeSu upalayanAd bhinnaviSayA eva / tena svabhAvasya eva sAdhyasAdhanabhAve
09304 api na sAdhyasAdhanasaMsargaH / tau na pratijJArthaikadezo  
09305 hetur iti / sa ca ayaM hetutvena apadizyamAnaH / 
09306 	upAdhibhedApekSo vA svabhAvaH kevalo athavA /  
09307 	ucyate sAdhyasiddhyarthaM nAze kAryatvasattvavat //186//
09308 apekSitaparavyApAro hi svabhAvaniSpattau bhAvaH kRtakaH / 
09309 tena iyaM kRtakazrutiH svabhAvAbhidhAyiny api paraupAdhim enam  
09310 AkSipati / etena pratyayabhedabheditvAdayo vyAkhyAtAH / evam  
09311 upAdhibhedApekSaH kvacit svabhAvo hetur ucyate kvacid anapekSaH
09312 sAmAnyena yathA anityatva eva sattvam kvacit svabhAvabhUta
09313 dharmavizeSaparigraheNa yathA tatra eva utpattiH / anayA  
09314 dizA anye api svabhAvahetupravibhAgA draSTavyAH /
09315 	sattAsvabhAvo hetuz cen na sattA sAdhyate katham /  
09316 	ananvayo hi bhedAnAM vyAhato hetusAdhyayoH //187//  
09317 yadi sattvam anityatve anyatra vA hetuH syAt sAdhyam api kasmAt 
09318 na iSyate / tat kila evaM prasAdhyamAnaM vizeSIbhavati /
09319 na ca vizeSaH sAdhAyituM zakyate ananvayAt / yathA Aha pramANa  
09320 vizeSayA ajJAnAd iti / so ayaM vizeSo na sAdhya eva vyAhanyate /  
09321 kiM tarhi / hetAv api tulyadoSatvAt / na hi hetur ananvayaH siddher 
09322 aGgaM tataH saMzayAt / na eSa doSaH / yasmAt /  
09401 	bhAvaupAdAnamAtre tu sAdhye sAmAnyadharmiNi /  
09402 	na kazcid arthaH siddhaH syAd aniSiddhaM ca tAdRzam //188//  
09403 na sarvathA sattAsAdhane vizeSaH sAdhito bhavati / bhAvamAtra 
09404 vizeSaNo asti kazcid dharmI iti prasAdhayato anirdiSTasvabhAvavizeSasya
09405 kasyacit sattAmAtre virodhAbhAvAt na iha sattAsAdhana 
09406 pratiSedhaH kiM tu sa tathA asti kazcid iti kaMcana asya bhedam 
09407 aparAmRzan bruvANaH kaM svArthaM puSNAti / tasmAd anena upAtta 
09408 bheda eva sAdhyaH /
09409 	upAttabhede sAdhye asmin bhavad hetur ananvayaH /  
09410 	sattAyAM tena sAdhyAyAM vizeSaH sAdhito bhavet //189//
09411 sa hi dharmI pradhAnalakSaNa eko nityaH sukhAdyAtmako anyo  
09412 vA iti yathAkathaMcid api vizeSitas tatsvabhAvaH prasAdhito bhavati /
09413 sa ca tathA na anveti / yad api sattAmAtram anveti na tena
09414 siddhena kiMcit / nanv evam agnyAdiSv api prasaGgaH / tatra api
09415 na agnisattAmAtre kazcid vivAdaH / viziSTAdhAravizeSasya tv abhimatasya
09416 ananvayAd asiddhiH / na vai sa AdhAras taM vizeSIkaroti / 
09417 tadayogavyavacchedena vizeSaNAd ity uktaM vakSyate ca / tasmAt
09418 tatra sAmAnyam eva sAdhyate tadayogavyavacchedena / 
09419 na tathA iha api kvacit sattAyAH sAdhanam / pradhAnAdizabdavAcyasya
09420 eva arthasya kvacid abhAvAn nirvizeSaNA eva sA / katham abhAvo 
09421 jJeyAbhidheyaprameyatvaiH so api siddha eva / tat kim idAnIM
09422 jJeyam asti iti siddhir astu / tathA api kiM siddhaM syAt / anyatra  
09423 tu tad eva agnisAmAnyaM tatra asiddham iti sAdhyate / nanu tatra api  
09424 tadayogavirahiNA sAmAnyena anvayo na siddha eva / na vai kazcit
09425 tathA abhUtena anvayaM karoti / pratipAdayatA hi paraM dhUmo
09426 agninAntarIyako darzanIyo yatra dhUmas tatra agnir iti / sa tathA agni
09427 mAtreNa vyAptaH siddho yatra eva svayaM dRzyate tatra eva agni
09501 buddhiM janayati / tatra ca sAdhyanirdezena na kiMcit / tatra  
09502 darzanasambandhAkhyAnamAtrAd iSTasiddheH / tadanirdeze ca  
09503 kathaM tadviziSTena anvayaH / tad ayam agninA avinAbhAvI 
09504 siddhaH / arthAd eva agnes tatpradezAyogaM vyavacchinatti iti sa
09505 tathA sAdhya ucyate / na punas tathA asya upanyAsapUrvako anvayaH
09506 sAdhyaukter iha anaGgatvAt / tatpUrvakatve vA kaH pratijJAM  
09507 sAdhanAd apAkaroti / tathA ca Aha / 
09508 liGgasya avyabhicAras tu dharmeNa anyatra darzyate /
09509 tatra prasiddhaM tadyuktaM dharmiNaM gamayiSyati /  
09510 tasmAn na agnyAdisAdhanavat sattAsAdhanam apy anavadyam iti / 
09511 	aparAmRSTatadbhede vastumAtre tu sAdhane /
09512 	tanmAtravyApinaH sAdhyasya anvayo na vihanyate //190// 
09513 sAdhane punaH sattve svabhAvavizeSAparigraheNa vastumAtra
09514 vyApini sAdhyadharme na anvayavyAghAtaH / na hi tatra avazyaM 
09515 vizeSaparigrahaH kAryaH / sanmAtrAzraye api sAdhanasAmarthyAt /
09516 na sAdhyatve / vaiphalyAt / api ca /
09517 	na asiddhe bhAvadharmo asti vyabhicAryubhayAzrayaH / 
09518 	dharmo viruddho abhAvasya sA sattA sAdhyate katham //191// 
09519 sattAyAM hi sAdhyAyAM / sarvas tadhetur na trayIM doSajAtim
09520 atipatati / asiddhiM vyabhicAraM virodhaM ca / tatra yadi bhAva  
09601 dharmo hetur ucyate / sa katham asiddhasattAke syAt / yo hi
09602 bhAvadharmaM tatra icchati / sa kathaM bhAvaM na icchet / svabhAva 
09603 eva hi kayAcid apekSayA dharma iti vyatirekI iva dharmiNo  
09604 nirdizyate / na hi dharmadharmivAcinoH zabdayor vAcye kazcid 
09605 vizeSo asti ity uktam etat / atha punar ubhayadharmaM brUyAt /
09606 anAzritavastuno aparyudAsena vyatirekamAtrasya abhAve apy avirodhAt /
09607 yathA na bhavati mUrta iti amUrtatvaM nirupAkhye api
09608 syAt / nirupAkhyasya abhAvAd na pratiSedhaviSayatA iti cet / tat kim
09609 idAnIM vidhiviSayo astu / tad api na iti cet / katham idAnIM na
09610 pratiSedhaviSayaH / vidhinivRttirUpatvAt pratiSedhasya /  tad 
09611 etad vyavacchedamAtraM dvayor api sambhavad vipakSapracAra  
09612 zaGkAvyavacchedena labhyaM gamakatvaM katham AtmasAt
09613 kuryAt / sa ca svayaM svavAcaubhayadharmatAM bruvANaH sato
09614 anyatra apy asya vRttiM bhASate sattAyAM ca avyabhicAram iti 
09615 kathaM na unmattaH / abhAvadharmaM tu bhAvamAtravyApino
09616 arthasya vyavacchedaM hetuM sattAyAM vadato asya viruddho
09617 hetuH syAt / tasya bhAve kvacid asambhavAt abhAve ca bhAva 
09618 vyavacchedasya bhAvAt / tad ayaM triprakAro api dharmaH sattA  
09619 sAdhane na hetulakSaNabhAk na ca anyA gatir asti / tasmAn na  
09620 sattA sAdhyate / sAdhanatve punar asyAH sAmAnyena tanmAtravyApini
09621 vastudharme siddhasattAke dharmiNi na asiddhiH / tena ca
09622 sAdhyadharmeNa vyAptir yadi kathaMcin nizcIyate na virodha  
09623 vyabhicArAv iti na ayaM prasaGgaH / anizcitAyAM tu vyAptau dharmi
09624 samAzraye vA tatsvabhAvatayA gamako na kazcid gamakaH /  
09625 ata eva svadharmeNa vyAptaH
09626 	siddhaH svabhAvo gamako /
09701 vAcyaH / na hi prakAzatayA prakAzayan pradIpas tadrUpApratipattau
09702 svAm arthakriyAM karoti /  
09703 	vyApakas tasya nizcitaH / 
09704 	gamyaH svabhAvas / 
09705 taddharmanizcayAd eva nizcito vyApakatvena tasya dharmiNo 
09706 dharmo gamyaH /  
09707 	tasya ayaM nivRttau vA nivartakaH //192//
09708 tasya vyApyasya dharmasya ayaM nivartako vyApako dharmaH 
09709 svayaM nirvartamAnaH / evaM hy ayam asya vyApakaH siddho 
09710 bhavati yady asya abhAve na bhavet / tad anena dvividhasya api  
09711 sAdhanaprayogasya gamakatAlakSaNam uktaM veditavyam / dvividho  
09712 hi prayogaH sAdharmyeNa vaidharmyeNa ca / yathA Ahur 
09713 eke anvayI vyatirekI ca iti / na anayor vastutaH kazcid bhedo anyatra
09714 prayogabhedAt / sAdharmyeNa api hi prayoge arthAd vaidharmya
09715 gatiH asati tasmin sAdhyena hetor anvayAbhAvAt / tathA  
09716 vaidharmye apy anvayagatiH asati tasmin sAdhyAbhAve hetvabhAvasya  
09717 sandehAd iti vistareNa vakSyAmaH /  
09718 	anityatve yathA kAryam akAryaM vA avinAzini / 
09719 anena udAharaNam anayor darzayati / tatra anvayI yat kiMcit kRtakaM  
09720 tat sarvam anityam yathA ghaTAdayaH zabdaz ca kRtaka
09721 iti kRtakatvasya apy anityatvena vyAptiM pradarzya zabdas ya 
09722 kRtakatve kathite sAmarthyAd eva anityaH zabda iti / tasmAn
09723 na avazyam iha pakSanirdeza iti / ayam anvayinaH prayogaH /  
09724 vyatireke api na anityatvAbhAve kRtakatvaM bhavati zabdaz ca  
09725 kRtaka iti / siddhatatsvabhAvatayA tadabhAve na bhavataH / 
09726 kRtakatvasya zabde ca bhAvakhyAtau tadAtmanaH sato bhAva iti
09727 sAmarthyAt siddheH pUrvavAn na pratijJAvacanam / anvayas tv
09801 arthApattyA siddhaH / na hy atadAtmaniyatasya tannivRttau
09802 nivRttiH / tasmAt tanniyamaM prasAdhya nivRttir vaktavyA /
09803 sA cet sidhyati tadAtmaniyamam arthAt sUcayati iti siddho anvayaH /  
09804 katham idAnIM kRtako avazyam anitya iti pratyetavyo
09805 yena evam ucyate / yasmAt / 
09806 	ahetutvAd vinAzasya svabhAvAd anubandhitA //193//
09807 na hi bhAvA vinazyantas tadbhAve hetum apekSante / svahetor 
09808 eva vinazvarANAM bhAvAt / tasmAd yaH kazcit kRtakaH sa prakRtyA  
09809 eva nazvaraH / tathA hi /
09810 	sApekSANAM hi bhAvAnAM na avazyaMbhAvitA IkSyate /
09811 nirapekSo bhAvo vinAze / sApekSatve hi ghaTAdInAM keSAMcin 
09812 nityatA api syAt / yena 
09813 	bAhulye api hi tadhetor bhavet kvacid asambhavaH //194//  
09814 yady api bahulaM vinAzakAraNAni santi teSAm api svapratyaya  
09815 adhInasaMnidhitvAn na avazyaM saMnidhAnam iti kazcin na vinazyed
09816 api / na hy avazyaM hetavaH phalavanto vaikalyapratibandha  
09817 sambhavAt /  
09818 	etena vyabhicAritvam uktaM kAryAvyavasthiteH / 
09819 	sarveSAM nAzahetUnAM hetumannAzavAdinAm //195//
09820 ity antarazlokaH / tad ayaM bhAvo anapekSas tadbhAvaM prati
09821 tadbhAvaniyato asambhavatpratibandhA iva kAraNasAmagrI sakalA
09822 kAryautpAdane / nanv anapekSANAm api keSAMcit kvacin na avazyaM  
09823 tadbhAvo bhUmibIjaudakasAmagryAm api kadAcid aGkurAnutpatteH /
09824 na tatra api santAnapariNAmApekSatvAt / na evaM bhAvasya kAcid  
09825 apekSA / tatra apy antyA kAraNasAmagrI yA avyavahitA kAryautpatteH
09826 sA phalavaty eva / sa eva ca tatra aGkurahetuH / anyas tu pUrvaH pariNAmas
09901 tadartha eva / na ca tAM tatra kazcit pratibandhuM samarthaH / 
09902 ekatra bhAve vikArAnutpatteH / utpattau vA ekatvahAneH /  
09903 tadAtmanaz ca apracyutasya tadutpAdanaM prati vaiguNyaM kAraNasya  
09904 akurvANasya pratibandhahetor apratibandhakatvAt / nanu  
09905 yavabIjAdayo api zAlyaGkure janye na sApekSAH / tadutpattipratyayAnAM
09906 kadAcit tatra api saMnidhAnAt / kathaM na sApekSAH /  
09907 yAvatA sa eva eSAM svabhAvo na asti yas tadutpAdanaH zAlibIjasya 
09908 iti tatsvabhAvApekSAH / evaM tarhi kRtakAnAm api keSAMcit
09909 satAM vA sa eva svabhAvo na asti yo vinazvaraH / tasmAt tatsvabhAvApekSatvAn
09910 na vinazvarAH / zAlibIjAdInAm api sa svabhAvaH
09911 svahetor iti yo na tadhetuH so atatsvabhAvaH syAt / niyatazaktiz 
09912 ca sa hetuH svarUpeNa pratIta eva / na ca svabhAvaniyamo  
09913 arthAnAm Akasmiko yuktaH / anapekSasya dezakAladravyaniyama 
09914 ayogAt tathA atra api niyamahetur vaktavyo yata ime kecin  
09915 nazvarAtmano jAtA na ca atra kazcin niyAmakaH svabhAvasya asti  
09916 sarvajanminAM vinAzasiddheH / janmisvabhAvo nAzI iti cet / na vai 
09917 janma nAzisvabhAvasya hetuH na ca ahetoH svabhAvaniyamaH /  
09918 tasmAn na atra kazcid hetoH svabhAvapravibhAgaH / tadabhAvAt 
09919 phalasya api na asti ity asamAnam / sA iyaM nirapekSatA vinAzasya  
09920 kvacit kadAcic ca bhAvavirodhinI tadabhAvaM svabhAvena 
09921 sAdhayati / yo hi svabhAvo nirapekSaH sa yadi kadAcid
09922 bhavet kvacid vA tatkAladravyApekSa iti nirapekSa eva na
09923 syAd ity uktam / sa tarhi nazvaraH svabhAvo nirapekSa iti ahetukaH
09924 syAt / na ahetukaH sattAhetor eva bhAvAt tathAutpatteH / sato  
09925 hi bhavatas tAdRzasya eva bhAvAt / na avazyaM sataH kutazcid  
09926 bhAva iti cet / AkasmikI tarhi sattA iti / na iyaM kasyacit kadAcit  
09927 kvacit viramet / tad hi kiMcid upalIyeta na vA yasya yatra kiMcit 
09928 pratibaddham apratibaddhaM vA / sA iyaM sattA apratibandhinI cet / 
10001 niyamavatI na syAt / tasmAn na iyam AkasmikI kvacit / kathaM  
10002 tarhi idAnIm ahetuko vinAza uktaH / jAtasya tadbhAve anyAnapekSaNAt /
10003 uktaM ca atra na vinAzo nAma anya eva kazcid bhAvAt svabhAva
10004 eva hi nAzaH sa eva hy ekakSaNasthAyI jAta iti / tam asya  
10005 mandAH svabhAvam UrdhvaM vyavasyanti na prAk darzane api  
10006 pATavAbhAvAd iti tadvazena pazcAd vyavasthApyate vikAradarzanena
10007 iva viSam ajJaiH / tad ayaM sattAvyatirekeNa na anyat kiMcid
10008 vinAzo apekSata iti tadvyApI / kathaM punar etad gamyate  
10009 nirapekSo vinAza iti /
10010 	asAmarthyAc ca tadhetor / 
10011 abhAvakAriNaH kriyApratiSedhAc ca iti cazabdAt / katham asAmarthyam /
10012 siddhe hi bhAve kArako na taM karoti / na apy anyakriyAyAM  
10013 tasya kiMcid iti / tadatadrUpAkaraNAc ca akiMcitkaro na apekSyata
10014 iti / kathaM kriyApratiSedho vinAza iti hi te bhAvAbhAvaM manyante /
10015 tad ayaM vinAzahetur abhAvaM karoti iti prAptam / tatra
10016 yady abhAvo nAma kazcit kAryaH syAt svabhAvaH sa eva bhAva iti
10017 na abhAvaH syAt / tasmAd abhAvaM karoti iti bhAvaM na karoti iti  
10018 kriyApratiSedho asya kRtaH syAt / tathA apy ayam akiMcitkaraH
10019 kim iti apekSyata iti siddhA vinAzaM praty anapekSayA bhAvasya /
10020 tasmAt 
10021 	bhavaty eSa svabhAvataH / 
10022 	yatra nAma bhavaty asmAd anyatra api svabhAvataH //196// 
10023 so ayaM kvacid bhavan dRSTo anapekSatvAt svabhAvata eva bhavati / 
10024 tathA anyatra api svabhAvabhAvI vizeSAbhAvAd iti /
10025 	yA kAcid bhAvaviSayA dvidhA eva anumitis tataH /
10026 	svasAdhye kAryabhAvAbhyAM sambandhaniyamAt tayoH //197// 
10028 tasmAt dviprakArA eva vastuviSayAnumitiH kAryaliGgA svabhAva
10101 liGgA ca / yathAsvaM vyApini sAdhye tayor eva pratibandhAt liGgayor  
10102 liGgini / yathoktaM prAk /  
10103 	pravRtter buddhipUrvatvAt tadbhAvAnupalambhane /  
10104 	pravartitavyaM na ity uktAnupalabdheH pramANatA //198// 
10105 tRtIyas tu hetur anupalabdhir avizeSeNa kvacid arthe gamaka ity 
10106 ucyate / sannizcayazabdavyavahArapratiSedhe hi sarvA eva anupalabdhir
10107 liGgam / sannizcayAd hi zabdavyavahArAH pravartante / 
10108 te pravRttir ity uktAH / tathA hy anupalabdhir eva asattvam 
10109 ity uktam prAk / tac ca pratipattRvazAn na vastuvazAt / tAvad hi 
10110 sa bhAvo asya na asti yAvad atra apratipattiH / satA api te na tadarthAkaraNAt /
10111 vastutas tv anupalabhyamAno na san na asan / satAm 
10112 api svabhAvAdiviprakarSAt kadAcid anupalambhAt tasya asatsv api
10113 tulyatvAt / tad etat pratipattuH pramANAbhAvAn nivRttaM sattvam  
10114 anupalabdhilakSaNaM svanimittAn zabdavyavahArAn nivartayati / 
10115 tena yA api iyam anupalabdhir upalabdhilakSaNaprAptAnAM 
10116 vastuto apy asattvarUpA apravRttiyogyatvAt / tasyA apy etat tulyaM  
10117 prAmANyam atra viSaye / asannizcayaphalA api sA / asattAyAm api iyaM
10118 pramANam eva / na hy asti sambhavo yad upalabdhiyogyaH
10119 sakaleSv anyeSu kAraNeSu san na upalabhyeta / na punaH pUrvA  
10120 asattAsAdhAnI / yasmAt 
10121 	zAstrAdhikArAsambaddhA bahavo arthA atIndriyAH /  
10122 	aliGgAz ca kathaM teSAm abhAvo anupalabdhitaH //199// 
10123 so ayam asattAM sAdhayann anupalabdhimAtreNa sarvArthAnAM
10201 pramANatrayanivRttyA sAdhayet / tatra na zAstranivRttir abhAvasAdhanI /
10202 tasya kvacid anadhikAre apravRtteH / zAstraM hi pravartamAnaM 
10203 kaMcit puruSArthasAdhanam upAyam Azritya pravartate /
10204 anyathA abaddhapralApasya aprAmANyAt / tatra ca prakaraNe  
10205 bahavo arthA na avazyaM nirdezyAH / yathA pratyAtmaniyatAH 
10206 kAzcana puruSANAM cetovRttayo aniyatanimittabhAvinyaH deza  
10207 kAlavyavahitA vA prakaraNAnupayogino dravyavizeSA na tAn
10208 zAstraM viSayIkaroti / na ca tathA viprakRSTeSu  svasAmarthyaupadhAnAv
10209 jJAnautpAdanazaktir asti / na ca avazyam eSAM kAryaupalambho 
10210 yena anumIyeran / na ca te pramANatrayanivRttAv api
10211 na santi iti zakyante vyavasAtum / tasmAn na sarvAnupalabdhiH
10212 sAdhanI nivRttinizcayasya / tad iyam /
10213 	sadasannizcayaphalA na iti syAd vA apramANatA /  
10214 na eva vA iyam anupalabdhiH pramANaM vyavasAyaphalatvAt pramANAnAm /
10215 na hi pravRttiniSedhe api iyaM niHzaGkaparicchedaM 
10216 cetaH karoti / saMzayAd api kvacil lokasya pravRtteH / tathAtve
10217 tan niravadyaM yadi nizcayapUrvaM vyavahared iti sA iyam apravRtti 
10218 phalA proktA / 
10219 	pramANam api kAcit syAl liGgAtizayabhAvinI //200//  
10220 atra na sarvAnupalabdhir apramANam / pramANam api kAcil liGga 
10221 vizeSabhAvinI yathAudAhRtA prAk / yat punar uktam apramANam anupalabdhir iti
10223 	svabhAvajJApakAjJAnasya ayaM nyAya udAhRtaH / 
10301 yasya kasyacit svabhAvo na upalabhyate dezAdiviprakarSAt na
10302 sa tadanupalambhamAtreNa asan nAma yathoktaM prAk / yo api  
10303 jJApakasya liGgasya abhAvAd atIndriyaH pratikSipyate arthaH svabhAva
10304 vizeSo vA yathA na asti viraktaM ceto devatAvizeSo vA 
10305 na asti dAnahiMsAviraticetanAnAm abhyudayahetutA iti pratyakSe
10306 apy arthe phalasya AnantaryAbhAvAd atatphalasAdharmyAd viparyasto  
10307 apavadeta api na tAvatA tadabhAva eva vyavahitAnAm
10308 api hetoH phalAnAm utpattidarzanAt mUSikAlarkaviSavikAravat / 
10309 tadbhAve virodhAbhAvAd atra anupalabdhimAtram apramANam /  
10310 bhAve kiM pramANam iti cet / ata eva saMzayo astu bhaved 
10311 vA pramANam ity apratikSepaH / tad atra keSAMcid arthAnAM
10312 svabhAvAnAM vA darzanapATavAbhAvAt kAraNAnAM kAryautpAdana
10313 niyamAbhAvAc ca bhavej jJApakAsiddhiH / neyatA tadabhAvaH /  
10314 punaH paryAyeNa keSAMcid abhivyakteH / 
10315 	kArye tu kArakAjJAnam abhAvasya eva sAdhakam //201// 
10316 svabhAvAbhAve sAdhye tadanupalambha eva apramANam ucyate /
10317 kArakAnupalambhas tu pramANam eva / na hy asti sambhavo yad  
10318 asati kAraNe kAryaM syAt / nanu kadAcit kAraNanAze api kAryasthitir
10319 dRSTA / na brUmaH kAraNasthitikAlabhAvI kAryam iti / heturahitA
10320 tu bhAvautpattir na asti ity ucyate / na ca tathA sthAyI bhAvas tadupAdAnaH /
10321 pAraMparyeNa tu santAnaupakArAt tatkAryavyapadezaH /  
10322 yady asya kathaMcid abhAvaH sidhyet tatphalaM na asti iti nizcIyate /  
10323 	svabhAvAnupalambhaz ca svabhAve arthasya liGgini /
10324 svabhAvAbhAva eva liGgini svabhAvAnupalambho api kazcit pramANam eva /
10325 eva yady anupalabhyamAno vyApakaH svabhAvo asya
10326 siddhiH syAt / yathA vRkSatvaM ziMzapAyAH /
10327 	tadabhAvaH pratIyeta hetunA yadi kenacit //202//  
10401 yady asya kAraNasya svabhAvasya vyApakasya vA abhAvaH kutazcid 
10402 gamakAd hetoH sidhyet / so ayaM asan eva svakAryaM  
10403 vyApyaM vA nivartayati / tadabhAvAsiddhau nivartye api saMzayAt /  
10404 katham idAnIM bhAvasya svayam anupalabdher abhAva  
10405 siddhiH /  
10406 	dRzyasya darzanAbhAvakAraNAsambhave sati /  
10407 	bhAvasya anupalabdhasya bhAvAbhAvaH pratIyate //203// 
10408 bhAvo hi yadi bhaved yathAsvaM grAhakeNa karaNena upalabhya
10409 eva bhavet / sa darzanapratibandhiSu vyavadhAnAdiSv asatsu upalabhya 
10410 eva / anupalabdhas tv asann iti nizcIyate / tAdRzaH sata 
10411 upalambhAvyabhicArAt / ayam eva hetur hetuvyApakayor abhAve  
10412 api veditavyaH / 
10413 	viruddhasya ca bhAvasya bhAve tadbhAvabAdhanAt /  
10414 	tadviruddhaupalabdhau syAd asattAyA vinizcayaH //204// 
10415 yo hi bhAvo yena saha na avatiSThate tadupAdAnayor anyonyavaiguNya  
10416 Azrayatvena ArambhavirodhAt tayor viruddhayor ekasya bhAve  
10417 apy anyAbhAvagatir bhavati yathoktaM prAk / idam anupalabdher  
10418 na pRthag vyavasthApyate / tata eva virodhagater virodhAc ca abhAva  
10419 sAdhanAt / bhavatu nAma evaMvidhAyA anupalabdher abhAva
10420 gatiH / sA punaH katham anumAnam / kathaM ca na syAt /  
10421 dRSTAntAnapekSaNAt / na hy asyAM kazcid dRSTAnto asti / kiM
10422 na nirupAkhyaM vyomakusumAdi dRSTAntaH / tad asat katham
10423 avagantavyaM yena evaM syAt / anupalabdher eva iti cet / tatra  
10424 katham adRSTAntikA asattAsiddhiH sadRSTAntatve vA anavasthAprasaGgaH /
10425 tathA ca apratipattiH / tasmAn nirupAkhyAbhAvasiddhivad  
10426 anyatra api dRSTAntAnapekSaNAd ananumAnam / zRNvann api devAnAMpriyo
10427 na avadhAraNapaTuH / nimittaM hy asaczabdavyavahArANAm 
10428 upalabhyAnupalabdhiH / sA svasaMnidhAnAt svanimittAn  
10429 etAn sAdhayati iti svanimittasAmagrIyogyasaMnidhAnaH sarvo atra  
10501 dRSTAntaH / asattA punar atra anupalabdhir eva / ata eva iyaM kAraNAt 
10502 kAryAnumAnalakSaNatvAt / svabhAvahetAv antarbhavati iti vakSyAmaH /
10503 saczabdavyavahArapratiSedhe api pramANanivRttyA  
10504 nimittavaikalyAbhAvino aGkurAdayo dRSTAntaH na kevalaM nirupAkhyam /
10505 nirupAkhye api iyam eva pravRttir niSidhyate / anupalabdhi  
10506 lakSaNA asattA siddhA eva / so ayaM mUDho nimittaM tadabhAvaM 
10507 vA abhyupagamya pravRttinivRttI vilomayan yathA abhyupagamaM 
10508 pratipAdyate nirupAkhyavad anyavad vA iti / sa eva tAvad
10509 upalabdhyabhAvaH kathaM siddha iti cet / etad uttaratra vakSyAmaH / 
10510 anyatra apy anumAne sAdhyadharmeNa vyAptaM sAdhanam
10511 icchan kim iti dRSTAntena pratyAyyo vyApyanirdezAd eva
10512 vyApnuvataH siddheH / nizcitArthasya api smRtyartho dRSTAnta iti  
10513 cet / tad itaratra api samAnaM / so ayam anyatra anupalambhamAtrAd 
10514 asadvyavahAraM pratipadyamAno api iha vyAmUDha iti 
10515 smAryate / atha yad idaM na santi pradhAnAdayo anupalabdher 
10516 iti / tatra katham asadvyavahAravidhiH sadvyavahAraniSedho 
10517 vA / kathaM ca na syAt / tadarthapratiSedhe dharmivAcino aprayogAd 
10518 abhidhAnasya nirviSayasya ca pratiSedhasya ayogAt / na eSa  
10519 doSaH / yasmAt /
10520 	anAdivAsanAudbhUtavikalpapariniSThitaH /
10521 	zabdArthas trividho dharmo bhAvAbhAvaubhayAzrayaH //205 //
10522 	tasmin bhAvAnupAdAne sAdhye asya anupalambhanam /  
10523 	tathA hetur na tasya eva abhAvaH zabdaprayogataH //206// 
10524 niveditam etat yathA na ete zabdAH svalakSaNaviSayA anAdi  
10525 vAsanAprabhavavikalpapratibhAsinam arthaM viSayatvena AtmasAt  
10526 kurvanti vaktuH zrotus ca tadvikalpabhAjo yathApratibhAsi
10527 vastupratipAdanasamIhAprayogAt tadAkAravikalpajananAc ca /
10601 na ca upAdAnakAryapratyayApratibhAsi rUpaM zakyaM tadviSayatvena
10602 adhyavasAtum / sa tu vikalpaH sadasadubhayapratyayAhita  
10603 vAsanAprabhava iti tatpratibhAsyAkArAdhyavasAyavazena
10604 ca bhAvAbhAvaubhayadharma ity ucyate / tad atra dharmiNi vyavasthitAH
10605 sadasattvaM cintayanti kim ayaM pradhAnazabdapratibhAsy
10606 artho bhAvaupAdAno na vA iti / tasya bhAvAnupAdAnatve  
10607 sAdhye sa eva pratyAtmavedyatvAd apratikSepArho artho dharmI /  
10608 na ca sa eva arthaH svalakSaNam iti zakyaM vaktum / asamprAptaniruddhayor
10609 apy arthayos tasya anapAyAt / vastuviparItAkAraniveziSv 
10610 api tIrthAnantarIyapratyayeSu bhAvAt /  
10611 	paramArthaekatAnatve zabdAnAm anibandhanA / 
10612 	na syAt pravRttir artheSu darzanAntarabhediSu //207//
10613 	atItAjAtayor vA api na ca syAd anRtArthatA / 
10614 	vAcaH kasyAzcid ity eSA bauddhArthaviSayA matA //208// 
10615 iti saMgrahazlokau / tasya ca yathA samIhitarUpAnupAdAnatve 
10616 sAdhye tathA anupalambho asya dharmo asti iti na sAdhanadharmAsiddhiH /
10617 na punar atra ayam eva zabdavikalpapratibhAsy artho apahnUyate /
10618 tasya buddhAv upasthApanAya zabdaprayogAt / tadabhAve  
10619 tadayogAt / api ca / 
10620 	zabdArthApahnave sAdhye dharmAdhAranirAkRteH /
10621 	na sAdhyaH samudAyaH syAt siddho dharmaz ca kevalaH //209// 
10623 yadi hi zabdArtha eva apohyeta pradhAnazabdavAcyasya dharmiNa 
10624 eva nirAkaraNAn nirAdhAraH sAdhyadharmaH syAt / tad ayam  
10625 AdhAravyavacchedAnapekSo na vivAdAzraya iti na upanyasanIya  
10626 eva syAt / kiM ca /
10627 	sadasatpakSabhedena zabdArthAnapavAdibhiH /
10628 	vastv eva cintyate hy atra pratibaddhaH phalaudayaH //210//
10629 	arthakriyAsamarthasya vicAraiH kiM tadarthinAm /  
10630 	SaNDhasya rUpavairUpye kAminyAH kiM parIkSayA //211//
10701 na hi zabdArtho asan san vA kaMcit puruSArtham uparuNaddhi  
10702 samAdadhAti vA / yathA abhinivezam atattvAt / yathA atattvaM ca asamIhitatvAt /
10703 tad ayaM pravartamAnaH sarvadA sadasaccintAyAm 
10704 avadhIritavikalpapratibhAso vastv eva adhiSThAnIkaroti yatra ayaM
10705 puruSArthaH pratibaddho yathA agnau zItapratIkArAdiH / na hy 
10706 atra zabdArthaH samarthas tadanubhavAptAv api tadabhAvAt /
10707 tad ayam arthakriyArthI tadasamarthaM prati dattAnuyogo bhavituM  
10708 na yuktaH / na hi vRSasyantI SaNDhasya rUpavairUpyaparIkSAyAm  
10709 avadhatte / yat punar etad uktaM kalpitasya anupalabdhir
10710 dharma iti tasya ko arthaH / 
10711 	zabdArthaH kalpanAjJAnaviSayatvena kalpitaH /
10712 	dharmo vastvAzrayAsiddhir asya ukto nyAyavAdinA //212// 
10713 kalpanAviSayatvAc zabdArtha eva kalpitaH / tasya vastv AzrayAnupalambho
10714 dharma ity abhiprAyaH / yad uktam na pramANa
10715 trayanivRttAv api bhAvAbhAvasiddhir iti / tan mA bhUd anyapramANa
10716 anivRttau nivRttiH / tayor asakalaviSayatvAd AgamaH punar  
10717 na kiMcin na vyApnoti / tannivRttiH kathaM na gamikA iti / uktam 
10718 atra na AgameSu sarvArthA upanibadhyante aprakaraNApannatvAd
10719 iti / api ca /
10720 	nAntarIyakatA abhAvAc zabdAnAM vastubhiH saha /
10721 	na arthasiddhis tatas te hi vaktRabhiprAyasUcakAH //213//
10722 na hi zabdA yathAbhAvaM vartante yatas tebhyo arthaprakRtir  
10723 nizcIyeta / te hi vaktur vivakSAvRttaya iti tannAntarIyakAs  
10724 tAm eva gamayeyuH / na ca puruSaicchAH sarvA yathArthabhAvinyaH /  
10725 na ca tadapratibaddhasvabhAvo bhAvo anyaM gamayati /
10801 yat tarhi idam AptavAdAvisaMvAdasAmAnyAd anumAnatA ity Agamasya
10802 prAmANyam anumAnatvam uktam tat kathaM / na ayaM puruSo anAzritya Agama
10803 AsituM samarthaH / atyakSaphalAnAM keSAMcit  
10804 pravRttinivRttyor mahAnuzaMsApApazravaNAt tadbhAve virodha 
10805 adarzanAc ca / tat sati pravartitavye varam evaM pravRtta 
10806 iti parIkSayA prAmANyam Aha / tac ca / 
10807 	sambaddhAnuguNaupAyaM puruSArthAbhidhAyakam /
10808 	parIkSAdhikRtaM vAkyam ato anadhikRtaM param //214//
10809 sambaddho vAkyAnAm ekArthaupasaMhAraupakAraH / na dazadADimAdi
10810 vAkyAnAm iva anupasaMhAra eva / anyathA vaktur vaiguNyam 
10811 udbhAvayet / azakyaupAyaphalAni ca zAstrANi phalArthI na Adriyeta
10812 vicArayitum apuruSArthaphalAni ca / viSazamanAya takSaka 
10813 phaNaratnAlaMkAraupadezavat kAkadantaparIkSAvac ca / tad viparyayeNa
10814 upasaMhAravat zakyaupAyaM puruSArthAbhidhAyi ca
10815 zAstraM parIkSyeta anyatra avadhAnasya eva ayuktatvAt / tad yadi na
10816 parIkSAyAM visaMvAdabhAk pravartamAnaH zobheta / kaH punar
10817 asya avisaMvAdaH / 
10818 	pratyakSeNa anumAnena dvividhena apy abAdhakam / 
10819 	dRSTAdRSTArthayor asya avisaMvAdas tadarthayoH //215//
10820 pratyakSeNa abAdhanaM pratyakSAbhimatAnAm arthAnAM tathAbhAvaH
10821 yathA nIlAdisukhaduHkhanimittaupalakSaNarAgAdibuddhInAm /
10822 atathAbhimatAnAM ca apratyakSatA yathA zabdAdirUpa  
10823 saMnivezinAM sukhAdInAM dravyakarmasAmAnyasaMyogAdInAM 
10824 ca / tathA anAgamApekSAnumAnaviSayAbhimatAnAM tathAbhAvaH
10825 yathA catUrNAm AryasatyAnAm / ananumeyAnAM tathAbhAvo  
10901 yathA AtmAdInAm / AgamApekSAnumAne api yathA rAgAdi 
10902 rUpaM tatprabhavaM ca adharmam abhyupagamya tatprahANAya 
10903 snAnAgnihotrAder anupadezaH / sa iyaM zakyaparicchedAzeSaviSaya
10904 vizuddhir avisaMvAdaH /  
10905 	AptavAdAvisaMvAdasAmAnyAd anumAnatA /  
10906 	buddher agatyA abhihitA parokSe apy asya gocare //216//
10907 tasya ca asya evaMbhUtasya AptabhUtasya avisaMvAdasAmAnyAd adRSTa
10908 vyabhicArasya pratyakSAnumAnAgamye apy arthe pratipattes  
10909 tadAzrayatvAt tadanyapratipattivad avisaMvAdo anumIyate /
10910 tataH zabdaprabhavA api satI na zAbdavad abhiprAyaM nivedayaty
10911 eva ity arthAvisaMvAdAd anumAnam api / athavA anyathA AptavAdasya
10912 avisaMvAdAd anumAnatvam ucyate /
10913 	heyaupAdeyatattvasya sa upAyasya prasiddhitaH /
10914 	pradhAnArthAvisaMvAdAd anumAnaM paratra vA //217// 
10915 heyaupAdeyatadupAyAnAM tadupadiSTAnAm avaiparItyam avisaMvAdaH /
10916 yathA catUrNAm AryasatyAnAM vakSyamANanItyA / tasya asya
10917 puruSArthaupayogino abhiyogArthasya avisaMvAdAd viSayAntare
10918 api tathAtvaupagamo na vipralambhAya anuparodhAt niSprayojana 
10919 vitathAbhidhAnavaiphalyAc ca vaktuH / tad etad agatyA ubhayathA 
10920 apy anumAnatvam Agamasya upavarNitam / varam AgamAt 
10921 pravRttAv evaM pravRttir iti / na khalv evam anumAnam anapAyam
10922 anAntarIyakatvAd artheSu zabdAnAm iti niveditam etat /
10923 	puruSAtizayApekSaM yathArtham apare viduH /
10924 yathArthadarzanAdiguNayuktaH puruSa Aptas tatpraNayanam avisaMvAda
11001 ity anye / 
11002 	iSTo ayam arthaH zakyeta jJAtuM so atizayo yadi //218// 
11003 sarve eva Agamam anAgamaM vA pravRttikAmo anveSate prekSA 
11004 pUrvakArI na vyasanena / api nAma anuSTheyam ato jJAtvA pravRtto
11005 arthavAn syAm iti / sa zakyadarzanAvisaMvAdapratyayena anyatra api 
11006 pravartate / evaMprAyatvAl lokavyavahArasya / puruSaparIkSayA
11007 tu pravRttAv apravRttir eva / tasya tathAbhUtasya jJAtum  
11008 azakyatvAt na aniSTeH / tAdRzAm avitathAbhidhAnAt / tathA hi / 
11009 	ayam evaM na vA ity anyadoSAnirdoSatA api vA /  
11010 	durlabhatvAt pramANAnAM durbodhA ity apare viduH //219//  
11011 caitasyebhyo hi guNadoSebhyaH puruSAH samyaGmithyApravRttaH 
11012 te ca atIndriyAH svaprabhavakAyavAgvyavahArAnumeyAH  
11013 syuH / vyavahArAz ca prAyazo buddhipUrvam anyathA api kartuM  
11014 zakyante puruSaicchAvRttitvAt teSAM ca citrAbhisandhitvAt / 
11015 tad ayaM liGgasaMkarAt katham anizcinvan pratipadyeta / atha
11016 kiM na eva sa tAdRzaH puruSo asti yo nirdoSaH /  
11017 	sarveSAM savipakSatvAn nirhrAsAtizayazritAm /  
11018 	sa AtmIbhAvAt tadabhyAsAd dhIyeran AzravAH kvacit //220 //
11019 sa tu prahINAzravo durjJAnaH / doSA hi nirhrAsAtizayadharmaNo
11020 vipakSAbhibhavautkarSApakarSaM sAdhayanti jvAlAAdivat / te hi  
11021 vikalpaprabhavAH saty apy upAdAne kasyacin manoguNasya abhyAsAd 
11022 apakarSiNaH / tatpATave niranvayavinAzadharmANaH  
11023 syuH / jvAlAAdivad eva / tena syAd api nirdoSaH / kathaM nirdoSo 
11024 nAma / yAvatA doSavipakSasAtmatve api doSasAtmano vipakSautpattivad
11025 yathApratyayaM doSautpattir api / na ayaM doSaH / yasmAt /
11027 	nirupadravabhUtArthasvabhAvasya viparyayaiH//
11028 	na bAdhA yatnavattve api buddhes tatpakSapAtataH //221 /  
11101 na hi svabhAvo ayatnena vinivartayituM zakyaH / zrotriya 
11102 kApAlikaghRNAvat / yatnaz ca prApyanivartyayoH svabhAvayor  
11103 guNadoSadarzanena kriyeta / tac ca vipakSasAtmanaH puruSasya 
11104 doSeSu na sambhavati / tasya nirupadravatvAt / azeSadoSahAneH
11105 paryavasthAnajanmapratibaddhaduHkhavivekAt prazamasukha  
11106 rasasya anudvejanAc ca / abhUtArthaM khalv apy upAdAna balabhAvi
11107 santAnasya viparyayaupAdAnAn na syAt / na tu bhUtArtham vastu
11108 balautpatteH / abhUtArthAz ca doSA na pratipakSasAtmyabAdhinaH /
11109 tasmAn na punar doSautpattiH / yatne api buddher guNapakSa
11110 pAtena pratipakSa eva yatnAdhAnAt parIkSAvato vizeSeNa aduSTa 
11111 AtmanaH / kaH punar eSAM doSANAM prabhavo yatpratipakSa 
11112 abhyAsAt prahIyante /  
11113 	sarvAsAM doSajAtInAM jAtiH satkAyadarzanAt// 
11114 	sA avidyA tatra tatsnehas tasmAd dveSAdisambhavaH //222 // 
11115 na hi na ahaM na mama iti pazyataM parigraham antareNa kvacit 
11116 snehaH / na ca ananurAgiNaH kvacid dveSaH / AtmAtmIyAnuparodhiny  
11117 uparodhapratighAtini ca tadabhAvAt / tasmAt samAnajAtIyAbhyAsajam
11118 Atmadarzanam AtmIyagrahaM prasUte / tau ca tatsnehaM
11119 sa ca dveSAdIni iti satkAyadarzanajAH sarvadoSAH / tad eva ca ajJAnam
11120 ity ucyate /  
11121 	moho nidAnaM doSANAm ata eva abhidhIyate//  
11122 	satkAyadRSTir anyatra tatprahANe prahANataH //223 //  
11123 mohaM doSanidAnam AhuH amUDhasya doSAnutpatteH punar 
11124 anyatra satkAyadRSTim / tac ca etat pradhAnanirdeze sati syAd
11125 anekajanmanAM doSANAm ekautpattivirodhAt / na ca dvayoH prAdhAnye
11201 ekaikanirdezaH parabhAgabhAk / ubhayathA apy ekasya  
11202 nirdeze na virodhaH / prAdhAnyaM punas tadupAdAnatvena / 
11203 tatprahANe doSANAM prahANAt / tasmAt sambhavati satkAyadarzana
11204 janmanAM doSANAM tatpratipakSanairAtmyadarzanAbhyAsAt 
11205 prahANam / sa tu kSINadoSo duranvayo yadupadezAd ayaM pratipadyeta / 
11206 mA bhUt puruSAzrayaM vacanam AgamaH praNetur
11207 duranvayatvAt / 
11208 	girAM mithyAtvahetUnAM doSANAM puruSAzrayAt// 
11209 	apauruSeyaM satyArtham iti kecit pracakSate //224 // 
11210 na khalu sarva eva AgamaH sambhAvyavipralambhaH / vipralambha
11211 hetUnAM doSANAM puruSAzrayAd apauruSeyaM satyArtham ity
11212 eke / kAraNAbhAvo hi kAryAbhAvaM sAdhayati iti / ya evaMvAdinas 
11213 tAn eva prati / 
11214 	girAM satyArthahetUnAM guNAnAM puruSAzrayAt//
11215 	apauruSeyaM mithyArthaM kiM na ity anye pracakSate //225 //  
11216 yathA rAgAdiparItaH puruSo mRSAvAdI dRSTas tathA dayAdharmatA
11217 AdiyuktaH satyavAk / tad yathA vacanasya puruSAzrayAn mithyArthatA
11218 tathA satyArthatA api iti / sa nivartamAnas tAm api nivartayati  
11219 ity AnarthakyaM syAd viparyayo vA / na hi zabdAH prakRtyA arthavantaH /
11220 samayAt tato arthakhyAteH kAyasaMjJAAdivat / aprAtikUlyaM
11221 tu yogyatA samaye tadicchApraNayanAt / nisargasiddheSv  
11222 icchAvazAt pratipAdanAyogAt / te anarthakAH puruSasaMskArAd 
11223 arthavantaH syuH / tatsaMskAryatA eva ca eSAM pauruSeyatA yuktA  
11224 na utpattiH / tata eva arthavipralambhAt / utpanno apy anyathA samito  
11225 na uparodhI tadanyapuruSadharmavat / tad ayaM nivartamAnaH 
11226 svakRtasamayasambhavAm arthapratibhAM nivartayati /  
11227 tat kutas tannivRttyA satyArthatA / atha punar utpattir eva pauruSeyatA /
11301 na samayAkhyAnam / 
11302 	arthajJApanahetur hi saMketaH puruSAzrayaH// 
11303 	girAm apauruSeyatve apy ato mithyAtvasambhavaH //226 //
11304 kiM hy asya apauruSeyatayA yato hi samayAd arthapratipattiH
11305 sa pauruSeyo vitatho api syAt / zIlasAdhanasvargavacanaM tad anyathA
11306 samayena viparyAsayet / tena ayathArtham api prakAzanasambhavAt
11307 sa eva doSaH /
11308 	sambandhApauruSeyatve syAt pratItir asaMvidaH//
11309 syAd etad akAryasambandhA eva zabdAH / na te artheSu puruSair  
11310 anyathA viparyasyante / tena adoSa iti / kim idAnIM saMketena /  
11311 sa hi sambandho yato arthapratItiH / sa ced apauruSeyo na ayaM 
11312 samayam apekSeta / apratItyAzrayo vA kathaM sambandhaH /  
11313 	saMketAt tadabhivyaktAv asamarthAnyakalpanA //227 //  
11314 na vai sambandho vidyamAno apy anabhivyaktaH pratItihetuH /
11315 saMketas tv enam abhivyanakti / sa tarhi siddhaupasthAyI kim  
11316 akAraNaM poSyate / nanv iyAn sambandhasya vyApAro yad artha 
11317 pratItijananam / tat samayena eva kRtam iti / na ayogye samayaH
11318 samartha iti yogyatA tatsambandhaz cet / tat kiM vai zabdaH 
11319 sambandho astu / samarthaM hi rUpaM zabdasya yogyatA kArya
11320 kAraNayogyatAvat / sA ced arthAntaraM kiM zabdasya iti sambandho
11321 vAcyaH / yogyatAupakAra iti cet / na / nityAyAH niratizayatvAt /
11322 tatra apy atiprasaGgAt upakArAsiddheH / yogyatAyAM ca svato yogyatve
11323 artha eva kiM na iSyate / samayas tarhi kathaM zabdArtha
11324 sambandhAH / puruSeSu vRtteH / na amizrANAM siddhAnAM kazcit
11325 sambandho abhedaprasaGgAt anapekSaNAc ca / arthavizeSasamIha 
11326 apreritA vAg ata idam iti viduSaH svanidAnAbhAsinam arthaM  
11327 sUcayati iti buddhirUpavAg vijJaptyor janyajanakabhAvaH sambandhaH  
11401 tataH zabdAt pratipattir avinAbhAvAt / tadAkhyAnaM  
11402 samayaH / tataH pratyAyakasambandhasiddheH sambandhAkhyAnAt /  
11403 na tu sa eva sambandhaH / astu vA anya eva nityaH sambandhaH /  
11404 tena 
11405 	girAm ekArthaniyame na syAd arthAntare gatiH// 
11406 na hi tena sambandhena asambaddhe arthe pratItir yuktA / tasya 
11407 vaiphalyaprasaGgAt / dRSTaz ca icchAvazAt kRtasamayaH sarvaH  
11408 sarvasya dIpakaH / 
11409 	anekArthAbhisambandhe viruddhavyaktisambhavaH //228 // 
11410 atha mA bhUd dRSTavirodha iti sarve sarvasya vAcakAH / tathA na  
11411 sarvaH sarvasAdhano asaMkarAt kAryakAraNatAyAH / tatra pratiniyata
11412 sAdhane abhimate arthe sarvasAdhyasAdhanasAdhAraNasya 
11413 zabdasyaiSTavyaktim eva samayakAraH karoti iti kuta etat so aniyato 
11414 niyamaM puruSAt pratipadyate / tadA /  
11415 	apauruSeyatAyAz ca vyarthA syAt parikalpanA// 
11416 api nAma asaMkIrNam arthaM jAnIyAm iti saMkarahetuH puruSaupAkIrNaH /
11417 tatra yAdRzAH puruSaiH kvacit prayuktAH saMkIryante 
11418 tAdRzA eva sarvasAdhAraNAH santaH kvacit tair viniyamitAs 
11419 tattvAparijJAnAt / prakRtyA eva vaidikA niyatA iti cet / na upadezam
11420 apekSeran na anyathA saMketena prakAzayeyuH vyAkhyAvikalpaz 
11421 ca na syAt / upadezasya ca iSTasaMvAdaH zakyavikalpe na asti iti
11422 vyarthA eva apauruSeyatA / 
11423 	vAcyaz ca hetur bhinnAnAM sambandhasya vyavasthiteH //229 // 
11425 arthA hi bAhyA na rUpaM zabdasya na zabdo arthAnAm / yena abhinna  
11426 AtmatayA vyavasthAbhede api nAntarIyakatA syAt / kRtakatva 
11427 anityatvavat / na apy ete vivakSAjanmAno dhvanayo ajanmAno  
11428 vA vivakSAvyaGgyAH na arthAyattAH / tataH katham idAnIM
11501 tatpratiniyamasaMsAdhyaM tadanvayaM sAdhayeyuH / na hy
11502 apratibaddhas tatsAdhana iti /
11503 	asaMskAryatayA pumbhiH sarvathA syAn nirarthatA//
11504 	saMskAraupagame mukhyaM gajasnAnam idaM bhavet //230//
11505 iti saMgrahazlokaH / api ca / zabdArthayoH sambandho nityo vA 
11506 syAd anityo vA / yady anityaH puruSaicchA vRttir avRttir vA / apuruSa
11507 adhInatve puruSANAM yathA abhiprAyaM dezAdiparAvRttyA tena  
11508 pratipAdanaM na syAt / icchAyAm apy anAyattasya kadAcid ayogAt /
11509 parvatAdivat / ayam eva nityatve api doSas tasya sthirarUpasya
11510 parAvRttyayogAd iti samaM sarvAvasthAne api iSTapratiniyamAbhAvAt /
11511 tato vizeSapratipattir na syAd iti pUrvavat prasaGgaH / 
11512 icchAvRttau ca pauruSeyatvam iti vipralambhAzaGkA / api ca /  
11513 	sambandhinAm anityatvAn na sambandhe asti nityatA /
11514 parAzrayo hi sambandho apratibandhe tayoH sambandhitA ayogAt /
11515 sa ca Azrayo anityaH / apAye asya sambandhasya apy apAyaH anyathA anAzritaH
11516 syAt / tato na nityaH / tadAzrayArthaz ca vaktavyaH / 
11517 nityasya anupakAryatvAt / anupakurvANaz ca anAzrayaH / jAter vAcyatvAd
11518 adoSa iti cet / na / tadvacane prayojanAbhAvAd iti nirloThitam 
11519 etat / sarvatra ca jAtyasambhavAd ayogo yAdRcchikeSu vyakti
11520 vAciSu sarvadA jAticodane vizeSAntaravyudAsena pravRttyayogAc
11521 ca / tasmAd anvayavyatirekiNo bhAvasya bhAvAbhAvau sambandhaH / 
11523 	arthair ataH sa zabdAnAM saMskAryaH puruSair dhiyA //231  /  
11524 tAv eva bhAvAbhAvAv AzrityAsaMsRSTAv api saMsRSTAv iva puruSasya
11601 vyavahArabhAvanAtaH pratibhAta iti pauruSeyo bhAvAnAM 
11602 saMzleSaH / kiM ca AzrayavinAzAn naSTe sambandhe sa zabdaH punar
11603 asambandhatvAn na apUrveNa yojyeta / utpannautpannAz ca bhAvAH
11604 sthitasambandhAbhAvAd asambandhino avAcyAH syuH / tatra
11605 api /
11606 	arthair eva saha utpAde /
11607 kalpyamAne
11608 	na svabhAvaviparyayaH /  
11609 	zabdeSu yuktaH / 
11610 atha mA bhUn naSTasambandhasya zabdasya arthAntare vaiguNyaM
11611 arthAnAM ca avAcyatA ity utpanno arthaH sambandhavAn yady utpadyeta  
11612 sa sambandha utpanno api na zabde syAt / tasya tena asambandhi 
11613 svabhAvasya svabhAvaviparyayam antareNa tadbhAvAyogAt /  
11614 arthena saha utpannasya anyataH siddhasya anupakAriNi zabde  
11615 asamAzrayAc ca / tasya api tadutpattisahakAritve samarthasya nityautpAdana
11616 prasaGgaH / anapekSatvAn nityasya anupakArAt / asAmarthye 
11617 api pazcAd api svabhAvAtyAgAd azaktiH /  
11618 	sambandhe na ayaM doSo vikalpite //232//  
11619 na hi bhAvazleSApekSI puruSabhAvanApratibhAsI tadapekSAlakSaNaH
11620 sambandhaH / so ayam nityAnAm apy aparAvartayan svabhAvaM 
11621 kutazcit svayam utprekSya ghaTayed iti te api tathA 
11622 syuH / na ca cyavanadharmANaH / yad uktam AzrayApAyena Azrita  
11623 sambandhavinAzAd anityaH sa iti tatra /  
11624 	nityatvAd AzrayApAye apy anAzo yadi jAtivat / 
11625 	nityeSv AzrayasAmarthyaM kiM yena iSTaH sa AzrayaH //233// 
11626 zrUyata etan nityA jAtir AzrayitA ca na apy AzrayeNa saha nazyati iti /
11627 kevalaM nityeSv AzrayasAmarthyaM na pazyAmaH yena asAv AzrayaH / 
11628 kRtasya karaNAbhAvAd akArakasya ca anapekSatvAt / vyaktir  
11629 upakAro jAteH sambandhasya ca AzrayAt tena Azraya iti cet
11701 	jJAnautpAdanahetUnAM sambandhAt sahakAriNAm /
11702 	tadutpAdanayogyatvena utpattir vyaktir iSyate //234// 
11703 	ghaTAdiSv api yuktijJair avizeSe avikAriNAm /
11704 	vyaJjakaiH svaiH kutaH ko artho vyaktAs tais te yato matAH //235 //
11706 sahakAriNaH sakAzAd upAdAnApekSAd jJAnajananayogyakSaNAntara
11707 utpattir eva ghaTAdInAm abhivyaktiH / anyathA anapekSya tadupakAraM
11708 jJAnautpAdanaprasaGgAt / sAmarthyakAriNaz ca janakatvAt /  
11709 tasya ca tadAtmakatvAt / arthAntaratve ca bhAvAnupakAra  
11710 prasaGgAt sAmarthyAc ca jJAnautpatter nityaM ghaTAdInAm
11711 agrahaNApatter anAlokApekSagrahaNaprasaGgAd anapekSAtmAnupakArAt /
11712 tad ime svaviSayajJAnajanane param apekSamANAs 
11713 tataH svabhAvAtizayaM svIkurvanti / tena asya te janyAH / jJeya  
11714 rUpAsAdhanAt tu jJAnavazena kAryAtizayavAcinA zabdena vizeSa
11715 khyAtyarthaM vyaGgyAH khyApyante / na evaM jAtisambandhAdayaH 
11716 kathaMcid apy anupakAryatvAd anupakAriNA vyaktA yujyate /
11717 	sambandhasya ca vastutve syAd bhedAd buddhicitratA / 
11718 sa ca ayaM sambandho vastu bhavan niyamena zabdArthAbhyAM 
11719 bhedAbhedau na ativartate / rUpaM hi vastu / tasya atattvam eva anyattvam 
11720 ity uktam / sa ca ayam aindriyaH san svabuddhau tadanyavivekina 
11721 apratibhAsamAno rUpeNa kathaM tathA syAt / dRzya avivekAdarzanayor
11722 vivekasattAviparyayAzrayatvAt / anyathA tatsthiter
11723 abhAvaprasaGgAt / atIndriyatvAd apratibhAse api indriyAdiSv iva adoSa
11724 iti cet / na / tato apratipattiprasaGgAt aprasiddhasya ajJApakatvAt / 
11725 saMnidhimAtreNa jJApane avyutpannAnAm api syAt /  
11726 na anumAnAt pratipattir liGgAbhAvAt dRSTAntAsiddhez ca tatra apy 
11801 atIndriyatvena sAdhanApekSaNAt / tulyam indriyAdiSv api iti cet / na /
11802 teSAm anyathAnumAnAt / jJAnaM hi keSucit satsu vyatirekAnvayavat
11803 tanmAtrAsambhavaM tadvyatiriktApekSAM ca sAdhayati / 
11804 tataH kAryadvAreNa indriyasiddhiH / na evaM sambandhasya / tasya asiddhau
11805 tatkAryasya eva jJAnasya abhAvAt / na hi tatra zabdarUpam 
11806 artho vA liGgaM tayoH sarvatra yogyatvAt / vizeSapratItisamAzrayasya
11807 apratyAyanAd apratItir asya / na hy asati sambandhavizeSe 
11808 sA yuktA / tasyAM vA animittAyAM tadvizeSapratItiniyamavad  
11809 arthapratipAdanam api zabdAnAm animittaM kiM na iSyate / tasmAt
11810 tat sadRzaM liGgaM sarvasambandhe tato avizeSeNa gamayet /
11811 tato avizeSeNa eva pratItiH syAt sarvasya ca / tasmAt sambandha 
11812 siddhyarthapratIter na kazcit sampradAyam apekSeta / sampradAya  
11813 sahitasya liGgatvam iti cet / tat kim anayA paraMparayA /
11814 sa eva sampradAyApekSo arthajJApanaM kiM na karoti / sa ca zabdo  
11815 yad abhiprAyaiH prayujyamAno dRSTo anyathA na dRSTo darzana  
11816 adarzanAbhyAM dhUmAdivat tatpratItiM janayati iti sa eva
11817 sambandho avinAbhAvAkhyaH / na ca atra anyasya sAmarthyaM pazyAmaH /
11818 na api siddhyupAyam / atha punar na zabdArthayor anya eva 
11819 sambandhaH /  
11820 	tAbhyAm abhede tAv eva na ato anyA vastuno gatiH //236//
11821 rUpabhedanibandhanatvAd vyavasthAntarasya tadrUpaM tad eva 
11822 syAt / dharmabhedas tu syAt pUrvauktakrameNa / sa ca aviruddha
11823 eva na vastubhedaH / na ca bhedAbhedau muktvA vastuno anyA  
11824 gatiH / tasya rUpalakSaNatvAd / rUpasya ca etadvikalpAnativRtteH / 
11825 api ca /
11826 	bhinnatvAd vasturUpasya sambandhaH kalpanAkRtaH /
11827 ity uktaM prAk / na hi zleSalakSaNaH sambandho azliSTeSu padArtheSu 
11828 sambhavati / na ca arthAntaram eSAM sambandhaH / yasmAt
11901 	saddravyaM syAt parAdhInaM sambandho anyasya vA katham //237 //
11903 na hi siddhaM sat param apekSate / na anapekSaH svatantraH sambandhaH /
11904 dravyam iti ca svabhAva ucyate / sa kathaM parabhAvasya  
11905 zleSaH syAt / na hi svabhAvAntarasattayA anyaH zliSTo nAma / 
11906 mA bhUd azliSTena zliSTena tu syAd iti cet / na / tasya eva
11907 tAbhyAM zleSAsiddheH sa eva asiddho yas tau zleSayet / tad ayam 
11908 atiprasaGgo yady arthAv arthAntareNa zliSyato vizeSAbhAvAt / kiM
11909 ca / 
11910 	varNA nirarthakAH santaH padAdi parikalpitam /
11911 	avastuni kathaM vRttiH sambandhasya asya vastunaH //238//
11912 vAcako hi vacanAGgena tadvAn syAt / santo apy avAcakA varNAH / 
11913 tan na teSu vAcyavAcakasambandhaH / tadvRttau svarUpahAni 
11914 prasaGgAt / kramavizeSeNa vAcakA varNA iti cet / na / kramasya anartha
11915 antaratvena abhedakatvAt / tadrUpasya kramAntare apy avizeSAt 
11916 tulyA syAt pratipattiH / arthAntaratvam api kramasya niSetsyAmaH / 
11917 tad asati varNAnAM vAcakatve padAdi vAcakaM syAt /
11918 tac ca na kiMcid vyatirekAvyatirekavirodhAt / tasmAd indriya  
11919 vijJAnavizeSAnubandhI sabhAgavAsanAupAdAna vikalpapratibhAsavibhramaH
11920 padam / vAkyaM ca ekAvabhAsi mithyA eva / ekAnekatva  
11921 ayogAt / na hy ekam / anekayA buddhyA krameNa grahaNAyogAt /  
11922 na ca tad ekayA grAhyaM varNAnukramagrahaNAt / ekavarNa  
11923 grahaNakAle ca anekabuddhivyatikramAt / kSaNikatvAt buddhInAm /
11924 kSaNasya ekaparamANuvyatikramakAlatvAt / Adhikye vibhAgavataH  
11925 paryavasAnAyogAt / anekANuvyatyayanimeSatulyakAlatvAd 
11926 antyavarNaparisamApteH / yathAnubhavaM smaraNAt  
11927 smRtir api tatkAla eva / anubhavasmaraNAnukramayor vizeSAnupalakSaNatvAc
11928 ca / na apy anekaM padAdi / abhedapratibhAsanAd buddhes 
11929 tadanekatvasya niSetsyamAnatvAc ca / tan na vastu / tasya  
11930 etadvikalpAnatikramAt / vastu ca sambandhaH / sa kathaM
12001 tadAzrayaH syAt / AzrayanIyAyogAt / anAzrito hy evaM syAt /  
12002 tathA ca asambandhaH / tasmAn na svAbhAvikaH zabdArthayoH  
12003 sambandhaH / tadabhiprAyasya prayogAd utpanno abhivyakto
12004 vA zabdo tadavyabhicArI iti tattvam asya sambandhaH / sA ca utpattir  
12005 abhivyaktir vA avyabhicArAzrayaH pauruSeyI iti pauruSeya eva  
12006 sambandhas taddvAreNa ca arthapratyAyane aniyamaH zabdAnAm 
12007 ity apauruSeyatve api sa eva vipralambhaH / 
12008 	apauruSeyatA api iSTA kartRRNAm asmRteH kila /
12009 yA api iyam apauruSeyatA vedavAkyAnAM kartur asmaraNAd varNyate /  
12011 	santy asya apy anuvaktAra iti dhig vyApakaM tamaH //239// 
12012 tasya eva tAvad IdRzaM prajJAskhalitaM kathaM vRttam iti savismaya
12013 anukampaM naz cetaH / tad apare apy anuvadanti iti nirdaya  
12014 AkrAntabhuvanaM dhig vyApakaM tamaH / kaH prANino hitaIpsA
12015 vipralabdhasya aparAdhaH / tathA hi smaranti saugatA mantrANAM 
12016 kartRRn aSTakAdIn / hiraNyagarbhaM ca kANAdAH / teSAM  
12017 sa mithyAvAda iti cet / ka idAnIm evaM pauruSeyo anyo api / 
12018 kumArasambhavAdiSv AtmAnam anyaM vA praNetAram upadizanto  
12019 yad evaM prativyUhyeran / tatra prativahane abhyupeta 
12020 bAdhA iti cet / nanv idam eva abhyupagamAGgam iti kasya bAdhA /  
12021 tat parasya api tulyam eva / tasya iSTatvAd adoSa iti cet / kuto asya iyam
12022 iSTir apramANikA prAg AsIt / akasmAd grAhI ca ayaM kiM punaH
12023 kvacit sAdhanam apekSate / yat pauruSeyApauruSeyacintayA 
12024 AtmAnam AyAsayati / tata eva iSTer anabhyupetabAdhAyAM tadanyasya
12025 api tulyam ity anupAlambhaH / anatizayadarzI ca ayaM 
12026 vAkyeSv evaMprakArANAm apauruSeyatvasAdhanAnAM kAryadharmANAM 
12101 vA kvacid atizayam abhyupeta ity apratyaya eva asya vRttiH / 
12102 dRzyante ca vicchinnakriyAsampradAyAH kRtakAz ca / tAn yatnavanta  
12103 upalabhanta iti cet / na / niyamAbhAvAt / anyatra anupalambhasya
12104 upalambhasya vA paraupadezAd apratyayAd anizcayArthatvAt /
12105 svayaMkRtAnAm apy apahnotRdarzanAt / niSThAgamanasya 
12106 azakyatvAt / 
12107 	yathA ayam anyato azrutvA na imaM varNapadakramam /  
12108 	vaktuM samarthaH puruSas tathA anyo api iti kazcana //240//
12109 tasya api tad eva uttaram evam apauruSeyatve api kim idAnIM pauruSeyam  
12110 ity Adi / tathA hi / 
12111 	anyo vA racito granthaH sampradAyAd Rte paraiH /
12112 	dRSTaH ko abhihito yena so apy evaM na anumIyate //241//
12113 na khalu kiMcid anyad apauruSeyatvAzrayo anyatra idAnIntanAnAm  
12114 anupadezapAThAzakteH / sA ca anyatra apy ekena racite granthe anyasya
12115 tulyA / tadanusAriNA sarvas tathA unneyo na vA kazcit / tasya 
12116 tathA aniSTatvAd ity AdAv apy uktam iSTes tadAzrayatvAd ity Adi / 
12117 api ca /
12118 	yajjAtIyo yataH siddhaH sa tasmAd agnikASThavat / 
12119 	adRSTahetur anyo apy aviziSTaH sampratIyate //242// 
12120 na adarzanAd hetur ahetuko nAma / adRSTahetavo api hi bhAvAs
12121 tadanyaiH svabhAvAbhedam anubhavantas tathAvidhAH samunnIyante /
12122 heturUpanivRttAv api tadrUpam anivRttaM kAryadharma  
12123 vyatikramAn na tataH syAd iti na kazcit tathA vacanIyaH / 
12124 rUpavizeSo vA tathA darzanIyo ya enaM hetum anuvidadhyAt / yena iSTa
12125 aniSTayor iSTaviparyayo na syAt / svabhAvanivRttez ca hetor  
12126 abhedane bhAvAnAM bhedaH syAd Akasmika iti na kvacid vinivarteta /
12127 tasmAd yaH svabhAvo yajjanmA dRSTaH so anyatra apy avibhajyamAnaH 
12128 svAtmanA tatkAryadharmatAM na ativartate agniindhanavat /
12201 	tatra apradarzya ye bhedaM kAryasAmAnyadarzanAt /  
12202 	hetavaH pravitanyante sarve te vyabhicAriNaH //243// 
12203 yathAdyo api pathikakRto agnir jvAlAntarapUrvako na araNinirmathana
12204 pUrvakaH / pathikAgnitvAt / anantarAgnivad iti / kathaM
12205 punaH pathikAgner vyabhicAraH / jvAlAudbhavasAmarthyaM hy
12206 Azritya dahanasya hetvantaraM pratikSipyate / yadi hi vinA jvAlayA
12207 syAd anyatra api syAd iti / tatra jvAlAitarajanmanor abAdhyabAdhakatve
12208 jvAlAprabhavatvam anyathA api syAd iti dharmayor  
12209 ekatra arthe sambhavAt sa pathikAgnir anyo vA artha ekapratiniyato
12210 na syAd ity AzaGkyate vyabhicAraH / so apy anyonyavyatirekI 
12211 dharmadvayAvatAro vastusAmAnye aviruddha ity ucyate na avasthA
12212 bhedini vizeSe / niSkalasya Atmanas tadatattvavirodhAt / na ca  
12213 jvAlAitarajanmanor bAdhyabAdhakatA pathikAgnau / tasya jvAlAprabhava
12214 vyatirekeNa asambhavAbhAvAt / evaMbhUtaH pathikAgnir
12215 jvAlAprabhava iti syAt / na sarvaH / tatra vizeSapratikSepasya 
12216 kartum azakyatvAt / sambhavadvizeSasya ca tAdavasthyAniyamAt /
12217 yad api vinA jvAlayA syAd anyatra api syAd iti / bhavaty eva / 
12218 yayA sAmagryA sambhavati sA yadi syAt asyAH sambhavaM pradarzya  
12219 tadabhAvaM pradarzayet tatra vA jvAlAm syAd etat /
12220 tasmAn na ekasya parapUrvakam adhyayanaM sarvasya tathAbhAvaM  
12221 sAdhayati / tasya anyathA asambhavAbhAvAt / tathAvidhasya 
12222 tu tatkriyApratibhArahitasya tathA syAd iti / tathAbhUtam  
12223 evaM vAcyaM syAt tad avizeSeNa sambhavadvizeSam ucyamAnaM 
12224 chAyAM na puSNAti / kathaM vizeSasya sambhavo yAvatA  
12225 teSAm api puruSANAm azaktir eva idAnIntanapuruSavat / atra api  
12226 zaktipuruSayor na kiMcid virodhadarzanam iti na aviruddhavidhir  
12227 anupalabdhiprayogo gamakaH / na hy atIndriyeSu virodhagatir
12228 asti ity uktam / na ca ayaM pUrvaprayogAd bhidyate / yadi puruSAH  
12229 zaktAH syur idAnIntanA api iti / vizeSAsambhava etat syAt /  
12301 sa ca duHsAdhaH / yatra ekasya azaktis tatra sarvapuruSANAm ity api
12302 pUrvavad vyabhicAri / bhAratAdiSv idAnIntanAnAm azaktAv  
12303 api kasyacit zaktisiddheH / tasmAt kAraNAni vivecayata artheSv  
12304 api tadatatpratibhaveSu svabhAvabhedo darzanIyaH / tadabhAve 
12305 sarvas tadAtmA na vA kazcit / na ca atra laukikavaidikayoH svabhAva 
12306 bhedaM pazyAmaH / asati tasmiMs tayoH sAmAnyasya eva 
12307 dRSTer ekasya kazcid dharmaM vivecayan tatsvabhAvasambhavinA 
12308 tena AzaGkyavyabhicAravAdaH kriyate / nanu vedAvedayos 
12309 tattvalakSaNo asty eva vizeSaH / satyam asti / na kevalaM 
12310 tayor eva / kiM tarhi / DiNDikapurANaitarayor api / na ca svaprakriyA
12311 bhedadIpano nAmabhedaH puruSakRtiM bAdhate / anyatra 
12312 api prasaGgAt / yadi tAdRzIM racanAM puruSAH kartuM na zaknuyuH
12313 kRtAM vA akRtasaMketo vivecayet vyaktam apauruSeyo
12314 vedaH / nanu na zaknuvanty eva puruSA mantrAn kartum / etad 
12315 uttaratra vicArayiSyAmaH / api ca / na mantro nAma anyad eva  
12316 kiMcit / kiM tarhi / satyatapaHprabhAvavatAM samIhitArthasAdhanaM
12317 vacanam / tad adyatve api puruSeSu dRzyata eva / yathAsvaM
12318 satyAdhiSThAnabalA viSadahanAdistambhanadarzanAt /
12319 zabarANAM ca keSAMcid adya api mantrakaraNAt / avaidikAnAM
12320 ca bauddhAdInAM mantrakalpAnAM darzanAt / teSAM ca puruSa
12321 kRteH / tatra apy apauruSeyatve katham idAnIm apauruSeyam
12322 avitatham / tathA hi bauddhaitarayor mantrakalpayor hiMsA
12323 maithunAtmadarzanAdayo anabhyudayahetavo anyathA ca varNyate /
12324 tat katham ekatra viruddhAbhidhAyi dvayam satyaM  
12325 syAt / tatra arthAntarakalpane tad anyatra api tulyam ity arthAnirNayAt
12326 kvacid apratipattiH / tathA ca sad apy anupayogam apauruSeyam / 
12327 bauddhAdInAm amantratve tadanyatra api koSapAnaM
12328 syAt karaNIyam / viSakarmAdhikRto bauddhA api dRzyante / tatra amantratvam
12401 api vipratiSiddham / mudrAmaNDaladhyAnair apy 
12402 anakSaraiH karmANi kriyante / na ca tAny apauruSeyANi nityAni
12403 yujyate / teSAM kriyAsambhave akSararacanAyAM kaH pratighAtaH
12404 puruSANAm / tasmAn na kiMcid azakyakriyam eSAm / katham  
12405 idAnIM satyaprabhavau mantrakalpau parasparavirodhinau / na 
12406 vai sarvatra tau satyaprabhavau / prabhAvayuktapuruSapratijJA
12407 lakSaNAv api tau staH / sa prabhAvo gatisiddhivizeSAbhyAm api 
12408 syAt / yadi pauruSeyA mantrAH kiM na sarve puruSA mantrakAriNaH /
12409 tatkriyAsAdhanavaikalyAt / yadi tAdRzaiH satyatapaHprabhRtibhir  
12410 yuktAH syuH kurvanty eva / api ca kAvyAni puruSaH karoti iti 
12411 sarvaH puruSaH kAvyakRt syAt / akaraNe vA na eva kazcit tadvad  
12412 ity apUrvA eSA vAcoyuktiH / satyaM mantrakriyAsAdhanavikalA 
12413 mantrAn na kurvate / tat tu kasyacit sAkalyaM na pazyAmaH /
12414 puruSANAM samAnadharmatvAt / uktam atra na mantro nAma anyad  
12415 eva kiMcit satyAdimatAM vacanasamayAd iti / tAni ca kvacit
12416 puruSeSu dRzyante / sarvapuruSAs tadrahitA ity api tatsambhava
12417 virodhAbhAvAd anirNayaH / na ca atyakSasvabhAveSv anupalabdhir
12418 nizcayahetuH / na ca smRtimatiprativedhasatyazaktayaH 
12419 sarvabhAvinyaH / tatsAdhanasampradAyabhedavad guNAntara  
12420 sAdhanAny api syuH / na api sann api sarvo draSTuM zakyaH / ata
12421 eva adRSTasya anapahnavaH / na api puruSeSu utpitsoH kasyacid guNasya  
12422 pratiroddhA / bAdhyAdRSTer bAdhyabAdhakabhAvAsiddheH / etena  
12423 sarvajJAnapratiSedhAdayo api nirvarNitauttarAH / tatra apy evaMbhUto
12424 yAdRzo ayam asambhavattatsAdhanasampradAyo na iti nyAyaH /
12425 na adRSTajJApako atatsvabhAva ity api / satAM api kAryAnArambha 
12426 sambhavAt / svabhAvaviprakarSeNa draSTum azakyatvAc ca /  
12427 tasmAd adhyayanam adhyayanAntarapUrvakam adhyayanAd iti 
12428 bhAratAdhyayane api bhAvAd vyabhicAri / nanu vedane vizeSaNAd
12429 adoSaH / kaH punar atizayo vedAdhyayanasya yad anyathA adhyetuM  
12430 na zakyate / na hi vizeSaNam aviruddhaM vipakSeNa asmAd
12431 hetuM vyAvartayati / aviruddhayor ekatra sambhavAt / idAnIntanAnAm  
12501 adhyayanAd iti cet / uktauttaram etat / adarzanAd 
12502 iti cet / idam api prativyUDham / na apy adarzanamAtram abhAvaM 
12503 gamayati iti vyabhicAra eva / tasmAn na vizeSaNam atizaya  
12504 bhAg ity anupAttasamam / yat kiMcid vedAdhyayanaM sarvaM 
12505 tadadhyayanAntarapUrvakam ity api vyAptir na sidhyati / 
12506 sarvasya tathAbhAvAsiddheH / yAdRzaM tu tannimittaM dRSTaM 
12507 tat tathA iti syAt / dRSTe vizeSe tannimittatayA tattyAgena sAmAnya
12508 grahaNaM vyabhicAry eva / hutAzanasaMsiddhau pANDudravyatvavat /  
12509 etena vacanAdayo rAgAdisAdhane pratyuktAH / astu 
12510 vedam adhyayanam adhyayanapUrvatAsAdhanam /  
12511 	sarvathA anAditA sidhyed evaM na apuruSAzrayaH /
12512 	tasmAd apauruSeyatve syAd anyo apy anarAzrayaH //244// 
12513 puruSa eva hi svayam abhyUhya parato vA AdhIyate / na eSAm avyApRta
12514 karaNAnAM svayaM zabdA dhvananti yena apauruSeyAH syuH /  
12515 api syur apauruSeyA yadi puruSANAm AdiH syAt / tadA apy anya 
12516 pUrvakaM na sidhyati / adhyApayitur abhAvAt / tat prathamo
12517 adhyetA kartA eva syAt / tad ayam anAdiH pUrvapUrvadarzana 
12518 pravRtto DimbhakapAMsukrIDAAdivat puruSavyavahAra iti syAt / 
12519 na apauruSeya iti / anAditvAd apauruSeyatve bahutaram idAnIm 
12520 apauruSeyam / tathA hi / 
12521 	mlecchAdivyavahArANAM nAstikyavacasAm api / 
12522 	anAditvAt tathAbhAvaH pUrvasaMskArasantateH //245// 
12523 mlecchavyavahArA api kecin mAtRvivAhAdayo madanautsavAdayaz  
12524 ca anAdayaH / nAstikyavacAMsi ca apUrvaparalokAdyapavAdIni / na
12525 hi tAny anAhitasaMskArAH paraiH pravartayanti / sva pratibhAracita
12526 samayAnAm api yathA azrutArthavikalpasaMhAreNa eva pravRtteH / 
12527 tat kiMcit kutazcid Agatam ity ekasya upadeSTuH prabandhena abhAvAd 
12528 aparapUrvakam ity ucyate / prAg eva yathAdarzana  
12529 pravRttayaH samyaGmithyApravRttayo lokavyavahArAH / nanv
12530 AdikalpikeSv adRSTA eva vyavahArAH pazcAt pravRttA iSyante / 
12601 na / teSAm apy anyasaMskArAhitAnAM yathApratyayaM prabodhAt /  
12602 bhavatu sarveSAm apauruSeyatvam iti cet /  
12603 	tAdRze apauruSeyatve kaH siddhe api guNo bhavet / 
12604 kAmam avisaMvAdakam ity apauruSeyatvam iSTam / tad visaMvAdakAnAm 
12605 api keSAMcid anAditvAd asti iti kim apauruSeyatvena /
12606 sati vA vedavAkyAnAm eva apauruSeyatve /  
12607 	arthasaMskArabhedAnAM darzanAt saMzayaH punaH //246// 
12608 yadi apauruSeyatve api pratiniyatAm eva tadarthapratibhAM 
12609 janayed AzvAsanaM syAt / yatheSTaM tu samAropApavAdAbhyAM  
12610 nairuktamImAMsakAdayo vedavAkyAni vizasanto dRzyante / na  
12611 ca te arthAs teSAM na saMghaTante / samayaprAdhAnyAd arthanivezasya
12612 ekasya api vAkyasya anekavikalpasambhavAt / prakRtipratyayAnAm  
12613 anekArthapAThAt / rUDher apy ekAntena ananumateH / arUDha  
12614 zabdabAhulyAt / tadarthasya puruSaupadezApekSaNAt / tadupadezasya
12615 tadicchAnuvRtter anirNaya eva vedavAkyArtheSu / 
12616 api ca / ayam apauruSeyatvaM sAdhayan varNAnAM vA sAdhayed  
12617 vAkyasya vA / tatra /
12618 	anyAvizeSAd varNAnAM sAdhane kiM phalaM bhavet /
12619 na hi lokavedayor nAnA varNAH / bhede api ca pratyabhijJAnAvizeSAt
12620 tata ekatvAsiddhiprasaGgAd bhedAnupalakSaNAc ca vaidika  
12621 varNAsiddhiH pratyabhijJAnAd apratipattiprasaGgAt / anabhyupagamAc 
12622 ca / teSAM ca apauruSeyatvasAdhane te tulyAH sarvatra  
12623 iti kim anena parizeSitam / tathA ca sarvo vyavahAro apauruSeyaH /
12624 na ca sarvo avitatha iti vyarthaH parizramaH / atha vAkyam
12625 apauruSeyam iSTam / 
12701 	vAkyaM na bhinnaM varNebhyo vidyate anupalambhanAt //247//
12703 na hi vayaM devadattAdipadavAkyeSu dakArAdipratibhAsaM muktvA
12704 anyaM pratibhAsaM buddheH pazyAmaH / dvitIyavarNapratibhAsavat /
12705 na ca apratibhAsamAnaM grahaNe grAhyatayA iSTam asty
12706 anyad vA iti zakyam avasAtum / AkArAntaravat / anyAsambhavi
12707 kAryaM gamakam iti cet / syAd etat yadi teSu varNeSu satsu api 
12708 tat kAryaM na syAt / na bhavati teSAm avizeSe api padavAkyAntare
12709 abhAvAd iti cet / na / teSAm avizeSAsiddheH / avizeSaH pratyabhijJAnAt
12710 siddha iti cet / na / tasya vyabhicArAd anidarzanatvAc  
12711 ca / varNAvizeSe api vAkyabhedAt pratipattibhedaH kAryabhedaH  
12712 syAt / sA ca vAkyAt / tac ca atIndriyam iti kutaH syAt / saMnidhAna
12713 mAtreNa janane avyutpannasya api syAt / tasmAn na vAkyaM  
12714 nAma kiMcid arthAntaraM varNebhyo yasya apauruSeyatvaM sAdhyeta /  
12715 tadabhAvAd vedAviziSTavarNApauruSeyatvam api prathama
12716 pakSe pratyuktam / api ca astv arthAntaraM vAkyam /  
12717 tad anekAvayavAtmakaM vA syAd anavayavaM vA /
12718 	anekAvayavAtmatve pRthak teSAM nirarthatA / 
12719 te api tasya bahavo avayavAH pRthak prakRtyA yady anarthakAH /
12720 	atadrUpe ca tAdrUpyaM kalpitaM siMhatAAdivat //248//
12721 arthavAn eva AtmA vAkyam / te ca avayavAH svayam anarthakAH / 
12722 teSu sa AtmA kalpanAsamAropitaH syAt / siMhatAAdivan mANavaka 
12723 AdiSv iti pauruSeya eva / atha mA bhUd eSa doSa iti pratyekaM 
12724 te avayavAH sA arthakA iSyante /
12801 	pratyekaM sA arthakatve api mithyAnekatvakalpanA / 
12802 	ekAvayavagatyA ca vAkyArthapratipad bhavet //249//  
12803 parisamAptArthaM hi zabdarUpaM vAkyam / te ca avayavAs tathAvidhAH
12804 pRthak pRthag iti pratyekaM te vAkyam / tathA ca na aneka  
12805 avayavaM vAkyam / ekAvayavapratipattyA ca vAkyArthapratipatter
12806 avayavAntarApekSA kAlakSepaz ca na syAt / tasya niSkala 
12807 AtmanaH kSaNena pratipatter ekajJAnautpattau niHzeSAvagamAt /
12808 anyathA ca ekatvavirodhAt /  
12809 	sakRc zrutau ca sarveSAM kAlabhedo na yujyate / 
12810 atha mA bhUd avayavAntarApratIkSaNena ekasmAd eva avayavAd 
12811 vAkyArthasiddher anekAvayavatvahAnir vAkyasya iti sakRt sarva
12812 avayavAnAM zravaNam iSyate / tadA api kAlakSepo na yukta eva / 
12813 ekAvayavapratipattikAla eva sarveSAM zravaNAt / kramazravaNe
12814 ca pRthag arthavatAm ekasmAd eva tadarthasiddher anyasya  
12815 vaiyarthyAt / sakRc zrutau ca pRthag artheSv adRSTasAmarthyAnAm
12816 arthavattA ca na sidhyati / sahiteSv arthadarzanAd adoSaH /
12817 na / pRthag asato rUpasya saMghAte apy asambhavAt / arthAntarAnutpattez
12818 ca / zabdautpAdavAdinas tAvad ayam adoSa eva / pRthag 
12819 asamarthAnAm apy avayavAnAm upakAravizeSAd atizayavatAM  
12820 kAryavizeSaupayogAt / pratyekaM tv avayaveSu samartheSu vyarthA
12821 syAd anyakalpanA / atha punar ekam eva anavayavaM vAkyam /
12822 tatra / 
12823 	ekatve api hy abhinnasya kramazo gatyasambhavAt //250//
12824 kAlabheda eva na yujyate / na hy ekasya krameNa pratipattir 
12825 yuktA / gRhItAgRhItayor abhedAt / krameNa ca vAkyapratipattir  
12826 dRSTA / sarvavAkyavyAhArazravaNasmaraNakAlasya anekakSaNanimeSa
12827 anukramaparisamApteH / varNarUpAsaMsparzinaz ca ekabuddhipratibhAsinaH
12828 zabdAtmano apratibhAsanAt / varNAnukramapratIteH /
12829 tadavizeSe apy anukramakRtatvAd vAkyabhedasya anukramavatI 
12901 vAkyapratItiH / varNAnukramaupakArAnapekSaNe tair yathA kathaMcit
12902 prayuktair api yat kiMcid vAkyaM pratIyeta vinA vA  
12903 varNaiH / tair anukramavadbhir akramasya upakArAyogAt / akrameNa
12904 ca vyAhartum azakyatvAt / gatyantarAbhAvAc ca / na eva 
12905 vAkye varNAH santi tad ekam eva zabdarUpaM vyaJjakAnukrama
12906 vazAd anukramavad varNavibhAgavac ca pratibhAti iti cet /  
12907 anukramavatA vyaJjakena akramasya vyaktiH pratyuktA / vyakta  
12908 avyaktavirodhAd avarNabhAge ca vAkye asakalazrAviNo asakalavAkya
12909 gatir na syAt ekasya zakalAbhAvAt sakalazrutir na vA kasyacit / 
12910 samastavarNasaMskAravatyA antyayA buddhyA vAkyAvadhAraNam
12911 ity api mithyA / tasya avarNarUpasaMsparzinaH kasyacit kadAcid 
12912 apratipatteH / varNAnAM ca akrameNa apratipatteH kuto akramam  
12913 ekabuddhigrAhyaM vAkyaM nAma / na ca antyavarNapratipatter
12914 Urdhvam anyam azakalaM zabdAtmAnam upalakSayAmaH / na api  
12915 svayam ayaM vaktA vibhAvayati / kevalam evaM yadi syAt sAdhu 
12916 me syAd iti kalyANakAmatAmUDhamatir antyAyAM buddhau samAptakalaH
12917 zabdo bhAti iti svapnAyate / na hi smaryamANAyor 
12918 api padavAkyayor varNAH kramavizeSam antareNa vibhAvyante / 
12919 akramAyAM buddhau paurvAparyAbhAvAt / teSAM tatkRtaH pada 
12920 vAkyabhedAnAM bhedo na syAt / na apy avarNakramam anyac 
12921 zabdarUpaM pazyAma ity uktam / sati vA tad anityaM vA syAn 
12922 nityaM vA / yadi / 
12923 	anityaM yatnasambhUtaM pauruSeyaM kathaM na tat /
12924 avazyaM hy anityam utpattimat kutazcid bhavati / tathA hy Akasmikatve
12925 sattvasya dezAdiniyamo na syAd ity uktam / tac ca yatna 
12926 preritAviguNakaraNAnAM dRSTam anyathA na iti / kAraNadharma
13001 darzanAt puruSavyApAra eva kAraNam / ataH pauruSeyaM syAt /
13002 	nityaupalabdhir nityatve apy anAvaraNasambhavAt //251// 
13003 atha tac zabdarUpaM nityaM syAd upalabhyasvabhAvaM ca /  
13004 sa ca tasya svabhAvaH kadAcin na apaiti iti nityam upalabhyeta /  
13005 evaM hi sa nityaH syAd yadi na kutazcit sAmarthyAt pracyavet /
13006 jJAnajananasAmarthyasya tadAtmakatvAt / arthAntaratvasya ca  
13007 prAg eva niSiddhatvAt / na api tasya upalabhyAtmanaH kiMcid upalambha
13008 AvaraNaM sambhavati / tasya sato api tadAtmanam akhaNDayataH
13009 sAmarthyatiraskArAyogAt / na hi tatra atizayam anutpAdayan
13010 kiMcitkaro nAma / akiMcitkaraz ca kaH kasya AvaraNam anyad  
13011 vA iti nirloThitaprAyam etat / kuDyAdayo ghaTAdInAM kam atizayam
13012 utpAdayanti khaNDayanti vA yena AvaraNam iSyante / na brUmas 
13013 te kaMcid atizAyayanti iti / api tu na sarve ghaTakSaNAH sarvasya indriya
13014 vijJAnahetavaH parasparasahitAs tu viSayaindriyAlokAH
13015 parasparato viziSTakSaNAntarautpAdAd vijJAnahetavaH / anupakAryasya
13016 apekSAyogAt / zaktasvabhAvasya nityaM jananam ajananaM  
13017 vA anyasya sarvadA syAd ity uktam / te ca avyavahitAH pratighAtinA  
13018 anyena anyonyasya upakAriNaH / avyavadhAnadezayogyatAsahakAritvAt 
13019 teSAm anyonyAtizayautpatteH / vyavadhAne tu hetvabhAvAt  
13020 samarthakSaNAntarAnutpatter jJAnAnutpattiH / tasmAt pUrvautpanna
13021 samarthanirodhAt sati kuDye anyasya utpitsoH kAraNAbhAvena anutpatteH
13022 kAraNavaikalyAj jJAnAnutpattir iti kuDyAdaya AvaraNam  
13023 ucyante / na prAg yogyasya pratibandhAt / tasya svabhAvAd apracyuteH /
13024 athavA sambhavaty api bhAvAnAM kSaNikAnAm anyonyaupakAro 
13025 acintyatvAd hetupratyayasAmarthyasya asarvavidA / tena
13026 yad indriyaviSayamadhyasthitam AvaraNaM tat tau vijJAnautpatti  
13027 vaiguNyatAratamyabhedena atizayayed api / AvaraNabhede na zabda  
13028 Adau zrutimAndyapATavadarzanAt / anyathA kvacid apy akiMcitkarasya
13101 	saMnidhAnasya apy asaMnidhAnatulyatvAt tasya idam ity /
13102 upasaMhAro vikalpanirmita eva syAt / na vastvAzrayaH / na ca
13103 samAropAnuvidhAyino arthakriyAH / na hi mANavako dahanaupacArAd  
13104 AdhIyate pAke / tasmAt satyAm api kalpanAyAm atatparAvRttayo  
13105 bhAvA yathAsvabhAvavRttaya eva syuH / tat saty apy
13106 AvaraNe vijJApayeyur eva indriyAdayaH / na ca tathA / tasmAt tena AdheyavizeSA
13107 iti gamyante / na khalv evaM nityAnAM zabdAnAM
13108 kasmiMzcit saty atizayahAnir utpattir vA / tad yadi teSAM jJAnajananaH
13109 svabhAvaH sarvasya sarvadA sarvANi svaviSayajJAnAni sakRj 
13110 janayeyuH / no cen na kadAcit kasyacit kiMcid ity ekAnta eSaH /
13111 	azrutir vikalatvAc cet kasyacit sahakAriNaH / 
13112 syAd etat / na AvaraNAn nityaM sarve zabdA na zrUyante / api tu 
13113 kiMcit teSAM pratipattau sahakAri pratiniyatam asti / tat kadAcit
13114 kasyacid bhavati iti tatkRtam eSAM kadAcit kvacic zravaNam iti / 
13115 	kAmam anyapratIkSA astu niyamas tu virudhyate //252//
13116 na vai vayaM kAraNAnAM sahakArINi pratikSipAmaH / kiM tv
13117 apekSanta eva kAraNAni tadavasthAupakAriNam / tato labhyasya atizayasya
13118 kAryaupayogAt / tathA zabdo api yadi kiMcid apekSya 
13119 kAryaM kuryAt karotu pUrvasvabhAvaniyata ity etan na syAt /  
13120 tasya pracyuteH / apekSyAc ca svabhAvAntarapratilambhAt /  
13121 na hy anupakAry apekSyata ity uktam etat / tadupakArasya
13122 ca arthAntaratve tasya iti sambandhAbhAvAdayo apy uktAH / tasya
13123 ca ajJeyatvam / upakArAd eva jJAnautpatteH / tasmAd eva zabdo na indriyaM
13124 na saMnikarSaM na AtmAnam anyad vA kiMcij jJAnautpattisamAzrayaM
13125 svavijJAnajanane apekSate / sarvasya tatra akiMcitkaratvAd / 
13126 api ca / ete zabdA vyApino vA syur avyApino vA /
13127 	sarvatra anupalambhaH syAt teSAm avyApitA yadi /  
13128 	kathaM ekadezavartitaM taczUnyadezasthita upalabdheta / aprApta
13201 grahaNapakSe adoSa iti cet / na / tatra api yogyadezasthititAratamya
13202 apekSaNAd ayaskAntAdivat / anyathA spaSTAspaSTazruti  
13203 bhedo na syAt / sati ca upalambhapratyaye sarvatra deze tulyam 
13204 upalabhyeran / tasmAn na avyApinaH /
13205 	sarveSAm upalambhaH syAd yugapad vyApitA yadi //253//
13206 na hi kazcic zabda kvacin na asti iti sarve yugapad upalabhyeran 
13207 sarvadezasthitaiz ca / yogyaindriyatvAd viSayasaMnidhAnAd apratibandhAc
13208 ca / 
13209 	saMskRtasya upalambhe ca kaH saMskartA avikAriNaH / 
13210 syAd etat / sann api na sarvaH zabda upalabhyate sarveNa / saMskRtasya
13211 saMskRtena eva upalambhAd iti / tatra na saMskRtasya upalambho  
13212 anAdheyavikArasya saMskArAyogAt /  
13213 	indriyasya hi saMskAraH zRNuyAn nikhilaM ca tat //254//
13214 tatra yadi saMskRtena upalambha ity asaMskRtena indriyo na upalabhate /
13215 yasya saMskAra indriyasya kRtaH sa sarvazabdAn yugapac zRNuyAd  
13216 iti prasaGgo anivRtta eva /  
13217 	saMskArabhedAd bhinnatvAd ekArthaniyamo yadi / 
13218 	anekazabdasaMghAte zrutiH kalakale katham //255//  
13219 atha api syAt / pratiniyatAH te saMskArAH zabdAnAm / tatra kenacit  
13220 saMskRtam indriyaM kasyacid eva grAhakam iti na yugapat sarva  
13221 zabdazrutir iti / saMskAravizeSAc zrutiniyame indriyANAm aneka 
13222 zabdasaMghAtasya kalakalasya zrutir na syAt / na hy ekaH
13223 zabdaH kalakalo nAma / bhinnasvabhAvAnAM yugapac zravaNAt /
13224 svabhAvabhedAzrayatvAc ca bhedavyavasthiteH / laghuvRtteH 
13225 sakRczrutir bhrAntir iti cet / vaMzAdisvaradhArAyAM gamakAvayava
13226 saMhArAt saMkulA pratipattiH syAt / vakSyate ca atra pratiSedhaH / 
13227 tasmAd ekagatizaktipratiniyamAd indriyasya anekAtmA 
13228 kalakalo na zrUyeta / 
13229 	dhvanayaH kevalaM tatra zrUyante cen na vAcakAH /  
13301 na vai kalakale varNapadavAkyAni zrUyante / dhvanInAm eva 
13302 kevalAnAM zravaNAt / vAcake ca pratiniyatazaktiindriyaM na 
13303 dhvaniSu / tatra /  
13304 	dhvanibhyo bhinnam asti iti zraddhA iyaM atibahv idam //256// 
13305 na hi vayaM dhvaniM zabdaM ca vAcakaM pRthagrUpam upalakSayAmaH /  
13306 ekam eva ekadA varNAnukramazravaNe zabdAtmAnaM  
13307 vyavasyAmaH / tat kathaM vyavasAyapUrvakaM vyavahAram 
13308 avyavasyantaH pravartayAmaH / tasmAd dhvanivizeSa eva iyaM 
13309 varNAdyAkhyaH / api ca /
13310 	sthiteSv anyeSu zabdeSu zrUyate vAcakaH katham / 
13311 na dhvanir ato bhinnas tena saha pRthag vA / na hi pratyakSe 
13312 arthe paraupadezo garIyAn / tad ayaM sthiteSv anyeSu vyAhartRSu 
13313 kevalam eva zabdaM zRNvaMs tadupalambhapratyayAnAM tadanya 
13314 niSpAdane sAmarthyAbhAvaM pratyeti / yadi hi samarthAH syus
13315 tatsAdhitaM tad upalabhyeta / tatsvabhAvA eva punaH pratyayAH 
13316 kathaM kalakale arthAntaram Arabheran / na hi kAraNAbhede  
13317 kAryabhedo yuktaH / bhedasya ahetukatvaprasaGgAd ity uktam / 
13318 na ca kalakale vAcako na zrUyate / padavAkyavicchedAnAm upalakSaNAt /
13320 	kathaM vA zaktiniyamAd bhinnadhvanigatir bhavet //257//
13321 tAni pratiniyatazaktIny api indriyANi nAnArUpAn pratizabdaniyatAn
13322 dhvanIn zRNvanti na tv evaM zabdAn iti kaH zabdeSv eSAM 
13401 nirvedaH / yad uktaM na dhvanayo bhedena vAcakebhyaH siddhA 
13402 iti kathaM na siddhAH / vacanAd arthapratipatteH / na hi dhvani 
13403 bhAgAd alpIyasaH zabdArthapratItiH / na ca so anyaM sama iti /  
13404 tad iyaM samastapadavAkyarUpasAdhyArthapratItir asamasta
13405 bhAgeSu dhvaniSu na sambhavati iti siddham akramasattvaM 
13406 zabdarUpam / kramavadbhAgaz ca dhvanir iti / tan na / akramasya 
13407 kramavadvyatirekiNaH prAg eva niSiddhatvAt / atiprasaGgaz ca evam / 
13408 karmabhAgAnAM pUrveNa aparasya apratisandhAnAt ekAMzAc 
13409 ca apratipatteH tadvyatirekI hastasaMjJAAdiSv arthapratipatti  
13410 hetuH samastarUpaH karmAtmAbhyupagantavyaH syAt / zabdavad 
13411 eva / kramabhAvina eva yathAsvaM karaNaprayogAd bhinnA  
13412 varNabhAgAH karmabhAgA vA krameNa vikalpaviSayA yathA 
13413 saMketam eva arthapratItiM janayanti iti nyAyyam / kiM ca /
13414 	dhvanayaH sammatA yais te doSaiH kair apy avAcakAH /  
13415 	dhvanibhir vyajyamAne asmin vAcake api kathaM na te //258//  
13416 kramautpAdibhir dhvanibhAgair vyaktaH kila vAcako vakti / tam  
13417 api te na eva sakRt prakAzayanti / kramabhAvAt / na apy eka eva 
13418 bhAgaH zabdaM vyanakti / tadanyavaiyarthyaprasaGgAt / ekavarNa  
13419 bhAgakAle ca samastarUpAnupalakSaNAt / tad ayam apratisaMhita 
13420 sakalaupalambho dhvanivad upalambhasAkalyasaMnidhAna
13421 arthaM kathaM sAdhayet / ko hi vizeSo atyantAnupalambhe 
13422 sadasator upalambhasAdhyeSv artheSu / na ca saMnidhimAtreNa 
13423 sAdhanaH / vyaktyapekSaNAt / sA ca iyaM kramabhAvinI sadasatos 
13424 tulyaupayoga iti dhvanibhir azakyasAdhanaM kAryam / tatra api
13425 tathA ity alaM anyena / tasmAn na varNeSu vAkye vA apauruSeyatA / 
13426 	varNAnupUrvI vAkyaM cen na varNAnAm abhedataH /
13501 na arthAntaram eva zabdarUpaM vAkyam apauruSeyam / kiM tarhi / 
13502 varNAnukramalakSaNaM hi no vAkyam / tad apauruSeyaM sAdhyam 
13503 iti cet / na / varNAnAm AnupUrvyA abhedAt / na iyam arthAntaraM
13504 varNebhyaH / dRzyAyAM vibhAgaupalambhAt / adRzyAyAM tato  
13505 apratipatter liGgAbhAvAt / bhedavatyAz ca AnupUrvyA abhAve varNa
13506 mAtram avaziSTaM sarvatra iti pUrvavat prasaGgaH / 
13507 	teSAM ca na vyavasthAnaM kramAntaravirodhataH //259//  
13508 yady akRtakAnupUrvI varNAnAM te ca na bahavaH samAnajAtIyAH 
13509 yena kecid vyavasthitakramAH syuH anye yathA iSTaparAvRttayaH  
13510 kiM tarhi eka eva trailokye akAras tathA gakAraH tadA agnir ity 
13511 eva syAt na gaganam iti akAragakArayoH pUrvAparabhAvasya 
13512 vyavasthitatvAt / kRtakAnAm api hetupariNAmaniyamavatAm
13513 azakyaH kramaviparyayaH kartum / yathA bIjAGkurapatrAdInAm  
13514 RtusaMvatsarAdInAM ca / kiM punar acalitAvasthAsvabhAvAnAm
13515 akRtakAnAM kathaMcit sthitAnAM pUrvAvasthAtyAgam antareNa 
13516 anyathA abhAvAyogAt / tyAge vA vinAzaprasaGgAt / vizeSeNa nityAyAm
13517 AnupUrvyAM / tad etat pratipadaM kramAnyatvaM varNa  
13518 anyatve apUrvautpAdAt varNabAhulyAd vA syAt / tac ca anabhimatam /
13519 api ca /
13520 	dezakAlakramAbhAvo vyAptinityatvavarNanAt / 
13521 sA ca iyam AnupUrvI vArNAnAM dezakRtA vA syAt / yathA pipIlikAnAM
13522 paGktau / kAlakRtA vA yathA bIjAGkurAdInAm / sA dvividhA api
13523 varNeSu na sambhavati / vyApter nityatvAc ca / anyonyadeza
13601 parihAreNa vRttir hi dezapaurvAparyam / tat sarvasya sarveNa  
13602 tulyadezatvAd varNeSu na sambhavati / vAtAtapavad AtmAdivac 
13603 ca / tathA kAlaparihAreNa vRttiH kAlapaurvAparyam / yadA eko
13604 na asti tadA anyasya bhAvAt / tad api nityeSu na sambhavati / sarvadA 
13605 sarvasya bhAvAt / na ca anyA gatir asti / tat kathaM varNapaurvAparyaM
13606 vAkyaM yad apauruSeyaM sAdhyeta /
13607 	anityAvyApitAyAM ca doSaH prAg eva kIrtitaH //260//  
13608 atha mA bhUd ayaM doSa ity anityAn avyApinaz ca varNAn icchet /
13609 tAv api pakSau prAg eva nirAkRtAv ity aparihAraH / 
13610 	vyaktikramo api vAkyaM na nityavyaktinirAkRtaH / 
13611 na varNAnAM rUpAnupUrvI vAkyam / kiM tarhi / tadvyakteH / sA 
13612 yathA svavarNAbhivyaktipratyayakramAd bhavantI kramayoginI
13613 iti tadA anupUrvI vAkyam ity api mithyA / tasyA nityeSu prAg eva
13614 niSiddhatvAt / kAryatAvizeSa eva sAkSAczaktyupadhAnena jJAna
13615 janakAnAM vyaktir ity AkhyAtam etat /
13616 	vyApArAd eva tatsiddheH karaNAnAM ca kAryatA //261//
13617 yat khalu rUpaM yata eva upalabhyate tasya tadupalabdhinAntarIyakAm 
13618 upalabdhim eva Azritya lokaH kAryatAM prajJApayati / sA  
13619 varNeSv apy asti / sa eva ca anyatra api tadAzrayaH / na vizeSaH / tat  
13620 kathaM tulye abhyupagamanibandhane na varNAH kAryAH / na  
13621 ca etad upalabdhyAzrayA kAryatAsthitiH / kiM tarhi / yat saty eva
13622 bhavati iti sattAAzrayA / sA sattA kutaH siddhA ye na kAryatAM sAdhayet /
13623 na hy asiddhAyAm asyAm evaM bhavati iti / tasmAt sattA  
13624 siddhis tatsAdhanI / sA ca upalabdhir eva / satyam evaM yadi tasya 
13625 prAG na sattA siddhA syAt / sA hi sattAsiddhir yA asiddhipUrvikA /
13701 nanu tad rUpam asiddham eva yat tathAbhUtavijJAnAvyavadhAna
13702 upayogi / siddham eva tad anyavaikalyAn na upayuktam iti 
13703 cet / katham idAnIm upayuktAnupayuktayor abhedaH / na api bhedaH
13704 zabdasvabhAvAsaMsparzI tasya eva atizayasya upayogasiddhes tasyA
13705 kAraNatvaprasaGgAt / yasya eva bhAve sAdhyasiddhis tad eva hi 
13706 tatra upayogi yuktam / tadatizayaupayoge apy asya tadvat prasaGgaH / 
13707 tasmAd atizeta eva avyavahitasAmarthyaupayogo avasthA
13708 bhedas tadanyam / anatizayasya apekSA ca prAg eva niSiddhA / 
13709 sa ca karaNavyApArAd eva siddha iti sarvakAryatulyadharmA /
13710 tasya tAdRzasya vyaktau sarvaM vyaGgyaM na vA kiMcid apy 
13711 avizeSAt / tathA hi /
13712 	svajJAnena anyadhIhetuH siddhe arthe vyaJjako mataH /  
13713 	yathA dIpo anyathA vA api ko vizeSo asya kArakAt //262// 
13714 svapratipattidvAreNa anyapratipattihetur loke vyaJjakaH siddhaH
13715 dIpAdivat sa cet prAk siddhaH syAt samAnajAtIyaupAdAnakSaNasiddheH /
13716 na tasya eva atizayasya jJAnahetoH tasya tatsAmagrIpratyayatvAt /  
13717 ye punar asiddhaupalambhanAH kArakA eva kulAlAdivad  
13718 ghaTAdau / pratyabhijJAnAdayo api siddhihetavo na hetulakSaNaM
13719 puSNanti / yad api kiMcid uttarA akArapratItir akArapratIteH pUrvAbhinna
13720 viSayA tadvad ity Adi / tad api na svalakSaNayor abheda 
13721 sAdhane samartham / tatsvabhAvasiddheH / sAmAnyena vacane  
13801 bhinnaviSayatvasya apy avirodhaH / ekaviSayayoz ca pratItyoH  
13802 pUrvAparabhAvAyogAt saMnihitAsaMnihitakAraNatvena utpAdAnutpAdAt
13803 saMnidhAne apy anutpannasya atatkAraNatvAt tayor 
13804 bhinnAkhilakAraNatvam / tatra ekAbhede api zaktasya apratIkSaNAd  
13805 yuktiviruddhaM pUrvAparayoH pratItyor ekaviSayatvam / pratIti
13806 pratibhAsasvabhAvabhede api nAmasAmyAd ekaviSayatvam apy
13807 ayuktam / ghaTAdiSv api prasaGgAt / tatra dRSTavirodhAd asAdhanatvam
13808 iti cet / iha api virodhAbhAvaH kena siddhaH / yAvat tathA abhidheyatA
13809 arthabhedena vyAptA na sAdhyate tAvat sandigdho 
13810 vyatirekaH / pratikaraNabhedaM ca bhinnasvabhAvaH zabdaH
13811 zrutau nivizamAno yadA ekaH sAdhyate kiM na ghaTAdayaH / tatra
13812 api zakyam evaM vyaJjakabhedAt pratibhAsabheda iti pratyavasthAtum / 
13813 api ca /
13814 	karaNAnAM samagrANAM vyApArAd upalabdhitaH /
13815 	niyamena ca kAryatvaM vyaJjake tadasambhavAt //263//  
13816 na hi kadAcid vyApRteSu karaNeSu zabdAnupalabdhiH / na ca avazyaM  
13817 vyaJjakavyApAro artham upalambhayati / kvacit prakAze  
13818 api ghaTAdyanupalabdheH / sA iyaM niyamena upalabdhis  tadvyApArAc
13819 zabdasya tadudbhave syAt / akartur vyApAre api tatsiddhyayogAt /
13820 vyApinityatvAd upalambha iti cet / ka idAnIM ghaTAdiSu
13821 samAzvAsaH / teSAM tathA aniSTer iti cet / zabdaH kim iSTas tatsamAna
13822 dharmA / na ca asya kazcid atizaya ity uktam / pratiSiddhe ca
13823 vyApinityatve / ghaTAdInAM vyaJjakAntarasambhavAd adoSaH /  
13824 prakAzo hy eSAM vyaJjakaH siddhaH / kulAlAdInAM vyaJjakatve 
13825 tAdRzA eva syuH / atizerate ca / tato vyaJjakAtizayAt kArakA eva / 
13826 upakArakasya gatyantarAbhAvAt / tad etac zabdeSv api tulyam / 
13827 tatra api indriyayogyadezatAAdibhyaH karaNAnAm atizayAt / ghaTAdi
13828 kArakadharmasya ca karaNeSu dRSTeH / tasya eva pradIpAder  
13829 viSayAntarasya ca kasyacid vyaJjakAntarAbhAvAt / tat kAraNAni  
13830 teSAM vyaJjakAni syuH / tasmAn na vyaktiH zabdasya / bhavantI
13901 vA karaNebhyo atizayattA vA zabdasya vyaktir AvaraNavigamo  
13902 vijJAnaM vA gatyantarAbhAvAt / tatra na atizayautpattir anityatA
13903 prasaGgAt / tasyAH pUrvApararUpahAnyupajananalakSaNatvAt /
13904 atha 
13905 	tadrUpAvaraNAnAM ca vyaktis te vigamo yadi /
13906 	abhAve karaNagrAmasAmarthyaM kiM nu tad bhavet //264// 
13907 na hy AvaraNasya akiMcitkarANi karaNAni samarthAni nAma / vigamaz 
13908 ca abhAvaH / na ca abhAvaH kAryam iti niveditam etat / na api zabdasya 
13909 nityasya kiMcid AvaraNam asAmarthyAd ity apy uktam / 
13910 tasmAn na AvaraNe karaNaupakSepaH / na apy eSAm asAmarthyam / tad
13911 vyApArAbhAve zabdAnupalabdheH / ato yuktam ete yac zabdAn
13912 kuryuH / anyathA / 
13913 	zabdAvizeSAd anyeSAm api vyaktiH prasajyate /  
13914 	tathA abhyupagame sarvakAraNAnAM nirarthatA //265//
13915 yadi sarvakAraNasamAnadharmANy api karaNAni vyaJjakAni na
13916 kiMcid idAnIM kAryaM syAt / na ca etad yuktam / sarvakAraNAnAm 
13917 AnarthakyaprasaGgAt / vastuno anAdheyavizeSatvAt / AvaraNAbhAvasya
13918 akAryatvAt / vastuvad eva jJAnasya api siddhatvAt / jJAnaM 
13919 prati kArakatve kasyacit tathAbhUtAnAm anyeSAm api tathAbhAva 
13920 prasaGgena sarvasya kAryatAprasaGgAt / tasmAd ayaM 
13921 kArakAbhimato arthakalApo na vyaktau na kriyAyAm iti vyartha 
13922 eva syAt / tathA ca idam anupakAryaupakArakaM nirIhaM jagat syAt / 
13923 zabdanityatve ca / 
13924 	sAdhanaM pratyabhijJAnaM satprayogAdi yan matam /  
13925 	anudAharaNaM sarvabhAvAnAM kSaNabhaGgataH //266//  
14001 kSaNabhaGgino hi sarvabhAvA vinAzasya akAraNatvAd ity uktaM  
14002 vakSyate ca / utpattimantaz ca parataH / sattAyA AkasmikatvAyogAt / 
14003 tan na idaM pratyabhijJAnaM satprayogAdikaM kvacid anveti
14004 sthiraekarUpe / aparAparasvabhAvaparAvRttiSv eva dIpAdiSu  
14005 dRSTam iti viruddham eva / na / abhinnajanmanaH sAdharmyavipralambhAd 
14006 bhrAntyA dIpAdiSu bhAvAt / abhinnajanma iti kena avaSTambhena 
14007 ucyate / tasya eva abhedasya sarvatra paurvAparyeNa 
14008 cintyatvAt / tathA bhedasya api iti cet / tena eva saMzayo astu / na ca 
14009 saMzayitAt siddhiH / vivekAdarzanAd ekatvam iti cet / na / jJAna 
14010 paurvAparyeNa sadasattvasiddheH / svabhAvavivekasadbhAvAt /  
14011 yady aparANi jJAnAni prAk saMnihitakAraNAni pUrvajJAnavaj jAtAni
14012 eva syuH / ajAtAni tu kAraNavaikalyaM sUcayanti / samarthasya  
14013 jananAd asamarthasya api punaH sAmarthyApratilambhAt /
14014 pratilambhe vA sthairyAyogAt / tad ayam satprayoga ity api jananam
14015 eva prayoktuH sAmarthyAt / svayaM samarthe tasya anupayogAt / 
14016 prayoga ity api iSTasAdhanasamarthautpAdanam eva samAna
14017 jAtIyaupAdAnApekSam anapekSaM vA vAsyAdiprayogavat karma
14018 Adiprayogavac ca kathyate / yo api manyate samakSe pratyabhijJAnaM 
14019 pratyakSam eva tataH pratyakSAd eva sthairyasiddhir
14020 iti / tad apy uttaratra niSetsyAmaH /
14021 	dUSyaH kuhetur anyo api / 
14022 na eva kazcid dharmo yaH samAnajAtIyam anveti / sarvadharmANAm  
14023 etadavasthatvAt / sarvasthairyapratijJAyAz ca yathA abhidhAnaM
14024 yuktivirodhAd anye api nityahetavo vAcyadoSAH /
14025 	buddher apuruSAzraye /  
14026 	bAdhA abhyupetapratyakSapratItAnumitaiH samam //267//
14027 yadi vyaktir buddhis tadA AnupUrvI vAdyam / tasyA apauruSeyatva 
14101 prasAdhane buddhInAM puruSaguNatvAbhyupagamAt sama yo  
14102 asya bAdhyate / pratyakSaM khalv apy etad yad imA buddhayaH 
14103 puruSasaMkhyAtebhyaH puruSaguNebhyo vA manaskArAdibhyo
14104 bhavanti iti / na ca kAryatA nAma anyA bhAvAbhAvavizeSAbhyAm / 
14105 sa ca bhAvaH pratyakSo abhAvo apy anupalabdhilakSaNaH pratyakSa 
14106 sAmarthyasiddha iti vakSyAmaH / tata eva puruSakAryatA  
14107 buddhInAm anumeyAnvayavyatirekaliGgatvAd asyAH / kiM ca /
14108 	AnupUrvyAz ca varNebhyo bhedaH sphoTena cintitaH / 
14109 	kalpanAAropitA sA syAt kathaM vA apuruSAzrayA //268//  
14110 varNavyatirekiNy AnupUrvI sphoTavicArAnukrameNa eva prativihitA /
14111 na api sA varNasvabhAvA / vastusvabhAvasya etadvikalpa anatikramAt /
14112 atadrUpeSu tadrUpasamAropapratibhAsinyA buddher
14113 ayaM vibhramaH syAd AnupUrvI iti / sA ca katham apauruSeyI / buddhi
14114 viThapanapratyupasthApanAt / api ca / Atyantikasya kasyacit  
14115 svabhAvasya abhAvAd bhavatA dhvaninA anAtyantikena bhavitavyam / 
14116 sa ca ahetuko anyahetuko vA nityaM bhaven na ca puruSa  
14117 vyApArAt / tasmAn na apauruSeyaH / katham idaM gamyate anAtyantiko
14118 dhvanir anyo vA bhAva iti /
14119 	sattAmAtrAnubandhitvAn nAzasya anityatA dhvaneH /
14120 na hi nAzo bhAvAnAM kutazcid bhavati / tad bhAvasvabhAvo bhavet /  
14121 bhAvasya eva svahetubhyas taddharmaNo bhAvAt / na ca bhAva
14122 vizeSasvabhAvaH / tasya niSetsyamAnatvAt / tasmAd bhAva
14123 mAtrasvabhAvaH syAt / tena zabdo anyo vA sattAbhAjanaH sarva
14124 eva bhAvo anAtyantika iti siddham / na siddham / tasya eva vinAzasya
14125 aparajanmAsiddheH / tathA hy agninA kASThaM daNDena  
14201 ghaTa iti vinAzahetavo bhAvAnAM dRzyante / anvayavyatirekAnuvidhAnaM
14202 hetutadvator lakSaNam AhuH / na / pUrvasya svarasa 
14203 nirodhe anyasya viziSTapratyayAzrayeNa vikRtasya utpatteH /  
14204 astu vA agniH kASThavinAzahetuH / sa vinAzo agnijanmA kiM kASTham
14205 eva Ahosvid arthAntaram / 
14206 	agner arthAntarautpattau bhavet kASThasya darzanam //269// 
14207 	avinAzAt /  
14208 kim ity arthAntarAd arthAntarajanmani kASTham abhUtaM nAma 
14209 na dRzyate vA / atiprasaGgo hy evaM syAt /
14210 	sa eva asya vinAza iti cet /
14211 yadi sa eva artho agnijanmA abhAvas tad idam abhUtatvAn na  dRzyata  
14212 iti / bhavatu tasya idaM nAma abhAva iti / tathA api 
14213 	katham / 
14214 anyo anyasya vinAzaH / na hi kasyacid arthasya nAmakaraNamAtreNa  
14215 kASThaM na dRzyata iti yuktam / na ca anyo anyasya vinAzaH /
14216 atiprasaGgAt / vizeSAbhAvAt tasya arthAntaratvena vastubhUtasya
14217 tadanyebhyaH / kASThe agnikRtaH svabhAvo vinAzo na sarva iti
14218 cet / kASThA iti kaH sambandhaH / AzrayAzrayi sambandhaz cet /  
14219 na / tasya niSetsyamAnatvAt / janyajanakabhAvAz cet / agner iti 
14220 kim / kASThAd eva bhAvAt / tadapekSAd utpatter adoSa iti cet /
14221 anatizayalAbhinaH kA apekSA / lAbhe vA aparakASThajanma syAt /  
14222 pUrvaM tv apracyutikAraNam / tathA eva dRzyata / tata eva agneH 
14223 pUrvavinAza iti cet / pUrveNa asya kaH sambandha iti sa eva prasaGgo  
14224 aparyavasAnaz ca / tad avazyaM vinAzasambandhayogyam uttaram 
14225 atizayaM pratyupakurvANo agnir apUrvam eva janayati iti  
14226 pUrvaM tadavasthaM dRzyeta / kASThavinAza iti ca kASThAbhAva  
14301 ucyate / na ca abhAvaH kAryaH / tatkArI ca akAraka eva ity anapekSaNIya
14302 ity uktam / svabhAvAbhAvasya ca tato bhede tato nivartamAnasya 
14303 bhAvasya svabhAva eva samarthitaH syAd iti katham 
14304 abhUto nAma / tasmAn na
14305 	anyo anyasya vinAzo astu kASThaM kasmAn na dRzyate //270//  
14306 ko ayam arthAntarabhAvakASThadarzanayor virodhaH /  
14307 	tatparigrahataz cen na tena anAvaraNaM yataH /
14308 yadi tena arthAntareNa parigRhItam iti kASThaM na dRzyeta tat
14309 kASThasya AvaraNam ity Apannam / na ca etad yuktam / AvaraNaM
14310 hi darzanaM bibadhnIyAn na abhighATAdIni dravyasAmarthyAni / 
14311 sarvapratibandhe ca nanv anena eva dravyaM vinAzitaM syAt /
14312 sarvazaktipracyAvanAt / punas tatra apy agnAv iva prasaGgAd
14313 anavasthA / apracyuteSu vA asya abhighATasAmarthyAdiSu satA vA  
14314 tena anyena kiM vinAzitam / yadi ca agnisamudbhavasya vinAzAkhyasya 
14315 arthasya parigrahAt kASThaM na dRSTam /
14316 	vinAzasya vinAzitvaM syAd utpattes tataH punaH //271//
14317 	kASThasya darzanaM /  
14318 avazyaM hy utpattimatA vinAzena vinaSTavyam / tasmin vinaSTe
14319 punaH kASThAdInAm unmajjanaM syAt /  
14320 	hantRghAte caitrApunarbhavaH /  
14321 	yathA atra apy evam iti ced dhantur na amaraNatvataH //272//  
14322 vinAzavinAze api ca vastunaH pratyApattiH / na hi hantari hate api 
14323 tadvataH pratyujjIvati iti cet / na / hantus tadghAtahetutvAt /
14324 na brUmo vinAzahetor agnidaNDAder nivRttau bhAvena bhavitavyam 
14325 iti / kiM tarhi / bhAvAbhAvasya atyantAnupalabdhilakSaNasya /
14326 tannivRttau kAnyA gatiH svabhAvasthiteH / hantA hi caitrasya  
14327 na nAzakalpaH / kiM tarhi / daNDAdikalpaH / nAzakalpaM
14401 hy asya maraNaM / tannivRttau ca syAd eva asya punarbhAvaH /  
14402 	ananyatve api nAzasya syAn nAzaH kASTham eva tu / 
14403 	tasya sattvAd ahetutvaM na ato anyA vidyate gatiH //273  /  
14404 anarthAntarabhUto vinAzaH kASThAt / tad eva tad bhavati / tac
14405 ca prAg eva asti iti kim atra sAmarthyaM vahnyAdInAm / tasmAt
14406 tadanupakArAt tena na apekSyante kathaMcit / na apy asya idam iti
14407 sambandham arhati / tasya upakAranibandhanatvAt / anyathA  atiprasaGgAt /
14408 pAraMparyeNa upakAre apy avazyam ayaM vikalpo anveti  
14409 sa kim upakAro arthAntaram Ahosvit tad eva iti / tadarthAntaratve
14410 api tasya iti punar upakAratvAdiparyanuyogas tadavasthA
14411 eva / tathA ananyatve / tasmAt sato rUpasya tattvAnyattvAvyatikramAt /
14412 upakArautpAdanasya ca rUpaniSpAdanalakSaNatvAt /
14413 tadatatkriyAvikalo na kartA eva iti na kasyacid hetur ahetuz ca
14414 na apekSate / tasmAt svayam ayaM bhAvas tatsvabhAva iti siddham /  
14415 	ahetutve api nAzasya nityatvAd bhAvanAzayoH// 
14416 	sahabhAvaprasaGgaz ced asato nityatA kutaH //274// 
14417 syAd etad yasya api vinAzo ahetukaH so avazyaM nitya iti bhAvas  
14418 tadabhAvalakSaNo vinAzaz ca saha syAtAm iti / na / tasya nityAnitya
14419 dharmAyogAt / na hy asaty ayaM vikalpaH sambhavati / 
14420 tayor vastudharmatvAt / vinAzasya ca akiMcitvAt / bhavato hi 
14421 kenacit sahabhAvaH syAt / na ca vinAzo bhavati / tasmAd adoSaH /
14422 	asattve abhAvanAzitvaprasaGgo api na yujyate / 
14423 	yasmAd bhAvasya nAzena na vinAzanam iSyate //275//
14424 katham asan vinAzo bhAvaM nAzayet / ato avinAzI bhAvaH syAd  
14425 ity aprasaGga eva / vinAzAd bhAvanAzAnabhyupagamAt / yo hi  
14501 vinAza iti kiMcin na ity Aha sa kathaM tato bhAvanAzam icchet /
14502 katham idAnIm asati vinAze bhAvo naSTo nAma / na hy asadvinAzA 
14503 naSTA gaNyante pratyutpannAvasthAyAm / na hi yo yena atadvAn  
14504 sa tena tathA vyapadizyate pratIyate vA / yathA azvo viSANena /
14505 na vai vinAzo na asty eva / sa tu na asti yo bhAvasya bhavanti / bhAva 
14506 eva tu kSaNasthitidharmA vinAzaH / tam asya svabhAvam uttara  
14507 kAlaM vibhAvayanto vinAzo asya bhUta iti yathApratIti vyapadizanti 
14508 ity uktam / na hi bhAvasya kiMcit kadAcid bhavati / sa eva  
14509 kevalaM svahetubhyas tathAbhUto bhavati / tan na kenacid bhavatA
14510 sa naSTaH / kiM tarhi / svabhAva eva asya yena sa naSTo nAma / 
14511 kathaM tarhi idAnIm ahetuko vinAzo bhavati ity ucyate / 
14512 	nazyan bhAvo aparApekSa iti tajjJApanAya sA /  
14513 	avasthA ahetur uktA asyA bhedam Aropya cetasA //276//  
14514 na bhAvo jAto aparasmAn nAzaM pratilabhate tathAbhUtasya eva 
14515 svayaM jAter ity aparApekSadharmAntarapratiSedhArthaM tat 
14516 svabhAvajJApanena arthAntaram iva dharmiNo dharmaM cetasA  
14517 vibhajya tanmAtrajijJAsAyAM svabhAva eva tathA ucyate / tad etan  
14518 mandabuddhayaH kvacit tathA darzanAd ghoSamAtravipralabdhA  
14519 nAzaM guNaM tasya ca bhAvam Aropya sahetukam ahetukaM
14520 vA apratiSThitatattvayA bhAvacintayA AtmAnam Akulayanti /  
14521 	svato api bhAve abhAvasya vikalpaz ced ayaM samaH / 
14522 nanv aparabhAvitve api vinAzasya svata eva bhAvasya bhavato
14523 ayaM tattvAnyattvavikalpas tulyaH / tadA kim arthAntarabhAve 
14524 bhAvo na dRzyate anarthAntaratve api tad eva tad bhavati /  
14525 tan na kiMcid asya jAtam iti kathaM vinaSTo nAma / nanv atra / 
14526 	na tasya kiMcid bhavati na bhavaty eva kevalam //277//
14527 ity uktam / na hy ayaM vinAzo anyo vA kazcid bhAvasya bhavati ity
14528 Aha / kiM tarhi / sa eva bhAvo na bhavati iti / yadi hi kasyacid bhAvaM
14529 brUyAt na bhAvo anena nivartitaH syAt / tathA ca bhAvanivRttau  
14530 prastutAyAm aprastutam eva uktaM syAt / na hi kasyacid  
14531 bhAvena bhAvo na bhUto nAma / tadA na bhUto yadi svayaM na 
14601 bhavet / na bhavati iti ca prasajyapratiSedha eSa na paryudAsaH /  
14602 anyathA iha api kasyacid bhAve na pratiSedhaparyudAsayo rUpa  
14603 bhedaH syAd ubhayatra vidheH prAdhAnyAt / evaM ca apratiSedhAt 
14604 kasyacit paryudAso api kvacin na syAt / yadi hi kiMcit kutazcin
14605 nivarteta tadA tadvyatireki saMspRzyeta / tatparyudAsena / tac  
14606 ca na asti / sarvatra nivRttir bhavati ity ukte kasyacid bhAvasya eva 
14607 pratItiH / tathA ca anena arthAntarabhAva eva uktaH syAt / na tayoH
14608 parasparaM vivekaH / aviveke ca na paryudAsaH / tad evaM vyatireka
14609 abhAvAd anvayo api na syAt / tasya ekasvabhAvasthitilakSaNatvAt / 
14610 tatsthitiz ca tadanyavyatireke sati syAt / sa ca na asti ity  apravRtti
14611 nivRttikaM jagat syAt / tasmAd yasya nAzo bhavati ity ucyate sa  
14612 svayam eva na bhavati ity uktaM syAt / na vai ghoSasAmyAd viSaya  
14613 antaradRSTo vidhiH sarvatra yojanAm arhati / na hi gardabha iti
14614 nAmakaraNAd bAleyadharmA manuSye api saMyojyAH / tathA na
14615 caitrasya putro bhavati ity atra dRSTo vidhir nAze api virodhAt / 
14616 evaM ca abhidhAne api prayojanam Aveditam eva / ataH  
14617 	bhAve hy eSa vikalpaH syAd vidher vastvanurodhataH /  
14618 bhAvo avazyaM bhavantam apekSate / sa ca svabhAva eva / niHsvabhAvasya 
14619 kvacid vyApAre samAvezAbhAvAt / vyApAra iti hi 
14620 tathAbhUtasvabhAvautpattiH / sA niHsvabhAvasya kathaM syAt /
14621 katham idAnIM bhavaty abhAvaH zazaviSANam ity AdivyavahAraH /  
14622 na vai zazaviSANaM kiMcid bhavati ity ucyate / api tv evam  
14623 asya na bhavati iti bhAvapratiSedha eva kriyate / api ca / vyavahartAra
14624 eva etad evaM vyApAravad iva samAropya Adarzayanti  
14625 prakaraNena kenacit / na tu tat tathA / sarvArthavivecanaM hi
14626 tatra tattvaM na kasyacit samAvezaH / na khalv evaM vinAzo vastuni
14627 tadbhAvAt / asAv api yadi vaktRbhir evaM khyApyate na 
14628 tu svayaM tathA tadA na kiMcid bhavati iti iSTam eva / tasmAt svayaM
14701 bhavan svabhAvo vikalpadvayaM na ativartate tattvam anyattvam  
14702 iti / atattvam eva svabhAvasya anyattvam / na hi rUparasayor  
14703 apy anyad eva parasparam anyattvam / svabhAvApratibandho
14704 anyattvam iti cet / ko ayaM pratibandho nAma yena sa ca na
14705 syAt / na anyasvabhAvaz ca / janma iti cet / sarvakAraNAnAM parasparam
14706 avAcyatA syAt / tathA ca sarvaH sarvasya kathaMcid upayogi
14707 iti na kazcit kutazcid anyaH syAt / evaM ca avAcyatA ity api kAryakAraNabhAva
14708 eva zabdAntareNa uktaH syAt / na arthabhedaH / svabhAvAnanugamanaM 
14709 tv anyattvaM brUmaH / sa ca svabhAvavatAM parasparam 
14710 asty eva ity anyattvam eva / na ca tajjanmalakSaNAt  
14711 svabhAvapratibandhAd anyaH pratibandho nAma / anAyattasya 
14712 vyabhicArAvirodhAt / tato dharmabhedAc ca anyattvam / jJAnakRtaH  
14713 pratibandha iti cet / syAd etat / yatpratipattinAntarIyakaM 
14714 yajjJAnaM tadgatau niyamena tatpratibhAsanAt tad atadrUpam  
14715 apy avAcyam iti / na / tasya niHsvabhAvatvAt svayam / sa eva hi  
14716 tasya svabhAvo yaH pratibhAti / svabhAvavattve asya tadvat pratibhAsa  
14717 prasaGgAt / apratibhAsamAnasya ca dRzyAbhAvAt / adRzyatve 
14718 api na tadrUpaM jJAnam iti kasya kim AyattA pratipattiH / na
14719 ca yad yad AyattApratItikaM tasya svabhAvapratibhAsa eva nazyati / 
14720 prakAzAyattapratItInAm iva nIlAdInAm / kA vA tasya pratyAsattiH /
14721 tatra yat tasminn anAtmarUpe pratIyamAne sa svayaM 
14722 pratyupatiSThate / atiprasaGgo hy evaM syAt / pratIyamAnasya  
14723 tadupAdAnatA iti cet / ko ayam upAdAnArthaH / na kAryakAraNabhAvo
14724 anabhyupagamAt / abhyupagame vA na kAryakAraNe anyonya  
14725 pratItipratyupasthApane / pratIter eva tannAntarIyakatA  
14726 pratyAsattir iti cet / nanu sa eva asati pratibandhe na yuktA ity ucyate /
14801 na akAryakAraNayoH kazcit pratibandha iti ca uktam / yat pratipatti  
14802 nAntarIyakaM yajjJAnam ity api tajjJAne sati syAt /  na hi yo
14803 vijJAne svarUpeNa asvarUpAsaMsargeNa na pratibhAsate tasya kiMcij
14804 jJAnam / tadabhAvAn na sidhyati avAcyatAlakSaNam artha
14805 rUpasya / tad bhavatA vastutas tattvAnyattvabhAjA bhavitavyam /  
14806 yasya tu vinazyato bhAvasya na kiMcid bhavati / tena  
14807 	na bhAvo bhavati ity uktam abhAvo bhavati ity api //278// 
14808 yad apy ayaM bhAvasya abhAvo bhavati ity Aha / tad api bhAvo 
14809 na bhavati ity eva uktaM bhavati / evaM hi sa nivartito bhavati /
14810 pratiSedhe vidher asambhavAt / tata eva asya vinAze na kazcid hetuH /
14811 tathA hi /
14812 	apekSyeta paraH kAryaM yadi vidyata kiMcana /  
14813 	yad akiMcitkaraM vastu kiM kenacid apekSyate //279//  
14814 sati hi kArye kAraNaM bhavati / na ca nazyato bhAvasya kiMcit 
14815 kAryam ity uktam / tasmAd yo nAma nAzahetuH sa bhAve na 
14816 kiMcit karoti ity akiMcitkaro na apekSyaNIyaH / tat kathaM idAnIm
14817 anutpannAtizayas tadavastha eva bhAvo naSTo nAma / 
14818 nanv atizayautpattAv api sa eva tasya atizaya utpanna iti kathaM  
14819 sa naSTo nAma / tena na ayaM tadavastho naSTo nAma / yena svayaM 
14820 na bhavati tena naSTaH / na arthAntarautpAdAd ity uktam / na hy
14821 atizayautpattyA svayaM na bhUto nAma / abhAvasya sarvAtizayaupAkhyA
14822 nivRttyA sarvabhAvadharmavivekalakSaNatvAt / bhAvasya 
14823 ca utpattisamAvezalakSaNatvAt / tasmAn na abhAve kasyacid bhAvaupakSepo
14824 anyasya /  
14825 	etena ahetukatve api hy abhUtvA nAzabhAvataH /
14826 	sattAnAzitvadoSasya pratyAkhyAtaM prasaJjanam //280// 
14827 yo api manyate ahetuke api vinAze abhUtvA asya bhAvAt sattAnityatvaM 
14901 ca durnivAram / abhUtvA bhavann ahetuko bhavati ity api 
14902 viruddham iti / so apy anena eva pratyAkhyAtaH / kasyacid bhAvAn abhyupagamAt /
14904 	yathA keSAMcid eva iSTaH pratigho janminAM tathA /
14905 	nAzaH svabhAvo bhAvAnAM na anutpattimatAM yadi //281//
14906 atha api syAd / bhavatu nAma svabhAva eSa bhAvAnAM ya ime kSaNa 
14907 sthitidharmANAH / sa tu utpattimatAm eva bhaviSyati / na hi svabhAva 
14908 iti sarvaH sarvasya svabhAvo bhavati pratighAtmatAvat / 
14909 satyam etat / tathA api /
14910 	svabhAvaniyamAd hetoH svabhAvaniyamaH phale / 
14911 	na anitye rUpabhedo asti bhedakAnAm abhAvataH //282//  
14912 na vai pratigho anyo vA svabhAvo akasmAt pratiniyamavAn / yAdRzI 
14913 tu svahetoH zaktisthitis tAdRzaM phalaM bhavati iti  hetusvabhAva
14914 niyamAt phalasvabhAvaniyamaH / Akasmikatve apy asya ukto 
14915 doSaH / pratighAtAtmatAhetusvabhAvapratiniyamavan na nazvara 
14916 jananapratiniyatasvabhAvaM bhAvaM pazyAmaH / yena tajjanmA
14917 tathA syAn na anyaH / sarvAkArajanmanAM vinAzadarzanAt / 
14918 nanv idam apy anizceyam eva sarvAkArajanmAno nazyanti iti /
14919 tAsAm aniHzeSadarzanAt / vicitrazaktayo hi sAmagryo dRzyante /
14920 tatra kAcit syAd api yA anazvarAtmAnaM janayet / na / jJeyAdhikArAt /
14921 ye kadAcit kvacit kenacij jJAtAH santo na jJAyante  
14922 teSAM sattAnubandhI nAza iti brUmaH / ta eva kRtakA anityAH
14923 sAdhyante / na hy ayaM sambhavo asti yat te jJAnajanana
14924 svabhAvAH punar anaSTA na janayeyur apekSeran vA param /  
14925 tajjananasvabhAvasya niSpatteH / na ca teSv anapekSeSu kasyacit  
14926 kadAcit kiMcij jJAnaM nivarteta / na ca evaMbhUtaM kiMcid asti / 
14927 sarvasya kenacit kadAcit jJAnAt / jJAnamAtrArthakriyAyAm apy 
15001 asAmarthye vastv eva na syAt / tathA hi tallakSaNaM vastv iti
15002 vakSyAmaH / tasya ca vinAzAvyabhicArAt sa sattAnubandhI / 
15003 	pratyAkhyeyA ata eva eSAM sambandhasya api nityatA / 
15004 ata eva yathA uktAd vastumAtrAnubandhAd vinAzasya zabdavat 
15005 sambandhanityatA api pratyAkhyeyA / yA ca zabdazaktir yogyatAAkhyA
15006 arthapratipattyA Azrayo jaiminIyair varNyate sa arthAntaram
15007 eva na bhavati / tathA hi / yogyatA iti rUpAtizaya eva bhAvAnAm ity
15008 AveditaM prAk / astu vA arthAntaram / tathA api
15009 	sambandhadoSaiH prAg uktaiH zabdazaktiz ca dUSitA //283// 
15010 ukto hi sambandhArthAntaravAde anekaprakAro doSaH / tena eva
15011 sA zabdazaktir dUSitA iti na punar ucyate / api ca /
15012 	na apauruSeyam ity eva yathA arthajJAnasAdhanam /
15013 	dRSTo anyathA api vahnyAdir aduSTaH puruSAgasA //284//  
15014 bhavantu nAma apauruSeyA vaidikAH zabdAH tathA api sambhAvyam 
15015 eva eSAm ayathArthajJAnahetutvam / na hi puruSadoSaupadhAnAd  
15016 eva artheSu jJAnavibhramaH / tadrahitAnAm api pradIpAdInAM nIlautpala
15017 AdiSu vitathajJAnajananAt / tad ime zabdAH saMskAranirapekSAH 
15018 prakRtyA ca artheSu pratibhAnahetavaH syuH / svabhAvavizeSAd  
15019 vahnyAdivat / vitathavyaktayaz ca niyamena eva niyama 
15020 kAraNAbhAvAd ayuktam iti cet / avitathavyaktiniyame kiM
15021 kAraNam / tasmAd yathArthavyaktiniyamavat prakRtyA ayathArtha 
15022 vyaktiniyamaH kiM na kalpyate / athavA vahnyAdivad eva artheSu 
15023 ubhayajJAnahetutvaM syAt / na hy apauruSeyA api vahnyAdaya
15024 ekatra yathArthajJAnahetavo api sarvatra tathA bhavanti iti /  
15025 tathA zabdAnAm apy apauruSeyatve apy ubhayaM syAd iti / bhavatu 
15026 vahnyAdInAm kRtakatvAd yathApratyayam anyatra anyathAtvam
15027 na punar nityeSu zabdeSv etad asti / nanv evaMvidho anyatra
15028 apy asty eva dharmas teSAm api saMketabalAd anyathAvRtteH
15029 kAryajananasvabhAvasthitau ca eSAM samayAder apekSaNIyasya abhAvAt /
15101 tataH pratItir artheSu sarvasya sarvadA syAt / na ca asti  
15102 tasmAn na zabdAH sthitasvabhAvA iti / api ca /  
15103 	na jJAnahetutA eva syAt tasminn akRtake mate / 
15104 	nityebhyo avastusAmarthyAn na hi janma asti kasyacit //285//
15105 yady akRtakaH zabdo na tasmAd artheSu pratItir eva syAt / pratIti 
15106 janmaitarakAlayos tulyarUpasya pratItijanmani sAmarthya 
15107 sambhAvanAyogAt / evam ayaM janako na evam iti vivecanIyasya  
15108 rUpabhedasya abhAvAt / na yAdRzo asya ajanakas tAdRza eva janako 
15109 yuktaH / anyApekSA api niSiddhA eva / tasmAn na nityAnAM kvacid 
15110 vijJAnajananasAmarthyam / kadAcid ajanane nityam ajanana prasaGgAt /
15111 kAryasAtatyAdarzanAc ca na te kathaMcit kartAra ity 
15112 uktaprAyam / yA apy etA nityAbhimateSv AkAzAdiSu pratipattayas 
15113 tA api na tatsvabhAvabhAvinyaH / tathA hi / 
15114 	vikalpavAsanAudbhUtAH samAropitagocarAH / 
15115 	jAyante buddhayas tatra kevalaM na arthagocarAH //286//  
15116 svalakSaNaviSayA hi buddhir niyamena tadyogyatAupasthApanAnuvidhAyinI
15117 iti tasmin saty asyAH kAraNe yogye sA bhavaty eva /  
15118 tad yadi nityAnAM padArthAnAM svalakSaNe kasyacit jJAnaM syAt
15119 sarvasya sarvadA syAt / kAryavizeSA hi vyaktayaH kathaMcit  
15120 kvacid upayujyamAnAs tadupajananayogyAtizayapratilambhahetum
15121 vastuvizeSam apekSanta iti yuktam / tathA akAryavizeSo nityo 
15122 bhAvaH kenacid gRhyamANas tatkAraNApekSo yadi grahaNam asya
15123 janayet yuktaM yat tena eva gRhyeta / tac ca sthitasvabhAvatvAn  
15124 na sambhavati iti / sarveNa samaM gRhyeta na vA kenacid iti  
15125 san kenacid dRSTo na nityaH kazcid atIndriyaH syAt / na ca etad asti /  
15126 tasmAd arthasAmarthyAnapekSAH samAropitagocarA Antaram eva upAdAnaM
15127 vikalpavAsanAprabodham Azritya bAhyArthazUnyA bhrAntaya 
15128 eva AkAzAdiSu sarvasya sambhavanti iti / tasmAn na aparAvRtti  
15201 dharmAH zabdAH / tattve vA kuta etad avitathArthapratItaya
15202 eva iti / na hy agnir himasya bheSajam ity AdiSu zItapratighAta
15203 sAmarthyaM lokaprasiddhaM agneH khyApyata iti sarvaM tathA
15204 bhavati / lokasya svaicchAkRtasaMketAnuvyavahArAt kim ayaM  
15205 lokaH svasaMketam anuvidadhad evaM pratyeti Ahosvic zabda  
15206 svabhAvasthiter iti sandehaH / lokaicchayA api parAvartyamAnAH  
15207 zabdAH punar anyatra anyathA dRzyanta iti lokaprasiddhyanuvidhAne 
15208 api sambhavaty eSAm anyathAbhAvaH / tasmAt kasyacid avaiparItya  
15209 darzane api sarveSAM tathAbhAvo na sidhyati / akRtakasvabhAvatve  
15210 api hy eSAM kazcin mithyArthaniyato api syAd iti svabhAva
15211 aparijJAnAt sarvatra saMzayaH syAt /  
15212 	mithyAtvaM kRtakeSv eva dRSTam ity akRtaM vacaH / 
15213 	satyArthaM vyatirekasya virodhivyApanAd yadi //287//
15214 atha yan mithyArthaM vacanaM tad akhilaM kRtakam iti hetuvyatirekeNa  
15215 sAdhyavyatirekasya vyApter anyatra asambhavAd akRtakaM 
15216 satyArtham iti syAd vinA apy anvayena / yo hi yena avyAptas 
15217 tatra tadvyatirekaH zaGkyeta / na ca viruddhayor ekatra sambhavo
15218 asti / asambhave ca vijAtIyasya gatyantarAbhAvAd iSTArthasiddhes
15219 tatsAdhanatvAc ca liGgasya vyartham anvayadarzanaM  
15220 vyatirekamAtreNa eva siddher iti / satyam etat yadi vipakSayor 
15221 vyApyavyApakabhAvaH sidhyet / sa tu na siddhaH / yasmAt
15222 	hetAv asambhave anukte bhAvas tasya api zaGkyate /  
15223 	viruddhAnAM padArthAnAm api vyApakadarzanAt //288// 
15224 yadi hetoH sAdhyavipakSe abhAvaH sidhyet tadA sAdhyavyatirekaM
15225 hetuvyatireko vyApnuyAt / na ca tatra tasya asambhave
15226 pramANaM pazyAmaH / na ca aviruddhavidhiH pratiSedhasAdhano 
15227 yuktaH / atiprasaGgAt / na ca ekatra dRSTasya anyatra asambhava 
15228 eva pRthag viruddhasahabhAvinAm api darzanAt / anityatvavat 
15229 prayatnAnantarIyakaitarayor / na ca tathAvidhasya adarzanAd asattvam 
15230 eva / yasmAt /  
15301 	na asattAsiddhir ity uktaM sarvato anupalambhanAt /
15302 	asiddhAyAm asattAyAM sandigdhA vyatirekitA //289// 
15303 na hy ayaM puruSamAtrakaH sarvaM draSTuM samartho yena asya 
15304 darzananivRttyA na tathA syAt / yasya hi jJAnaM jJeyasattAM  
15305 na vyabhicarati sa evaM bruvANaH zobheta adarzanAn na asti iti /  
15306 tad ime svabhAvadezakAlaviprakarSeNa santo apy anupalakSyAH
15307 syuH / tathA hi / ko atyantaparokSe arthe saMvAdanam itarad 
15308 vA sarvadarzI vacanasya akRtakasya itarasya vA vibhAvayituM samarthaH /
15309 pratipAditaM ca etat kvacit tathA dRSTAnAm apy arthAnAM 
15310 punaH kathaMcid anyathAbhAvaH yathA kvacid deze madhurANi 
15311 nimbaphalAni saMskAravizeSAd AmalakIphalAni ca na 
15312 ca idAnIm ataddarzinA tAni pratikSeptavyAny eva iti / tasmAd akRtakaM
15313 ca syAd syAn mithyArthaM ca iti na virodhaM pazyAmaH / na  
15314 hi iyam anupalabdhir adRzyAtmasv abhAvasAdhikA ity uktam / tena 
15315 yat kiMcin mithyArthaM tat sarvaM pauruSeyam ity anizcayAd
15316 avyAptiH / tathA hi /
15317 	anvayo vyatireko vA sattvaM vA sAdhyadharmiNi /
15318 	tannizcayaphalair jJAnaiH sidhyanti yadi sAdhanam //290//  
15319 tathA ca uktam / ya eva tu ubhayanizcitavAcI sa sAdhanam dUSaNam
15320 vA na anyataraprasiddhasandigdhavAcI punaH sAdhanApekSatvAd  
15321 iti / ko hy adRSTavirodhasya sambhavaM pratyAcakSIta / tad ayaM  
15322 vyatirekaH saMzayAd asAdhanam / api ca /  
15323 	yatra sAdhyavipakSasya varNyate vyatirekitA /
15324 	sa eva asya sapakSaH syAt sarvo hetur ato anvayI //291//  
15401 yat kiMcin mithyArthaM tat sarvaM pauruSeyam iti hetuvipakSeNa  
15402 sAdhyavipakSasya vyAptis tadabhAve abhAvasiddhau syAt / tadabhAve
15403 bhavatas tena vyAptyayogAt / yA eva ca vijAtIyayor vyAvRttisiddhiH
15404 sA eva anvayasthitir vipakSavyavacchedalakSaNatvAt sAdhyasya /
15405 pratiSedhadvayAc ca vidhisiddhir iti kA ananvayAvyatirekavyApti  
15406 siddhiH / tan na kazcid hetur ananvayo nAma / ekavyavacchedasya
15407 tadvijAtIyasiddhinAntarIyakatvAt / anityanirAtmatAAdivyavacchede
15408 api tathA syAt / na / vyatirekavyavacchedasya bhAvarUpatvAt /
15409 na bhAvarUpavyavacchede bhAvAnuSaGgaH / abhAvavyavacchedas 
15410 tu niyamena bhAvaupasthApanaH / bhAvAbhAvayor anyonyaviveka
15411 rUpatvAt / abhAvarUpas tu vyatirekaH sa ca vyatiricyamAno  
15412 bhAvam upasthApayati na evaM nairAtmyAdayaH svabhAvavizeSAt / 
15413 kriyAbhogAdhiSThAnAsvatantro hy AtmA nirAtmA /  tatsvAtantrya
15414 lakSaNatvAd AtmanaH / tad rUpaM nairAtmyaM na AtmanivRtti 
15415 mAtram anyathA nirupAkhye kRtakatvAdyayogAn na tato nairAtmya
15416 siddhiH syAt / Atmavyavacchedena nirAtmano bhAvasya parAmarzAd  
15417 adoSa iti cet / paryudAsena vastusaMsparzAd tad eva  
15418 vasturUpaM nairAtmyam AyAtam / yasya api na abhAvarUpo vyatirekas 
15419 tasya bhAvarUpavyavacchede na bhAvasiddhiH syAd iti 
15420 na anvayAnuSaGgaH / tathA anairAtmye api na bhAvasiddhiH / yathA
15421 na idaM nirAtmakaM jIvaccharIraM prANAdimattvAd iti vipakSayor 
15422 vyApyavyApakabhAvacintAyAm aprANAdimattva eva nairAtmyaM 
15423 dRSTam / tadabhAve ca na asti iti svayaM na bhavad api prANAdInAM 
15424 na Atmani siddhim upasthApayati / tathA sAdhye api prANAdibhir 
15425 vyudasyamAnaM syAt kevalam / nairAtmye abhAvAt prANa
15426 Adayas tannirasanA na AtmaupasthApanAH / tatra bhAvAsiddheH /
15427 na ca nairAtmyanivRttyA AtmasiddhiH syAt / vipakSavyatirekadarzane
15428 api prasaGgAt / tan na viparyayavyAptir vyatirekAsiddhau /  
15429 tatsiddhir eva ca anvayasiddhiH / asiddhau vA tadvyatirekanivRtti
15430 siddhAv api tadasiddhir iti sAdhye api prasaGgaH / tan na  ananvayA
15431 vyatirekavyAptiH / mithyArthatAyAs tu pauruSeyatvena vyAptyA  
15501 apauruSeyAn nivRttAv api na satyArthatvam / prakArAntarasambhavAt /
15502 dvairAzye tu zabdAnAm etat syAd ekanivRttau gatyantara 
15503 abhAvAt / te tv anarthakA api syur iti na iSTasiddhiH /  arthapratIter
15504 na anarthakA iti cet / eSa puruSavyApAraH syAt / arthAntaravikalpavat /
15505 yathA atadarthatve api bharataurvazyAdicaritAdikam artham  
15506 anye anyathA vyAcakSate / tadanusAreNa ca keSAMcit pratItiH / 
15507 tathA ayam anarthakeSv arthavikalpaH puruSakRtaH na tu zabda 
15508 svabhAvakRtaH puruSaupadezApekSaNAt arthAntaravad eva / 
15509 na hi prakRtyA prakAzanAs tam apekSante vahnyAdayaH / puruSas  
15510 tu svasamayavyApAram AcakSANa upadizati iti nyAyyam /  
15511 puruSasamitanisargasiddhayor upadezApekSaNAvizeSAd anyavizeSAbhAvAc
15512 ca eko naisargiko anyas tu pauruSeya iti duravasAnam / asti  
15513 vizeSaH pramANasaMvAda iti cet / etad uttaratra niSetsyAmaH na asty
15514 atyantaparokSe arthe pramANAntaravRttir iti / samAnadharmiNor 
15515 arthayoH pramANasaMvAdamAtravizeSAd ekatra apauruSeyatve bahutaram 
15516 idAnIm apauruSeyam / santi puruSakRtAny api vAkyAni kAnicid  
15517 evaMvidhAni iti teSv api prasaGgaH / tadvad eSAm apy abhimata 
15518 arthavattA pauruSeyI ca syAt pramANAnurodhinI ca / api 
15519 ca idaM mantrA apauruSeyAz ca iti vyAhataM pazyAmaH / tathA hi /  
15520 	samayatve hi mantrANAM kasyacit kAryasAdhanam /  
15521 	yuktam / yady ete mantrAH kasyacit samayo yathA matpraNItam / 
15522 etad abhimatArthaupanibandhanaM vAkyam evaM niyuJjAnam
15523 anena arthena yojayAmi iti parArthaparatAnurodhena anyato vA kutazcid
15524 hetoH kRtaH syAt tadA mantraprayogAt kadAcit artha
15525 niSpattir yuktA kavisamayAd iva pAThakAnAm /
15526 	atha api bhAvazaktiH syAd anyatra apy avizeSataH //292// 
15527 na vai puruSasamayAn mantrebhyo arthasiddhiH / kiM tarhi / bhAva  
15528 svabhAva eSa yad ime kathaMcin niyuktAH phaladAH / tat tarhi  
15601 rUpaM varNAnAM sarvatra aviziSTam iti yathAkathaMcit prayuktAd 
15602 api phalaM syAt / varNA eva hi mantro na anyat kiMcit / tatkramo
15603 mantra iti cet / 
15604 	kramasya arthAntaratvaM ca pUrvam eva nirAkRtam /  
15605 na varNavyatirikto anyaH krama iti niveditam etat / avyatireke  
15606 ca varNA eva mantras / te ca aviziSTAH sarvatra iti sarvadA phaladAH  
15607 syuH / upaplavas tv alpIyaso api kramasya bhraMzAd dRSTaH / 
15608 kasyacid anuSThAnAd devatAsaMnidher asAkalyena virAdhanAc 
15609 ca / sarvabhraMze tu kasyacid eva samayasya ananuSThAnAd asaMnidher 
15610 na arthAnarthau / kiM ca kramasya arthAntaratve anarthAntaratve 
15611 vA varNAtmanas tatkramAtmano vA mantrasya arthahetor
15612 akRtakatvAn nityasya nityaM saMnidhAnam iti  
15613 	nityaM tadarthasiddhiH syAd /
15614 yato hi bhAvazakteH phalautpattiH sA avikalA iti na phalavaikalyaM  
15615 syAt / na hi kAraNasAkalye kAryavaikalyaM yuktam / tasya
15616 akAraNatvaprasaGgAt / na kevalAn mantraprayogAd iSTasiddhis
15617 tasya vidhAnApekSatvAd iti cet /  
15618 	asAmarthyam apekSaNe //293// 
15619 yadi mantrA vidhAnAd anyato vA kaMcit svabhAvAtizayam AsAdayeyuH 
15620 sa tatra samartho apekSyaH syAt / na ca nityeSv etad
15621 asti ity uktam / tat kim ayam asamartho apekSyata ity anapekSAH 
15622 sadA kuryur na vA kadAcid anatizayAt /  
15623 	sarvasya sAdhanaM te syur bhAvazaktir yad IdRzI /
15624 	prayoktRbhedApekSA ca na asaMskAryasya yujyate //294// 
15625 yadi bhAvazaktayA eva mantrAH siddhipradA na te kaMcit parihareyur  
15626 yajamAnam anyaM vA / na hy anyaM prati svabhAvo atadbhAvo
15627 bhavati / tasya tena anapakarSaNAd anyena ca anutkarSaNAt / 
15701 kenacit saha kAryakAraNabhAvAyogAt / pratyAsattiviprakarSa 
15702 abhAvAt / ata eva asya asaMskAryatvAt prayoktA api na asti / ataH 
15703 prayoktA phalam aznuvIta / 
15704 	saMskAryasya api bhAvasya vastubhedo hi bhedakaH / 
15705 	prayoktRbhedAn niyamaH zaktau na samaye bhavet //295// 
15706 AdheyavizeSA hy anityA bhAvAH / tadhetoH svabhAvabhede tataH
15707 samAsAditAtizayatvAd anyatra anyathA syuH / na abhede / kAraNAvizeSe
15708 kAryAvizeSAt / vizeSe tasya ahetukatvaprasaGgAd ity ukta  
15709 prAyam / tad ime mantrAH svabhAvAtizayAt phaladAyinaH kAryA
15710 api na zUdrAdiprayoge apy anyathA syuH / zUdraviprAbhidhAnayoH
15711 puruSayoH svabhAvAbhedAt / na hi puruSaicchAnuvidhAyino 
15712 nAmavyavahArabhedAt svabhAvabhedAnubandhinAm arthAnAm anyathAtvam
15713 asti / tayor jAtibheda iti cet / sa khalv AkRtiguNazakti
15714 bhede dRSTo gavAzvavat / anupadezaM ca enaM lokaH pratipadyate / 
15715 na tadvad anayoH kaMcid api guNaM viniyataM pazyAmaH / apazyantaz  
15716 ca kathaM bhedaM pratipadyemahi / yo apy ayaM nAma 
15717 bhedAnvayo loke pratItibhedaH so asaty api jAtibhede vyApAra  
15718 vizeSAnuSThAnAd anvayAc ca syAd vaidyavaNigvyapadezAdivat /
15719 tad ime aviziSTena prayujyamAnA mantrAs tato aviziSTam eva 
15720 svabhAvam AsAdayanti / tena avizeSeNa eva phaladAH syuH / yadA 
15721 tu samayAd ebhyaH phalaM tadA ayam adoSaH / samayakArasya  
15722 ruceH phalautpattiniyamAt / svabhAvavRttayo hi bhAvAs tanmukhena  
15723 prasaGgam arhanti / na puruSaicchAvRttayas teSAM 
15724 yathAkathaMcid vRtteH / yad api prayoktA phalam aznuta iti  
15725 prayogaM samIhitArthayogyasya utpAdanaM santAnapariNAmanaM
15726 vA pazyAmaH / tad ubhayaM vizeSajanmani syAt / anyathA / 
15727 	anAdheyavizeSANAM kiM kurvANaH prayojakaH / 
15728 yena tataH kazcit phalam aznute anyo na /
15729 	prayogo yady abhivyaktiH sA prAg eva nirAkRtA //296//
15801 na hi nityAnAM kAcid abhivyaktir ity uktam yato abhivyaJjakaH 
15802 prayoktA syAt /
15803 	vyaktiz ca buddhiH sA yasmAt sa phalair yadi yujyate / 
15804 	syAc zrotaH phalasambandho vaktA hi vyaktikAraNam //297// 
15806 na hi zabdasya anyataH svarUpapariNAmo vyaktir na apy AvaraNavigamanam /
15807 kiM tu tadviSayA pratItir azrUyamANe avyaktavyapadezAt /  
15808 tatra yadi buddhihetur vaktA syAt tat tulyaM zrotary api iti  
15809 so api phalaM vaktRvad aznuvIta / na hi vaktuH kazcid anyas 
15810 tadbhAvo anyatra tadbuddhihetutvAt / paraupAdhibuddhiH  zrotur 
15811 na vaktur iti vizeSa iti cet / kaH punar upayogo vaktuH zrotari  
15812 yena upAdhir iSyate / tataH zabdazrutir iti cet / nanu tad eva idaM  
15813 paryanuyujyate kathaM tata iti / asambandhAt / viSayaupanayanAd  
15814 ayam asya zrAvakaH syAt / tac ca na zakyam tasya kathaMcid  
15815 apy apariNAmAt / indriyasaMskArAdayo apy uktAH / mAM 
15816 zrAvayaty ahaM zrAvayAmi iti tayoH pratyayAd vaktRzrotRbheda  
15817 iti cet / anupakAryaupakArakAd bhrAntimAtrAt tadbhAve atiprasaGgAt /
15818 anyatra api bhrAntyA pratyayadarzanAt / sarvathA upakArAbhAve 
15819 ca tathA pratyayo na yuktaH / sarveSAM parasparam evaM
15820 prasaGgAt / bhrAntir api kutazcid upakAre sati kayAcit  pratyAsattyA
15821 anyatra bhavati / sa apy atyantAnupakAre na syAt / tasmAd
15822 vaktRzrotror vyaktihetutve avizeSAt tulyaH phalasambandhaH
15823 syAt / api ca / 
15824 	anabhivyaktazabdAnAM karaNAnAM prayojanam /
15825 	manojapo vA vyarthaH syAc zabdo hi zrotragocaraH //298//
15826 zrotragrahaNalakSaNaH zabdaH / tadatikrame atiprasaGgAt / nanv  
15827 evaM sAmAnye api prasaGgaH / na brUmaH zabda eva iti / zabdas tv 
15828 avazyaM tallakSaNas tasya lakSaNAntarAbhAvAt /  tatra yadi
15829 zabdAtmanAM mantrANAM vyaktihetuH prayoktA anabhivyakta 
15901 zrutiviSayANAM karaNAnAM prayoktA jApI na mantraphalena 
15902 yujyate na api manasA japan / na hi tadA zrotreNa kaMcid arthaM
15903 vibhAvayAmaH / na ca azabdAtmA mantraH / 
15904 	pAraMparyeNa tajjatvAt tadvyaktiH sA api cen matiH / 
15905 na hi manasA dhyAyato api mantrAbhAsA buddhiH zabdazravaNAd
15906 Rte / tataH zabdaprabhavAt sA api zabdavyaktir eva / anavasthA evaM  
15907 syAt / zabdArthavikalpAnAm api paraMparayA prasUtir asti iti / 
15908 	te / 
15909 api  
15910 	tathA syus tadarthA ced asiddhaM kalpanAnvayAt //299//  
15911 na brUmaH sarvA zabdaprabhavA buddhis tadvyaktir iti / yA tu  
15912 tadviSayA sA tasya vyaktir iti / manovikalpasya tadviSayatvam  
15913 asiddham / na hi svalakSaNe vikalpAnAM vRttir iti nivedayiSyAmaH /  
15914 te hi yathAsvam AntarAd vikalpavAsanAprabodhAd anapekSita  
15915 bAhyArthaupanidhayo bhavanti / bAhyApAyAnAgame api bhAvAt / 
15916 na hi yo yasya sattAupadhAnaM na apekSate sa tasya hetuH / ahetuz
15917 ca kathaM viSayaH / tasmAn na manovikalpaH zabdavyaktir 
15918 yatas tadvAn prayoktA syAt / tatprasUtA tadviSayA buddhis tadvyaktiH /
15919 tadAzrayaH prayoktA iti / atra apy uktaM zrotary api prasaGga 
15920 iti / tajjJAne ca prayoge zabdaH puruSe vyApriyate / tasya jJAna  
15921 jananAt / na puruSaH zabde / tadAtmany anupakArAt / atha ca  
15922 puruSaH zabdAnAM prayoktA ity alaukiko ayaM vyavahAraH /  
15923 sarvathA zabdasvabhAvAnAM mantrANAM prayogAt phalAvAptau 
15924 vyartho manojapo vikalpasya zabdarUpAsaMsparzAt / 
15925 	svasAmAnyasvabhAvAnAm ekabhAvavivakSayA /  
15926 	ukteH samayakArANAm avirodho na vastuni //300// 
15927 samayakAras tu svalakSaNam indriyaviSayaM sAmAnyalakSaNaM 
15928 ca vikalpapratibhAsaM yathAvyavahAraM saMvRtyA saMkalayya  
15929 samayam Arocayet yathAsamayaM ca arthaM niSpAdayed iti 
16001 na manojapAdAu doSaH / vastusvabhAvAt tu phalAv AptAv atat 
16002 svabhAvasaMsparze na syAt / yad uktaM na varNebhyo anyA kAcid  
16003 AnupUrvI iti tatra /
16004 	AnupUrvyAm asatyAM syAt saro rasa iti zrutau /
16005 	na kAryabheda iti ced / 
16006 na hi saro rasa ity AdipAdeSu kazcid varNabhedo na ca varNavyatiriktam
16007 anyad yataH kAryabhedaH syAt / bhinnAM ca tayoH  
16008 pratibhAM pazyAmaH AnupUrvIm eva ca atulyAm / na ca kAraNAbhede  
16009 kAryabhedo yuktaH / tasmAd asti sA bhedavatI yato ayam
16010 pratItibhedaH / satyam  
16011 	asti sA puruSAzrayA //301// 
16012 tathA hi /
16013 	yo yadvarNasamutthAnajJAnajAj jJAnato dhvaniH /
16014 	jAyate tadupAdhiH sa zrutyA samavasIyate //302// 
16015 	tajjJAnajanitajJAnaH sa zrutAv apaTuzrutiH /
16016 	apekSya tatsmRtiM pazcAd Adhatte smRtim Atmani //303//  
16017 	ity eSA pauruSeyy eva tadhetugrAhicetasAm /  
16018 	kAryakAraNatA varNeSv AnupUrvI iti kathyat //304// 
16019 cittasamutthAnA hi vAgvijJaptir varNapadavAkyAbhidhAnA / tatra
16020 sakArasamutthApanacetasA samanantarapratyayena akarautthApana  
16021 cittam utthApyate / tathA rephAkAravisarjanIyautthApanAni pUrva
16022 pUrvapratyayAni / tad ime anyAnyahetavo varNAH svakAraNa
16101 AnupUrvIjanmAnaH / zrutikAle api yadA mandacAriNaH pUrva
16102 varNajJAnasahakAripratyayApekSAH svajJAnaM janayanti tadA
16103 pUrvavarNasmaraNApekSA eva smRtim upalIyante / sa eSa varNAnAM
16104 bhinnakAryakAraNabhAvapratyayanirvRttidharmA bhinna  
16105 nirvartanadharmA ca svabhAvaH puruSasaMskArabhedabhinnaH  
16106 krama ity ucyate /  
16107 	anyad eva tato rUpaM tad varNAnAM padaM padam / 
16108 	kartRsaMskArato bhinnaM sahitaM kAryabhedakRt //305//  
16109 tasmAn na khalv eka eva padeSu varNAnAM svabhAvaH kartRcitta  
16110 saMskArabhedena bhedAt / sa ca parasparasahitaH kAryabheda
16111 hetuH / 
16112 	sA ca AnupUrvI varNAnAM pravRttA racanAkRtaH / 
16113 	icchA aviruddhasiddhInAM sthitakramavirodhataH //306//
16114 kAryakAraNabhUtapratyayautpannasvabhAvavizeSo varNAnAm AnupUrvI
16115 ity uktam / sA ca puruSavitarkavicArakRtA iti na sthitakramA
16116 varNAH / icchA aviruddhasiddhikramatvAt / kramavizeSAnukramavat /
16117 na hi sthitakramANAM dezakAlayor himavadvindhyamalaya 
16118 AdInAM bIjAGkurAdInAM ca svaicchayA kramaracanA zakyate 
16119 kartum / tata eva puruSadharmasaMkhyAte vikalpAnukrame sati 
16120 bhAvAd asati ca abhAvAt /
16121 	kAryakAraNatAsiddheH puMsAM varNakramasya ca /  
16122 	sarvo varNakramaH pumbhyo dahanaindhanayuktivat //307// 
16123 sati indhane dAhavRtter asaty abhAvAd adRSTaindhano api dahano 
16124 na anindhanas tasya dezakAlaniyamAyogAt / niyame ca tasya  eva indhanatvAd
16125 dahanaupAdAnalakSaNatvAd indhanasya / tathA ayam  
16201 api varNAnukramaH puruSavikalpaM yadi na apekSeta nirAlambanaH  
16202 svayaM prakAzeta / yatne api na zakyeta / atatprabhavAt /  
16203 kvacic zaktau sarvas tathA syAt / vizeSAbhAvAt / tadbhAvabhAvino
16204 atadviziSTasya ca atatkRtau sarvatra kAryakAraNabhAvaz ca
16205 nirAkRtaH syAt / anvayavyatirekalakSaNatvAt tasya /  lakSaNAntaraM 
16206 vA vaktavyam / sarve api ghaTAdayo bhAvAH kRtrimA akRtrimA  
16207 prasajanti / tatra apy evaM vikalpanAyAH sambhavAt / vizeSAbhAvAc  
16208 ca / tAn api hi parakriyAdarzanapUrvakam eva anyaH karoty avidita
16209 kartAraz ca kecid iti sarveSAM keSAMcid vA akriyAbhinivezo astu /  
16210 tasmAt sarvA eva iyaM varNAnupUrvI prasiddhakAryakAraNabhAva  
16211 vastudharmAnatikramAt puruSakRtA / ata eva /
16212 	asAdhAraNatA siddhA mantrAkhyakramakAriNAm /  
16213 	puMsAM jJAnaprabhAvAbhyAm anyeSAM tadabhAvataH //308//
16214 ayaM kramo varNAnAM viSanirghAtAdisamartho na anya iti yady 
16215 anyo api jAnIyat taM tathA eva pratipadyeta / na ca evam / tasmAd  
16216 ayam anukramaH svabhAvato api kAryakRt kaizcid eva vijJAta ity  
16217 asti parokSArthadarzI puruSaH / na hy ayam arthaH samartho
16218 na ayam iti zakyam unnetum asaMkIrNasya liGgavizeSasya asiddheH / 
16219 pratyakSayor apy anupadiSTayor mantrAmantrayor aparijJAnAt / 
16220 upadeze api kathaMcit svabhAvavivekApratIter anyatra kArya
16221 saMvAdAt tasya ca karaNAt prAg draSTum azakyatvAt / na
16222 ca ayam anukramaH svabhAvataH kArakaH kasyacid Azusiddher 
16223 anyasya cirAd aparasya vratacaraNAdyapekSaNAt / ekasmAd api 
16224 karmaNaH kayozcid arthAnarthasandarzanAt / vahatAm api
16225 mantrANAM punaH kvacid visaMvAdAt / na hy ayaM prakAraH  
16226 svabhAve yuktaH / svabhAvasya sarvatra avizeSAt / puruSas  
16227 tu svaicchApravRttiH sattvasabhAgatAAdivazAt sA iva avizeSAd vA
16301 kaMcid anugRhNAti na aparam iti yuktam / vratacaryAbhraMza  
16302 AdinA dharmAdharmaupacaye dharmAdharmAtmanor vA prakRtyA
16303 siddhyasiddhI iti cet / na / dharmaviruddhAnAm api kraurya 
16304 steyamaithunahInakarmAdibahulAnAM vratAnAM DAkinIbhaginI 
16305 tantrAdiSu darzanAt / taiz ca siddhivizeSAt / na ca evaMvidho 
16306 dharmasvabhAva iti ca yathAvasaraM nivedayiSyAmaH / maitrI 
16307 zaucadharmaparAyaNAnAM ca tannimittam eva kasyAzcid siddher  
16308 asiddher viparyaye ca punaH siddheH / na ca ekarUpAt karmaNaH 
16309 sa tadvirodhI dharmo yukto adharmaz ca / katham idAnIM dharma  
16310 phalam iSTam adharmAtmano vratAder aznute / na vai tasya eva  
16311 tad iSTaM phalaM vratAder vipAko api tu pUrvasya karmaNaH / 
16312 brahmahatyAdezAnuSThAnAd grAmapratilambhavat / tasya tv
16313 adharmAtmano vratasya AgAmi phalam aniSTam / sa tu mantrAdi
16314 prayogas tasya iSTaphalasya karmaNaH kathaMcid upakArAt pAcakaz 
16315 citratvAd upakArakazakteH / puruSavizeSAzrayavipAka 
16316 dharmA sa dharmas tena kRtaH sa tathA tadArAdhanena phalati iti / 
16317 tatprayogaupakAravipAkadharmaNaH kRtatvAt tatphalasya karmaNaH /  
16318 vinA api puruSeNa tadupakArAt phalam iti cet / na / puruSAkAra
16319 svabhAvacaryAdhimuktivaiyarthyaprasaGgAt / tasya apy upakAratve
16320 siddhaH puruSavizeSo asAdhAraNaguNaH / tadadhimukter  
16321 eva hi viSakarmAdikaraNAt / tasmAt ca mantrAH puruSapranItA 
16322 api tadupayoganirapekSAH svabhAvena phaladAH /  
16323 	ye api tantravidaH kecin mantrAn kAMzcana kurvate / 
16324 	prabhuprabhAvas teSAM sa taduktanyAyavRttitaH //309//
16325 rathyApuruSA api kecana tantrajJAH svayaMkRtair mantraiH  
16326 kiMcit karma kurvanti / tathA anyo apy anatizayaz ca kartA ca  
16327 mantrANAm iti / na / teSAM prabhAvavatA eva adhiSThAnAt / tatkRtaM
16328 hi te samayam anupAlayantas tadupadezena ca vartamAnAH samarthAH /
16329 tatsamayaupadezanirapekSANAm asAmarthyAt / tatra api
16401 tadAkAradhyAnAder eva prayogAt / tasmAt tadadhiSThAnam eva 
16402 tat tAdRzam ity unneyam / api ca / so api tAdRzaH prabhAvavAn
16403 eva ananyasAdhAraNazaktitvAd iti puruSavizeSa eva samarthitaH 
16404 	kRtakAH pauruSeyAz ca vAcyA mantrAH phalaIpsunA / 
16405 na hy akRtakAnAM prayogaH sambhavati na ca aprayuktebhyaH
16406 phalam iti prayogAt phalam icchatA kRtakA mantrA vAcyAH pauruSeyAz 
16407 ca / puruSAdhiSThAnam antareNa anyato asambhavatphalAnAM
16408 phaladarzanAt / kRtasamayakAvyAdivat /
16409 	azaktisAdhanaM puMsAm anena eva nirAkRtam //310// 
16410 pratipAditA hi puruSakRtAs tadadhiSThAnAc ca phaladA mantrAH /
16411 tad asti kazcid atizayavAn iti tatpratikSepasAdhanAny api prativyUDhAni /
16413 	buddhiindriyauktipuMstvAdi sAdhanaM yat tu varNyate /
16414 	pramANAbhaM yathArthA asti na hi zeSavato gatiH //311//  
16415 yat tu buddhiindriyavacanayogAt puMstvAd iti puruSAtizayapratikSepa
16416 sAdhanaM tat tv agamakam eva / pratikSepasAmAnya
16417 sAdhanayor asambhavAt / na hy atIndriyeSv ataddarzinaH pratikSepaH 
16418 sambhavati / satAm apy eSAm ajJAnAt / ata eva virodhAsiddheH /
16419 avirodhinA ca saha sambhavAvirodhAd ity apy uktam / 
16420 na api itarasAmAnyasiddhir vizeSAsambhavasya jJAtum azakyatvAt /  
16421 IdRzeSu ca anupalabdher hetutvapratikSepAt / puMstvAdisAmye 
16422 api kasyacid vizeSasya darzanAt / sambhavadvizeSe ca  sAmyAsiddhir
16423 ity uktam / tasmAc zeSavad anumAnam etat / vyatirekasya  
16424 sandehAd asamartham adarzane api vipakSavRtteH / api  
16425 ca evaMvAdino jaiminIyAH svam eva vAdaM svavAcA vidhurayanti /  
16426 tathA hi /
16501 	artho ayaM na ayam artho na iti zabdA vadanti na /
16502 	kalpyo ayam arthaH puruSais te ca rAgAdisaMyutAH //312// 
16503 	tatra ekas tattvavin na anya iti bhedaz ca kiMkRtaH /
16504 	tadvat puMstve katham api jJAnI kazcit kathaM na vaH //313// 
16506 na khalv ete vaidikAH zabdA evaM vikrozanti eta bhavanto brAhmaNA  
16507 ayam asmAkam artho grAhyo na anya iti kevalam anabhivyakta 
16508 arthavizeSasaMsargAH zrutim abhipatanti / tatra ekaH puruSaH  
16509 kaMcid arthaM kalpayaty anyo aparam / na ca zabdAnAM  
16510 kazcit svabhAvapratiniyamo yena ekam artham anurundhate na aparam /  
16511 kevalaM samayavazAt taM tam Avizanto dRzyante / teSAm  
16512 aviditArthaniyamAnAm atyakSAvezAd avidvAn eva doSaupaplavaH  
16513 kazcit tattvaM vyAcaSTe na apara iti na nyAyyam / atha
16514 kutazcid atizayAd buddhiindriyAdInAM sa eva vetti na aparaH / tasya
16515 kuto ayam atIndriyajJAnAtizayaH / tathA anyo api draSTA dezakAla
16516 svabhAvaviprakRSTAnAm arthAnAM kim asambhavI dRSTaH / na  
16517 hi tatpratikSepasAdhanAni kAnicid yAni na enam upalIyante /  
16518 yathA ayaM tatsAdhanasambhave apy asya vizeSas tathA anyasya api  
16519 syAd ity anabhiniveza eva yuktaH /
16520 	yasya pramANasaMvAdi vacanaM so arthavid yadi / 
16521 	na hy atyantaparokSeSu pramANasya asti sambhavaH //314//
16522 syAd etan na vayaM puruSaprAmANyAt kasyacid vyAkhyAnam 
16523 abhiniviSTAH / kiM tarhi / pramANAntarasaMvAdAd / bahuSv api 
16524 vyAkhyAtRSu yaH pramANaM pratyakSAdikaM saMsyandayati so
16525 anumanyate / tan na / atIndriyeSv adRSTAdiSu pramANAntarAvRtteH /
16526 tadasambhavAd eva hy Agamas tatpratipattyartham upayAcyate /  
16527 anyathA saty api tasmin pramANAntarAvRttAv apratipatteH / tataz 
16601 ca kevalAd arthapratipatter asAdhanam eva AgamaH syAt / kevalAd anyato
16602 apy atIndriyeSv apratipattir iti cet / katham atIndriyaz ca nAma 
16603 pratyakSAdiviSayaz ca / te punaH svaviSaye apy Agamam apekSya eva
16604 sAdhakAz cet / anAgamAd dhUmAder agnyAdipratyayo na syAt / 
16605 na vai pravRtta Agame pramANAntaram anviSyate kiM tarhi
16606 sa eva AgamapravRttir na jJAyata iti cet / svayaM samarthasya prasAdhane
16607 asya tadAgamaupadhAnaM kam atizayaM puSNAti / asamarthaM 
16608 tv AgamapravRttim api na eva sAdhayiSyati / sA ca atIndriya
16609 arthasambaddhA AgamapravRttir atIndriyA katham anyena
16610 siddhA / anyac ca evam AgamalakSaNaM syAt / tathA hi /
16611 	yasya pramANasaMvAdi vacanaM tatkRtaM vacaH /  
16612 	sa Agama iti prAptaM nirarthA apauruSeyatA //315//
16613 tulye apy AgamavAde pramANabalAd Agamasya kvacid Agamatve
16614 pramANasaMvAdo vacanAnAm AgamalakSaNaM syAt / na apuruSa
16615 kriyA / tasyAH sarvArtheSu tulyatve api pramANAbAdhanAt pratipatteH /
16516 tadbhAve apy anyatra pramANAsaMvAdiny aniSTatvAt /  
16617 kiM ca /
16618 	yady atyantaparokSe arthe anAgamajJAnasambhavaH / 
16619 	atIndriyArthavit kazcid asti ity abhimataM bhavet //316// 
16620 yady AgamAnapekSaM jJAnAyAthAtathyaM puruSasya iSyate parokSe  
16621 arthe santi puruSA atIndriyArthadRza iti iSTaM syAt / pratyakSa
16622 pUrvakAnAM pramANAnAm ataddarzane asambhavAt / pratyakSa
16624 pramANam AgamaH / pramANAntaravRttis tu pratyakSam
16625 anvAkarSati iti na puruSAtizayo nivAryaH syAt / tasmAn na asty atIndriyeSu
16626 pramANAntaravRttiH / ata eva Agamasya arthavizeSavRtter
16627 aparijJAnAd ayaM jaiminir anyo vA  
16628 	svayaM rAgAdimAn na arthaM vetti vedasya na anyataH /  
16701 	na vedayati vedo api vedArthasya kuto gatiH //317// 
16702 sarva eva hi puruSo anatikrAntadoSaviplavas tam atIndriyam
16703 arthavizeSapratiniyamaM vyAkhyAtA na svayaM vetti / na apy enam
16704 anyo vedayati / tasya api tulyaprasaGgatvAt / na hy andhena AkRzyamANo 
16705 andhaH panthAnaM pratipadyate / na api svayaM vedaH 
16706 svArthaM vivRNoti / upadezavaiyarthyaprasaGgAt / tad ayam aparijJAta
16707 arthaH zabdagaDur evaM zalyabhUto asaddarzanasnAyu 
16708 vinibaddho duruddharo duHkham Asayati / 
16709 	tena agnihotraM juhuyAt svargakAma iti zrutau /  
16710 	khAdec zvamAMsam ity eSa na artha ity atra kA pramA //318//
16711 kvacid apy arthe pratyAsattiviprakarSarahitasya agnihotraM  juhuyAt
16712 svargakAma ity AdivAkyasya bhUtavizeSe yathAbhimataM 
16713 ghRtAdi prakSiped ity ayam arthaH na punaH zvamAMsaM khAded 
16714 iti na atizayaM pazyAmaH / nanv ayaM sarvatra samAnaH prasaGgaH
16715 parokSadaizikAnAM vacanAnAm arthaM yathAbhiprAyam idAnIntanAH 
16716 kiM samanuyanti Ahosvid viparyayam iti / na / upadeSTuH
16717 svAbhiprAyaprakAzanena sampradAyasambhavAt / na hy 
16718 ayam adaizikAnAM zabdAnAM sambhavati / lokapratyAyanAbhiprAyaz
16719 ca bruvANo lokasaMketaprasiddhim anupAlayati iti tato api 
16720 tadarthasiddhiH syAt na apauruSeyANAM zabdAnAm / tatra kasyacit 
16721 samIhAbhAvAt / api ca nyAyam eva anupAlayantaH paNDitA heyaupAdeya
16722 tadAzrayeSu saMghaTante na tu pravAdamAtreNa iti na samAnaH
16723 prasaGgaH / tac ca yathAvasaraM pratipAdayiSyAmaH / nanu
16724 kazcil lokasaMnivezAdir ayuktiviSayo api sambhAvanIyapuruSa
16725 vacanAd arthaH pratipadyate / na / apratyayAt / na hi kvacid
16726 askhalita iti sarvaM tathA / vyabhicAradarzanAt / tat pravRtter 
16801 avisaMvAdena vyAptyasiddhez ca / agatyA ca idam AgamalakSaNam 
16802 iSTam / na ato nizcayaH / tan na pramANam Agama ity apy uktam /
16803 apauruSeyAnAM zabdAnAm arthajJAnaM na sampradAyAn  
16804 na yukter na lokAd iti tatra apratipattir nyAyyA / tatra api /
16805 	prasiddho lokavAdaz cet /  
16806 pratipattihetuH /  
16807 	tatra ko atIndriyArthadRk / 
16808 	anekArtheSu zabdeSu yena artho ayaM vivecitaH //319//  
16809 na hy ayaM lokavyavahAro apauruSeyAc zabdArthasambandhAt /
16810 kiM tarhi / samayAt / svazAstrakArasamayAt pANinIyAdivyavahAravat /
16811 upadezApekSaNAt / na hy apauruSeye tasminn upadezo
16812 yuktaH / tasya kenacid ajJAnAt / atIndriyatvAt / aindriyakatve  
16813 svayaM pratipattiprasaGgAt / rUpAdivat / upadeze ca puruSANAM 
16814 svatantrANAM yathAtattvam upadezena avisaMvAdasya asiddher anAzvAsaH /
16815 vedavat tadvyAkhyAnam apy apauruSeyaM sampradAyAvicchedAd
16816 Agatam tato arthasiddhir iti cet / tasya api zabdAtmakatve  
16817 tulyaH paryanuyogaH katham asya artho vidita iti / puruSo  
16818 hi svayaM samitAnAM zabdAnAm arthaM zRGgagrAhikayA api tAvad
16819 abudhaM bodhayed ity asti pauruSeyAnAM zabdAnAm arthagatAv 
16820 upAyaH / apauruSeyas tu zabdo na evaM karoti / na ca asya kazcit
16821 kvacit sambandhaniyamaM jJAtum Iza ity apratipattir eva 
16822 tadarthasya / api ca vedas tadvyAkhyAnaM vA puruSeNa puruSAya upadizyamAnam
16823 anaSTasampradAyam eva anuvartata ity atra api samayaH 
16824 zaraNam / AgamabhraMzakAriNAm AhopuruSikayA taddarzana  
16825 vidveSeNa vA tatpratipannakhalIkArAya dhUrtavyasanena anyato
16826 vA kutazcit kAraNAd anyathA racanAdarzanAt / api ca atra  
16827 bhavAn svam eva mukhavarNaM svavAdAnurAgAn nUnaM vismRtavAn
16901 puruSo rAgAdibhir upapluto anRtam api brUyAd iti na asya
16902 vacanaM pramANam iti / tad iha api kiM na pratyavekSyate  
16903 sambhavati na vA iti / sa eva upadizann upaplavAd vedaM vedArthaM  
16904 vA anyathA apy upadized iti / zrUyante hi kaizcit puruSair utsannauddhRtAni
16905 zAkhAntarANi / idAnIm api kAnicid viralAdhyetRkANi /  
16906 tadvat pracurAdhyetRkANAm api kasmiMzcit kAle kathaMcit  
16907 saMhArasambhavAt / punaH sambhAvitapuruSapratyayAt pracuratA  
16908 upagamanasambhAvanAsambhavAt / teSAM ca punaH pratAnayitRRNAM
16909 kadAcid adhItavismRtAdhyanAnAm anyeSAM vA sambhAvanA  
16910 bhraMzabhayAdinA anyathA upadezasambhavAt / tatpratyayAc 
16911 ca tadbhaktAnAm avicAreNa pratipatter bahuSv adhyetRSu 
16912 sambhAvitAt puruSAd bahulaM pratipattidarzanAt / tato api kathaMcid  
16913 vipralambhasambhavAt / kiM ca / parimitavyAkhyAtR 
16914 puruSaparaMparAm eva ca atra bhavatAm api zRNumaH / tatra  
16915 kazcid dviSTAjJadhUrtAnAm anyatamaH syAd api ity anAzvAsaH / 
16916 tasmAn na apauruSeyAd vyAkhyAnAn na api sAmayikAl lokavyavahArAd
16917 vedArthasiddhiH / asAmayikatve api nAnArthAnAM zabdAnAM 
16918 vyavahAre darzanAt kasyacid aprasiddhArthasya aprasiddhasya vA 
16919 punar vyutpattidarzanena sarvatra tadAzaGkAnivRtteH / sarveSAM  
16920 yathArthaniyoge apy avaiguNyena yathAsamayaM pratIti 
16921 jananAt / iSTAniSTayor avizeSAt / aviziSTAnAM sarvArtheSv ekam artham
16922 atyakSasaMyogam anatyakSadarzini puruSasAmAnye ko vivecayed 
16923 yato lokAt pratItiH syAt / api ca / svayam apy ayaM na  
16924 sarvatra prasiddhim anusarati / yasmAt /
16925 	svargaurvazyAdizabdaz ca dRSTo arUDhArthavAcakaH / 
16926 anena eva nirvarNyamAnaH / manuSyAtizAyipuruSavizeSaniketo atimAnuSa
16927 sukhAdhiSThAno nAnAupakaraNaH svargaH tannivAsiny  
16928 apsarA urvazI nAma iti lokavAdaH / tam anAdRtyAnyAm eva arthakalpanAm
16929 ayaM kurvANaH zabdAntareSu kathaM prasiddhiM pramANayet / 
16930 tatra avirodhAd abhyupagama iti cet / na / atra apy atIndriye  
16931 virodhAsiddheH / anyatra apy avirodhasya duranvayatvAt / viruddhAm 
16932 apy agnihotrAt svargAvAptiM mAndyAd ayaM na lakSayed 
16933 api / virodhAvirodhau ca bAdhakasAdhakapramANavRttI / te ca atyakSe
17001 na abhimate / tat kathaM tadvazAt pratItiH / na ca vacana 
17002 vRtter eva avirodho anyatra api prasaGgAt / apauruSeya Agamas tasya  
17003 pravAdAd arthasiddhiH / tatra punar virodhacintAyAm anAzvAsa
17004 Agame syAt / saty api tasminn atathAbhAvAd arthasya apramANa 
17005 vRtter anyasya api zaGkanIyatvAt / yad uktam agnihotraM  
17006 juhuyAt svargakAma ity atra zvamAMsabhakSaNadezanAvikalpo 
17007 bhavatv iti sa na bhavati / pradezAntareSu tathA tasya carcanAt /  
17008 na / tasya arthAparijJAnAt / pradezAntareSv api tathAvidhArtha  
17009 kalpanAyA anivAryatvAt / yadi hi kvacid viditArtho ayam apauruSeyaH 
17010 zabdarAziH syAt tadA tato arthapratItiH syAt / te tu bAhulye
17011 apy andhA eva sarva iti yatheSTam praNIyante / tasmAt / 
17012 	zabdAntareSu tAdRkSu tAdRzy eva astu kalpanA //320//
17013 yAdRzy agnihotraM juhuyAt svargakAma ity asya vAkyasya / api 
17014 ca
17015 	prasiddhiz ca nRNAM vAdaH pramANaM sa ca na iSyate /  
17016 	tataz ca bhUyo arthagatiH kim etad dviSTakAmitam //321//
17017 na prasiddhir nAma anyA anyatra janapravAdAt / te ca sarve janA  
17018 rAgAdyavidyAparItatvAd asambhAvanIyayAthAtathyavacanAH /
17019 tad eSAM pravAdo na pramANam / na hi kasyacid api samyakpratipatter
17020 abhAve bAhulyam arthavad bhavati / pArazIkamAtRmithyA 
17021 cAravat / teSAm eva puruSANAM vacanAt punaH parokSArthasampratipattir
17101 iti kathaM tad eva yugapad dveSyaM ca kAmyaM ca /
17102 	atha prasiddhim ullaGghya kalpane na nibandhanam / 
17103 	prasiddher apramANatvAt tatgrahe kiM nibandhanam //322// 
17105 prAptapratilomanena anyatra pravRttir guNadoSasandarzanena yuktA  
17106 iti prasiddher anvaya iti cet / na / prApteH pramANavRttilakSaNatvAt /
17107 yat kiMcana grahaNaM hi prasiddhim apramANayatas 
17108 tanmukhena pratItiH / nyAyAt prAptipratiSedhAt / tulyA svapara  
17109 vikalpayor ubhayathA api vRttir iti kaH prasiddhAv anurodhaH / 
17110 api ca iyam
17111 	utpAditA prasiddhyA eva zaGkA zabdArthanizcaye /  
17112 	yasmAn nAnArthavRttitvaM zabdAnAM tatra dRzyate //323//
17113 na prasiddher ekArthanizcayaH zabdAnAM tata eva zaGkAutpatteH / 
17114 nAnArthA hi zabdA loke dRzyante / lokavAdaz ca pratItiH / ata eva
17115 nAnArthA iti tata ekArthaniyamo na yuktaH /  
17116 	anyathA asambhavAbhAvAn nAnAzakteH svayaM dhvaneH / 
17117 	avazyaM zaGkayA bhAvyaM niyAmakam apazyatAm //324//  
17118 ity antarazlokaH / tasmAd aviditArthavibhAgeSu zabdeSv ekam 
17119 artham atyakSasaMyogam anAlambanasamAropaM vinizcitya vyAcakSANo 
17120 jaiminis tadvyAjena svam eva matam Aha iti na tIrthakara 
17121 antarAd asya vizeSaM pazyAmaH / tathA hi / tadarthavacanavyApAra 
17122 zUnyasya tatsamAropeNa abhidhAnaM na svavacanam atizete /
17123 tatkAriNA kevalaM mithyAvinItatA eva AtmanaH samuddyotitA syAt  
17124 tathA hi /
17125 	eSa sthANur ayaM mArga iti vakti iti kazcana /
17126 	anyaH svayaM bravImi iti tayor bhedaH parikSyatAm //325// 
17201 nirabhiprAyavyApAravacane sthANau samAropyaupadizataH svatantrasya
17202 vA svayaM vacanaupagame na kazcid vizeSo anyatra jaDasya 
17203 pratipattimAndyAt / api ca ekArthaniyame saty enaM jaiminir
17204 jAnIyAt / sa eva zabdasya 
17205 	sarvatra yogyasya ekArthadyotane niyamaH kutaH /
17206 na hi zabdasya kazcid arthaH svabhAvaniyataH sarvatra yogyatvAt /
17207 ayogyatve ca tadapracyuter avidheyasya puruSANAM kvacid 
17208 upanayanApanayanAsambhavAt /  
17209 	jJAtA vA atIndriyAH kena vivakSAvacanAd Rte //326// 
17210 puruSapraNIte hi zabde kayAcid vivakSayA sa tAM kadAcit kvacin
17211 nivedayed api iti vivakSApUrvakAnAM zabdAnAm arthaniyamaH
17212 pratIyeta api / apauruSeye tu vidyamAno apy arthaniyamaH kathaM  
17213 vijJeyaH / svabhAvabhedasya abhAvAt / sati vA pratyakSasya svayaM 
17214 pratItiprasaGgAt / apratyakSe api kenacij jJAtum azakyatvAt / 
17215 na ca asti kazcid vizeSaH / sarvazabdA hi sarvArtha pratyAsattiviprakarSa
17216 rahitAH / tatas teSAm
17217 	vivakSA niyame hetuH saMketas tatprakAzanaH /
17218 	apauruSeye sA na asti tasya sA ekArthatA kutaH //327//  
17219 vivakSayA hi zabdo arthe niyamyate na svabhAvataH tasya  
17220 kvacid apratibandhena sarvatra tulyatvAt / yatra api pratibandhas  
17221 tadabhidhAnaniyamAbhAvAt / sarvazabdaiH karaNAnAm abhidhAna
17222 prasaGgAt / tasmAd vivakSAprakAzanAya abhiprAya nivedanalakSaNaH 
17223 saMketaH kriyate / apauruSeye tu na vivakSA na saMketaH
17224 kasyacid abhiprAyAbhAvAd iti na niyamo na tajjJAnam /  
17225 	svabhAvaniyame anyatra na yojyeta tayA punaH /  
17226 yadi saMketanirapekSaH svabhAvata eva artheSu zabdo nilInaH
17227 syAt uktam atra apratibandhAd aniyama iti / api ca / svAbhAvike 
17228 vAcyavAcakabhAve na punar vivakSayA yatheSTaM niyujyeta / 
17301 	saMketaz ca nirarthaH syAd /
17302 na hi svabhAvabheda indriyagamyaH svapratItau paribhASAAdikam 
17303 apekSate / nIlAdibhedavat / tadapekSapratItayas tu na vastusvabhAvAH /
17304 kiM tarhi / sAmayikA rAjacihnAdivat / yaz ca sAmayikaH  
17305 sa svabhAvaniyato ayuktas tasya icchAvRtteH / ata eva saMketAt 
17306 svabhAvavizeSasya 
17307 	vyaktau ca niyamaH kutaH //328// 
17308 svaicchAvRttiH saMketaH sa iha eva kartuM zakyate na anyatra iti 
17309 na uparodho asti / sa ca puruSaiH svaicchayA kriyamANas tam eva  
17310 svabhAvaM vyanakti na anyam iti na niyamo asti / 
17311 	yatra svAtantryam icchAyA niyamo nAma tatra kaH / 
17312 	dyotayet tena saMketo na iSTAm eva asya yogyatAm //329//  
17313 iti antarazlokaH /  
17314 	yasmAt kila IdRzaM satyaM yathA agniH zItanodanaH / 
17315 	vAkyaM vedaekadezatvAd anyad apy aparo abravIt //330//  
17316 anyas tv apauruSeyam AgamalakSaNaM parityajya anyathA prAmANyaM
17317 vedasya sAdhayitukAmaH prAha avitathAni vedavAkyAni 
17318 yatra apratipattir vedaekadezatvAt yathA agnir himasya  
17319 bheSajam ity AdivAkyam iti / tasya idam /
17320 	rasavat tulyarUpatvAd ekabhANDe ca pAkavat /  
17321 	zeSavad vyabhicAritvAt kSiptaM nyAyavidA IdRzam //331//
17322 svayam IdRzam AcAryeNa anumAnaM naiyAyikazeSavadanumAna 
17323 vyabhicAram udbhAvayatA tulyarUpatayA phalAnAM tulyarasa  
17324 sAdhanavad ekasthAlyantargamAd dRSTavad adRSTataNDulapAka
17325 sAdhanavac ca asAdhanam uktam / tadasAdhanatvanyAyaz ca
17326 pUrvam eva uktaH / uktaM ca idam AgamalakSaNam asmAbhiH tat tu  
17327 sarvasya zakyavicArasya viSayasya yathAsvaM pramANena vidhi 
17328 pratiSedhavizuddhau nAntarIyakatvAbhAve api zabdAnAm artheSu
17401 varaM saMzayitasya vRttiH tatra kadAcid avisaMvAdasambhavAt  
17402 na tv anyatra dRSTapramANauparodhasya puruSasya pravRttir iti /  
17403 yaH punaH prAkRtaviSayasya vahneH zItapratighAtasAmarthyasya 
17404 abhidhAnaM satyArtham upadarzya sarvaM satyArtham Aha
17405 zAstraM zakyaparicchede api viSaye pramANavirodhAd bahutaram
17406 ayuktam api / 
17407 	nityasya puMsaH kartRtvaM nityAn bhAvAn atIndriyAn /
17408 	aindriyAn viSamaM hetuM bhAvAnAM viSamAM sthitim //332//  
17409 	nivRttiM ca pramANAbhyAm anyad vA vyastagocaram /
17410 	viruddham AgamApekSeNa anumAnena vA vadat //333//
17411 	virodham asamAdhAya zAstrArthaM ca apradarzya saH / 
17412 	satyArthaM pratijAnAno jayed dhArSTyena bandhakIm //334//
17414 apracyutAnutpannapUrvApararUpaH pumAn kartA krameNa karmaNAM  
17415 karmaphalAnAM ca bhoktA samavAyikAraNAdhiSThAnabhAvAdinA  
17416 ity Aha vedaH tac ca ayuktam ity AveditaprAyam nityatvaM ca  
17417 keSAMcid bhAvAnAm akSaNikasya vastudharmAtikramAd ayuktam
17418 apratyakSAny eva hi sAmAnyAdIni pratyakSAni janmasthiti
17419 nivRttIz ca viSamAH padArthAnAm anAdheyavizeSasya prAg akartuH  
17420 parApekSayA janakatvam niSpatter akAryarUpasya Azrayavazena 
17421 sthAnam kAraNAc ca vinAza ity Adikam anyad api pratyakSa
17422 anumAnAbhyAM prasiddhiviparyayam AgamAzrayeNa ca anumAnena 
17423 bAdhitam agnihotrAdeH pApazodhanasAmarthyAdikam / tasya  
17424 evaMvAdino vedasya sarvatra zAstrazarIre pramANavirodham
17425 apratisamAdhAya sambandhAnuguNaupAyapuruSArthAbhidhAnAni ca  
17426 zAstradharmAn apradarzyAtyantaprasiddhaviSayasatyAbhidhAna 
17427 mAtreNa prajJAprakarSaduravagahagahane api niratyayatAM sAdhayitukAmo
17428 bandhakIm api prAgalbhyena vijayate / kAcit kila  
17429 bandhakI svayaM svAminA vipratipattisthAne dRSTvA upAlabdhA / sA 
17430 taM pratyuvAca / pazyata mAtaH puruSasya vaiparItyaM / mayi
17431 dharmapatnyAM pratyayam akRtvA AtmIyayor na itarAbhidhAnayor
17432 jalabudbudayoH / karoti / tena jaratkANena grAmyakASThahArakeNa
17501 prArthitA api na saMgatA / rUpaguNAnurAgena kila mantrimukhya
17502 dArakaM kAmaye aham iti / evaMjAtIyakam etad api vahneH
17503 zItapratIkAravacanena dRSTapramANavirodhasya apy atyantaparokSe
17504 arthe avisaMvAdAnumAnam /
17505 	sidhyed pramANaM yady evam apramANam atha iha kim / 
17506 	na hy ekaM na asti satyArthaM puruSe bahubhASiNi //335//
17507 yathA idam atiduSkaram atyantasatyAbhidhAnaM tathA atyantAsatya  
17508 abhidhAnam api / tatra ekasya vacanasya kathaMcit saMvAdena  aviziSTasya
17509 tadvacanarAzes tathAbhAve na kazcit puruSo anAptaH
17510 syAt / api ca / 
17511 	na ayaM svabhAvaH kAryaM vA vastUnAM vaktari dhvaniH / 
17512 	na ca tadvyatiriktasya vidyate avyabhicAritA //336//  
17513 na tAvad etad vacanaM vAcyAnAM svabhAvaH / na apy eSAM kAryam
17514 tadabhAve api vaktur vivakSAmAtreNa bhAvAt / na ca anyaH  
17515 kazcit kasyacid avyabhicArI / vyabhicAre ca tato anyathA api 
17516 tatsambhavAt tadbhAvAt tatpratItir ayuktA /  
17517 	pravRttir vAcakAnAM ca vAcyadRSTikRtA iti cet /
17518 syAd etat kAryam eva vacanaM vAcakasya vAcyadarzanavRtter /
17519 evaM sati /
17520 	parasparaviruddhArthA katham ekatra sA bhavet //337//  
17521 yady eSa pratiniyamo vAcyaM vastv antareNa zabdo na pravartata  
17522 iti / bhinneSu pravAdeSv ekatra vastuni viruddhasvabhAvaupasaMhAreNa 
17523 vacanavRttir na syAt / na hy ayaM sambhavo asti ekaH  
17524 zabdo niSparyAyaM nityaz ca syAd anityaz ca iti /  
17525 	vastubhir na AgamAs tena kathaMcin nAntarIyakAH /  
17526 	pratipattuH prasidhyanti kutas tebhyo arthanizcayaH //338//
17527 AgamaM pramANaM tadAdarzitArthapratipattaye ajJo janaH samanveSate 
17528 samadhigatayAthAtathyAnAm upadezAnapekSaNAt / ajJasya 
17529 ca atIndriyaguNapuruSavivecane asAmarthyAt / vacanAnAM
17601 samIhitArthasattAm antareNa api vRttiM pazyato bhavitavyam eva adRSTa
17602 vyabhicAravacasAm api puruSANAM vAci zaGkayA kiM yathArthA 
17603 na vA iti / tena na yuktam anena kasyacid vacanena kiMcin 
17604 nizcetum / 
17605 	tasmAn na tannivRttyA api bhAvAbhAvaH prasidhyati /  
17606 yad uktaM sarvaviSayatvAd Agamasya sati vastuny avisaMvAdena 
17607 asya vRttes tannivRttilakSaNAnupalabdhir abhAvaM sAdhayati 
17608 iti tad asya sarvaviSayatve api vastvantareNa AvRttau syAt / tac
17609 ca na asti / tataH pratipattukAmasya asiddhir ity uktam / 
17610 	tena asaMnizcayaphalA anupalabdhir na sidhyate //339// 
17611 tasmAn na pramANatrayanivRttAv api viprakRSTeSv abhAvanizcayaH /
17613 	vedaprAmANyaM kasyacit kartRvAdaH / 
17614 	snAne dharmaicchA jAtivAdAvalepaH /  
17615 	santApArambhaH pApahAnAya ca iti /  
17616 	dhvastaprajJAne paJca liGgAni jADye //340//  
17617 iti pramANavArttike prathamaH paricchedaH // //  